\documentclass[11pt,openany,UTF8]{ctexbook}
\usepackage{amsmath,amssymb}
\usepackage{array,booktabs,longtable}
\usepackage{fancyvrb}
\usepackage{xcolor}
\usepackage[a4paper,top=2.4cm,bottom=2.4cm,left=2.3cm,right=2.3cm,headsep=0.6cm]{geometry}
\usepackage{fancyhdr}
\usepackage{needspace}
\usepackage{enumitem}
\usepackage[hidelinks,bookmarksnumbered=true]{hyperref}
\usepackage{graphicx}
\providecommand{\iddots}{\reflectbox{$\ddots$}}
\providecommand{\xlongequal}[1]{\xrightarrow{\ #1\ }}
\providecommand{\rank}{\operatorname{rank}}
\providecommand{\tr}{\operatorname{tr}}


% ---- fonts
\setCJKmainfont{Songti SC}[BoldFont=Heiti SC]
\xeCJKsetup{CJKmath=true}
\xeCJKDeclareCharClass{CJK}{"2500 -> "257F, "25A0 -> "25FF, "2160 -> "2183, "2190 -> "21FF}
\setCJKsansfont{Heiti SC}
\setCJKmonofont{Songti SC}
\setmainfont{Times New Roman}

% ---- colours (restrained)
\definecolor{accent}{HTML}{1A1A1A}
\definecolor{keybar}{HTML}{1F4E79}
\definecolor{warnbar}{HTML}{9C2A20}
\definecolor{exambar}{HTML}{4A4A4A}
\definecolor{notebar}{HTML}{2E6B4F}
\definecolor{codebg}{HTML}{F4F4F4}
\definecolor{rulesep}{HTML}{BFBFBF}

% ---- table column type
\newcolumntype{L}[1]{>{\raggedright\arraybackslash}p{#1}}
\newcommand{\ccaret}{\ifmmode\hat{}\else\textasciicircum{}\fi}
\newcommand{\circled}[1]{\textcircled{\scriptsize\,#1}}
\newcommand{\ctilde}{\ifmmode\tilde{}\else\textasciitilde{}\fi}

% ---- code blocks
\DefineVerbatimEnvironment{codebox}{Verbatim}%
  {frame=single,framerule=0.3pt,rulecolor=\color{rulesep},%
   framesep=6pt,fontsize=\small,baselinestretch=0.98,%
   xleftmargin=4pt,xrightmargin=4pt}

% ---- callout environments (breakable: only a bar + rule, no box)
\newcommand{\callouthead}[2]{%
  \par\addvspace{7pt}\noindent
  {\color{#1}\rule[-0.15em]{3pt}{1.05em}}\hspace{0.5em}%
  {\bfseries\color{#1}#2}\par\nobreak\vspace{-0.35em}%
  {\color{rulesep}\hrule height 0.4pt}\nobreak\vspace{0.45em}%
  \begingroup\small\setlength{\parindent}{0pt}%
  \leftskip=1.1em\rightskip=0.6em\noindent\ignorespaces}
\newcommand{\calloutend}{\par\endgroup\addvspace{7pt}}

\newenvironment{callkey}[1]{\callouthead{keybar}{#1}}{\calloutend}
\newenvironment{callwarn}[1]{\callouthead{warnbar}{#1}}{\calloutend}
\newenvironment{callexam}[1]{\callouthead{exambar}{#1}}{\calloutend}
\newenvironment{callnote}[1]{\callouthead{notebar}{#1}}{\calloutend}

% ---- headings
\ctexset{
  chapter = {format = \raggedright\Large\bfseries, name = {第,章}, number = \chinese{chapter},
             aftername = {\quad}, beforeskip = 0pt, afterskip = 18pt, pagestyle = plain},
  section = {format = \raggedright\large\bfseries, name = {,、}, number = \chinese{section},
             aftername = {}, beforeskip = 16pt, afterskip = 8pt},
  subsection = {format = \raggedright\normalsize\bfseries, name = {（,）}, number = \chinese{subsection},
             aftername = {}, beforeskip = 11pt, afterskip = 6pt},
  subsubsection = {format = \raggedright\normalsize\itshape, name = {,.}, number = \arabic{subsubsection},
             aftername = {}, beforeskip = 9pt, afterskip = 4pt},
}
\setcounter{secnumdepth}{3}
\setcounter{tocdepth}{2}
\setlist[itemize]{topsep=3pt,itemsep=1.5pt,parsep=0pt,leftmargin=1.5em}
\setlist[enumerate]{topsep=3pt,itemsep=1.5pt,parsep=0pt,leftmargin=1.8em}
\linespread{1.16}
\tolerance=3000
\emergencystretch=3em
\hbadness=10000
\vbadness=10000
\setlength{\parskip}{3.2pt}
\setlength{\parindent}{0pt}
\renewcommand{\arraystretch}{1.25}

% ---- header / footer
\pagestyle{fancy}
\fancyhf{}
\fancyhead[L]{\small\color{exambar}考研数学二 · 高等数学知识点手册}
\fancyhead[R]{\small\color{exambar}\leftmark}
\fancyfoot[C]{\small\thepage}
\renewcommand{\headrulewidth}{0.3pt}
\renewcommand{\headrule}{\hbox to\headwidth{\color{rulesep}\leaders\hrule height \headrulewidth\hfill}}
\fancypagestyle{plain}{\fancyhf{}\fancyfoot[C]{\small\thepage}\renewcommand{\headrulewidth}{0pt}}

\begin{document}


\begin{titlepage}
\centering
\vspace*{2.6cm}
{\Huge\bfseries 考研数学二\par}
\vspace{0.5cm}
{\Huge\bfseries 高等数学知识点手册\par}
\vspace{0.9cm}
{\Large 按《考研数学二考试大纲》编排 · 五星重难点分级\par}
\vspace{1.8cm}
{\normalsize
\begin{tabular}{rl}
考查科目 & 高等数学（数学二不考概率统计） \\
分值占比 & 约 80\%（150 分中约 120 分） \\
本册章节 & 5 章，覆盖大纲高等数学全部条目 \\
分级方式 & $\bigstar$ 五级：考频 $\times$ 难度 $\times$ 分值 \\
主要素材 & 高等数学大讲义 · 喻老强化讲义 · 中值定理专题 \\
 & 考研数学大观知识点 · 反常积分专题 · 数学二真题
\end{tabular}\par}
\vfill
{\small 本手册严格按《考研数学二考试大纲》顺序编排，每个考点标注五星重难点等级，供复习检索与精力分配使用。\par}
\vspace{1.0cm}
\end{titlepage}


\frontmatter
\pagestyle{fancy}
\chapter*{使用说明}
\addcontentsline{toc}{chapter}{使用说明}
\markboth{使用说明}{使用说明}
\section*{一、这本手册是什么}

这是一本\emph{严格按《考研数学二考试大纲》顺序}编排的高等数学知识点手册。它不替代教材与习题册，
解决的是三件事：

\begin{enumerate}
  \item \textbf{对照考纲查漏}：大纲上出现的每一条，本册都写到；大纲外的内容一律不收，避免把时间花在
        数一专属的级数、空间解析几何、三重积分与曲线曲面积分上。
  \item \textbf{按星级分配精力}：每个考点都标了五级星标，先攻$\bigstar\bigstar\bigstar\bigstar\bigstar$，再补$\bigstar\bigstar\bigstar\bigstar$，不要平均用力。
  \item \textbf{考前快速回顾}：公式、条件、易错点集中成块，配合每章末尾的易错清单，可在两小时内过完一章。
\end{enumerate}

\section*{二、五星分级标准}

星标取\textbf{考频 $\times$ 难度 $\times$ 分值}的综合，不是单纯的"重要程度"：

\begin{center}
\begin{tabular}{@{}clll@{}}
\toprule
\textbf{星标} & \textbf{名称} & \textbf{含义} & \textbf{复习要求} \\
\midrule
$\bigstar\bigstar\bigstar\bigstar\bigstar$ & 决胜考点 & 几乎年年必考，且是难题/易错点 & 闭卷默写 + 限时算对 \\
$\bigstar\bigstar\bigstar\bigstar$ & 高频重点 & 反复考，套路固定但计算量不小 & 熟练到条件反射 \\
$\bigstar\bigstar\bigstar$ & 常考点 & 隔年出现或以小题形式考 & 理解 + 会用 \\
$\bigstar\bigstar$ & 一般考点 & 低频，或只是解题的基础工具 & 会算即可 \\
$\bigstar$ & 了解性考点 & 大纲要求"了解"，考到即送分 & 记住结论 \\
\bottomrule
\end{tabular}
\end{center}

每章开头的"星级总表"给出该章全部考点的星级、大纲要求层次、难度与常见题型，可以先看表再决定精读哪几节。

\section*{三、试卷结构与分值}

数学二满分 150 分，考试时间 180 分钟，闭卷笔试。内容结构：\textbf{高等数学约 80\%}、
\textbf{线性代数约 20\%}，不考概率论与数理统计。题型结构：

\begin{center}
\begin{tabular}{@{}lll@{}}
\toprule
\textbf{题型} & \textbf{题量} & \textbf{分值} \\
\midrule
单项选择题 & 10 小题 & 每小题 5 分，共 50 分 \\
填空题 & 6 小题 & 每小题 5 分，共 30 分 \\
解答题（包括证明题） & 6 小题 & 共 70 分 \\
\bottomrule
\end{tabular}
\end{center}

由此可得一条实用结论：\textbf{选择填空共 80 分，占比过半}，而它们考的是基本概念、基本计算与常见陷阱。
本册中标注$\bigstar\bigstar\bigstar$以上、题型为"选择/填空"的考点，是性价比最高的得分点。

\section*{四、数学二不考什么（去糟粕清单）}

以下内容属于数一/数三或大纲外，本册一律不收，复习时可以直接跳过：

\begin{itemize}
  \item 向量代数与空间解析几何（平面、直线、曲面、曲线）
  \item 无穷级数（数项级数、幂级数、傅里叶级数）
  \item 三重积分、曲线积分、曲面积分、高斯公式、斯托克斯公式、场论
  \item 含参变量积分、$\Gamma$ 函数等分析类拓展
  \item 微分方程的欧拉方程、方程组、级数解法、拉普拉斯变换、差分方程
  \item 概率论与数理统计（数学二完全不考）
\end{itemize}

在线性代数部分，数学二与数一范围基本一致（行列式、矩阵、向量、线性方程组、特征值与特征向量、
二次型六章全考），但\textbf{不考}线性空间与坐标变换、线性变换、若尔当标准形等更深的定理。

\section*{五、怎么用：三轮法}

\begin{enumerate}
  \item \textbf{第一轮（打基础）}：按章节顺序读，重点看"定义与结论"和公式的\textbf{适用前提}，
        每个$\bigstar\bigstar\bigstar\bigstar$以上的考点都动手算一遍例题。
  \item \textbf{第二轮（补漏洞）}：只看每章的"星级总表"和"易错清单"，把$\bigstar\bigstar$以下的快速扫过，
        把$\bigstar\bigstar\bigstar\bigstar$以上的重新推一遍；错题回到对应考点的warning框。
  \item \textbf{第三轮（考前）}：只看$\bigstar\bigstar\bigstar\bigstar\bigstar$考点与公式，配合本册的公式速查部分，
        限时做真题，检查是否还有"知道套路但算错"的考点。
\end{enumerate}

\section*{六、使用提醒}

\begin{itemize}
  \item 本册按大纲顺序编排，因此\textbf{章号即大纲条目号}（如"一元函数微分学"对应大纲第二部分）。
  \item 手册中的"真题考法"只描述\textbf{出题角度}，不编造年份与题号；确需核实的年份，请以真题册为准。
  \item 所有公式都以考纲与统考习惯写法为准；不同教材符号不一致处，正文中已注明"统考按 \dots 理解"。
  \item 手册是复习的\textbf{索引与骨架}，不能替代刷题。看完一个考点，请立刻用对应真题或习题验证。
\end{itemize}

\cleardoublepage
\tableofcontents
\mainmatter

\chapter{函数极限连续}
\markboth{函数极限连续}{函数极限连续}
\subsection*{本章考纲要求}
\addcontentsline{toc}{subsection}{本章考纲要求}

\begin{quote}
\textbf{考试内容}

\begin{enumerate}
\item 函数的概念及表示法
\item 函数的有界性、单调性、周期性和奇偶性
\item 复合函数、反函数、分段函数和隐函数
\item 基本初等函数的性质及其图形
\item 初等函数
\item 函数关系的建立
\item 数列极限与函数极限的定义及其性质
\item 函数的左极限和右极限
\item 无穷小量和无穷大量的概念及其关系
\item 无穷小量的性质及无穷小量的比较
\item 极限的四则运算
\item 极限存在的两个准则：单调有界准则和夹逼准则
\item 两个重要极限：$\lim_{x\to 0}\frac{\sin x}{x}=1$，$\lim_{x\to\infty}\left(1+\frac1x\right)^x=e$
\item 函数连续的概念
\item 函数间断点的类型
\item 初等函数的连续性
\item 闭区间上连续函数的性质
\end{enumerate}
\textbf{考试要求}

\begin{enumerate}
\item 理解函数的概念，掌握函数的表示法，会建立应用问题的函数关系。
\item 了解函数的有界性、单调性、周期性和奇偶性。
\item 理解复合函数及分段函数的概念，了解反函数及隐函数的概念。
\item 掌握基本初等函数的性质及其图形，了解初等函数的概念。
\item 理解极限的概念，理解函数左极限与右极限的概念以及函数极限存在与左极限、右极限之间的关系。
\item 掌握极限的性质及四则运算法则。
\item 掌握极限存在的两个准则，并会利用它们求极限，掌握利用两个重要极限求极限的方法。
\item 理解无穷小量、无穷大量的概念，掌握无穷小量的比较方法，会用等价无穷小量求极限。
\item 理解函数连续性的概念（含左连续与右连续），会判别函数间断点的类型。
\item 了解连续函数的性质和初等函数的连续性，理解闭区间上连续函数的性质（有界性、最大值和最小值定理、介值定理），并会应用这些性质。
\end{enumerate}
\end{quote}
\subsection*{本章星级总表}
\addcontentsline{toc}{subsection}{本章星级总表}

\begingroup\small\setlength{\tabcolsep}{4.0pt}
\begin{longtable}{@{}L{0.4484\linewidth}L{0.1114\linewidth}L{0.1114\linewidth}L{0.1114\linewidth}L{0.1114\linewidth}@{}}
\toprule
\textbf{考点} & \textbf{星级} & \textbf{考纲要求} & \textbf{难度} & \textbf{常见题型} \\
\midrule
\endhead
函数的概念、表示法与函数关系的建立 & $\bigstar$ & 理解/掌握 & 基础 & 小题 \\
函数的有界性、单调性、周期性、奇偶性 & $\bigstar$$\bigstar$ & 了解 & 基础 & 选择 \\
复合函数、反函数、分段函数、隐函数 & $\bigstar$$\bigstar$ & 理解/了解 & 基础 & 小题 \\
基本初等函数的性质与图形、初等函数 & $\bigstar$$\bigstar$ & 掌握/了解 & 基础 & 小题 \\
数列极限的定义与性质 & $\bigstar$$\bigstar$$\bigstar$ & 理解 & 中等 & 选择 \\
函数极限的定义、左极限与右极限、极限存在的充要条件 & $\bigstar$$\bigstar$ & 理解 & 基础 & 选择 \\
极限的性质、四则运算与"存在性传播"规律 & $\bigstar$$\bigstar$$\bigstar$ & 掌握 & 中等 & 选择 \\
无穷小与无穷大的概念及关系 & $\bigstar$$\bigstar$$\bigstar$ & 理解 & 基础 & 选择、填空 \\
无穷小阶的比较与确定（含求参数） & $\bigstar$$\bigstar$$\bigstar$$\bigstar$ & 掌握 & 中等 & 选择、填空 \\
等价无穷小替换及其适用条件 & $\bigstar$$\bigstar$$\bigstar$$\bigstar$$\bigstar$ & 掌握 & 中等 & 选择、填空、解答 \\
极限存在的两个准则：夹逼准则与单调有界准则 & $\bigstar$$\bigstar$$\bigstar$$\bigstar$ & 掌握 & 难 & 解答、证明 \\
两个重要极限 & $\bigstar$$\bigstar$$\bigstar$$\bigstar$ & 掌握 & 基础 & 填空、解答 \\
七类未定式与极限计算总思路 & $\bigstar$$\bigstar$$\bigstar$$\bigstar$$\bigstar$ & 掌握 & 中等 & 解答 \\
数列极限的计算：n 项和与 n 项积、递推数列 & $\bigstar$$\bigstar$$\bigstar$$\bigstar$ & 掌握 & 难 & 解答 \\
泰勒/等价无穷小在极限中的联合运用 & $\bigstar$$\bigstar$$\bigstar$$\bigstar$$\bigstar$ & 掌握 & 难 & 解答 \\
已知极限确定参数、求函数表达式 & $\bigstar$$\bigstar$$\bigstar$$\bigstar$ & 掌握 & 中等 & 填空、解答 \\
函数的连续性与左、右连续 & $\bigstar$$\bigstar$$\bigstar$ & 理解 & 基础 & 选择 \\
间断点及其类型（第一类/第二类） & $\bigstar$$\bigstar$$\bigstar$ & 会判别 & 中等 & 选择、填空 \\
连续函数的性质与初等函数的连续性 & $\bigstar$$\bigstar$ & 了解 & 基础 & 选择 \\
闭区间上连续函数的性质（有界性、最值定理、介值定理、零点定理） & $\bigstar$$\bigstar$$\bigstar$ & 理解 & 中等 & 选择、证明 \\
\bottomrule
\end{longtable}
\endgroup

\section{函数}

\subsection{函数的概念与表示}

\subsubsection{函数的概念、表示法与函数关系的建立 $\bigstar$}

\textbf{定义}：设 $D$ 是非空实数集，若存在对应法则 $f$，使对每个 $x\in D$ 都有唯一确定的实数 $y$ 与之对应，则称 $f$ 为定义在 $D$ 上的函数，记作 $y=f(x)$。其中 $D$ 称为定义域，$R_f=\{f(x)\mid x\in D\}$ 称为值域。

三个要素：\textbf{定义域、对应法则、值域}（其中值域由前两者确定）。判定两个函数是否为同一函数，只看定义域与对应法则是否都相同，与自变量、因变量用什么字母无关。

\textbf{求定义域的优先顺序（"定义域优先原则"）}：

\begin{itemize}
\item 分式的分母不为零；
\item 偶次根式的被开方数非负；
\item 对数的真数大于零（底数大于零且不等于 1）；
\item $\arcsin u$、$\arccos u$ 要求 $\vert u\vert\le 1$；
\item $\tan u$ 要求 $u\ne k\pi+\frac{\pi}{2}$；
\item 多个条件同时出现时取交集。
\end{itemize}
对于复合函数 $f(g(x))$，定义域由 $x\in D_g$ 且 $g(x)\in D_f$ 共同决定，即"由外向内解不等式链"。

\textbf{解题套路（求表达式）}

\begin{enumerate}
\item 已知 $f(g(x))$ 的表达式求 $f(x)$：令 $t=g(x)$，反解 $x=g^{-1}(t)$ 代入，再把 $t$ 换成 $x$；
\item 已知 $f(x)$ 满足某个含 $f(-x)$、$f(\frac1x)$ 的方程，用"换元构造联立方程"消去多余未知函数；
\item 求分段函数在分段点处的值，必须严格按分界条件代入对应分支。
\end{enumerate}
\textbf{例}：求 $f(x)=\frac{\sqrt{4-x^2}}{\ln(x+1)}$ 的定义域。 需要 $4-x^2\ge 0\Rightarrow -2\le x\le 2$；$x+1>0\Rightarrow x>-1$；$\ln(x+1)\ne 0\Rightarrow x\ne 0$。 三者取交，定义域为 $(-1,0)\cup(0,2]$。

\textbf{关于"函数关系的建立"（应用问题）}

考纲要求"会建立应用问题的函数关系"，落到题目上就是两类工作：一是由几何关系（周长、面积、体积、相似三角形）写出目标量关于自变量的表达式；二是由物理或经济关系（路程与时间、成本与产量）写出函数关系。写法固定为三步：

\begin{enumerate}
\item 设自变量（通常是"所求最值的那个量"或题中给定变化的量），并\textbf{写出它的实际取值范围}；
\item 用几何/物理关系把其余量都化成自变量的表达式；
\item 写出目标函数并注明定义域。
\end{enumerate}
\textbf{例}：用长为 $L$ 的篱笆围成一个靠墙的矩形（墙的一边不用篱笆），设垂直于墙的边长为 $x$，写出矩形面积 $S$ 与 $x$ 的函数关系。 平行于墙的边长为 $L-2x$，故 $S(x)=x(L-2x)$，定义域为 $0<x<\frac L2$（同时要求 $L-2x>0$）。\textbf{定义域不能省略}，否则后续求最值时会出现不合理的驻点。

\textbf{求复合函数定义域的三种题型}

\begingroup\small\setlength{\tabcolsep}{4.0pt}
\begin{longtable}{@{}L{0.1894\linewidth}L{0.2367\linewidth}L{0.1894\linewidth}L{0.2959\linewidth}@{}}
\toprule
\textbf{题型} & \textbf{已知} & \textbf{求} & \textbf{做法} \\
\midrule
\endhead
正向 & $f(x)$ 的定义域 $D_f$ & $f(g(x))$ 的定义域 & 解 $g(x)\in D_f$ \\
反向 & $f(g(x))$ 的定义域 $D$ & $f(x)$ 的定义域 & 求 $g(x)$ 在 $x\in D$ 上的值域 \\
构造 & $f(g(x))$ 的定义域 $D_1$ & $f(h(x))$ 的定义域 & 先用反向求出 $D_f$，再解 $h(x)\in D_f$ \\
\bottomrule
\end{longtable}
\endgroup

\textbf{例}：设 $f(x-1)$ 的定义域为 $[-1,3]$，求 $f(3x)$ 的定义域。 $x\in[-1,3]\Rightarrow x-1\in[-2,2]$，故 $D_f=[-2,2]$。再令 $3x\in[-2,2]$，得 $x\in\left[-\frac23,\frac23\right]$。

\begin{callwarn}{易错}
定义域必须写成\textbf{区间或集合}，不能只写不等式；由 $\ln(x+1)\ne 0$ 只能推出 $x\ne 0$（不是 $x>0$），漏掉这一条是本考点最常见的失分点。

\end{callwarn}

\begin{callexam}{真题考法}
通常以填空或选择的"送分小题"出现，考"已知 $f(g(x))$ 的表达式反求 $f(x)$"或"求复合函数的定义域"。陷阱集中在：反解时忘记限制 $t$ 的取值范围；把 $\ln$ 的真数条件与函数值条件混为一谈。

\end{callexam}

\subsection{函数的四大特性}

\subsubsection{函数的有界性、单调性、周期性、奇偶性 $\bigstar$$\bigstar$}

\textbf{有界性}：若存在 $M>0$，使对一切 $x\in I$ 有 $\vert f(x)\vert\le M$，则称 $f$ 在 $I$ 上有界；否则无界。判断无界的标准手法是\textbf{找一列点} $\{x_n\}\subset I$ 使 $\vert f(x_n)\vert\to+\infty$。

\[f \text{ 在 } I \text{ 上有界} \iff \exists M>0,\ \forall x\in I,\ \vert f(x)\vert\le M .\]

常用放缩不等式（证有界、做夹逼都要用）：

\[\vert\sin x\vert\le 1,\quad \vert\cos x\vert\le 1,\quad \vert\arctan x\vert<\frac{\pi}{2},\quad \frac{\vert x\vert}{1+x^2}\le\frac12 .\]

\textbf{单调性}：对区间 $I$ 上任意的 $x_1<x_2$，若 $f(x_1)\le f(x_2)$ 则称单调增加；若 $f(x_1)<f(x_2)$ 则称严格单调增加。\textbf{注意}：本考点只要求"了解"，判别单调性的主要工具（导数符号法）放在一元微分学一章，此处只需会按定义验证。

\textbf{奇偶性}：定义域必须关于原点对称。

\begin{itemize}
\item 偶函数：$f(-x)=f(x)$，图形关于 $y$ 轴对称；
\item 奇函数：$f(-x)=-f(x)$，图形关于原点对称。
\end{itemize}
奇偶性的运算规律（"奇奇得偶"等）：

\begingroup\small\setlength{\tabcolsep}{4.0pt}
\begin{longtable}{@{}L{0.1462\linewidth}L{0.1462\linewidth}L{0.1462\linewidth}L{0.1462\linewidth}L{0.1462\linewidth}L{0.1462\linewidth}@{}}
\toprule
\textbf{运算} & \textbf{奇 $\pm$ 奇} & \textbf{偶 $\pm$ 偶} & \textbf{奇 $\times$ 奇} & \textbf{偶 $\times$ 偶} & \textbf{奇 $\times$ 偶} \\
\midrule
\endhead
结果 & 奇 & 偶 & 偶 & 偶 & 奇 \\
\bottomrule
\end{longtable}
\endgroup

复合函数的奇偶性判据：\textbf{内偶则偶，内奇同外}——即若 $g$ 为偶函数，则 $f(g(x))$ 为偶函数；若 $g$ 为奇函数，则 $f(g(x))$ 的奇偶性与 $f$ 相同。

\textbf{周期性}：若存在 $T>0$ 使 $f(x+T)=f(x)$ 对定义域内一切 $x$ 成立，则称 $f$ 为周期函数。常用结论：若 $f$ 以 $T$ 为周期，则 $f(ax+b)$ 以 $\frac{T}{\vert a\vert}$ 为周期；$\sin x,\cos x$ 的周期为 $2\pi$，$\tan x$ 的周期为 $\pi$。

\textbf{例}：判断 $f(x)=\ln\left(x+\sqrt{1+x^2}\right)$ 的奇偶性。 $f(-x)=\ln\left(-x+\sqrt{1+x^2}\right)=\ln\frac{1}{x+\sqrt{1+x^2}}=-\ln\left(x+\sqrt{1+x^2}\right)=-f(x)$，故为奇函数。这个函数是反双曲正弦函数 $\operatorname{arsinh}x$，是必须记住的"奇函数加无穷阶可导"的典型代表。

\textbf{奇偶性与周期性、有界性的常用组合结论}

\begin{itemize}
\item 若 $f$ 是周期为 $T$ 的连续函数，则 $f$ 在 $\mathbb{R}$ 上有界（只需在一个周期 $[0,T]$ 上用闭区间连续函数的有界性定理）；
\item 若 $f$ 为奇函数且在 $x=0$ 处有定义，则必有 $f(0)=0$；
\item 偶函数的图形关于 $y$ 轴对称，因此在 $\int_{-a}^{a}$ 上偶函数积分为 $2\int_0^a$，奇函数积分为 $0$，这条技巧在一元积分学中反复使用；
\item $f(x)+f(-x)$ 必为偶函数，$f(x)-f(-x)$ 必为奇函数，故任何定义域对称的函数都可分解为偶函数与奇函数之和：
\end{itemize}
\[f(x)=\underbrace{\frac{f(x)+f(-x)}{2}}_{\text{偶部分}}+\underbrace{\frac{f(x)-f(-x)}{2}}_{\text{奇部分}} .\]

\textbf{判断有界的常用手法}

\begin{enumerate}
\item \textbf{直接放缩}：把函数拆成有界因子的乘积，如 $\left\vert\frac{\sin x}{x}\right\vert\le 1\ (x\ne 0)$；
\item \textbf{极限控制法}：若 $\lim_{x\to a^+}f(x)$ 与 $\lim_{x\to b^-}f(x)$ 都存在，且 $f$ 在 $(a,b)$ 内连续，则 $f$ 在 $(a,b)$ 内有界（口诀："连续 + 两端极限存在 $\Rightarrow$ 有界"）；
\item \textbf{导数有界法}：若 $f'$ 在区间 $I$ 上存在且有界，则 $f$ 在 $I$ 上有界（由拉格朗日中值定理 $f(x)-f(x_0)=f'(\xi)(x-x_0)$ 放缩得到）；
\item \textbf{证明无界}：构造一列点 $x_n$ 使 $\vert f(x_n)\vert\to+\infty$。如 $f(x)=x\sin x$，取 $x_n=2n\pi+\frac{\pi}{2}$，则 $f(x_n)=2n\pi+\frac{\pi}{2}\to+\infty$，故无界。
\end{enumerate}
\begin{callwarn}{易错}
\begin{enumerate}
\item 判断奇偶性前\textbf{必须先检查定义域是否关于原点对称}，例如 $f(x)=x^2,\ x\in[0,1]$ 既非奇也非偶；
\item 偶函数若存在反函数，必为单射——所以"非平凡的偶函数没有反函数"；
\item 有界与"存在最大值"不是一回事：$f(x)=\arctan x$ 有界但在 $\mathbb{R}$ 上取不到最值。
\end{enumerate}
\end{callwarn}

\begin{callexam}{真题考法}
多以选择题出现，把奇偶性、周期性、有界性与原函数（变限积分）问题串起来考，例如"设 $f$ 为连续奇函数，问 $F(x)=\int_0^x f(t)\,dt$ 的奇偶性"。关键是记住：\textbf{连续奇函数的原函数是偶函数，连续偶函数的原函数中只有 $F(x)=\int_0^x f(t)\,dt$ 为奇函数}（因为此时 $F(0)=0$）。

\end{callexam}

\subsubsection{复合函数、反函数、分段函数、隐函数 $\bigstar$$\bigstar$}

\textbf{复合函数}：设 $y=f(u)$，$u=g(x)$，当 $g$ 的值域与 $f$ 的定义域有非空交集时，可定义复合函数 $y=f(g(x))$。复合顺序不可交换，$f(g(x))$ 与 $g(f(x))$ 一般不同。

\textbf{反函数}：若 $y=f(x)$ 在 $D$ 上严格单调，则存在反函数 $x=f^{-1}(y)$，且 $y=f(x)$ 与 $y=f^{-1}(x)$ 的图形关于直线 $y=x$ 对称。

\begin{itemize}
\item 求反函数的步骤：由 $y=f(x)$ 解出 $x=\varphi(y)$，再把 $x,y$ 互换写成 $y=f^{-1}(x)$，并注明其定义域（即原函数的值域）；
\item 奇函数的反函数（若存在）仍为奇函数。
\end{itemize}
\textbf{分段函数}：不同区间用不同表达式表示的函数，\textbf{分段点必须单独讨论}。求分段点的函数值、极限、连续性、可导性时，一律回到定义，分别从左、右两侧处理。

\textbf{隐函数}：由方程 $F(x,y)=0$ 确定的函数关系。本考点只要求"了解"概念，隐函数求导法放在微分学一章。

\textbf{例}：设 $f(x-1)=x^2-3x+2$，求 $f(x)$。 令 $t=x-1$，则 $x=t+1$，$f(t)=(t+1)^2-3(t+1)+2=t^2-t$，故 $f(x)=x^2-x$。

\begin{callwarn}{易错}
分段函数在分段点处不能直接代入某一分支求值，必须先看分段点归属哪个区间（必要时用极限与函数值分别计算）。另外，分段函数求复合时要把内层的取值范围代回外层，逐段确定表达式，不能只换主分支。

\end{callwarn}

\subsection{基本初等函数与初等函数}

\subsubsection{基本初等函数的性质与图形、初等函数 $\bigstar$$\bigstar$}

五类基本初等函数：\textbf{幂函数 $x^\mu$、指数函数 $a^x\ (a>0,a\ne 1)$、对数函数 $\log_a x$、三角函数、反三角函数}。由基本初等函数经过有限次四则运算与复合，且能用一个式子表示的函数称为\textbf{初等函数}。

必须熟记的图形特征与关键值：

\begingroup\small\setlength{\tabcolsep}{4.0pt}
\begin{longtable}{@{}L{0.2430\linewidth}L{0.2603\linewidth}L{0.4252\linewidth}@{}}
\toprule
\textbf{函数} & \textbf{定义域} & \textbf{关键性质} \\
\midrule
\endhead
$x^\mu$ & 视 $\mu$ 而定，均过 $(1,1)$ & $\mu>0$ 时在 $(0,+\infty)$ 单调增 \\
$a^x$ & $\mathbb{R}$，过 $(0,1)$ & $a>1$ 单调增，$0<a<1$ 单调减 \\
$\log_a x$ & $(0,+\infty)$，过 $(1,0)$ & 与 $a^x$ 互为反函数 \\
$\arcsin x$ & $[-1,1]$，值域 $\left[-\frac{\pi}{2},\frac{\pi}{2}\right]$ & 单调增，奇函数 \\
$\arccos x$ & $[-1,1]$，值域 $[0,\pi]$ & 单调减 \\
$\arctan x$ & $\mathbb{R}$，值域 $\left(-\frac{\pi}{2},\frac{\pi}{2}\right)$ & 单调增，奇函数，有界 \\
\bottomrule
\end{longtable}
\endgroup

\textbf{必须记牢的两个恒等式}：

\[\arcsin x+\arccos x=\frac{\pi}{2}\ (x\in[-1,1]),\qquad \arctan x+\operatorname{arccot} x=\frac{\pi}{2}\ (x\in\mathbb{R}).\]

\textbf{需要用图形记牢的三组结论}

\begin{enumerate}
\item \textbf{幂函数与指数、对数的增长对比}：当 $x\to+\infty$ 时，$\ln x\ll x^{\alpha}\ll a^{x}$（$\alpha>0$，$a>1$）。这条结论是"抓大头"和"无穷大比较"的全部依据。
\item \textbf{反三角函数的全局形态}：$\arcsin x$ 在 $[-1,1]$ 上从 $-\frac{\pi}{2}$ 单调升到 $\frac{\pi}{2}$；$\arctan x$ 在 $\mathbb{R}$ 上有两条水平渐近线 $y=\pm\frac{\pi}{2}$，这是"$\arctan x$ 有界但取不到端点值"的图形依据。
\item \textbf{分段定义的幂函数}：$x^{a}$ 的图形随 $a$ 的符号变化很大，$a<0$ 时在 $x\to 0^+$ 处趋于 $+\infty$，$0<a<1$ 时在 $x=0$ 处不可导（此结论在微分学中使用）。
\end{enumerate}
\textbf{常用初等不等式的图像记忆（做夹逼与证明有界必备）}

\[\vert\sin x\vert\le\vert x\vert\ (x\in\mathbb{R}),\qquad e^{x}\ge 1+x\ (x\in\mathbb{R}),\qquad \frac{x}{1+x}\le\ln(1+x)\le x\ (x>-1),\]

\[\sin x\le x\le\tan x\ \left(0\le x<\frac{\pi}{2}\right),\qquad \cos x\ge 1-\frac{x^{2}}{2}\ (x\in\mathbb{R}).\]

\begin{callwarn}{易错}
$\sqrt{x^2}=\vert x\vert$，不是 $x$；$e^{\ln x}=x$ 只在 $x>0$ 时成立，而 $\ln e^x=x$ 对一切 $x$ 成立。求极限时把 $\sqrt{x^2}$ 直接写成 $x$ 是高频错误（在 $x\to-\infty$ 时符号恰好相反）。

\end{callwarn}

\section{极限}

\subsection{极限的概念与性质}

\subsubsection{数列极限的定义与性质 $\bigstar$$\bigstar$$\bigstar$}

\textbf{定义（$\varepsilon$-$N$ 语言）}：设 $\{x_n\}$ 为数列，$a$ 为常数。若对任意给定的 $\varepsilon>0$，总存在正整数 $N$，使当 $n>N$ 时恒有 $\vert x_n-a\vert<\varepsilon$，则称数列 $\{x_n\}$ 收敛于 $a$，记为 $\lim_{n\to\infty}x_n=a$。用"通俗语言"说：\textbf{当 $n$ 充分大时，$x_n$ 与 $a$ 要多近有多近。}

用定义证 $\lim_{n\to\infty}x_n=a$ 的四步：

\begin{enumerate}
\item 写出 $\vert x_n-a\vert$ 的表达式；
\item 适当放缩，把它放大成便于解不等式的形式（如 $\frac{1}{n}$、$\frac{C}{n}$）；
\item 解不等式 $\vert x_n-a\vert<\varepsilon$，求出 $n$ 的范围 $n>N(\varepsilon)$；
\item 取 $N=\left[\ \right]+1$（取整加一），再按定义顺序书写。
\end{enumerate}
\textbf{收敛数列的基本性质}：

\begin{itemize}
\item \textbf{唯一性}：极限存在则唯一；
\item \textbf{有界性}：收敛数列必有界（反之不成立，如 $\{(-1)^n\}$）；
\item \textbf{保号性}：若 $\lim_{n\to\infty}x_n=a>0$，则存在 $N$，当 $n>N$ 时 $x_n>0$；反之，若 $x_n\ge 0$ 且极限存在，则 $a\ge 0$；
\item \textbf{子列性质}：$\{x_n\}$ 收敛于 $a$ $\iff$ 它的任一子列都收敛于 $a$。
\end{itemize}
由此得到两个最常用的\textbf{判发散}手法：若某个子列极限不存在，则原数列极限不存在；若有两个子列极限存在但不相等，则原数列极限不存在。例如 $\{(-1)^n\}$ 的奇子列趋于 $-1$、偶子列趋于 $1$，故发散。

数列极限存在的充要条件（奇偶子列形式）：

\[\lim_{n\to\infty}x_n=A \iff \lim_{n\to\infty}x_{2n}=\lim_{n\to\infty}x_{2n-1}=A .\]

\textbf{用定义证明时的放缩技巧（"放大法"）}

直接解 $\vert x_n-a\vert<\varepsilon$ 往往很难，标准做法是\textbf{把 $\vert x_n-a\vert$ 放大成一个更容易处理的形式}，只要放大后的量能任意小，原量也就能任意小。常用放大手段：

\begin{itemize}
\item 分母放大：$\frac{1}{n^{2}+k}\le\frac1{n^{2}}$，$\frac{n}{n^{2}+k}\le\frac{1}{n}$；
\item 分子放大：$\vert\sin n\vert\le 1$，$\vert(-1)^{n}\vert=1$；
\item 分段放大：当 $n\ge 2$ 时 $\frac{1}{n}\le\frac{1}{2}$，把含 $n$ 的因子拆成"前有限项 + $n$ 很大时的项"分别控制。
\end{itemize}
\textbf{例}：用定义证明 $\lim_{n\to\infty}\frac{\sin n}{n}=0$。 任给 $\varepsilon>0$，取 $N=\left[\frac1\varepsilon\right]+1$。当 $n>N$ 时，

\[\left\vert\frac{\sin n}{n}-0\right\vert=\frac{\vert\sin n\vert}{n}\le\frac1n<\varepsilon .\]

故 $\lim_{n\to\infty}\frac{\sin n}{n}=0$。

\begin{callwarn}{易错}
\begin{enumerate}
\item "数列有界"推不出"数列收敛"，只能由收敛推有界；
\item 保号性中，\textbf{由极限大于零只能推出$n$ 充分大后数列大于零}，不能推出所有项都大于零；反向由 $x_n>0$ 只能得到 $a\ge 0$，不能得到 $a>0$；
\item 用定义证明时，$N$ 必须明显依赖于 $\varepsilon$，写"取 $N$ 充分大"是不合格的。
\end{enumerate}
\end{callwarn}

\begin{callexam}{真题考法}
选择题常考"由 $\lim x_n=a$ 能推出什么"：等价变形有 $\lim x_{n+1}=a$、$\lim x_{n+k}=a$（$k$ 固定）、$\lim\vert x_n\vert=\vert a\vert$（反之不成立）。也常用反例法——命题对"$f$ 单调有界但 $f(x_n)$ 未必收敛"的判错（取 $f(x)=\arctan x$、$x_n=n$）。

\end{callexam}

\subsubsection{函数极限的定义、左极限与右极限、极限存在的充要条件 $\bigstar$$\bigstar$}

\textbf{六个极限过程}：$x\to x_0$、$x\to x_0^+$、$x\to x_0^-$、$x\to\infty$、$x\to+\infty$、$x\to-\infty$。其中

\[\lim_{x\to x_0}f(x)=A \iff \forall\varepsilon>0,\ \exists\delta>0,\ \text{当 } 0<\vert x-x_0\vert<\delta \text{ 时},\ \vert f(x)-A\vert<\varepsilon .\]

注意 $0<\vert x-x_0\vert$ 表示\textbf{去心邻域}，即 $f$ 在 $x_0$ 处有无定义、取值多少都不影响极限。

\textbf{极限存在的充要条件（左右极限）}：

\[\lim_{x\to x_0}f(x) \text{ 存在} \iff \lim_{x\to x_0^-}f(x)=\lim_{x\to x_0^+}f(x)=A .\]

对无穷远过程：$\lim_{x\to\infty}f(x)=A \iff \lim_{x\to-\infty}f(x)=\lim_{x\to+\infty}f(x)=A$。

\textbf{海涅定理（函数极限与数列极限的关系）}：若 $\lim_{x\to x_0}f(x)=A$，$\{x_n\}$ 是 $f$ 定义域内满足 $x_n\ne x_0$ 且 $x_n\to x_0$ 的任一数列，则 $\lim_{n\to\infty}f(x_n)=A$。常用其\textbf{逆否命题}：若能找到两个趋近于 $x_0$ 的数列使 $f(x_n)$ 的极限不同（或不存在），则 $\lim_{x\to x_0}f(x)$ 不存在。

\begin{callwarn}{易错}
含 $\vert x-a\vert$、$e^{1/x}$、$\arctan\frac1x$、$\frac{1}{x-a}$ 的函数在 $x=a$ 处\textbf{必须分左右讨论}：$\lim_{x\to 0^+}e^{1/x}=+\infty$，而 $\lim_{x\to 0^-}e^{1/x}=0$，两侧完全不同。另外，极限存在与函数在该点是否有定义无关。

\end{callwarn}

\begin{callexam}{真题考法}
选择题常以"下列命题中正确的是"的形式考查概念：把 $f(x_0)$ 与 $\lim_{x\to x_0}f(x)$ 存在的因果关系、左右极限与 $\vert x-a\vert$ 的关系、$x\to\infty$ 与 $x\to\pm\infty$ 的关系混在一起设置陷阱。看到"极限存在"要先问"在哪个过程下、是双侧还是单侧"。

\end{callexam}

\subsubsection{极限的性质、四则运算与"存在性传播"规律 $\bigstar$$\bigstar$$\bigstar$}

设在同一极限过程中 $\lim f(x)=A$，$\lim g(x)=B$，则

\[\lim[f(x)\pm g(x)]=A\pm B,\quad \lim[f(x)g(x)]=AB,\quad \lim\frac{f(x)}{g(x)}=\frac{A}{B}\ (B\ne 0).\]

\textbf{使用前提}：拆开后的每一部分极限都必须存在；商的法则还要求分母极限不为零。幂与根式：$\lim[f(x)]^n=A^n$，当 $A\ge 0$ 时 $\lim\sqrt[n]{f(x)}=\sqrt[n]{A}$。

\textbf{极限的性质}：唯一性、局部有界性、局部保号性、保序性。

\begin{itemize}
\item \textbf{局部有界性}：若 $\lim_{x\to x_0}f(x)=A$，则存在 $\delta>0$ 与 $M>0$，当 $0<\vert x-x_0\vert<\delta$ 时 $\vert f(x)\vert\le M$；
\item \textbf{局部保号性}：若 $A>0$，则存在 $\delta>0$，当 $0<\vert x-x_0\vert<\delta$ 时 $f(x)>0$；
\item \textbf{保序性}：若在去心邻域内 $f(x)\ge g(x)$ 且两侧极限存在，则 $A\ge B$（不等号在取极限后可能变成等号）。
\end{itemize}
\textbf{"存在性传播"规律（务必背下来）}：

\begingroup\small\setlength{\tabcolsep}{4.0pt}
\begin{longtable}{@{}L{0.5629\linewidth}L{0.3828\linewidth}@{}}
\toprule
\textbf{运算} & \textbf{结论} \\
\midrule
\endhead
存在 $\pm$ 存在 & 存在 \\
存在 $\pm$ 不存在 & 不存在 \\
不存在 $\pm$ 不存在 & 不确定 \\
存在 $\times$ 存在 & 存在 \\
存在且非零 $\times$ 不存在 & 不存在 \\
存在且为零 $\times$ 不存在 & 不确定 \\
不存在 $\times$ 不存在 & 不确定 \\
\bottomrule
\end{longtable}
\endgroup

\textbf{两个高频结论}（来自商法则的反向使用）：

\begin{itemize}
\item 若 $\lim\frac{f(x)}{g(x)}=A$ 存在，且 $\lim g(x)=0$，则 $\lim f(x)=0$；
\item 若 $\lim\frac{f(x)}{g(x)}=A\ne 0$，且 $\lim f(x)=0$，则 $\lim g(x)=0$。
\end{itemize}
\textbf{复合函数的极限运算法则}

设 $y=f[g(x)]$ 在 $x_0$ 的某去心邻域内有定义，若

\[\lim_{x\to x_0}g(x)=u_0,\qquad \lim_{u\to u_0}f(u)=A,\]

且存在 $\delta_0>0$，当 $0<\vert x-x_0\vert<\delta_0$ 时 $g(x)\ne u_0$，则

\[\lim_{x\to x_0}f[g(x)]=\lim_{u\to u_0}f(u)=A .\]

其中的条件"$g(x)\ne u_0$"是关键：它保证 $x\to x_0$ 时内层的取值确实能遍历 $u_0$ 的去心邻域。若内层函数恰为常数 $u_0$，结论就不成立——这是概念题的常考细节。特别地，当外层 $f$ 在 $u_0$ 处连续时，条件"$g(x)\ne u_0$"可以去掉，直接有 $\lim f(g(x))=f(\lim g(x))$。

\textbf{例}：已知 $\lim_{x\to 1}[f(x)+g(x)]=3$，$\lim_{x\to 1}[f(x)-g(x)]=1$，求 $\lim_{x\to 1}f(x)g(x)$。 两式相加得 $\lim 2f(x)=4$，故 $\lim f(x)=2$；相减得 $\lim 2g(x)=2$，故 $\lim g(x)=1$。于是 $\lim f(x)g(x)=2\times 1=2$。

\begin{callwarn}{易错}
\begin{enumerate}
\item 最常见的错误是\textbf{随手拆极限}：在 $\lim[f(x)+g(x)]$ 中若 $\lim f(x)$ 不存在，绝不能拆成两项分别求；
\item "存在且为零 $\times$ 不存在"是不确定——如 $x\to 0$ 时 $x\cdot\frac{1}{x^2}$ 中 $\frac{1}{x^2}$ 不存在，但乘积仍不存在；而 $x\cdot\sin\frac1x\to 0$ 存在，所以不能一概而论；
\item 由 $f(x)\ge g(x)$ 取极限只能得到 $A\ge B$，不能得到 $A>B$（如 $f(x)=1+\frac1x$，$g(x)=1$）。
\end{enumerate}
\end{callwarn}

\begin{callexam}{真题考法}
选择题的高频考法：给出 $\lim f(x)$、$\lim g(x)$ 的存在性，让判断 $\lim[f(x)g(x)]$、$\lim\frac{f(x)}{g(x)}$ 是否存在；或给出"$\lim\frac{f(x)}{x}=2$"这类商极限，要求用两个高频结论反推 $f(x)$ 的极限、甚至反求 $f(0)$ 与 $f'(0)$。\textbf{看到"分母趋于零而分式极限存在"，第一反应就是"分子也趋于零"。}

\end{callexam}

\subsection{无穷小与无穷大}

\subsubsection{无穷小与无穷大的概念及关系 $\bigstar$$\bigstar$$\bigstar$}

\textbf{无穷小}：若在某个极限过程中 $\lim f(x)=0$，则称 $f(x)$ 是该过程下的无穷小量。\textbf{无穷小是一个函数（变量），不是"一个很小的数"}——说"$10^{-100}$ 是无穷小"是错的。

\textbf{无穷小的性质}：

\begin{itemize}
\item 有限个无穷小的和、积仍是无穷小；
\item \textbf{有界函数与无穷小之积仍是无穷小}（这一条最常用，如 $\lim_{x\to 0}x\sin\frac1x=0$）；
\item 无穷小一定有界；有界变量不一定是无穷小；
\item $\lim f(x)=A \iff f(x)=A+\alpha(x)$，其中 $\alpha(x)$ 是同一过程下的无穷小。
\end{itemize}
\textbf{无穷大}：若对任意 $M>0$，存在 $X>0$（或 $\delta>0$），使当 $\vert x\vert>X$（或 $0<\vert x-x_0\vert<\delta$）时 $\vert f(x)\vert>M$，则称 $f(x)$ 为该过程下的无穷大，记作 $\lim f(x)=\infty$。

\textbf{无穷小与无穷大的关系}（同一过程下）：

\[\lim f(x)=\infty \iff \lim\frac{1}{f(x)}=0\quad(\text{要求 } f(x)\ne 0).\]

\textbf{无穷大与无界的区别（高频概念题）}：

\begin{itemize}
\item \textbf{无界}：只要\textbf{存在某个点} $x_0$，使 $\vert f(x_0)\vert$ 可以任意大（即存在一列点使函数值趋于无穷）；
\item \textbf{无穷大}：要求\textbf{所有充分接近该点的点}的函数值都要多大有多大。
\end{itemize}
所以"无穷大 $\Rightarrow$ 无界"，但"无界 $\not\Rightarrow$ 无穷大"。反例：$f(x)=x\sin x$ 在 $x\to+\infty$ 时无界（取 $x=2k\pi+\frac{\pi}{2}$），但不是无穷大（取 $x=2k\pi$ 时函数值为 $0$）。

\textbf{"伪无穷大"}：有些函数看起来在不断变大，但在某些点列上取不到无穷。判断 $\lim_{x\to+\infty}f(x)=\infty$ 是否成立，必须检查\textbf{所有}趋于 $+\infty$ 的路径，只要找到一条路径使 $\vert f(x)\vert$ 不趋于无穷，结论就不成立。$x\sin x$、$x(2+\sin x)$ 都是这类"伪无穷大"。

\textbf{无穷小的比较与无穷大的比较是对偶的}：$\alpha=o(\beta)$（$\alpha$ 比 $\beta$ 高阶的无穷小）等价于 $\frac1\alpha$ 是比 $\frac1\beta$ \textbf{高阶}的无穷大（在两者均不为零、且同号趋于无穷的前提下）。所以熟记 $\ln x\ll x^{\alpha}\ll a^{x}$ 就同时掌握了两个方向的比较。

\textbf{无穷大的阶（当 $x\to+\infty$ 时，从低到高）}：

\[\ln x \ll x^{\alpha}\ (\alpha>0)\ll a^{x}\ (a>1).\]

数列形式（当 $n\to\infty$ 时，从低到高）：

\[\ln n \ll n^{k}\ (k>0)\ll a^{n}\ (a>1)\ll n!\ll n^{n}.\]

\begin{callwarn}{易错}
\begin{enumerate}
\item 无穷小与无穷大都必须\textbf{指明极限过程}（$x\to 0$、$x\to\infty$ 等），脱离过程谈无穷小没有意义；
\item "$0$"是无穷小，但"很小的数"不是无穷小；
\item 无穷大是\textbf{极限不存在}的一种情形（趋势为无穷），不能参与通常的四则运算，$\infty-\infty$、$\frac{\infty}{\infty}$、$0\cdot\infty$ 都是未定式。
\end{enumerate}
\end{callwarn}

\begin{callexam}{真题考法}
选择题常考"无界是否等价于无穷大"（答案：不等价，必须举 $x\sin x$ 这类反例）；填空题常考"$x\to\infty$ 时比较 $\ln x$、$x^{100}$、$e^{0.01x}$ 的阶"。结论记牢：\textbf{指数增长永远快于多项式，多项式永远快于对数}。

\end{callexam}

\subsubsection{无穷小阶的比较与确定（含求参数）$\bigstar$$\bigstar$$\bigstar$$\bigstar$}

设同一极限过程下 $\lim\alpha(x)=0$，$\lim\beta(x)=0$，且 $\beta(x)\ne 0$。

\[\lim\frac{\alpha}{\beta}=
\begin{cases}
0, & \text{则 } \alpha \text{ 是比 } \beta \text{ 高阶的无穷小，记 } \alpha=o(\beta);\\
c\ne 0, & \text{则 } \alpha \text{ 与 } \beta \text{ 是同阶无穷小};\\
1, & \text{则 } \alpha \text{ 与 } \beta \text{ 是等价无穷小，记 } \alpha\sim\beta;\\
\infty, & \text{则 } \alpha \text{ 是比 } \beta \text{ 低阶的无穷小}.
\end{cases}\]

\textbf{$k$ 阶无穷小}：若 $\lim\frac{\alpha(x)}{[x-x_0]^{k}}=c\ne 0$（$c$ 为有限非零常数），则称 $\alpha(x)$ 是关于 $x-x_0$ 的 $k$ 阶无穷小。

\textbf{确定阶数的套路}

\begin{enumerate}
\item \textbf{定义法}：设 $\alpha\sim c(x-x_0)^k$，计算 $\lim\frac{\alpha(x)}{(x-x_0)^k}$，令其等于非零有限常数，解出 $k$ 与 $c$；
\item \textbf{泰勒展开法}：把 $\alpha(x)$ 展开成幂函数的多项式（带余项），\textbf{消去相消项后的第一个非零项的幂次就是阶数}；
\item \textbf{已知 $o(\ )$ 关系求参数}：先由"分子极限必须为零"定出部分参数，再用泰勒展开比较同次幂系数。
\end{enumerate}
\textbf{高阶无穷小的运算律}（写成小 $o$ 的形式）：

\[o(x^m)\pm o(x^n)=o(x^{\min(m,n)}),\qquad o(x^m)\cdot o(x^n)=o(x^{m+n}),\qquad k\,o(x^n)=o(x^n).\]

一句话：\textbf{低阶吸收高阶}，即 $x^2+x^5\sim x^2$（$x\to 0$）。

\textbf{例}：当 $x\to 0$ 时，判断 $\cos x-e^{-x^2/2}$ 是 $x$ 的几阶无穷小。

\[\cos x=1-\frac{x^2}{2}+\frac{x^4}{24}+o(x^4),\qquad e^{-x^2/2}=1-\frac{x^2}{2}+\frac{x^4}{8}+o(x^4).\]

相减得 $\cos x-e^{-x^2/2}=\left(\frac{1}{24}-\frac{1}{8}\right)x^4+o(x^4)=-\frac{1}{12}x^4+o(x^4)$，故为 $4$ 阶无穷小。

\begin{callwarn}{易错}
\begin{enumerate}
\item 比较阶数必须\textbf{在同一极限过程下}进行，$x\to 0$ 与 $x\to\infty$ 的结论完全不同；
\item 写 $o(x^m)\pm o(x^n)=o(x^{\min(m,n)})$ 时不要颠倒——是\textbf{低阶（幂次小）的留下}；
\item 用"定义法"求阶数时，$k$ 必须是使极限为非零\textbf{有限}常数的那个值，取 $k$ 偏大或偏小都会得到 $0$ 或 $\infty$。
\end{enumerate}
\end{callwarn}

\begin{callexam}{真题考法}
填空题常直接问"当 $x\to 0$ 时，$f(x)$ 是 $x$ 的几阶无穷小"；选择或解答题常考"已知 $f(x)$ 是 $x$ 的 $k$ 阶无穷小，求其中的参数"。陷阱在于把等价无穷小直接代入\textbf{加减式}得到错误的阶数（如由 $\sin x\sim x$ 误判 $\sin x-x$ 是 $1$ 阶）。

\end{callexam}

\subsection{极限的计算工具}

\subsubsection{等价无穷小替换及其适用条件 $\bigstar$$\bigstar$$\bigstar$$\bigstar$$\bigstar$}

\textbf{常用等价无穷小（$x\to 0$ 时，务必默写）}：

\[\sin x\sim x,\quad \tan x\sim x,\quad \arcsin x\sim x,\quad \arctan x\sim x,\quad e^{x}-1\sim x,\quad \ln(1+x)\sim x,\]

\[a^{x}-1\sim x\ln a\ (a>0,\ a\ne 1),\quad 1-\cos x\sim\frac{x^2}{2},\quad (1+x)^{\alpha}-1\sim\alpha x\ (\alpha\ne 0),\]

\[\ln\!\left(x+\sqrt{1+x^2}\right)\sim x,\qquad (1+x)^{x}-1\sim x^{2}.\]

\textbf{差函数型（高频，必须记牢，符号最容易错）}：

\[x-\sin x\sim\frac{x^{3}}{6},\qquad \arcsin x-x\sim\frac{x^{3}}{6},\qquad \tan x-x\sim\frac{x^{3}}{3},\qquad x-\arctan x\sim\frac{x^{3}}{3},\]

\[x-\ln(1+x)\sim\frac{x^{2}}{2},\qquad e^{x}-1-x\sim\frac{x^{2}}{2}.\]

\begin{callkey}{必背：等价无穷小的"使用边界"}
\textbf{只有处于乘除（因子）位置的无穷小才能直接替换}；加减位置的无穷小必须整体替换或用泰勒。

\[\lim_{x\to 0}\frac{\sin x(1+\cos x)}{x}=\lim_{x\to 0}\frac{x(1+\cos x)}{x}=2\quad(\text{替换正确})\]

\[\lim_{x\to 0}\frac{\tan x-\sin x}{x^{3}}=\frac12\ \ne\ 0\quad(\text{若换成 } x \text{ 则错})\]

\end{callkey}

\textbf{替换的适用条件（两大核心问题）}

\begin{enumerate}
\item \textbf{何时能换}：当无穷小处于因子位置（乘除）时可以直接换；处于加减位置时一般不能单独换。若一定要在加减中处理，必须\textbf{整体替换}并保证精度——例如 $\lim_{x\to 0}\frac{\tan x-\sin x}{x^3}$ 中若把 $\tan x$、$\sin x$ 都换成 $x$，会得到错误答案 $0$；正确答案是 $\frac12$。
\item \textbf{换成谁}：必须换成与它\textbf{等价}的$(\ )$，并注意复合函数的内层不要单独换。对 $\alpha(x)\to 0$，若 $f(u)\sim u$（$u\to 0$），则 $f(\alpha(x))\sim\alpha(x)$——这叫"整体代换"，是正确用法。
\end{enumerate}
\textbf{极限计算的一般流程}

\[化简 \to 提取非零因子 \to 等价替换 \to 泰勒展开 \to 洛必达法则.\]

\textbf{例}：求 $\lim_{x\to 0}\frac{\sin x(1+\cos x)}{x}$。 $\sin x$ 处于分子因子位置，可直接换成 $x$：

\[\lim_{x\to 0}\frac{x(1+\cos x)}{x}=\lim_{x\to 0}(1+\cos x)=2 .\]

\textbf{等价替换的三种常用"变形后再替换"套路}

\begin{enumerate}
\item \textbf{加减项创造条件}：在分子上 $+1-1$ 或 $+x-x$，凑出可替换的结构。例如
\end{enumerate}
\[\lim_{x\to 0}\frac{e^{x}-1-x}{x^{2}}=\lim_{x\to 0}\frac{(e^{x}-1)-x}{x^{2}}=\lim_{x\to 0}\frac{x-x}{x^{2}}\ \text{（错！）}\]

必须改用泰勒：$e^{x}-1-x\sim\frac{x^{2}}{2}$，故极限为 $\frac12$。这个例子说明"加减中间项"不能只靠一阶等价。

\begin{enumerate}
\item \textbf{变限积分的整体等价}：若 $f(t)\sim g(t)$（$t\to 0$），则
\end{enumerate}
\[\int_0^{\varphi(x)}f(t)\,dt \sim \int_0^{\varphi(x)}g(t)\,dt \quad\bigl(\varphi(x)\to 0\bigr),\]

即"被积函数等价 $\Rightarrow$ 变限积分等价"。例如 $\lim_{x\to 0^+}\frac{\int_0^{\sin x}\ln(1+t)\,dt}{x^{2}}$：由 $\ln(1+t)\sim t$ 得积分 $\sim\int_0^{\sin x}t\,dt=\frac{\sin^{2}x}{2}\sim\frac{x^{2}}{2}$，故极限为 $\frac12$。

\begin{enumerate}
\item \textbf{对数中"先合并再替换"}：$\ln\frac{A}{B}=\ln A-\ln B$ 后往往能凑出差函数型，如
\end{enumerate}
\[\lim_{x\to 0}\frac{\ln(1+x)-x}{x^{2}}=\lim_{x\to 0}\frac{-\frac{x^{2}}{2}+o(x^{2})}{x^{2}}=-\frac12 .\]

\textbf{例}：求 $\lim_{x\to 0}\left(\frac{a^{x}+b^{x}+c^{x}}{3}\right)^{\frac1x}$（$a,b,c>0$）。 这是 $1^\infty$ 型。底数减 $1$：

\[\frac{a^{x}+b^{x}+c^{x}-3}{3}=\frac{(a^{x}-1)+(b^{x}-1)+(c^{x}-1)}{3}\sim\frac{x\ln a+x\ln b+x\ln c}{3}=\frac{x\ln(abc)}{3}.\]

故 $A=\lim_{x\to 0}\frac1x\cdot\frac{x\ln(abc)}{3}=\frac13\ln(abc)$，原极限为 $e^{\frac13\ln(abc)}=\sqrt[3]{abc}$，即三项的几何平均值。这个结论对 $n$ 项同样成立：$a_1,\dots,a_n>0$ 时 $\lim_{x\to 0}\left(\frac{a_1^{x}+\cdots+a_n^{x}}{n}\right)^{\frac1x}=\sqrt[n]{a_1a_2\cdots a_n}$。

\begin{callwarn}{易错}
\begin{enumerate}
\item \textbf{加减中替换}是本章最高频的错误来源：$\lim_{x\to 0}\frac{\tan x-\sin x}{x^3}$ 用"都换成 $x$"得到 $0$，与正确答案 $\frac12$ 相差甚远；
\item 复合函数的内层不能换：$\sin(\sin x)$ 不能写成 $\sin x$ 就算完，要继续换成 $x$；而 $\ln(\sin x)$ 在 $x\to 0^+$ 请先用 $\sin x\sim x$ 化为 $\ln x$，不要反过来；
\item 差函数型的\textbf{符号与系数}要逐个核对，$x-\sin x\sim\frac{x^3}{6}$，而 $\sin x-x\sim-\frac{x^3}{6}$；$\tan x-x\sim\frac{x^3}{3}$，而 $x-\tan x\sim-\frac{x^3}{3}$；
\item 替换前先"提取极限非零因子"：$\lim_{x\to 0}\frac{e^x\sin x}{x(1+x)}$ 中 $e^x\to 1$、$\frac{1}{1+x}\to 1$ 可以直接提出去，让式子变简单。
\end{enumerate}
\end{callwarn}

\begin{callexam}{真题考法}
解答题第一题几乎必考，标准解法就是"化简 $\to$ 等价替换 $\to$ 泰勒"三步。真题的陷阱设置有两种：一是把待替换的无穷小放在加减位置，迫使考生改用泰勒或通分；二是把等价无穷小藏在变限积分里，此时要求"被积函数比值极限为 $1$"才可作整体等价。\textbf{只要出现加减相消，就老老实实写泰勒展开。}

\end{callexam}

\subsubsection{极限存在的两个准则：夹逼准则与单调有界准则 $\bigstar$$\bigstar$$\bigstar$$\bigstar$}

\textbf{夹逼准则}

\begin{itemize}
\item 数列情形：若存在 $n_0$，当 $n>n_0$ 时 $x_n\le y_n\le z_n$，且 $\lim_{n\to\infty}x_n=\lim_{n\to\infty}z_n=a$，则 $\lim_{n\to\infty}y_n=a$。
\item 函数情形：若在 $x_0$ 的某去心邻域内有 $g(x)\le f(x)\le h(x)$，且 $\lim_{x\to x_0}g(x)=\lim_{x\to x_0}h(x)=A$，则 $\lim_{x\to x_0}f(x)=A$。
\end{itemize}
使用的关键是\textbf{"两边同时逼近"}：一边缩（取和式的最小项或积的最小项），一边放（取最大项），必须使两边极限相同。常用放缩工具：

\begin{itemize}
\item 分数放缩：分子越大、分母越小，分数越大；
\item 抓大头：$n+\frac1i$ 抓 $n$，$n+\sqrt i$ 抓 $n$，$n^2+n$ 抓 $n^2$；
\item 均值不等式：$\frac{n}{\sum\frac{1}{a_k}}\le\sqrt[n]{a_1a_2\cdots a_n}\le\frac{a_1+\cdots+a_n}{n}$（正数）；
\item 常见不等式：$\sin x\le x\ (x\ge 0)$，$e^{x}\ge 1+x$，$\frac{x}{1+x}\le\ln(1+x)\le x\ (x>-1)$。
\end{itemize}
\textbf{单调有界准则}：单调增加且有上界的数列必收敛；单调减少且有下界的数列必收敛。

处理递推数列 $x_{n+1}=f(x_n)$ 的四步：

\begin{enumerate}
\item \textbf{证有界}：用数学归纳法（先猜出上界，再验证 $x_{k+1}\le M$）；
\item \textbf{证单调}：比较 $x_{n+1}-x_n$ 与 $0$（或比较 $\frac{x_{n+1}}{x_n}$ 与 $1$）；也可用 $f'(x)>0$ 且 $x_1<x_2$ 推单调增；
\item \textbf{设极限为 $A$}：在递推式两端取极限得 $A=f(A)$；
\item \textbf{取舍}：若方程有多个根，用单调性与有界范围排除。
\end{enumerate}
\textbf{例}：求 $\lim_{n\to\infty}\sqrt[n]{1+2^{n}+3^{n}}$。 设 $y_n=\sqrt[n]{1+2^n+3^n}$，则 $3=\sqrt[n]{3^n}\le y_n\le\sqrt[n]{3\cdot 3^n}=3\sqrt[n]{3}$，而 $\lim_{n\to\infty}3\sqrt[n]{3}=3$，由夹逼准则得极限为 $3$。

\textbf{例}：设 $x_1=\sqrt2$，$x_{n+1}=\sqrt{2+x_n}$，求 $\lim_{n\to\infty}x_n$。 先用归纳法证 $x_n<2$：$x_1<2$；若 $x_k<2$，则 $x_{k+1}=\sqrt{2+x_k}<\sqrt4=2$。 再证单调增：$x_{n+1}-x_n=\sqrt{2+x_n}-x_n$，因 $x_n<2$ 有 $2+x_n>x_n^2$，故 $\sqrt{2+x_n}>x_n$。 由单调有界准则极限存在，设为 $A$，两端取极限得 $A=\sqrt{2+A}$，即 $A^2-A-2=0$，$A=2$（舍去 $A=-1$）。

\begin{callwarn}{易错}
\begin{enumerate}
\item 夹逼准则中两边的极限\textbf{必须相同}，只放缩不验证同步等于没有做；
\item 对 $n$ 项和用夹逼时，若分母中变化的量级与主体量级\textbf{相同}（如 $\frac{1}{n^2+k^2}$），夹逼会把误差放大到不可忽略，这时应改用定积分定义；
\item 递推数列\textbf{必须先用单调有界准则证明极限存在，才能设极限为 $A$ 并"两端取极限"}，否则是逻辑循环；
\item 数列单调性的常用结论：$n^n\gg n!\gg a^n\ (a>1)\gg n^k\gg\ln n$。
\end{enumerate}
\end{callwarn}

\begin{callexam}{真题考法}
解答题（含证明）常考两类：一是"$n$ 次根号下 $n$ 项和"（放缩到最大项），二是"递推数列求极限"（先证收敛再解方程）。设问常分两小问：第一问证明极限存在，第二问求值或进一步求 $\lim nx_n^2$ 这类衍生极限（后者要重新回到 $\frac{1^\infty}{0/0}$ 的处理）。

\end{callexam}

\subsubsection{两个重要极限 $\bigstar$$\bigstar$$\bigstar$$\bigstar$}

\textbf{第一重要极限}：

\[\lim_{x\to 0}\frac{\sin x}{x}=1 .\]

其"推广形式"（$\square\to 0$ 时 $\sin\square\sim\square$）：

\[\lim_{x\to 0}\frac{\tan x}{x}=1,\qquad \lim_{x\to 0}\frac{\arcsin x}{x}=1,\qquad \lim_{x\to 0}\frac{1-\cos x}{x^{2}}=\frac12 .\]

\textbf{第二重要极限}：

\[\lim_{x\to\infty}\left(1+\frac1x\right)^{x}=e,\qquad \lim_{x\to 0}(1+x)^{\frac1x}=e .\]

推广为"$\square$ 型"：当 $\square\to 0$ 且 $\square\ne 0$ 时，$\lim(1+\square)^{\frac{1}{\square}}=e$。由此得到三条常用等价关系：

\[\ln(1+x)\sim x,\qquad e^{x}-1\sim x,\qquad \lim_{n\to\infty}\left(1+\frac{x}{n}\right)^{n}=e^{x}\ (\text{即 } 1^\infty \text{ 型通法}).\]

\textbf{$1^{\infty}$ 型的通用算法}：若 $\lim u(x)=1$，$\lim v(x)=\infty$，令

\[A=\lim v(x)\bigl[u(x)-1\bigr],\qquad \text{则 } \lim u(x)^{v(x)}=e^{A}.\]

推导只需一步：$u^{v}=e^{v\ln u}$，而 $\ln u=\ln[1+(u-1)]\sim u-1$。

\textbf{例}：求 $\lim_{x\to 0}\left(\cos x\right)^{\frac{1}{x^{2}}}$。 这是 $1^\infty$ 型。取 $A=\lim_{x\to 0}\frac{\cos x-1}{x^{2}}=-\frac12$，故原极限为 $e^{-1/2}=\frac{1}{\sqrt e}$。

\begin{callwarn}{易错}
\begin{enumerate}
\item 第一重要极限的标准形式必须凑成 $\frac{\sin\square}{\square}$，且 $\square\to 0$：如 $\lim_{x\to 0}\frac{\sin 5x}{\tan 3x}=\frac{5}{3}$，系数来自"内层系数之比"，不要漏；
\item $1^\infty$ 型\textbf{只能}用 $e^{A}$ 处理，直接写"$1^\infty=1$"或"$=e$"都是错的；
\item 用 $u^{v}=e^{v\ln u}$ 处理幂指函数时，$\ln u$ 的等价替换必须\textbf{整体}做，即 $\ln u=\ln[1+(u-1)]\sim u-1$，而不是 $\ln u\sim u$；
\item 第二重要极限的 $x\to 0$ 形式是 $(1+x)^{1/x}$，不要与 $x\to\infty$ 的 $(1+\frac1x)^{x}$ 混着抄错指数。
\end{enumerate}
\end{callwarn}

\begin{callexam}{真题考法}
填空题多为"$\lim_{x\to 0}\frac{\sin ax}{\tan bx}$"这类标准小题；解答题的常见形态是 $1^\infty$ 型，且往往与"脱极限号求参数"结合，例如由 $\lim_{x\to 0}\left(\frac{1+\tan x}{1+\sin x}\right)^{\frac{1}{x^{3}}}$ 的类型判断，进而定出参数。陷阱在于底数趋于 $1$ 但\textbf{加减项需要先通分或提公因子}才能看出 $1^\infty$ 结构。

\end{callexam}

\subsubsection{七类未定式与极限计算总思路 $\bigstar$$\bigstar$$\bigstar$$\bigstar$$\bigstar$}

极限计算的所有题目，最终都归结为下表七种。\textbf{识别类型是第一步}，类型认错，方法全错。

\begingroup\small\setlength{\tabcolsep}{4.0pt}
\begin{longtable}{@{}L{0.1801\linewidth}L{0.4277\linewidth}L{0.3208\linewidth}@{}}
\toprule
\textbf{未定式} & \textbf{首选方法} & \textbf{补充手段} \\
\midrule
\endhead
$\frac00$ & 等价无穷小替换、泰勒展开 & 洛必达、因式分解约分、有理化 \\
$\frac{\infty}{\infty}$ & 抓大头（分子分母同除最高阶项） & 洛必达 \\
$0\cdot\infty$ & 同除简单因式，化为 $\frac00$ 或 $\frac{\infty}{\infty}$ & 倒代换 $x=\frac1t$ \\
$\infty-\infty$ & 通分（含分数）、有理化（含根式） & 倒代换、提公因子 \\
$0^{0}$ & 化幂指函数 $u^{v}=e^{v\ln u}$ & 再用 $0\cdot\infty$ 的手法 \\
$\infty^{0}$ & 化幂指函数 $u^{v}=e^{v\ln u}$ & 先对底数取对数简化 \\
$1^{\infty}$ & $e^{A}$，其中 $A=\lim v(u-1)$ & 加减项先等价替换 \\
\bottomrule
\end{longtable}
\endgroup

\textbf{通用操作顺序（"三大思路"：化简、替换、展开）}

\begin{enumerate}
\item \textbf{化简}：约分、通分、提取极限非零因子、三角恒等变形、对数公式、换元（含倒代换）；
\item \textbf{判断类型}：代入求值不行，看是七类中的哪一类；
\item \textbf{替换}：乘除位置用等价无穷小；
\item \textbf{展开/洛必达}：加减相消或有根式时用泰勒；分母为一次式或含变限积分时优先洛必达；
\item \textbf{回代验算}：把结果代入原式或取特殊值检验数量级是否正确。
\end{enumerate}
\textbf{"改加减为乘除"的具体手法}

\begin{itemize}
\item \textbf{通分}：$\lim_{x\to 0}\left(\frac{1}{\sin x}-\frac1x\right)$，通分后 $\frac{x-\sin x}{x\sin x}\sim\frac{x^3/6}{x^2}=\frac{x}{6}\to 0$；
\item \textbf{有理化}：$\lim_{x\to+\infty}\left(\sqrt{x^{2}+x}-x\right)$，分子有理化得 $\frac{x}{\sqrt{x^2+x}+x}\to\frac12$；三次根号用立方差公式 $a^3-b^3=(a-b)(a^2+ab+b^2)$；
\item \textbf{倒代换}：$x\to\infty$ 与 $x\to 0$ 互换时令 $t=\frac1x$，便于产生分母、便于通分；
\item \textbf{提公因子}：指数函数相减时提出后面（或前面）一项，形如 $e^{a}-e^{b}=e^{b}(e^{a-b}-1)$；
\item \textbf{幂指函数恒等变形}：$u^{v}=e^{v\ln u}$。
\end{itemize}
\textbf{例}：求 $\lim_{x\to 0}\frac{\sqrt{1+\tan x}-\sqrt{1+\sin x}}{x\ln(1+x)-x^{2}}$。 分母 $x\ln(1+x)-x^2=x[\ln(1+x)-x]\sim x\cdot\left(-\frac{x^{2}}{2}\right)=-\frac{x^{3}}{2}$。 分子有理化：$\frac{\tan x-\sin x}{\sqrt{1+\tan x}+\sqrt{1+\sin x}}$，分母趋于 $2$，故分子 $\sim\frac{\tan x-\sin x}{2}\sim\frac{x^{3}}{4}$。 原极限 $=\frac{x^{3}/4}{-x^{3}/2}=-\frac12$。

\textbf{洛必达法则的适用条件与三条禁忌}

设 $f$、$g$ 在 $x_0$ 的去心邻域内可导，$g'(x)\ne 0$，$\lim\frac{f(x)}{g(x)}$ 为 $\frac00$ 或 $\frac{\infty}{\infty}$ 型，且 $\lim\frac{f'(x)}{g'(x)}=A$（$A$ 可以是有限数，也可以是 $\infty$），则 $\lim\frac{f(x)}{g(x)}=A$。

\begin{itemize}
\item \textbf{禁忌一}：不是 $\frac00$ 或 $\frac{\infty}{\infty}$ 型时绝不能洛。例如 $\lim_{x\to 0}\frac{1-\cos x}{x^{2}}$ 是 $\frac00$ 型可以洛；但 $\lim_{x\to 0}\frac{x}{1+\cos x}$ 不是未定式，洛必达会得到错误结果。
\item \textbf{禁忌二}：若 $\lim\frac{f'(x)}{g'(x)}$ 不存在（也不为无穷），不能反推原极限不存在。例如 $\lim_{x\to+\infty}\frac{x+\sin x}{x}=1$ 存在，但 $\lim(1+\cos x)$ 不存在——洛必达"失效"不等于"极限不存在"。此时应改用其他手法：$\frac{x+\sin x}{x}=1+\frac{\sin x}{x}\to 1$。
\item \textbf{禁忌三}：不要反复洛到"越洛越复杂"。当分母次数高、分子含变限积分或含根式时，先化简、等价替换或泰勒展开，再决定是否洛。
\end{itemize}
\textbf{洛必达的使用时机（什么时候优先用）}

\begin{enumerate}
\item 分母为一次式（如 $\lim_{x\to 0}\frac{\int_0^x\sin t^{2}\,dt}{x^{3}}$ 中分母是幂函数）：洛一次后变简单；
\item 含变限积分：$\frac{d}{dx}\int_a^{x}f(t)\,dt=f(x)$，洛必达是把积分"拿掉"的最快通道（要求被积函数连续）；
\item 分子分母含 $e^{x}$、$\ln x$ 等求导后形式变简单的情形。
\end{enumerate}
\textbf{例}：求 $\lim_{x\to 0}\frac{\int_0^{x}\ln(1+t)\,dt}{x^{2}}$。 $\frac00$ 型，洛必达一次：$\lim_{x\to 0}\frac{\ln(1+x)}{2x}=\lim_{x\to 0}\frac{x}{2x}=\frac12$。

\begin{callwarn}{易错}
\begin{enumerate}
\item \textbf{$\frac{\infty}{\infty}$ 抓大头的适用边界}：只对多项式/幂函数型因式有效。对复合结构（如 $\sqrt{x^2+x}$）要先把根号处理掉（同乘共轭式）或倒代换，不能"每个因式保留最高阶"了事；
\item $0\cdot\infty$ 同除时，要除\textbf{容易求导}的那个因式，否则洛必达后式子会更复杂；
\item 洛必达法则只适用于 $\frac00$ 与 $\frac{\infty}{\infty}$，且要检查"分子分母在去心邻域内可导、分母导数不为零"；用一次洛必达后类型若已消失，就不要再洛；
\item $1^\infty$ 型中，底数往往是通分/提公因子后的形式，直接看原式容易误判为 $\frac{\infty}{\infty}$。
\end{enumerate}
\end{callwarn}

\begin{callexam}{真题考法}
解答题第一题（约 10 分）几乎是"七类未定式的组合"：常见形态是分母为 $x^{3}$ 或 $x^{4}$ 的 $\frac00$ 型，需要先把分子有理化/通分，再用泰勒展开到与分母同阶。常见陷阱是把 $\frac{\infty}{\infty}$ 的抓大头误用到含根式的加减结构上；或者在 $x\to-\infty$ 时忘记 $\sqrt{x^2}=-x$ 的符号。

\end{callexam}

\subsubsection{泰勒/等价无穷小在极限中的联合运用 $\bigstar$$\bigstar$$\bigstar$$\bigstar$$\bigstar$}

\textbf{泰勒公式（带皮亚诺余项）}：若 $f$ 在 $x_0$ 处 $n$ 阶可导，则

\[f(x)=\sum_{k=0}^{n}\frac{f^{(k)}(x_0)}{k!}(x-x_0)^{k}+o\bigl((x-x_0)^{n}\bigr).\]

\textbf{必须背熟的麦克劳林展开（$x\to 0$）}：

\[e^{x}=1+x+\frac{x^{2}}{2!}+\frac{x^{3}}{3!}+o(x^{3}),\qquad \ln(1+x)=x-\frac{x^{2}}{2}+\frac{x^{3}}{3}+o(x^{3}),\]

\[\sin x=x-\frac{x^{3}}{6}+o(x^{3}),\qquad \cos x=1-\frac{x^{2}}{2}+\frac{x^{4}}{24}+o(x^{4}),\]

\[\tan x=x+\frac{x^{3}}{3}+o(x^{3}),\qquad \arcsin x=x+\frac{x^{3}}{6}+o(x^{3}),\qquad \arctan x=x-\frac{x^{3}}{3}+o(x^{3}),\]

\[(1+x)^{\alpha}=1+\alpha x+\frac{\alpha(\alpha-1)}{2!}x^{2}+o(x^{2}),\qquad \frac{1}{1-x}=1+x+x^{2}+o(x^{2}).\]

\textbf{"展到第几阶"的判定规则（核心技巧）}

\begin{itemize}
\item 若分母的阶数已确定（如分母为 $x^{3}$），则分子各项\textbf{按分母的阶数展开}即可；
\item 若分母阶数不确定（含参数），则\textbf{试展到分子、分母相消后第一个没有抵消的幂次}为止。
\end{itemize}
\textbf{例}：求 $\lim_{x\to 0}\frac{\sin x-\tan x}{x^{3}}$。

\[\sin x=x-\frac{x^{3}}{6}+o(x^{3}),\qquad \tan x=x+\frac{x^{3}}{3}+o(x^{3}),\]

相减得 $\sin x-\tan x=-\frac{x^{3}}{2}+o(x^{3})$，故原极限 $=-\frac12$。

\textbf{例}：已知 $\lim_{x\to 0}\frac{e^{x}+ax^{2}+bx}{x^{2}}=1$，求 $a,b$。 分母为 $x^{2}$，分子展开到 $2$ 阶：

\[e^{x}+ax^{2}+bx=1+(1+b)x+\left(\frac12+a\right)x^{2}+o(x^{2}).\]

分母趋于零而极限存在且等于 $1$，故分子的 $1$ 次项必须为零：$1+b=0$，得 $b=-1$；此时极限为 $\frac12+a=1$，得 $a=\frac12$。

\begin{callwarn}{易错}
\begin{enumerate}
\item \textbf{泰勒与等价无穷小不能混着乱用}：等价无穷小只是泰勒展开的一阶近似，分子里凡是会出现加减相消的项，一律不得用一阶等价替代；
\item \textbf{"泰勒嵌套泰勒"}：复合函数 $f(g(x))$ 中，外层展开的阶数要按内层 $g(x)$ 的阶数折算。例如 $e^{\sin x}$ 展到 $x^3$ 项时，$e^{u}=1+u+\frac{u^2}{2}+\frac{u^3}{6}$，需代入 $u=\sin x=x-\frac{x^3}{6}+\cdots$ 后再合并同次幂，不能只保留 $u$ 的一次项；
\item 展开\textbf{不足}是最大陷阱：只写到 $x^{2}$ 项却去比 $x^{3}$ 阶，会得到错误的系数；
\item 洛必达的次数有上限：$f$ 在点处 $n$ 阶可导时最多洛 $n-1$ 次，最后一步要回到导数定义；而 $f$ 有 $n$ 阶连续导数时可以洛 $n$ 次。
\end{enumerate}
\end{callwarn}

\begin{callexam}{真题考法}
解答题的核心考点。典型设问分两种：一是"求极限"（分母多为 $x^{3}$，考查差函数型等价与泰勒展开的配合）；二是"已知极限值反求参数 $a,b$ 或阶数 $n$"（必须把分子展开到分母同阶，再令低次项系数为零）。\textbf{只要题目中出现 $\sin x-\tan x$、$x-\ln(1+x)$、$\arcsin x-x$ 这类"差函数型"结构，就直接用差函数等价或泰勒，不要尝试洛必达。}

\end{callexam}

\subsubsection{已知极限确定参数、求函数表达式 $\bigstar$$\bigstar$$\bigstar$$\bigstar$}

\textbf{能反推出什么}

\begingroup\small\setlength{\tabcolsep}{4.0pt}
\begin{longtable}{@{}L{0.3513\linewidth}L{0.5944\linewidth}@{}}
\toprule
\textbf{已知条件} & \textbf{可推出} \\
\midrule
\endhead
$\lim\limits_{x\to x_0}\frac{f(x)}{g(x)}=A$ 存在，$\lim g(x)=0$ & $\lim f(x)=0$（分母推分子） \\
$\lim\limits_{x\to 0}\frac{f(x)}{x}=A$，$f$ 在 $0$ 处连续 & $f(0)=0$，且 $f'(0)=A$（导数定义） \\
$\lim\limits_{x\to x_0}f(x)=A>0$ & 存在去心邻域使 $f(x)>0$（保号性） \\
$\lim\limits_{x\to x_0}f(x)=A$ & $f(x)=A+\alpha(x)$，$\alpha(x)$ 为无穷小（脱极限号） \\
\bottomrule
\end{longtable}
\endgroup

\textbf{"脱极限号"的关键}：脱去极限号后必须补上\textbf{高阶无穷小项}。例如由 $\lim_{x\to 0}\frac{f(x)}{x}=2$ 只能写 $f(x)=2x+o(x)$，不能写 $f(x)=2x$。

\textbf{求参数的套路}

\begin{enumerate}
\item \textbf{定阶}：把等式整理成 $\frac{P(x)}{Q(x)}\to A$ 的形式，先由"分母趋于零 $\Rightarrow$ 分子趋于零"定出第一组条件；
\item \textbf{展开比较}：把分子泰勒展开，令\textbf{低于分母阶数的项的系数全为零}，令与分母同阶项的系数等于 $A$；
\item \textbf{验算}：把求得的参数代回原式，确认极限确实存在且等于给定值。
\end{enumerate}
\textbf{例}：已知 $\lim_{x\to 0}\frac{\ln(1+x)+ax+bx^{2}}{x^{2}}=2$，求 $a,b$。 分子展开：$\ln(1+x)+ax+bx^{2}=\left(a+1\right)x+\left(b-\frac12\right)x^{2}+o(x^{2})$。 需 $a+1=0$，即 $a=-1$；此时极限为 $b-\frac12=2$，得 $b=\frac52$。代回验证：$\frac{\ln(1+x)-x+\frac52x^{2}}{x^{2}}\to-\frac12+\frac52=2$，正确。

\textbf{例}：设 $\lim_{x\to 0}\frac{f(x)}{x}=2$，求 $\lim_{x\to 0}\frac{f(3x)}{\sin 2x}$。 由条件得 $f(3x)=2\cdot 3x+o(x)=6x+o(x)$，而 $\sin 2x\sim 2x$，故原式 $=\frac{6x}{2x}=3$。

\begin{callwarn}{易错}
\begin{enumerate}
\item 由"分母趋于 $0$ 且分式极限存在"推"分子趋于 $0$"时，\textbf{必须强调分式极限存在}，否则不能推；
\item 脱极限号必须补 $o(\ )$ 项：由 $\lim\frac{f(x)}{x}=2$ 直接写 $f(x)=2x$ 会丢掉高阶信息，遇到 $f(3x)$ 这类变形就出错；
\item 比较系数时只令"低次项为零"是不够的，还要检查\textbf{分母阶数与保留项阶数是否匹配}，否则会漏掉一组解。
\end{enumerate}
\end{callwarn}

\begin{callexam}{真题考法}
常以填空或解答的第一问出现，形式固定："已知 $\lim_{x\to 0}\frac{\cdots}{x^{n}}=A$，求 $a,b$"。答案的判定完全依赖泰勒展开的系数比较。陷阱是把 $a$、$b$ 求反（把 $b=-1$、$a=\frac12$ 写颠倒），以及忘记验证极限确实存在。

\end{callexam}

\subsection{数列极限的计算}

\subsubsection{数列极限的计算：n 项和与 n 项积、递推数列 $\bigstar$$\bigstar$$\bigstar$$\bigstar$}

数列极限有四种典型结构，\textbf{先识别结构再选工具}：

\begingroup\small\setlength{\tabcolsep}{4.0pt}
\begin{longtable}{@{}L{0.5075\linewidth}L{0.4382\linewidth}@{}}
\toprule
\textbf{结构} & \textbf{主要工具} \\
\midrule
\endhead
$n$ 项和，分母主体量级与变化量级\textbf{相同} & 定积分定义 \\
$n$ 项和，分母主体量级\textbf{高于}变化量级 & 夹逼准则 \\
$n$ 项积 & 取对数化为 $n$ 项和 \\
递推 $x_{n+1}=f(x_n)$ & 单调有界准则 + 解方程 \\
$n$ 次根号下 $n$ 项和 & 夹逼（抓最大项） \\
\bottomrule
\end{longtable}
\endgroup

\textbf{一、定积分定义求 $n$ 项和}

条件：区间 $[0,1]$、$n$ 等分、取端点。对应关系为 $\lim\sum\to\int_0^1$，$\frac1n\to dx$，$\frac{i}{n}\to x$。

\[\lim_{n\to\infty}\frac1n\sum_{i=1}^{n}f\!\left(\frac{i}{n}\right)=\int_0^1 f(x)\,dx .\]

操作步骤：把求和式改写成"$\frac1n$ 乘以某个关于 $\frac in$ 的式子"（哪里出现 $i$ 就在哪里改写为 $i$），提出 $\frac1n$，再改写为定积分。

\textbf{例}：求 $\lim_{n\to\infty}\sum_{k=1}^{n}\frac{n}{n^{2}+k^{2}}$。

\[\sum_{k=1}^{n}\frac{n}{n^{2}+k^{2}}=\frac1n\sum_{k=1}^{n}\frac{1}{1+\left(\frac kn\right)^{2}}\to\int_0^1\frac{dx}{1+x^{2}}=\arctan 1=\frac{\pi}{4}.\]

\textbf{注意两个推广}：$\circled{1}$ 若分的不是 $n$ 等分（差一项），则"多了拿出来、少了补上"，转化为 $n$ 等分（有限项极限为零）；$\circled{2}$ 若 $i$ 不从 $1$ 到 $n$，用下标代换移回来。

\textbf{二、夹逼准则求 $n$ 项和}

\textbf{例}：求 $\lim_{n\to\infty}\sum_{k=1}^{n}\frac{n}{n^{2}+k}$。 每一项都介于 $\frac{n}{n^{2}+n}$ 与 $\frac{n}{n^{2}+1}$ 之间，故

\[\frac{n^{2}}{n^{2}+n}\le\sum_{k=1}^{n}\frac{n}{n^{2}+k}\le\frac{n^{2}}{n^{2}+1},\]

两边极限都为 $1$，故原极限为 $1$。

\textbf{三、$n$ 项积：取对数}

\textbf{例}：求 $\lim_{n\to\infty}\left(\frac{n+1}{n}\cdot\frac{n+2}{n}\cdots\frac{n+n}{n}\right)^{\frac1n}$。 取对数得 $\frac1n\sum_{k=1}^{n}\ln\left(1+\frac kn\right)$，这是定积分结构，极限为 $\int_0^1\ln(1+x)\,dx=2\ln 2-1$，故原极限为 $e^{2\ln 2-1}=\frac{4}{e}$。

\textbf{四、$n$ 次根号下 $n$ 项和（"最大的赢"）}

步骤：找出最大项；右边放（都变成最大项，即 $n$ 倍最大项）；左边缩（只留最大项）。

\textbf{例}：$\lim_{n\to\infty}\sqrt[n]{1^{n}+2^{n}+3^{n}}=3$，因为

\[\sqrt[n]{3^{n}}\le\sqrt[n]{1+2^{n}+3^{n}}\le\sqrt[n]{3\cdot 3^{n}}=3\sqrt[n]{3},\qquad \lim_{n\to\infty}3\sqrt[n]{3}=3 .\]

再如 $\lim_{n\to\infty}\sqrt[n]{n}=1$：两端不好直接夹，改为取对数，由 $\frac{\ln n}{n}\to 0$ 得 $\sqrt[n]{n}=e^{\frac{\ln n}{n}}\to e^{0}=1$。

\textbf{注意}：用夹逼处理 $n$ 次根号时，若最大项 $M_n$ 随 $n$ 变化，必须先验证 $\lim M_n$ 存在，再让两端都收敛到这个值；若最大项本身不收敛，此法失效。

\begin{callkey}{必背：求极限的"总纲"}
\textbf{看到极限题，先判类型，再选工具。}

\begin{enumerate}
\item 先化简（通分、有理化、提非零因子、换元、取对数）；
\item 判类型：是七类未定式中的哪一类；
\item $\frac00$ 优先等价替换，加减相消或分母阶数高时改用泰勒；
\item $\frac{\infty}{\infty}$ 抓大头；含根式加减先有理化或倒代换；
\item $1^{\infty}$ 一律写成 $e^{A}$，$A=\lim v(u-1)$；
\item 数列题先看结构：量级相同用定积分定义，量级不同用夹逼，递推用单调有界。
\end{enumerate}
\end{callkey}

\textbf{五、递推数列}：见"单调有界准则"考点，四步是"证有界 $\to$ 证单调 $\to$ 设极限解方程 $\to$ 取舍"。

\textbf{六、数列极限的其他常用转化}

\begin{enumerate}
\item \textbf{数列未定式转函数未定式（海涅定理）}：$n\to\infty$ 型的数列极限，把 $n$ 改写成 $x$、把 $\frac1n$ 改写成 $\frac1x$（或直接令 $\frac1n=t$），转化为 $x\to+\infty$（或 $t\to 0$）的函数极限后再用洛必达、泰勒等工具。例如
\end{enumerate}
\[\lim_{n\to\infty}n\left(e^{\frac1n}-1\right)=\lim_{t\to 0^+}\frac{e^{t}-1}{t}=1 .\]

\begin{enumerate}
\item \textbf{"先求和再取极限"}：若通项是等比或可裂项的形式，先算部分和 $S_n$ 再求极限。例如 $\lim_{n\to\infty}\sum_{k=1}^{n}\frac{1}{k(k+1)}=\lim_{n\to\infty}\left(1-\frac{1}{n+1}\right)=1$。
\item \textbf{夹逼 + 定积分定义}：当和式不能直接写成 $\frac1n\sum f\left(\frac in\right)$ 时，先放缩处理分母，再用夹逼准则，或放缩后套定积分定义。
\item \textbf{递推数列的"衍生极限"}：第一问证出 $\lim x_n=0$ 后，第二问常要求 $\lim nx_n$ 或 $\lim\frac{x_{n+1}}{x_n}$，此时要把 $n$ 与 $x_n$ 的关系用递推式展开。典型手法是\textbf{取倒数造出可求和的差式}：若 $x_{n+1}=x_n-x_n^{2}$，则
\end{enumerate}
\[\frac{1}{x_{n+1}}-\frac{1}{x_n}=\frac{x_n-x_{n+1}}{x_nx_{n+1}}=\frac{x_n^{2}}{x_nx_{n+1}}=\frac{x_n}{x_{n+1}}\to 1,\]

由此可推出 $\frac{1}{x_n}$ 与 $n$ 同阶，即 $nx_n\to 1$。\textbf{不能只停留在"极限为 $0$"就交卷。}

\textbf{例}：求 $\lim_{n\to\infty}n\left(\sqrt[n]{a}-1\right)$（$a>0$）。 令 $t=\frac1n\to 0^+$，则原式 $=\lim_{t\to 0^+}\frac{a^{t}-1}{t}=\ln a$。

\begin{callwarn}{易错}
\begin{enumerate}
\item \textbf{判别量级是选对工具的前提}：$\frac{1}{n^2+k}$ 中分母主体 $n^2$ 远大于变化量级 $n$，用夹逼；而 $\frac{1}{n^2+k^2}$ 中两者同量级，必须用定积分定义；
\item 定积分定义要求 $\frac1n$ 恰好出现一次，若原式提出的是 $\frac{1}{n^{2}}$ 而不是 $\frac1n$，先检查是否漏项或需要凑；
\item 取对数时注意 $\ln$ 的运算律只对\textbf{正项}成立，使用前说明各项为正；
\item $n$ 项积取对数后是 $n$ 项和，若变成"$\frac{1}{n}\sum\ln(\cdots)$"就是定积分结构，若变成"$\sum$ 通项与 $n$ 无关"就不是。
\end{enumerate}
\end{callwarn}

\begin{callexam}{真题考法}
解答题常考 $\lim_{n\to\infty}\frac1n\sum f\left(\frac in\right)$ 形式的定积分定义题，以及"分母不同、不可直接求和"的夹逼题（例如带 $\sin\frac{k\pi}{n}$ 的和式）。陷阱在于把量级判断反了：本该夹逼的用了定积分定义（两边极限不相等），或本该用定积分定义的硬套夹逼（误差放大）。

\end{callexam}

\section{连续}

\subsection{连续与间断}

\subsubsection{函数的连续性与左、右连续 $\bigstar$$\bigstar$$\bigstar$}

\textbf{在一点连续}：设 $f$ 在 $x_0$ 的某邻域内有定义，若

\[\lim_{x\to x_0}f(x)=f(x_0),\]

则称 $f$ 在 $x_0$ 处连续。等价的增量形式：$\lim_{\Delta x\to 0}\Delta y=0$，其中 $\Delta y=f(x_0+\Delta x)-f(x_0)$。

\textbf{连续的充要条件（三段式，必须逐项验证）}：

\[\lim_{x\to x_0^-}f(x)=\lim_{x\to x_0^+}f(x)=f(x_0).\]

\textbf{左连续与右连续}：$\lim_{x\to x_0^-}f(x)=f(x_0)$ 称为在 $x_0$ 处左连续；$\lim_{x\to x_0^+}f(x)=f(x_0)$ 称为右连续。函数在 $x_0$ 处连续 $\iff$ 在 $x_0$ 处既左连续又右连续。

\textbf{区间连续}：在开区间 $(a,b)$ 内每点连续，称为在 $(a,b)$ 内连续；若还在左端点 $a$ 右连续、在右端点 $b$ 左连续，则称在闭区间 $[a,b]$ 上连续。

\textbf{例}：讨论 $f(x)=\begin{cases}\dfrac{\sin x}{x}, & x\ne 0,\\ 1, & x=0\end{cases}$ 在 $x=0$ 处的连续性。 $\lim_{x\to 0}\frac{\sin x}{x}=1=f(0)$，故在 $x=0$ 处连续。

\textbf{讨论分段函数的连续性（标准四步）}

\begin{enumerate}
\item 在非分段点处：直接说明"由初等函数的连续性，$f$ 在……上连续"；
\item 在分段点 $x_0$ 处：分别求 $f(x_0^-)$、$f(x_0^+)$、$f(x_0)$；
\item 令三者相等，解出使函数连续的参数条件（若题目含参数）；
\item 写出最终的连续区间，注意\textbf{区间端点处只要求单侧连续}。
\end{enumerate}
\textbf{例}：设 $f(x)=\begin{cases}\dfrac{e^{x}-1}{x}, & x>0,\\ a, & x=0,\\ x+b, & x<0,\end{cases}$ 在 $x=0$ 处连续，求 $a,b$。 $\lim_{x\to 0^+}\frac{e^{x}-1}{x}=1$，$\lim_{x\to 0^-}(x+b)=b$，$f(0)=a$。 连续性要求 $a=b=1$。注意此题只要求"在 0 处连续"，$a$ 与 $b$ 必须都等于 $1$ 缺一不可。

\begin{callwarn}{易错}
\begin{enumerate}
\item 讨论连续性时\textbf{不能漏掉"函数值 $f(x_0)$ 存在"这一项}：$\lim_{x\to x_0}f(x)$ 存在但 $f(x_0)$ 无定义，则函数在该点不连续（且为可去间断点）；
\item 分段函数在分段点处必须\textbf{分左右验证}，直接代公式是错的；
\item 函数在一点连续是"局部"性质，讨论区间连续时要逐点说明（分段点单独讨论）。
\end{enumerate}
\end{callwarn}

\begin{callexam}{真题考法}
选择题常以"下列函数在 $x=0$ 处连续的是"或"$f$ 在 $x_0$ 处连续的充要条件是"出现，把"极限存在""有定义""左右极限存在"等条件做排列组合设置陷阱。\textbf{正确选项必是三条同时满足。}

\end{callexam}

\subsubsection{间断点及其类型（第一类/第二类）$\bigstar$$\bigstar$$\bigstar$}

\textbf{间断点}：若 $f$ 在 $x_0$ 的某去心邻域内有定义，但在 $x_0$ 处不连续，则称 $x_0$ 为 $f$ 的间断点。候选点来自：无定义点、分段函数的分界点、分母零点。

\textbf{分类（按左右极限的存在性）}

\begingroup\small\setlength{\tabcolsep}{4.0pt}
\begin{longtable}{@{}L{0.3150\linewidth}L{0.4477\linewidth}L{0.1658\linewidth}@{}}
\toprule
\textbf{类型} & \textbf{判据} & \textbf{典型例子} \\
\midrule
\endhead
第一类：可去间断点 & 左右极限都存在且相等，但不等于 $f(x_0)$ 或 $f(x_0)$ 无定义 & $\frac{\sin x}{x}$ 在 $x=0$ \\
第一类：跳跃间断点 & 左右极限都存在但不相等 & $\lfloor x\rfloor$ 在 $x=0$；$\frac{\vert x\vert}{x}$ 在 $x=0$ \\
第二类：无穷间断点 & 至少一侧极限为 $\infty$ & $\frac1x$ 在 $x=0$ \\
第二类：振荡间断点 & 极限不存在且不趋于 $\infty$ & $\sin\frac1x$ 在 $x=0$ \\
\bottomrule
\end{longtable}
\endgroup

\textbf{判定流程}

\begin{enumerate}
\item 求出 $f$ 的定义域，找出所有无定义点与分段点，作为候选间断点；
\item 对每个候选点，分别计算 $\lim_{x\to x_0^-}f(x)$ 与 $\lim_{x\to x_0^+}f(x)$；
\item 两侧极限均存在：相等则为可去间断点，不等则为跳跃间断点；
\item 至少一侧不存在（含趋于无穷）：为第二类，再区分是无穷还是振荡。
\end{enumerate}
\textbf{可去间断点的"修补"}：若 $x_0$ 是可去间断点，令 $\tilde f(x_0)=\lim_{x\to x_0}f(x)$，函数在该点即变为连续。

\textbf{例}：求 $f(x)=\frac{x^{2}-1}{x^{2}-3x+2}$ 的间断点并分类。 分母 $x^{2}-3x+2=(x-1)(x-2)$，候选点 $x=1,2$。 $x=1$：约分后 $f(x)=\frac{x+1}{x-2}$（$x\ne 1$），$\lim_{x\to 1}f(x)=\frac{2}{-1}=-2$ 存在，为\textbf{可去间断点}。 $x=2$：$\lim_{x\to 2}\frac{x+1}{x-2}=\infty$，为\textbf{无穷间断点}（第二类）。

\textbf{三类"高危函数"的左右极限（必须分侧计算）}

\begingroup\small\setlength{\tabcolsep}{4.0pt}
\begin{longtable}{@{}L{0.2144\linewidth}L{0.2144\linewidth}L{0.2144\linewidth}L{0.2680\linewidth}@{}}
\toprule
\textbf{函数} & \textbf{$x\to x_0^-$} & \textbf{$x\to x_0^+$} & \textbf{间断类型} \\
\midrule
\endhead
$e^{\frac{1}{x-x_0}}$ & $0$ & $+\infty$ & 跳跃或第二类 \\
$\arctan\frac{1}{x-x_0}$ & $-\frac{\pi}{2}$ & $\frac{\pi}{2}$ & 跳跃（第一类） \\
$\frac{1}{1-e^{\frac{x}{x-x_0}}}$ & 需分别计算 & 需分别计算 & 常为跳跃 \\
$\vert x-x_0\vert$ 相关的分式 & 代入 $x<x_0$ & 代入 $x>x_0$ & 依结果定 \\
\bottomrule
\end{longtable}
\endgroup

\textbf{例}：讨论 $f(x)=\arctan\frac1x$ 在 $x=0$ 处的间断类型。 $\lim_{x\to 0^-}\arctan\frac1x=-\frac{\pi}{2}$，$\lim_{x\to 0^+}\arctan\frac1x=\frac{\pi}{2}$，两侧极限都存在但不相等，故 $x=0$ 是\textbf{跳跃间断点}（第一类）。

\textbf{例}：讨论 $f(x)=e^{\frac1x}$ 在 $x=0$ 处的间断类型。 $\lim_{x\to 0^-}e^{1/x}=0$（因 $\frac1x\to-\infty$），$\lim_{x\to 0^+}e^{1/x}=+\infty$。左极限存在、右极限为无穷，属\textbf{第二类间断点}（无穷间断）。注意这里不能只说"无穷间断"，因为左极限是存在的。

\begin{callwarn}{易错}
\begin{enumerate}
\item 判断间断点前\textbf{必须先求出极限}，不能只看"分母为零"就直接判为无穷间断——如 $\frac{x^2-1}{x-1}$ 在 $x=1$ 处是可去间断；
\item 遇 $e^{1/x}$、$\arctan\frac1x$、$\frac{1}{1-e^{x/(x-1)}}$ 时必须\textbf{分左右计算}，两侧结果往往不同（常见为跳跃间断）；
\item 判断"第一类"的判据是\textbf{左右极限都存在}，而不管是否等于函数值；即便函数在该点无定义，只要左右极限存在且相等，仍是可去间断点。
\end{enumerate}
\end{callwarn}

\begin{callexam}{真题考法}
选择题/填空题的高频考点，标准问法是"求函数的间断点并判别其类型"。陷阱集中在三类函数：含 $\vert x-a\vert$ 的（左右分段）、含 $e^{1/(x-a)}$ 的（左右极限一个是 $\infty$ 一个是 $0$）、含 $\arctan\frac{1}{x-a}$ 的（左右极限分别是 $\pm\frac{\pi}{2}$）。

\end{callexam}

\subsection{连续函数的性质}

\subsubsection{连续函数的性质与初等函数的连续性 $\bigstar$$\bigstar$}

\textbf{连续函数的运算性质}：若 $f$、$g$ 在 $x_0$ 处连续，则 $f\pm g$、$f\cdot g$、$\frac fg$（$g(x_0)\ne 0$）在 $x_0$ 处都连续。

\textbf{复合函数的连续性}：若 $u=g(x)$ 在 $x_0$ 处连续，$y=f(u)$ 在 $u_0=g(x_0)$ 处连续，则 $f(g(x))$ 在 $x_0$ 处连续。更有用的是\textbf{极限与函数符号可交换}：

\[\lim_{x\to x_0}f\bigl(g(x)\bigr)=f\!\left(\lim_{x\to x_0}g(x)\right).\]

\textbf{反函数的连续性}：严格单调的连续函数，其反函数在其值域上连续。

\textbf{初等函数的连续性}：基本初等函数在其\textbf{定义域}内连续；初等函数在其\textbf{定义区间}内连续。这里的"定义区间"指定义域内的区间，某些孤立点（如 $\sqrt{\cos x-1}$ 只在 $x=2k\pi$ 有定义的孤立点）不属于定义区间。

\textbf{例}：求 $\lim_{x\to 0}\ln\frac{\sin x}{x}$。 因 $\ln$ 在 $1$ 处连续，故 $\lim_{x\to 0}\ln\frac{\sin x}{x}=\ln\left(\lim_{x\to 0}\frac{\sin x}{x}\right)=\ln 1=0$。

\textbf{"$f$ 在 $x_0$ 处连续"的三个等价刻画（概念题必背）}

\[\lim_{x\to x_0}f(x)=f(x_0) \iff \lim_{\Delta x\to 0}\Delta y=0 \iff \lim_{x\to x_0^-}f(x)=\lim_{x\to x_0^+}f(x)=f(x_0).\]

\textbf{连续与可导的关系（本章只作结论性了解，展开在微分学一章）}

\begin{itemize}
\item 可导 $\Rightarrow$ 连续，连续 $\not\Rightarrow$ 可导（如 $\vert x\vert$ 在 $x=0$ 处连续但不可导）；
\item 连续是"曲线不断开"，可导是"曲线无尖角、切线不垂直于 $x$ 轴"；
\item 初等函数在其定义区间内连续，因此讨论初等函数的连续性归结为"求定义域 + 单独查分段点"。
\end{itemize}
\begin{callwarn}{易错}
\begin{enumerate}
\item "初等函数在其定义域内连续"的准确说法是"在\textbf{定义区间}内连续"，定义域中孤立的点处不讨论连续性；
\item 用"极限与函数符号交换"时，外层函数必须在对应点\textbf{连续}，如 $\lim_{x\to 0}\ln\frac{\sin x}{x}$ 中 $\ln$ 在 $1$ 处连续才可交换；
\item 不要与可导性混淆：连续不一定可导（如 $\vert x\vert$ 在 $0$ 处）。
\end{enumerate}
\end{callwarn}

\begin{callexam}{真题考法}
多以小题形式考查"求函数 $\lim_{n\to\infty}f_n(x)$ 的连续区间"或"下列函数在何处连续"，本质是把初等函数连续性与极限计算结合。陷阱是漏掉定义域中的孤立点与分段点。

\end{callexam}

\subsubsection{闭区间上连续函数的性质 $\bigstar$$\bigstar$$\bigstar$}

设 $f(x)$ 在闭区间 $[a,b]$ 上连续（记 $f\in C[a,b]$）。

\textbf{有界性定理}：$f$ 在 $[a,b]$ 上有界，即存在 $M>0$，使对一切 $x\in[a,b]$ 有 $\vert f(x)\vert\le M$。

\textbf{最值定理}：$f$ 在 $[a,b]$ 上必取得最大值 $M$ 与最小值 $m$，即存在 $\xi,\eta\in[a,b]$ 使

\[f(\xi)=m=\min_{x\in[a,b]}f(x),\qquad f(\eta)=M=\max_{x\in[a,b]}f(x).\]

\textbf{介值定理}：设 $f$ 在 $[a,b]$ 上连续，$m\le c\le M$，则存在 $\xi\in[a,b]$ 使 $f(\xi)=c$。特别地，若 $f(a)\ne f(b)$，则对介于 $f(a)$ 与 $f(b)$ 之间的任何实数 $c$，存在 $\xi\in(a,b)$ 使 $f(\xi)=c$。

\textbf{零点定理}（介值定理的推论）：若 $f\in C[a,b]$ 且 $f(a)\cdot f(b)<0$，则存在 $\xi\in(a,b)$ 使 $f(\xi)=0$。

\textbf{三点必须记牢}

\begin{itemize}
\item \textbf{闭区间条件不可少}：$f(x)=\frac1x$ 在 $(0,1]$ 上连续但无界、无最大值；
\item \textbf{连续性条件不可少}：$f(x)=\begin{cases}-1, & x<0,\\ 1, & x\ge 0\end{cases}$ 在 $[-1,1]$ 上无零点，虽然两端异号；
\item 介值定理中的 $\xi$ \textbf{不唯一}（如常值函数）。
\end{itemize}
\textbf{使用零点定理证题的套路}

\begin{enumerate}
\item 把待证等式移项整理成 $F(x)=0$ 的形式（构造辅助函数，常见是把常数移到一边）；
\item 指出 $F$ 在某个闭区间上连续；
\item 计算两端点函数值，验证异号（$F(a)F(b)<0$）；
\item 由零点定理得存在 $\xi$，再把结论翻译回原命题。
\end{enumerate}
\textbf{证明零点唯一}：在"存在"之后追加一步——证明 $f$ 在该区间上\textbf{严格单调}（通常用 $f'(x)>0$ 或 $f'(x)<0$）。

\textbf{例}：证明方程 $x^{3}-3x+1=0$ 在 $(0,1)$ 内至少有一个实根。 令 $f(x)=x^{3}-3x+1$，$f$ 在 $[0,1]$ 上连续，$f(0)=1>0$，$f(1)=-1<0$。由零点定理，存在 $\xi\in(0,1)$ 使 $f(\xi)=0$。

\textbf{例}：证明方程 $x=2\ln x+1$ 在 $(2,e^{2})$ 内至少有一个实根。 令 $f(x)=x-2\ln x-1$，$f$ 在 $[2,e^{2}]$ 上连续。 $f(2)=2-2\ln 2-1=1-2\ln 2<0$（因 $\ln 2>\frac12$）；$f(e^{2})=e^{2}-4-1=e^{2}-5>0$（因 $e^{2}\approx 7.39$）。 由零点定理，存在 $\xi\in(2,e^{2})$ 使 $f(\xi)=0$，即原方程在该区间内至少有一实根。

\textbf{"存在"之外再加"唯一"的标准写法}

若要证"有且仅有一个根"，在零点定理之后追加一段：

\begin{quote}
又对 $x\in(a,b)$ 有 $f'(x)>0$（或 $<0$），故 $f$ 在 $[a,b]$ 上严格单调，零点至多一个；结合前面的存在性，零点恰有一个。

\end{quote}
\textbf{证明"存在 $\xi$ 使 $f(\xi)=c$"的辅助函数构造表}

\begingroup\small\setlength{\tabcolsep}{4.0pt}
\begin{longtable}{@{}L{0.4728\linewidth}L{0.4728\linewidth}@{}}
\toprule
\textbf{待证形式} & \textbf{构造的辅助函数} \\
\midrule
\endhead
$f(\xi)=c$ & $F(x)=f(x)-c$ \\
$f(\xi)=\xi$ & $F(x)=f(x)-x$ \\
$f(\xi)=k\xi+b$ & $F(x)=f(x)-kx-b$ \\
$\int_a^{\xi}f(t)\,dt=\int_{\xi}^{b}f(t)\,dt$ & $F(x)=\int_a^{x}f(t)\,dt-\int_x^{b}f(t)\,dt$ \\
$f(\xi)=f(\xi+T)$ & $F(x)=f(x+T)-f(x)$ \\
\bottomrule
\end{longtable}
\endgroup

\begin{callwarn}{易错}
\begin{enumerate}
\item 用零点定理只能得到"零点\textbf{存在}"，不能得到"唯一"；要证唯一必须补上单调性；
\item "端点异号"不是零点的必要条件：$f(x)=x^{2}$ 在 $[-1,1]$ 上有零点但两端同号；
\item 定理要求的是\textbf{闭区间}上的连续函数，使用前必须先说明区间与连续性；
\item 用介值定理证 $f(\xi)=c$ 时，要先把原式整理成"$f(x)-c$"的形式，构造辅助函数是最容易漏的一步。
\end{enumerate}
\end{callwarn}

\begin{callexam}{真题考法}
证明题的高频工具，常见两种设问：一是"证明方程 $f(x)=0$ 在 $(a,b)$ 内有且仅有一个实根"（存在用零点定理，唯一用单调性）；二是与中值定理结合，先由介值定理取到某个中间值，再用拉格朗日或罗尔定理。陷阱是忘记指明"闭区间上连续"这一前提。

\end{callexam}

\subsection*{本章易错清单}
\addcontentsline{toc}{subsection}{本章易错清单}

\begin{enumerate}
\item \textbf{加减中滥用等价无穷小}：$\lim_{x\to 0}\frac{\tan x-\sin x}{x^{3}}$ 中把 $\tan x$、$\sin x$ 都换成 $x$ 得到 $0$；正确答案是 $\frac12$。等价替换只在乘除（因子）位置可直接进行。
\item \textbf{差函数型等价的符号与系数记错}：$x-\sin x\sim\frac{x^{3}}{6}$、$\tan x-x\sim\frac{x^{3}}{3}$、$x-\ln(1+x)\sim\frac{x^{2}}{2}$，符号与系数必须逐个核对。
\item \textbf{随手拆极限}：在 $\lim[f(x)+g(x)]$ 中，若 $\lim f(x)$ 不存在就绝不能拆成两项分别求；"存在且为零 $\times$ 不存在"是不确定，不能直接判为不存在。
\item \textbf{忽略定义域}：判断奇偶性前不检查定义域是否关于原点对称；求复合函数定义域时漏掉内层值域的限制；忘记 $\ln$ 的真数大于零。
\item \textbf{混淆无界与无穷大}：无穷大必无界，无界不一定是无穷大（$x\sin x$ 在 $x\to+\infty$ 时无界但非无穷大）。
\item \textbf{$\sqrt{x^{2}}=\vert x\vert$ 写错}：在 $x\to-\infty$ 时把 $\sqrt{x^{2}}$ 直接写成 $x$，符号恰好相反。
\item \textbf{洛必达法则的适用条件与次数}：只能在 $\frac00$、$\frac{\infty}{\infty}$ 型使用；$f$ 在点处 $n$ 阶可导时最多洛 $n-1$ 次，最后一步要回到导数定义。
\item \textbf{$\frac{\infty}{\infty}$ 抓大头用错对象}：对 $\sqrt{x^{2}+x}-x$ 这类含根式加减的结构，必须先有理化或倒代换，不能"每个因式保留最高阶"了事。
\item \textbf{$1^{\infty}$ 型处理不当}：必须用 $e^{A}$（$A=\lim v(u-1)$），不能写"$1^{\infty}=1$"或直接"$=e$"；幂指函数一律化 $u^{v}=e^{v\ln u}$。
\item \textbf{脱极限号不补高阶无穷小}：由 $\lim_{x\to 0}\frac{f(x)}{x}=2$ 应写 $f(x)=2x+o(x)$，写成 $f(x)=2x$ 会在后续变形中出错。
\item \textbf{间断点只看分母}：必须先算极限，$\frac{x^{2}-1}{x-1}$ 在 $x=1$ 处是可去间断而非无穷间断；含 $e^{1/x}$、$\arctan\frac1x$ 的必须分左右计算。
\item \textbf{数列极限中"先证存在再设极限"的顺序颠倒}：递推数列必须先用单调有界准则（或夹逼准则）证明极限存在，才能设 $\lim x_n=A$ 并在递推式两端取极限。
\item \textbf{零点定理只能证存在}：要证"有且仅有"必须补充单调性论证；且不能忘记验证"闭区间上连续"。
\item \textbf{连续的三段式验证漏项}：$f$ 在 $x_0$ 连续要求"左极限 $=$ 右极限 $=f(x_0)$"三条同时成立，漏掉哪一条都不行。
\end{enumerate}
\chapter{一元函数微分学}
\markboth{一元函数微分学}{一元函数微分学}
\subsection*{本章考纲要求}
\addcontentsline{toc}{subsection}{本章考纲要求}

\begin{quote}
\textbf{考试内容}

\begin{enumerate}
\item 导数和微分的概念
\item 导数的几何意义和物理意义
\item 函数的可导性与连续性之间的关系
\item 平面曲线的切线和法线
\item 导数和微分的四则运算
\item 基本初等函数的导数
\item 复合函数、反函数和隐函数的微分法
\item 高阶导数
\item 一阶微分形式的不变性
\item 微分中值定理
\item 洛必达（L'Hospital）法则
\item 函数单调性的判别
\item 函数的极值
\item 函数图形的凹凸性、拐点及渐近线
\item 函数图形的描绘
\item 函数的最大值与最小值
\item 弧微分
\item 曲率的概念
\item 曲率圆与曲率半径
\end{enumerate}
\textbf{考试要求}

\begin{enumerate}
\item 理解导数和微分的概念，理解导数与微分的关系，理解导数的几何意义，会求平面曲线的切线方程和法线方程，了解导数的物理意义，会用导数描述一些物理量，理解函数的可导性与连续性之间的关系。
\item 掌握导数的四则运算法则和复合函数的求导法则，掌握基本初等函数的导数公式。了解微分的四则运算法则和一阶微分形式的不变性，会求函数的微分。
\item 了解高阶导数的概念，会求简单函数的高阶导数。
\item 会求分段函数的导数，会求隐函数和由参数方程所确定的函数以及反函数的导数。
\item 理解并会用罗尔（Rolle）定理、拉格朗日（Lagrange）中值定理和泰勒（Taylor）定理，了解并会用柯西（Cauchy）中值定理。
\item 掌握用洛必达法则求未定式极限的方法。
\item 理解函数的极值概念，掌握用导数判断函数的单调性和求函数极值的方法，掌握函数的最大值和最小值的求法及其应用。
\item 会用导数判断函数图形的凹凸性（注：在区间 $(a,b)$ 内，设函数 $f(x)$ 具有二阶导数。当 $f''(x)>0$ 时，$f(x)$ 的图形是凹的；当 $f''(x)<0$ 时，$f(x)$ 的图形是凸的），会求函数图形的拐点以及水平、铅直和斜渐近线，会描绘函数的图形。
\item 了解曲率、曲率圆与曲率半径的概念，会计算曲率和曲率半径。
\end{enumerate}
\end{quote}
\subsection*{本章星级总表}
\addcontentsline{toc}{subsection}{本章星级总表}

\begingroup\small\setlength{\tabcolsep}{4.0pt}
\begin{longtable}{@{}L{0.4327\linewidth}L{0.1154\linewidth}L{0.1154\linewidth}L{0.1154\linewidth}L{0.1154\linewidth}@{}}
\toprule
\textbf{考点} & \textbf{星级} & \textbf{考纲要求} & \textbf{难度} & \textbf{常见题型} \\
\midrule
\endhead
导数的定义与用定义判定可导性（分段点、绝对值、抽象函数） & $\bigstar$$\bigstar$$\bigstar$$\bigstar$$\bigstar$ & 理解 & 难 & 选择、解答 \\
可导与连续的关系、左导数与右导数 & $\bigstar$$\bigstar$$\bigstar$$\bigstar$ & 理解 & 中等 & 选择 \\
导数的几何意义、切线与法线；导数的物理意义 & $\bigstar$$\bigstar$ & 理解/了解 & 基础 & 填空 \\
导数的四则运算与基本初等函数导数公式 & $\bigstar$$\bigstar$$\bigstar$$\bigstar$ & 掌握 & 基础 & 选择、填空 \\
复合函数、反函数、隐函数、参数方程求导 & $\bigstar$$\bigstar$$\bigstar$$\bigstar$ & 掌握/会 & 中等 & 填空、解答 \\
高阶导数与莱布尼茨公式 & $\bigstar$$\bigstar$$\bigstar$$\bigstar$ & 了解/会 & 中等 & 填空、解答 \\
微分的概念、微分的四则运算与一阶微分形式不变性 & $\bigstar$$\bigstar$ & 理解/了解 & 中等 & 选择 \\
可微与可导的关系、微分概念题 & $\bigstar$$\bigstar$$\bigstar$ & 理解 & 中等 & 选择 \\
罗尔定理、拉格朗日中值定理及其应用 & $\bigstar$$\bigstar$$\bigstar$$\bigstar$$\bigstar$ & 理解并会用 & 难 & 证明 \\
柯西中值定理（了解并会用） & $\bigstar$$\bigstar$$\bigstar$ & 了解并会用 & 难 & 证明 \\
泰勒定理与泰勒公式的应用（含中值定理证明中的泰勒法） & $\bigstar$$\bigstar$$\bigstar$$\bigstar$$\bigstar$ & 理解并会用 & 难 & 证明、解答 \\
洛必达法则及其适用条件与失效情形 & $\bigstar$$\bigstar$$\bigstar$$\bigstar$ & 掌握 & 中等 & 选择、解答 \\
函数单调性的判别 & $\bigstar$$\bigstar$$\bigstar$$\bigstar$ & 掌握 & 基础 & 解答 \\
函数的极值（含不可导点、驻点判定） & $\bigstar$$\bigstar$$\bigstar$$\bigstar$ & 理解/掌握 & 中等 & 解答 \\
最大值、最小值及其应用 & $\bigstar$$\bigstar$$\bigstar$ & 掌握 & 中等 & 解答、应用 \\
凹凸性、拐点、渐近线 & $\bigstar$$\bigstar$$\bigstar$ & 会 & 中等 & 选择、填空 \\
函数图形的描绘 & $\bigstar$$\bigstar$ & 会 & 中等 & 解答 \\
曲率、曲率圆与曲率半径的概念与计算 & $\bigstar$$\bigstar$ & 了解/会 & 基础 & 填空 \\
中值定理证明题的辅助函数构造 & $\bigstar$$\bigstar$$\bigstar$$\bigstar$$\bigstar$ & 会用 & 难 & 证明 \\
不等式证明与放缩技巧 & $\bigstar$$\bigstar$$\bigstar$$\bigstar$ & 会用 & 难 & 证明 \\
方程根的个数与零点存在性讨论 & $\bigstar$$\bigstar$$\bigstar$$\bigstar$ & 会用 & 难 & 解答 \\
分段函数的导数与可导性综合 & $\bigstar$$\bigstar$$\bigstar$$\bigstar$ & 会 & 中等 & 解答 \\
\bottomrule
\end{longtable}
\endgroup

\section{导数与微分的概念}

\subsection{导数的定义与可导性判定}

\subsubsection{导数的定义与用定义判定可导性（分段点、绝对值、抽象函数） $\bigstar$$\bigstar$$\bigstar$$\bigstar$$\bigstar$}

设 $y=f(x)$ 在 $x_0$ 的某邻域内有定义，给自变量一个增量 $\Delta x$（$x_0+\Delta x$ 仍在该邻域内），函数增量 $\Delta y=f(x_0+\Delta x)-f(x_0)$。若极限 $\lim_{\Delta x\to 0}\frac{\Delta y}{\Delta x}$ 存在，则称 $f(x)$ 在 $x_0$ 处\textbf{可导}，该极限称为 $f(x)$ 在 $x_0$ 处的\textbf{导数}，记作 $f'(x_0)$。若该极限不存在，则称 $f$ 在 $x_0$ 处不可导。

导数的三种等价写法，做题时哪个方便用哪个：

\[f'(x_0)=\lim_{\Delta x\to 0}\frac{f(x_0+\Delta x)-f(x_0)}{\Delta x}
=\lim_{h\to 0}\frac{f(x_0+h)-f(x_0)}{h}
=\lim_{x\to x_0}\frac{f(x)-f(x_0)}{x-x_0}.\]

\textbf{用定义求导的三步}：

\begin{enumerate}
\item 写增量比 $\frac{f(x_0+\Delta x)-f(x_0)}{\Delta x}$；
\item 化简（通分、有理化、提公因式、用三角恒等式）；
\item 令 $\Delta x\to 0$ 取极限。
\end{enumerate}
\textbf{抽象函数的常见变形}：设 $f$ 在 $x_0$ 处可导，则对任意常数 $a,b$ 有

\[\lim_{h\to 0}\frac{f(x_0+ah)-f(x_0+bh)}{h}=(a-b)f'(x_0).\]

推导方法：把分子拆成 $[f(x_0+ah)-f(x_0)]-[f(x_0+bh)-f(x_0)]$，再分别凑成定义式（各除以 $ah$、$bh$）。同理

\[\lim_{h\to 0}\frac{f(x_0+h)-f(x_0-h)}{2h}=f'(x_0),\qquad
\lim_{h\to 0}\frac{f(x_0+h)-f(x_0)}{h}=f'(x_0).\]

\textbf{关键限制}：分母中的变量必须与分子中的增量严格对应；若出现 $h\to 0^+$ 或分母阶数不匹配（如 $\frac{f(x_0+h)-f(x_0)}{h^2}$），就\textbf{得不到} $f'(x_0)$。

\textbf{例}：设 $f(x)=(x-1)(x-2)\cdots(x-n)$，求 $f'(1)$。 因为 $f(1)=0$，用定义：

\[f'(1)=\lim_{x\to 1}\frac{f(x)-f(1)}{x-1}
=\lim_{x\to 1}(x-2)(x-3)\cdots(x-n)=(-1)(-2)\cdots(1-n)=(-1)^{n-1}(n-1)!.\]

\textbf{例}：设 $f$ 在 $x_0=2$ 处可导，$f'(2)=3$，求 $\lim_{h\to 0}\frac{f(2+3h)-f(2-5h)}{h}$。 由上面的公式，系数 $a=3,b=-5$，故极限为 $(3-(-5))f'(2)=8\times 3=24$。

\textbf{例（抽象函数）}：设 $f(x)$ 在 $x=0$ 处可导，$f'(0)=1$，且对任意 $x,y$ 有 $f(x+y)=f(x)+f(y)$。求 $f'(x)$。 先取 $x=y=0$：$f(0)=2f(0)$，得 $f(0)=0$。再由定义，对任意固定的 $x$：

\[f'(x)=\lim_{h\to 0}\frac{f(x+h)-f(x)}{h}=\lim_{h\to 0}\frac{f(x)+f(h)-f(x)}{h}=\lim_{h\to 0}\frac{f(h)-f(0)}{h}=f'(0)=1.\]

即 $f'(x)=1$ 对一切 $x$ 成立。这个套路（\textbf{抽象函数求导一律回到定义，并利用函数方程把不同点转化成同一点}）在真题中反复出现。

\textbf{这一考点的三类题型}（识别题型即可立刻选工具）：

\begin{enumerate}
\item \textbf{分段点处}：先验连续，再用左右导数定义；
\item \textbf{含绝对值或 $\max,\min$}：先化为分段函数或分类讨论，再逐段处理；
\item \textbf{抽象函数}：只能用定义，并借助函数方程、奇偶性、周期性把待求极限凑成 $\frac{f(x_0+\Delta x)-f(x_0)}{\Delta x}$。
\end{enumerate}
\begin{callwarn}{易错}
\begin{itemize}
\item $f'(x_0)$ 与 $[f(x_0)]'$ 完全不同：后者是常数的导数，恒为 $0$。用定义求导时不能先把 $x_0$ 代进去。
\item 定义式中的极限是\textbf{双侧}的；只算 $\Delta x\to 0^+$ 得到的是右导数，不能推出可导。
\item 分母必须与分子的自变量增量一致，凑定义时"分子加一项减一项"要配平系数。
\item 若 $f$ 在 $x_0$ 处不连续，则必不可导，且此时用定义算出的极限一定不存在，不必再分批左右。
\end{itemize}
\end{callwarn}

\begin{callexam}{真题考法}
选择题的典型问法：给出一个形如 $\lim\limits_{h\to 0}\frac{f(x_0+ah)-f(x_0+bh)}{h}$ 或 $\lim\limits_{x\to x_0}\frac{f(x)-f(x_0)}{x-x_0}$ 的极限值，问 $f$ 在 $x_0$ 处是否可导、$f'(x_0)$ 是多少。陷阱有三个：分子是不同点的函数值之差、分母阶数与分子不匹配、把函数的连续性误当可导性。判可导或求导数值只能依赖定义，不要试图用求导公式。

\end{callexam}

\subsubsection{可导与连续的关系、左导数与右导数 $\bigstar$$\bigstar$$\bigstar$$\bigstar$}

\textbf{左、右导数}：

\[f'_-(x_0)=\lim_{\Delta x\to 0^-}\frac{f(x_0+\Delta x)-f(x_0)}{\Delta x},
\qquad
f'_+(x_0)=\lim_{\Delta x\to 0^+}\frac{f(x_0+\Delta x)-f(x_0)}{\Delta x}.\]

\textbf{可导的充要条件}：$f$ 在 $x_0$ 处可导等价于 $f'_-(x_0)$ 与 $f'_+(x_0)$ 都存在且相等，此时 $f'(x_0)=f'_-(x_0)=f'_+(x_0)$。

\textbf{可导与连续的关系}：可导必连续，连续不一定可导。这条单向关系是选择题的常客。原因是

\[f(x)-f(x_0)=\frac{f(x)-f(x_0)}{x-x_0}\cdot(x-x_0)\to f'(x_0)\cdot 0=0\quad(x\to x_0).\]

\textbf{连续但不可导的典型函数}（要能背出反例）：

\begin{itemize}
\item $f(x)=|x|$ 在 $x=0$：$f'_-(0)=-1\ne f'_+(0)=1$，尖点型；
\item $f(x)=\sqrt[3]{x}$ 在 $x=0$：$f'(0)=\lim_{h\to 0}\frac{\sqrt[3]{h}}{h}=\lim_{h\to 0}h^{-2/3}=+\infty$，垂直切线型；
\item $f(x)=x\sin\frac1x\ (x\ne 0)$，$f(0)=0$：$\frac{f(h)-f(0)}{h}=\sin\frac1h$ 无极限，振荡型。
\end{itemize}
\textbf{不可导的四种典型形态}：间断点、尖点（左右导数存在但不相等）、垂直切线（导数趋于无穷）、振荡（导数极限不存在）。

\textbf{导函数自身的性质}（选择题常考概念判断）：

\begin{itemize}
\item 导函数 $f'(x)$ \textbf{不一定连续}。典型例子是 $f(x)=x^2\sin\frac1x\ (x\ne0)$，$f(0)=0$：$f'(0)=0$ 存在，但当 $x\to0$ 时 $f'(x)=2x\sin\frac1x-\cos\frac1x$ 无极限，即 $f'$ 在 $0$ 处不连续——所以"$f$ 可导"推不出"$f'$ 连续"。
\item 导函数具有\textbf{介值性}（不属于大纲要求，只需知道"$f'$ 不一定连续、也不一定单调"即可）：若 $f$ 在 $[a,b]$ 上可导，则 $f'$ 能取到 $f'(a)$ 与 $f'(b)$ 之间的一切值。数二考纲内的证明题一律用罗尔、拉格朗日、柯西、泰勒四条定理，不需要此结论。
\item 若 $f'(x_0)>0$，只能说明 $f$ 在该点附近沿 $x$ 增大方向上升，\textbf{不能}由此断定 $f$ 在 $x_0$ 的某个邻域内单调递增（反例：$f(x)=x+2x^2\sin\frac1x$ 类函数）。
\end{itemize}
\textbf{$|f(x)|$ 的可导性判定（高频结论）}：设 $f$ 在 $x_0$ 处可导。

\begin{itemize}
\item 若 $f(x_0)\ne 0$：$|f(x)|$ 在 $x_0$ 处可导（因为在 $x_0$ 附近 $f$ 不变号，$|f|=\pm f$）；
\item 若 $f(x_0)=0$：$|f(x)|$ 在 $x_0$ 处可导的充要条件是 $f'(x_0)=0$，此时导数为 $0$。
\end{itemize}
\textbf{例}：$g(x)=(x^2-1)|x-1|$，讨论 $g$ 在 $x=1$ 处的可导性。 $g(1)=0$。$\frac{g(1+h)-g(1)}{h}=\frac{((1+h)^2-1)|h|}{h}=((1+h)^2-1)\frac{|h|}{h}=(2h+h^2)\frac{|h|}{h}$。 右导数：$h\to 0^+$ 时为 $(2h+h^2)\to 0$；左导数：$h\to 0^-$ 时为 $-(2h+h^2)\to 0$。左右相等，故 $g'(1)=0$。（也可用结论：此处 $f(x)=x^2-1$ 满足 $f(1)=0$ 且 $f'(1)=2\ne0$——注意这是把 $g$ 写成 $(x^2-1)\vert x-1\vert$，属于乘积形式而非 $|f|$ 形式，所以不能用上述结论，必须用定义。）

\begin{callwarn}{易错}
\begin{itemize}
\item "连续"与"可导"不能互推；把"可导必连续"记反是最常见的低级失分。
\item 求左右导数时必须回到定义，不能先求导函数再代入（分段点两侧表达式不同）。
\item 对 $|f(x)|$ 用结论前先确认形式：若是 $|f(x)|$ 本身且 $f(x_0)=0$，结论是"可导等价于 $f'(x_0)=0$"；若是 $g(x)|x-x_0|$ 这类乘积，要改用定义。
\item 连续不一定可导，但"可导"一定不能只看左右极限相等（那是函数极限，不是导数）。
\end{itemize}
\end{callwarn}

\begin{callexam}{真题考法}
选择题标准形式有三类：$\circled{1}$ 判断给定分段函数在某点"连续但不可导"；$\circled{2}$ 已知 $f$ 在 $x_0$ 可导，问 $|f(x)|$、$f(x)|x-x_0|$、$f(x)g(x)$ 的可导性；$\circled{3}$ 已知 $f$ 可导、$g$ 不可导，问 $f+g$、$fg$ 的可导性（结论取决于 $f(x_0)$ 是否为 $0$）。作答时要分情况讨论，不能一概而论。

\end{callexam}

\subsection{导数与微分的几何、物理意义}

\subsubsection{导数的几何意义、切线与法线；导数的物理意义 $\bigstar$$\bigstar$}

$f'(x_0)$ 就是曲线 $y=f(x)$ 在点 $(x_0,f(x_0))$ 处切线的斜率。由此得到

\textbf{切线方程}：$y-f(x_0)=f'(x_0)(x-x_0)$。 \textbf{法线方程}（当 $f'(x_0)\ne 0$）：$y-f(x_0)=-\frac{1}{f'(x_0)}(x-x_0)$。

两个特殊情形必须单独记：

\begin{itemize}
\item 若 $f'(x_0)=0$，切线为水平线 $y=f(x_0)$，法线为铅直线 $x=x_0$；
\item 若 $f'(x_0)$ 不存在但曲线有垂直切线（如 $y=\sqrt[3]{x}$ 在原点），切线为 $x=x_0$，法线为 $y=f(x_0)$。
\end{itemize}
\textbf{各类曲线求切线的方法}：

\begin{enumerate}
\item \textbf{显函数} $y=f(x)$：直接求 $f'(x_0)$；
\item \textbf{隐函数} $F(x,y)=0$：两边对 $x$ 求导解出 $y'$，或用 $y'=-\frac{F'_x}{F'_y}$；
\item \textbf{参数方程} $x=\varphi(t),y=\psi(t)$：斜率 $\frac{dy}{dx}=\frac{\psi'(t_0)}{\varphi'(t_0)}$，切点取 $(\varphi(t_0),\psi(t_0))$；
\item \textbf{极坐标} $r=r(\theta)$：先化为参数方程 $x=r\cos\theta,\ y=r\sin\theta$，再按参数方程处理。
\end{enumerate}
\textbf{导数的物理意义}：位置 $s(t)$ 的导数 $v(t)=s'(t)$ 是瞬时速度；速度的导数 $a(t)=v'(t)=s''(t)$ 是加速度。相关变化率公式

\[\frac{dy}{dt}=\frac{dy}{dx}\cdot\frac{dx}{dt}.\]

\textbf{例}：求曲线 $y=x^3-3x$ 在 $x=1$ 处的切线方程与法线方程。 $y'=3x^2-3$，$y'(1)=0$，切点为 $(1,-2)$。切线为 $y=-2$；法线为 $x=1$。

\textbf{例}：求过点 $(0,0)$ 且与曲线 $y=x^3-3x$ 相切的直线。 设切点 $(t,t^3-3t)$，斜率 $k=3t^2-3$。切线过原点，故 $t^3-3t=(3t^2-3)t$，即 $2t^3=0$，得 $t=0$。切点为原点，切线 $y=-3x$。

\textbf{例（相关变化率）}：一气球半径以 $0.2\ \mathrm{m/s}$ 的速率增大，求半径等于 $5\ \mathrm{m}$ 时体积的变化率。 设半径为 $r$，体积 $V=\frac43\pi r^3$，两边对 $t$ 求导：

\[\frac{dV}{dt}=4\pi r^2\frac{dr}{dt}.\]

代入 $r=5$、$\frac{dr}{dt}=0.2$，得 $\frac{dV}{dt}=4\pi\cdot 25\cdot 0.2=20\pi\ (\mathrm{m^3/s})$。

\textbf{例（相关变化率）}：圆锥形水池深 $10\ \mathrm{m}$、顶面半径 $5\ \mathrm{m}$，以 $3\ \mathrm{m^3/min}$ 注水，求水深 $5\ \mathrm{m}$ 时水面上升的速率。 设水深 $h$，水面半径 $r=\frac{h}{2}$（由相似比 $\frac{r}{h}=\frac{5}{10}$），故 $V=\frac13\pi r^2h=\frac{\pi h^3}{12}$，从而 $\frac{dV}{dt}=\frac{\pi h^2}{4}\frac{dh}{dt}$。代入 $\frac{dV}{dt}=3$、$h=5$，得 $\frac{dh}{dt}=\frac{12}{\pi\cdot 25}=\frac{12}{25\pi}\ (\mathrm{m/min})$。

\textbf{要点}：相关变化率问题的固定做法是——$\circled{1}$ 找出各变量之间的几何（或物理）关系式；$\circled{2}$ 两边对 $t$ 求导，把未知的变化率解出来；$\circled{3}$ 再代入具体数值。\textbf{必须先求导后代值}。

\begin{callwarn}{易错}
\begin{itemize}
\item "在点 $P$ 处的切线"与"过点 $P$ 的切线"是两回事：前者 $P$ 是切点，后者要设切点 $(t,f(t))$ 再解方程，可能有多条。
\item 法线斜率是切线斜率的负倒数，只在切线斜率非零时有意义。
\item 求"曲线与直线相切"的参数时，要同时满足\textbf{函数值相等}与\textbf{导数值相等}两个方程。
\item 切线截距型问题（切线与两坐标轴围成的面积）要先写出截距再算，别用错点。
\end{itemize}
\end{callwarn}

\begin{callexam}{真题考法}
填空题多考"求切线/法线方程""求公切线""求切点坐标"；解答题中会出现"已知两曲线相切求参数""曲线与直线相切"等综合问法。陷阱集中在"过某点的切线要分切点讨论，常有两解"，以及参数方程求切线时忘记把 $t_0$ 换成切点坐标。

\end{callexam}

\subsubsection{微分的概念、微分的四则运算与一阶微分形式不变性 $\bigstar$$\bigstar$}

\textbf{定义}：设 $y=f(x)$ 在 $x_0$ 的某邻域内有定义，若函数增量可表示为

\[\Delta y=f(x_0+\Delta x)-f(x_0)=A\Delta x+o(\Delta x)\quad(\Delta x\to 0),\]

其中 $A$ 是不依赖于 $\Delta x$ 的常数，则称 $f$ 在 $x_0$ 处\textbf{可微}，$A\Delta x$ 称为 $f$ 在 $x_0$ 处的\textbf{微分}，记作 $dy=A\Delta x=A\,dx$。

注意 $dx=\Delta x$ 只对\textbf{自变量}成立。

\textbf{核心关系}：一元函数 $f$ 在 $x_0$ 处可微的充要条件是 $f$ 在 $x_0$ 处可导，且 $A=f'(x_0)$。于是

\[dy=f'(x_0)\,dx,\qquad dy=f'(x)\,dx.\]

\textbf{几何意义}：$dy$ 是曲线在 $(x_0,f(x_0))$ 处切线上对应于增量 $\Delta x$ 的纵坐标增量，而 $\Delta y=dy+o(\Delta x)$。所以当 $|\Delta x|$ 很小时 $\Delta y\approx dy$，得到近似公式

\[f(x_0+\Delta x)\approx f(x_0)+f'(x_0)\Delta x,\]

常用取值 $x_0$ 使 $f(x_0)$、$f'(x_0)$ 都是好算的数。

\textbf{四则运算}（$u,v$ 可微）：

\[d(u\pm v)=du\pm dv,\qquad d(uv)=v\,du+u\,dv,\qquad d\!\left(\frac uv\right)=\frac{v\,du-u\,dv}{v^2}\ (v\ne0).\]

\textbf{一阶微分形式不变性}：无论 $u$ 是自变量还是中间变量，$y=f(u)$ 恒有

\[dy=f'(u)\,du.\]

这正是复合函数求导法则的微分形式：$y=f(u),u=\varphi(x)$ 时 $dy=f'(u)\varphi'(x)dx=f'(u)du$。

\textbf{求微分的四步}：$\circled{1}$ 求导数 $f'(x)$；$\circled{2}$ 写 $dy=f'(x)dx$；$\circled{3}$ 需要时代入 $x=a$ 得 $dy|_{x=a}=f'(a)dx$；$\circled{4}$ 用 $\Delta y\approx dy$ 做近似。

\textbf{例}：$y=\ln(1+2x)$，求 $dy$。$y'=\frac{2}{1+2x}$，故 $dy=\frac{2}{1+2x}dx$。

\textbf{例（近似计算）}：估算 $\sqrt[3]{8.02}$。 取 $f(x)=\sqrt[3]{x}$，$x_0=8$，$\Delta x=0.02$。$f'(x)=\frac{1}{3}x^{-2/3}$，$f'(8)=\frac{1}{12}$，

\[\sqrt[3]{8.02}\approx 2+\frac{1}{12}\times 0.02\approx 2.00167.\]

\textbf{例（一阶微分形式不变性）}：$y=\sin(e^{2x})$，求 $dy$。 把 $e^{2x}$ 看作中间变量 $u$：$dy=\cos u\,du=\cos(e^{2x})\cdot 2e^{2x}dx=2e^{2x}\cos(e^{2x})dx$。

\begin{callwarn}{易错}
\begin{itemize}
\item $\Delta y$ 与 $dy$ 一般不相等，$\Delta y=dy+o(\Delta x)$；只有当 $f$ 是一次函数时二者才恒等。
\item $dx=\Delta x$ 只是对自变量成立，对因变量 $dy\ne\Delta y$。
\item 用微分做近似时 $\Delta x$ 要小，且 $x_0$ 必须选在函数值与导数都容易计算的点。
\item 求微分不要漏写 $dx$；填空题只写 $f'(x)$ 会被判为导数而非微分。
\end{itemize}
\end{callwarn}

\begin{callexam}{真题考法}
选择题常问"$f$ 在 $x_0$ 可微的充要条件"，干扰项写"$\Delta y\sim A\Delta x$""$dy=f'(x_0)$"（少了 $dx$）等。填空题给一个函数求 $dy$ 或 $dy|_{x=a}$，偶尔结合近似计算（如估算 $\sqrt[3]{8.02}$、$\ln 1.01$）。

\end{callexam}

\subsubsection{可微与可导的关系、微分概念题 $\bigstar$$\bigstar$$\bigstar$}

三条关系必须记成一条链：\textbf{可微与可导等价，可导必连续}。写成一句话就是：

\begin{quote}
$f$ 在 $x_0$ 处可微 $\iff$ $f$ 在 $x_0$ 处可导 $\Longrightarrow$ $f$ 在 $x_0$ 处连续。

\end{quote}
其中"可微与可导等价"只对\textbf{一元函数}成立；"可导必连续、连续不一定可导"是单向的。

\textbf{判断可微的两个等价说法}：

\[\lim_{\Delta x\to 0}\frac{\Delta y-f'(x_0)\Delta x}{\Delta x}=0
\iff \Delta y=f'(x_0)\Delta x+o(\Delta x).\]

\textbf{概念题常见考法}：给出 $f$ 在 $x_0$ 处连续的某个条件（如 $\lim_{x\to x_0}\frac{f(x)-f(x_0)}{x-x_0}$ 存在、$\lim_{h\to 0}\frac{f(x_0+h)-f(x_0-h)}{h}$ 存在），让考生判断可导性与微分存在性。要抓住两点：

\begin{itemize}
\item $\lim_{h\to 0}\frac{f(x_0+h)-f(x_0-h)}{h}$ 存在\textbf{推不出} $f'(x_0)$ 存在，因为 $h$ 与 $-h$ 处的函数值"两边一起动"，可以恰好抵消（例如 $f(x)=|x|$ 在 $0$ 处该极限为 $0$，但 $f'(0)$ 不存在）；
\item $\lim_{h\to 0}\frac{f(x_0+h)-f(x_0)}{h}$ 存在才是 $f'(x_0)$ 存在。
\end{itemize}
\textbf{例}：设 $f(x)=|x|$，判断 $f$ 在 $0$ 处是否可微。 $\Delta y=|0+\Delta x|-0=|\Delta x|$。若可微，则 $A=f'(0)$ 应存在，但 $f'_-(0)=-1$、$f'_+(0)=1$ 不等，故不可导，从而不可微。

\textbf{例}：设 $f(x)$ 在 $x=0$ 处连续，且 $\lim_{x\to 0}\frac{f(x)}{x}=1$，问 $f$ 在 $0$ 处是否可微。 由 $\frac{f(x)}{x}\to 1$ 得 $f(x)\to 0$，又 $f$ 连续，故 $f(0)=0$。于是 $\lim_{x\to 0}\frac{f(x)-f(0)}{x-0}=1$，即 $f'(0)=1$ 存在，所以 $f$ 在 $0$ 处可微，且 $df=f'(0)dx=dx$。

\textbf{微分在近似与误差估计中的应用}：由于 $\Delta y=dy+o(\Delta x)$，当 $|\Delta x|$ 很小时可用 $dy$ 代替 $\Delta y$：

\begin{itemize}
\item 函数值的近似：$f(x_0+\Delta x)\approx f(x_0)+f'(x_0)\Delta x$；
\item 绝对误差：$|\Delta y-dy|$ 是比 $|\Delta x|$ 高阶的无穷小，故误差随 $\Delta x$ 的减小而迅速减小；
\item 相对误差：$\frac{\Delta y}{y}\approx\frac{dy}{y}=\frac{f'(x_0)\Delta x}{f(x_0)}$，常用于"测量 $x$ 时产生 $\Delta x$ 的误差，求 $y$ 的相对误差"这类题。
\end{itemize}
\textbf{例}：测得正方体棱长 $x=10\ \mathrm{cm}$，测量误差 $|\Delta x|\le0.05\ \mathrm{cm}$，求体积 $V=x^3$ 的绝对误差与相对误差限。 $dV=3x^2dx$，故 $|\Delta V|\approx 3\times 100\times 0.05=15\ (\mathrm{cm^3})$；相对误差限约为 $\frac{15}{1000}=1.5\%$。

\begin{callwarn}{易错}
\begin{itemize}
\item 一元函数中"可导"与"可微"完全等价，不要写成"可导是可微的充分不必要条件"。
\item 判断连续只需验 $\lim_{x\to x_0}f(x)=f(x_0)$；判断可导必须验左右导数；判断可微在一元情形直接等价于可导。
\item 对称形式的极限 $\lim_{h\to 0}\frac{f(x_0+h)-f(x_0-h)}{2h}$ 存在\textbf{不能}推出可导（反例 $|x|$），要格外小心。
\end{itemize}
\end{callwarn}

\begin{callexam}{真题考法}
选择题把"连续""可导""可微""极限存在"做成四个选项，唯一确定的方向是"可导 $\Rightarrow$ 连续"与"一元中可导等价于可微"。另一类是给分段函数问分段点的可微性，本质上是问可导性，用左右导数判定即可。

\end{callexam}

\subsubsection{分段函数的导数与可导性综合 $\bigstar$$\bigstar$$\bigstar$$\bigstar$}

分段函数求导是解答题的常见起点，核心是\textbf{分段点单独处理}。

\textbf{固定流程}：

\begin{enumerate}
\item \textbf{先验连续性}：算 $\lim_{x\to x_0^-}f(x)$、$\lim_{x\to x_0^+}f(x)$，若二者不相等或与 $f(x_0)$ 不等，则必不可导，直接结束；
\item \textbf{在分段点用定义}分别求 $f'_-(x_0)$、$f'_+(x_0)$，不能直接对两段表达式求导后代入；
\item 若 $f'_-(x_0)=f'_+(x_0)$，则 $f'(x_0)$ 取该公共值；否则 $f$ 在 $x_0$ 处不可导；
\item 在非分段点用求导公式求导，最后把 $f'(x)$ 写成分段函数，并标明分段点的取值。
\end{enumerate}
\textbf{快捷方法}：若 $f$ 在 $x_0$ 的两侧分别由连续可导的 $g(x)$（$x<x_0$）、$h(x)$（$x>x_0$）给出，且 $\lim_{x\to x_0^-}g'(x)$、$\lim_{x\to x_0^+}h'(x)$ 都存在，则

\[f'_-(x_0)=\lim_{x\to x_0^-}g'(x),\qquad f'_+(x_0)=\lim_{x\to x_0^+}h'(x).\]

这能大幅减少计算量。但要注意它只是\textbf{充分条件}：

\textbf{例}：$f(x)=\begin{cases}x^2\sin\frac1x, & x\ne 0,\\ 0, & x=0.\end{cases}$ 判断 $f'(0)$ 是否存在。

\[f'(0)=\lim_{x\to 0}\frac{x^2\sin\frac1x-0}{x}=\lim_{x\to 0}x\sin\frac1x=0\]

（无穷小乘有界量），故 $f'(0)=0$ 存在。

而当 $x\ne0$ 时 $f'(x)=2x\sin\frac1x-\cos\frac1x$，$\lim_{x\to 0}f'(x)$ 不存在。这说明\textbf{导函数在分段点可以不连续}，但 $f'(0)$ 确实存在——分段点必须回到定义。

\textbf{例（含参求参数）}：设 $f(x)=\begin{cases}x^2, & x\le 1,\\ ax+b, & x>1,\end{cases}$ 在 $x=1$ 处可导，求 $a,b$。 连续性：$\lim_{x\to1^-}f=1$，故 $a+b=1$。 左右导数：$f'_-(1)=\lim_{h\to0^-}\frac{(1+h)^2-1}{h}=2$；$f'_+(1)=\lim_{h\to0^+}\frac{a(1+h)+b-1}{h}=\frac{ah+a+b-1}{h}=a$。故 $a=2$，再由 $a+b=1$ 得 $b=-1$。

\begin{callwarn}{易错}
\begin{itemize}
\item 在分段点直接用导函数求值（先求 $f'$ 再代 $x_0$）是最常见的错误，必须用定义。
\item 若 $f$ 在 $x_0$ 处不连续则必不可导，先验连续性可以省去大量计算。
\item 分段点的归属（$\le$ 还是 $<$）要与题目一致，写 $f'(x)$ 时别漏掉 $(x_0,f'(x_0))$ 这一点。
\item 含绝对值函数的可导性讨论，本质是把绝对值化为分段函数后再按上述流程处理。
\end{itemize}
\end{callwarn}

\begin{callexam}{真题考法}
解答题常给含绝对值或含 $\sin\frac1x$、$e^{-1/x^2}$ 型的分段函数，要求写出 $f'(x)$ 并讨论分段点的可导性；另一类是"已知 $f$ 在分段点可导求参数"，需要同时用连续与左右导数相等两个方程。

\end{callexam}

\section{求导方法与高阶导数}

\subsection{基本求导法则}

\subsubsection{导数的四则运算与基本初等函数导数公式 $\bigstar$$\bigstar$$\bigstar$$\bigstar$}

\textbf{四则运算法则}：

\[(u\pm v)'=u'\pm v',\qquad (uv)'=u'v+uv',\qquad \left(\frac uv\right)'=\frac{u'v-uv'}{v^2}\ (v\ne0).\]

\textbf{基本初等函数导数公式}（必须闭卷默写）：

\begingroup\small\setlength{\tabcolsep}{4.0pt}
\begin{longtable}{@{}L{0.2278\linewidth}L{0.2278\linewidth}L{0.2278\linewidth}L{0.2278\linewidth}@{}}
\toprule
\textbf{原函数} & \textbf{导数} & \textbf{原函数} & \textbf{导数} \\
\midrule
\endhead
$C$（常数） & $0$ & $\sin x$ & $\cos x$ \\
$x^\alpha$ & $\alpha x^{\alpha-1}$ & $\cos x$ & $-\sin x$ \\
$a^x$ & $a^x\ln a$ & $\tan x$ & $\sec^2x$ \\
$e^x$ & $e^x$ & $\cot x$ & $-\csc^2x$ \\
$\log_ax$ & $\frac{1}{x\ln a}$ & $\sec x$ & $\sec x\tan x$ \\
$\ln x$ & $\frac1x$ & $\csc x$ & $-\csc x\cot x$ \\
$\arcsin x$ & $\frac{1}{\sqrt{1-x^2}}$ & $\arctan x$ & $\frac{1}{1+x^2}$ \\
$\arccos x$ & $-\frac{1}{\sqrt{1-x^2}}$ & $\mathrm{arccot}\,x$ & $-\frac{1}{1+x^2}$ \\
\bottomrule
\end{longtable}
\endgroup

\begin{callkey}{必背}
三组最容易写错的公式，闭卷默写三遍：

\[\left(\frac{1}{x}\right)'=-\frac{1}{x^2},\qquad (\sqrt{x})'=\frac{1}{2\sqrt{x}},\qquad (\ln|x|)'=\frac{1}{x}\]

\[(\arcsin x)'=\frac{1}{\sqrt{1-x^2}},\qquad (\arctan x)'=\frac{1}{1+x^2}\]

\[(u^v)'=u^v\left(v'\ln u+\frac{vu'}{u}\right)\quad(u>0)\]

\end{callkey}

\textbf{求导的三条总原则}：

\begin{enumerate}
\item \textbf{先化简再求导}：能通分、约分、用对数性质拆开、用倍角公式降幂的，先化简，可以省掉大量运算；
\item \textbf{认清结构}：四则运算、复合、隐式、参数式分别用对应法则，链式法则优先；
\item \textbf{由外向内}：复合函数逐层剥开，且不要漏掉任何一层的导数。
\end{enumerate}
\textbf{幂指函数} $u(x)^{v(x)}$（底数与指数都含变量）：不能用幂函数或指数函数的公式。标准做法是恒等变形

\[u^v=e^{v\ln u}\ \Longrightarrow\ (u^v)'=u^v\left(v'\ln u+\frac{vu'}{u}\right).\]

\textbf{对数求导法}：适用于"多个因式的积、商、开方"以及幂指函数。两边取对数把乘除降为加减再求导，可以避免商的法则的繁琐计算。

\textbf{例}：$y=\sqrt{\dfrac{(x+2)(3-x)^4}{(x+1)^5}}$，求 $y'$。 两边取对数：

\[\ln y=\frac12\ln(x+2)+4\ln(3-x)-5\ln(x+1),\]

两端对 $x$ 求导：$\frac{y'}{y}=\frac{1}{2(x+2)}-\frac{4}{3-x}-\frac{5}{x+1}$，于是

\[y'=\sqrt{\frac{(x+2)(3-x)^4}{(x+1)^5}}\left[\frac{1}{2(x+2)}-\frac{4}{3-x}-\frac{5}{x+1}\right].\]

\textbf{例}：$y=x^{\sin x}\ (x>0)$，求 $y'$。 取对数：$\ln y=\sin x\ln x$，故

\[y'=x^{\sin x}\left(\cos x\ln x+\frac{\sin x}{x}\right).\]

\begin{callwarn}{易错}
\begin{itemize}
\item $(\ln x)'=\frac1x$，而 $(\ln|x|)'=\frac1x\ (x\ne0)$，两者形式相同但定义域不同。
\item $(\arcsin x)'$ 与 $(\arccos x)'$ 只差一个负号，最易写错。
\item 商的导数分子是 $u'v-uv'$，顺序颠倒会整体错号。
\item 幂指函数直接套 $(x^\alpha)'=\alpha x^{\alpha-1}$ 是高频错误，指数含变量时绝不能这样用。
\end{itemize}
\end{callwarn}

\begin{callexam}{真题考法}
基础题，多出现在填空与选择题：直接考公式与四则运算，或给出含指数、对数、三角的组合函数求 $f'(x_0)$。陷阱是幂指函数、分段函数、含绝对值函数直接套公式。

\end{callexam}

\subsubsection{复合函数、反函数、隐函数、参数方程求导 $\bigstar$$\bigstar$$\bigstar$$\bigstar$}

\textbf{(1) 复合函数（链式法则）}：$y=f[\varphi(x)]$，则

\[\frac{dy}{dx}=f'[\varphi(x)]\,\varphi'(x),\qquad dy=f'[\varphi(x)]\,\varphi'(x)\,dx.\]

必须记住 $\{f[\varphi(x)]\}'\ne f'[\varphi(x)]$（右端漏掉了内层导数 $\varphi'(x)$）。

\textbf{(2) 反函数求导}：设 $y=f(x)$ 与 $x=g(y)$ 互为反函数，$f'(x)\ne0$，则

\[g'(y)=\frac{1}{f'(x)}=\frac{1}{f'[g(y)]},\qquad g''(y)=-\frac{f''(x)}{[f'(x)]^3}.\]

\textbf{(3) 隐函数求导}：由 $F(x,y)=0$ 确定 $y=y(x)$，两边对 $x$ 求导，遇 $y$ 按复合函数处理（乘 $y'$），再解出 $y'$。也可用

\[y'=-\frac{F'_x}{F'_y}.\]

求二阶导时，把 $y'$ 的表达式再对 $x$ 求导，$y$ 仍视为 $x$ 的函数，最后把已求出的 $y'$ 代入化简。

\textbf{(4) 参数方程求导}：$x=\varphi(t),\ y=\psi(t)$，则

\[\frac{dy}{dx}=\frac{\psi'(t)}{\varphi'(t)},\qquad
\frac{d^2y}{dx^2}=\frac{\psi''(t)\varphi'(t)-\psi'(t)\varphi''(t)}{[\varphi'(t)]^3}.\]

二阶导的记忆方法：$\frac{d^2y}{dx^2}=\frac{d}{dt}\!\left(\frac{dy}{dx}\right)\Big/\frac{dx}{dt}$，\textbf{不要忘记再除以 $\varphi'(t)$}。

\textbf{(5) 极坐标曲线}：$r=r(\theta)$ 化为参数方程 $x=r(\theta)\cos\theta,\ y=r(\theta)\sin\theta$，以 $\theta$ 为参数：

\[\frac{dy}{dx}=\frac{r'(\theta)\sin\theta+r(\theta)\cos\theta}{r'(\theta)\cos\theta-r(\theta)\sin\theta}.\]

\textbf{例}：设 $x=t^2,\ y=\ln(1+t)$，求 $t=1$ 处的 $\frac{dy}{dx}$ 与 $\frac{d^2y}{dx^2}$。

\[\frac{dy}{dx}=\frac{\frac{1}{1+t}}{2t}=\frac{1}{2t(1+t)},\qquad \frac{dy}{dx}\Big|_{t=1}=\frac14 .\]

\[\frac{d^2y}{dx^2}=\frac{d}{dt}\!\left(\frac{1}{2t(1+t)}\right)\Big/2t
=\frac{-\frac{2(1+2t)}{4t^2(1+t)^2}}{2t}
=-\frac{1+2t}{4t^3(1+t)^2},\qquad \frac{d^2y}{dx^2}\Big|_{t=1}=-\frac{3}{16}.\]

Let me double check: $\frac{d}{dt}\left(\frac{1}{2t(1+t)}\right) = \frac{d}{dt}\left(\frac{1}{2(t+t^2)}\right) = -\frac{1+2t}{2(t+t^2)^2}$. Divide by $\frac{dx}{dt}=2t$: $-\frac{1+2t}{4t(t+t^2)^2} = -\frac{1+2t}{4t\cdot t^2(1+t)^2}=-\frac{1+2t}{4t^3(1+t)^2}$. $\checkmark$ At t=1: $-3/(4\cdot4)=-3/16$ $\checkmark$.

\textbf{例（隐函数）}：由 $e^y+xy=e$ 确定 $y=y(x)$，求 $y''(0)$。 $x=0$ 时 $e^y=e$，故 $y(0)=1$。两边对 $x$ 求导：

\[e^y y'+y+xy'=0.\]

代入 $x=0,y=1$：$ey'+1=0$，得 $y'(0)=-\frac1e$。再对 $x$ 求导：

\[e^y(y')^2+e^yy''+2y'+xy''=0.\]

代入 $x=0,y=1,y'=-\frac1e$：$e\cdot\frac{1}{e^2}+ey''-\frac2e=0$，得 $y''(0)=\frac{1}{e^2}$。

\textbf{例（反函数）}：设 $y=f(x)$ 单调可导，$f(1)=2$，$f'(1)=3$，求反函数 $x=g(y)$ 在 $y=2$ 处的导数。 $g'(2)=\frac{1}{f'(1)}=\frac13$。

\textbf{例（隐函数求导的完整流程）}：设 $\sin(xy)+x-y=0$ 确定 $y=y(x)$，求 $y'$。 两边对 $x$ 求导（注意 $y$ 是 $x$ 的函数）：

\[\cos(xy)\,(y+xy')+1-y'=0 .\]

解出：$y'=\frac{y\cos(xy)+1}{1-x\cos(xy)}$，只要分母不为零即可。

\textbf{对数求导法的适用条件}（用之前先判断）：

\begin{enumerate}
\item 目标函数是多个因式的积、商、乘方、开方（此时直接求导很繁琐）；
\item 函数是幂指函数 $u(x)^{v(x)}$；
\item 用对数求导法时，若函数可能取负值，应改为对 $|y|$ 取对数，或先讨论使 $y>0$ 的区间。
\end{enumerate}
\textbf{参数方程求导的两个易错细节}：

\begin{itemize}
\item 求 $\frac{dy}{dx}$ 时先把 $\psi'(t)$、$\varphi'(t)$ 都算出来再作商，避免中途代入 $t_0$ 造成混乱；
\item 求 $\frac{d^2y}{dx^2}$ 时，把 $\frac{dy}{dx}$ 看成 $t$ 的函数，先关于 $t$ 求导，再除以 $\varphi'(t)$。若题目只要某点处的值，也要先化简再代 $t_0$。
\end{itemize}
\begin{callwarn}{易错}
\begin{itemize}
\item 参数方程的二阶导最常错在"只把 $\frac{dy}{dx}$ 对 $t$ 求导就当成 $\frac{d^2y}{dx^2}$"，漏除了 $\frac{dx}{dt}$。
\item 隐函数求二阶导时忘记 $y$ 是 $x$ 的函数，导致含 $y'$ 的项被漏掉。
\item 反函数求导的分母是 $f'(x)$ 在\textbf{对应点}的值，写成 $\frac{1}{f'(y)}$ 是错的。
\item 复合函数求导时漏掉内层导数，尤其是内层为 $\ln$、$\sqrt{\ }$、反三角函数时。
\item 对数求导法在 $y<0$ 时应对 $|y|$ 取对数，或说明 $y>0$ 的区间。
\end{itemize}
\end{callwarn}

\begin{callexam}{真题考法}
填空题的常客：参数方程的 $\frac{dy}{dx}$、$\frac{d^2y}{dx^2}$，隐函数的 $y'$、$y''$，反函数的导数，幂指函数与对数求导法。解答题常与"求切线方程""求极值"结合。陷阱就是上面列的符号与漏项问题。

\end{callexam}

\subsection{高阶导数}

\subsubsection{高阶导数与莱布尼茨公式 $\bigstar$$\bigstar$$\bigstar$$\bigstar$}

\textbf{定义}：$f^{(n)}(x)=[f^{(n-1)}(x)]'$（$n\ge2$），记作 $y^{(n)}$ 或 $\frac{d^ny}{dx^n}$。二阶导数记 $f''$ 或 $y''$。

\textbf{求高阶导数的五条路}：

\begin{enumerate}
\item \textbf{逐次求导找规律}（再用数学归纳法证明）；
\item \textbf{公式法}（见下表）；
\item \textbf{莱布尼茨公式}（乘积形式）；
\item \textbf{有理分式先拆部分分式}，化为 $\frac{1}{ax+b}$ 型的和；
\item \textbf{泰勒展开反求} $f^{(n)}(0)$：若 $f(x)=a_0+a_1x+\cdots+a_nx^n+\cdots$，则 $f^{(n)}(0)=n!\,a_n$。
\end{enumerate}
\textbf{常用 $n$ 阶导数公式}：

\[\begin{aligned}
(e^{ax})^{(n)}&=a^ne^{ax}, & (a^x)^{(n)}&=a^x(\ln a)^n,\\
(\sin x)^{(n)}&=\sin\!\left(x+\frac{n\pi}{2}\right), & (\cos x)^{(n)}&=\cos\!\left(x+\frac{n\pi}{2}\right),\\
(\sin(ax+b))^{(n)}&=a^n\sin\!\left(ax+b+\frac{n\pi}{2}\right), & (\cos(ax+b))^{(n)}&=a^n\cos\!\left(ax+b+\frac{n\pi}{2}\right),\\
\left(\frac{1}{ax+b}\right)^{(n)}&=\frac{(-1)^na^nn!}{(ax+b)^{n+1}}, & (\ln(1+x))^{(n)}&=\frac{(-1)^{n-1}(n-1)!}{(1+x)^n},\\
(x^m)^{(n)}&=m(m-1)\cdots(m-n+1)x^{m-n}, & (x^n)^{(n)}&=n!.
\end{aligned}\]

\textbf{莱布尼茨公式}：

\[(uv)^{(n)}=\sum_{k=0}^{n}C_n^k\,u^{(n-k)}v^{(k)}
=u^{(n)}v+nu^{(n-1)}v'+\frac{n(n-1)}{2}u^{(n-2)}v''+\cdots+uv^{(n)}.\]

\textbf{使用要点}：把其中因式取为\textbf{多项式}（求若干次导后为 $0$，求和式自动截断），另一个取 $e^{ax}$ 或 $\sin,\cos$ 等公式已知的函数。

\textbf{例}：设 $y=x^2e^{2x}$，求 $y^{(20)}$。 取 $v=x^2$（三次以上导数为 0），$u=e^{2x}$：

\[y^{(20)}=C_{20}^0x^2(e^{2x})^{(20)}+C_{20}^1(2x)(e^{2x})^{(19)}+C_{20}^2\cdot 2\cdot(e^{2x})^{(18)}
=e^{2x}\left(2^{20}x^2+40\cdot 2^{19}x+380\cdot 2^{18}\right).\]

\textbf{例（奇偶性）}：设 $f(x)=\arctan x$，求 $f^{(100)}(0)$。 $f$ 是奇函数（$\arctan(-x)=-\arctan x$），奇函数求偶数阶导数仍为奇函数，故 $f^{(100)}(x)$ 是奇函数，$f^{(100)}(0)=0$。

\textbf{奇偶性规律}（可省大量计算）：

\begin{itemize}
\item 可导奇函数的导函数是偶函数，可导偶函数的导函数是奇函数；
\item 于是奇函数在 $0$ 处的\textbf{偶数阶}导数为 $0$，偶函数在 $0$ 处的\textbf{奇数阶}导数为 $0$；
\item 若 $f$ 是奇函数且 $f^{(n)}(0)$ 存在，则 $n$ 为偶数时 $f^{(n)}(0)=0$。
\end{itemize}
\textbf{例（有理分式）}：求 $y=\frac{1}{x^2-5x+6}$ 的 $n$ 阶导数。 拆部分分式：$\frac{1}{(x-2)(x-3)}=\frac{1}{x-3}-\frac{1}{x-2}$，故

\[y^{(n)}=(-1)^nn!\left[\frac{1}{(x-3)^{n+1}}-\frac{1}{(x-2)^{n+1}}\right].\]

\textbf{例（莱布尼茨公式取项的技巧）}：设 $f(x)=x^2\ln(1+x)$，求 $f^{(10)}(0)$。 写 $f=uv$，$u=x^2$，$v=\ln(1+x)$。因为 $u$ 的三阶以上导数恒为 $0$，故莱布尼茨公式中只有 $k=8,9,10$ 三项非零：

\[f^{(10)}=C_{10}^{8}u''v^{(8)}+C_{10}^{9}u'v^{(9)}+C_{10}^{10}uv^{(10)} .\]

在 $x=0$ 处 $u(0)=0$，$u'(0)=0$，$u''(0)=2$；又 $v^{(8)}(0)=(-1)^{7}\frac{7!}{1^8}=-5040$，故

\[f^{(10)}(0)=45\times 2\times(-5040)=-453600 .\]

\textbf{例（泰勒展开法）}：设 $f(x)=\frac{1}{1+x}$，求 $f^{(n)}(0)$。 按 $(1+x)^{-1}$ 的展开式，$f(x)=1-x+x^2-\cdots+(-1)^nx^n+\cdots$，所以 $x^n$ 的系数为 $(-1)^n$，于是 $f^{(n)}(0)=n!\,(-1)^n$。这与公式 $\left(\frac{1}{ax+b}\right)^{(n)}$（取 $a=1,b=1$）一致。

\begin{callwarn}{易错}
\begin{itemize}
\item 莱布尼茨公式的系数是组合数 $C_n^k$，不是 $n^k$，也不是二项式系数外的其他东西。
\item $(\sin x)^{(n)}$ 中的相位是 $\frac{n\pi}{2}$，不要写成 $\frac{\pi}{2}$。
\item 用泰勒展开反求 $f^{(n)}(0)$ 时，要保证展开式是按 $x$ 的正整数次幂写出的，且 $f^{(n)}(0)=n!\times$ 该幂次项的系数。
\item 有理分式一定要先拆，否则求高阶导数无从下手。
\item 用奇偶性只能得到\textbf{特定点（0 处）}某些阶导数为 0，不能得到整个导函数为 0。
\end{itemize}
\end{callwarn}

\begin{callexam}{真题考法}
填空题高频考 $f^{(n)}(0)$ 或 $f^{(n)}(x_0)$，主套路是"奇偶性 + 泰勒展开 + 莱布尼茨公式"三选一；解答题偶有要求证明递推关系式（如 $x^2y''+\cdots$ 型恒等式），做法是逐次求导后代入。

\end{callexam}

\section{微分中值定理与证明题}

\subsection{中值定理基础}

\subsubsection{罗尔定理、拉格朗日中值定理及其应用 $\bigstar$$\bigstar$$\bigstar$$\bigstar$$\bigstar$}

\textbf{费马引理}：设 $f$ 在 $x_0$ 的某邻域内有定义，在 $x_0$ 处可导，且 $x_0$ 是 $f$ 的极值点，则 $f'(x_0)=0$。 几何意义：可导函数在极值点处的切线平行于 $x$ 轴。注意这是\textbf{必要}条件：$f'(x_0)=0$ 推不出 $x_0$ 是极值点（如 $x^3$ 在 $0$ 处）。端点处的极值不适用（必须在开区间内部）。

\textbf{罗尔定理}：设函数 $f$ 满足

\begin{enumerate}
\item 在闭区间 $[a,b]$ 上连续；
\item 在开区间 $(a,b)$ 内可导；
\item $f(a)=f(b)$，
\end{enumerate}
则至少存在一点 $\xi\in(a,b)$，使 $f'(\xi)=0$。

\textbf{拉格朗日中值定理}：设函数 $f$ 满足

\begin{enumerate}
\item 在闭区间 $[a,b]$ 上连续；
\item 在开区间 $(a,b)$ 内可导，
\end{enumerate}
则至少存在一点 $\xi\in(a,b)$，使

\[f(b)-f(a)=f'(\xi)(b-a)\]

或写成比值形式

\[f'(\xi)=\frac{f(b)-f(a)}{b-a}\qquad(a<\xi<b).\]

等价形式：$f(x_0+\Delta x)-f(x_0)=f'(x_0+\theta\Delta x)\Delta x\ (0<\theta<1)$，这个形式在证明"与 $\theta$ 有关"的题时常用。

几何意义：连续光滑的曲线上至少存在一点，其切线平行于端点连线。

\textbf{两个重要推论}：

\begin{itemize}
\item 若 $f$ 在区间 $I$ 上连续、在 $I$ 内可导且 $f'(x)\equiv0$，则 $f$ 在 $I$ 上为常数；
\item 若 $f,g$ 在 $I$ 上连续、在 $I$ 内可导且 $f'(x)\equiv g'(x)$，则 $f(x)=g(x)+C$。
\end{itemize}
\textbf{证明拉格朗日定理的方法（必须会写）}：构造

\[F(x)=f(x)-f(a)-\frac{f(b)-f(a)}{b-a}(x-a),\]

则 $F$ 在 $[a,b]$ 上连续、$(a,b)$ 内可导，且 $F(a)=F(b)=0$，用罗尔定理即得。

\textbf{应用一：证明恒等式}。要证 $f(x)=g(x)$，先证 $f'(x)-g'(x)\equiv0$，得 $f-g\equiv C$，再取一个特殊的 $x$ 定出 $C=0$。

\textbf{应用二：证明不等式}。把不等式整理成 $\frac{f(b)-f(a)}{b-a}$ 的形式，再用 $f'(\xi)$ 在 $[a,b]$ 上的范围放缩：

\[\vert f(b)-f(a)\vert=\vert f'(\xi)\vert\,\vert b-a\vert\le M\vert b-a\vert,\qquad M=\max_{[a,b]}\vert f'\vert .\]

\textbf{应用三：证明"存在 $\xi$ 使 $f'(\xi)=k$"}。构造 $F(x)=f(x)-kx$，对 $F$ 用罗尔定理。

\textbf{例（证明不等式）}：证明当 $b>a>0$ 时 $\frac{b-a}{b}<\ln\frac ba<\frac{b-a}{a}$。 取 $f(x)=\ln x$，在 $[a,b]$ 上用拉格朗日中值定理：

\[\ln\frac ba=\ln b-\ln a=\frac{b-a}{\xi},\qquad a<\xi<b .\]

由 $\frac{1}{b}<\frac{1}{\xi}<\frac{1}{a}$ 两端同乘 $b-a>0$，即得所证。

\textbf{例（证明根的存在）}：设 $f(x)=(x-1)(x-2)(x-3)$，证明 $f'(x)=0$ 在 $(1,3)$ 内至少有两个实根。 $f(1)=f(2)=f(3)=0$。在 $[1,2]$ 上用罗尔定理得 $\xi_1\in(1,2)$ 使 $f'(\xi_1)=0$；在 $[2,3]$ 上用罗尔定理得 $\xi_2\in(2,3)$ 使 $f'(\xi_2)=0$。由于 $\xi_1<2<\xi_2$，两个根互异。

\textbf{反复罗尔的推广（导数零点定理）}：若 $f$ 在 $[a,b]$ 上有 $n+1$ 个零点且在 $[a,b]$ 上有 $n$ 阶导数，则 $f^{(n)}(x)$ 在 $(a,b)$ 内至少有一个零点。这是处理"含高阶导数结论"的基本工具。

\textbf{应用四：证明恒等式}。要证 $f(x)=g(x)$，只需证 $(f-g)'\equiv0$ 再用一点定常数。

\textbf{例}：证明 $\arcsin x+\arccos x=\frac{\pi}{2}$（$x\in[-1,1]$）。 令 $F(x)=\arcsin x+\arccos x$，则

\[F'(x)=\frac{1}{\sqrt{1-x^2}}-\frac{1}{\sqrt{1-x^2}}=0\quad(-1<x<1),\]

故 $F$ 在 $(-1,1)$ 内为常数。取 $x=0$ 得 $F(0)=0+\frac\pi2=\frac\pi2$，由连续性该式在端点也成立。

\textbf{应用五：用拉格朗日中值定理求极限}（重要考法）。当极限式中出现\textbf{同一个函数在两点处的差}时，用 $f(a)-f(b)=f'(\xi)(a-b)$，再由 $\xi$ 介于两点之间确定 $f'(\xi)$ 的极限。

\textbf{例}：求 $\lim\limits_{x\to 0}\frac{e^{x}-e^{\sin x}}{x-\sin x}$。 取 $f(t)=e^t$，对两点 $x$ 与 $\sin x$（当 $x\to0$ 时二者都趋于 $0$）用拉格朗日中值定理：

\[e^{x}-e^{\sin x}=e^{\xi}(x-\sin x),\]

其中 $\xi$ 介于 $x$ 与 $\sin x$ 之间，故 $x\to0$ 时 $\xi\to0$。于是

\[\lim_{x\to 0}\frac{e^{x}-e^{\sin x}}{x-\sin x}=\lim_{\xi\to 0}e^{\xi}=1 .\]

\textbf{例}：证明对任意实数 $a,b$ 有 $|\arctan b-\arctan a|\le|b-a|$。 取 $f(x)=\arctan x$，若 $a=b$ 显然成立；若 $a\ne b$，由拉格朗日中值定理 $|\arctan b-\arctan a|=\frac{|b-a|}{1+\xi^2}\le|b-a|$，因为 $\frac{1}{1+\xi^2}\le1$。

\begin{callwarn}{易错}
\begin{itemize}
\item 罗尔定理的三个条件缺一不可：闭区间不连续、开区间内不可导、端点不等高，都会使结论失效。答题时必须逐条交代。
\item 题目要求"存在 $\xi\in(a,b)$"，不能写成闭区间 $[a,b]$。
\item 拉格朗日定理要求的是"闭区间连续 + 开区间可导"，不要写成"闭区间可导"。
\item 用拉格朗日证明不等式时，放缩方向取决于 $f'$ 在 $[a,b]$ 上的单调性或上下界，先判断再放缩。
\item 罗尔定理只保证 $\xi$ 存在，位置不确定；若题目给出具体函数，可解出 $\xi$ 来检验。
\end{itemize}
\end{callwarn}

\begin{callexam}{真题考法}
证明题的第一问常是"证明存在 $\xi$ 使 $f'(\xi)=0$"，直接对所给两个（或三个）等值点用罗尔定理；第二问升级为含 $f'(\xi)+g(\xi)f(\xi)=0$ 的形式，需要构造辅助函数。选择题则考"条件与结论的对应"，例如问四个选项中哪个条件能推出结论。

\end{callexam}

\subsubsection{柯西中值定理（了解并会用） $\bigstar$$\bigstar$$\bigstar$}

\textbf{定理}：设函数 $f(x),g(x)$ 满足

\begin{enumerate}
\item 在闭区间 $[a,b]$ 上连续；
\item 在开区间 $(a,b)$ 内可导；
\item 对任意 $x\in(a,b)$，$g'(x)\ne0$，
\end{enumerate}
则至少存在一点 $\xi\in(a,b)$，使

\[\frac{f(b)-f(a)}{g(b)-g(a)}=\frac{f'(\xi)}{g'(\xi)}.\]

条件 3 同时保证了 $g(b)\ne g(a)$（否则由罗尔定理会得出 $g'$ 在 $(a,b)$ 内有零点，与条件矛盾），因此分母不为零。

\textbf{证明（要会写）}：构造函数

\[F(x)=f(x)-f(a)-\frac{f(b)-f(a)}{g(b)-g(a)}\big[g(x)-g(a)\big],\]

则 $F(a)=0$，且 $F(b)=f(b)-f(a)-\frac{f(b)-f(a)}{g(b)-g(a)}(g(b)-g(a))=0$。对 $F$ 用罗尔定理，得 $F'(\xi)=0$，即

\[f'(\xi)-\frac{f(b)-f(a)}{g(b)-g(a)}g'(\xi)=0,\]

整理即得结论。

\textbf{与其他定理的关系}：取 $g(x)=x$ 即退化为拉格朗日中值定理；若再补上 $f(a)=f(b)$ 就退化为罗尔定理。所以柯西中值定理是一般形式。

\textbf{使用信号}：结论中出现\textbf{两个不同函数}的增量之比，或结论可整理成 $\frac{f'(\xi)}{g'(\xi)}=$ 常数（端点的量）的形式。若只有一个函数出现，用拉格朗日。

\textbf{例}：设 $0<a<b$，$f$ 在 $[a,b]$ 上连续、在 $(a,b)$ 内可导，证明存在 $\xi\in(a,b)$ 使

\[2\xi\big[f(b)-f(a)\big]=(b^2-a^2)f'(\xi).\]

分析：把结论整理成 $\frac{f(b)-f(a)}{b^2-a^2}=\frac{f'(\xi)}{2\xi}$，分母正是 $(x^2)'$ 在 $\xi$ 处的值，故取 $g(x)=x^2$（在 $(a,b)$ 上 $g'(x)=2x\ne0$，因为 $0<a<b$）。用柯西中值定理即得。

\textbf{例}：设 $f$ 在 $[a,b]$ 上连续、在 $(a,b)$ 内可导，$0<a<b$，证明存在 $\xi\in(a,b)$ 使

\[f(b)-f(a)=\xi f'(\xi)\ln\frac ba .\]

分析：把结论整理成 $\frac{f(b)-f(a)}{\ln b-\ln a}=\frac{f'(\xi)}{\frac1\xi}$，右端分母正是 $(\ln x)'$ 在 $\xi$ 处的值，故取 $g(x)=\ln x$，它在 $(a,b)$ 上满足 $g'(x)=\frac1x\ne0$（因为 $a>0$）。由柯西中值定理

\[\frac{f(b)-f(a)}{\ln b-\ln a}=\frac{f'(\xi)}{\frac1\xi}=\xi f'(\xi),\]

两边同乘 $\ln\frac ba$ 即得所证。

\textbf{识别柯西的三个信号}：$\circled{1}$ 结论中出现两个不同函数（或一个函数与 $x^2$、$e^x$、$\ln x$、$\frac1x$ 等）的差之比；$\circled{2}$ 分母恰好像是某个常见函数的导数；$\circled{3}$ 结论中出现了 $\ln\frac ba$、$e^b-e^a$、$b^2-a^2$ 这类"函数的增量"。

\begin{callwarn}{易错}
\begin{itemize}
\item 柯西定理分子分母的 $\xi$ 必须是\textbf{同一个}点，不能分子取 $\xi_1$、分母取 $\xi_2$。
\item 忘记验证 $g'(x)\ne0$；若 $g'$ 在区间内有零点（如 $g(x)=x^2$ 且区间含 0），条件不满足，需要换个 $g$ 或分段处理。
\item 构造的辅助函数必须逐条验证罗尔定理条件，尤其 $F(b)=0$ 要写清楚。
\item 柯西中值定理不能由"对 $f$ 和 $g$ 分别用拉格朗日再相除"得到，因为两个 $\xi$ 未必相同。
\end{itemize}
\end{callwarn}

\begin{callexam}{真题考法}
数二对柯西中值定理的要求是"了解并会用"，通常在证明题中作为工具：结论里出现两个函数的差之比，或分母出现 $e^x$、$\ln x$、$x^2$ 等特殊函数时，就是柯西的信号。考查形式多为"证明存在 $\xi$ 使某个含两个函数的等式成立"。

\end{callexam}

\subsubsection{泰勒定理与泰勒公式的应用（含中值定理证明中的泰勒法） $\bigstar$$\bigstar$$\bigstar$$\bigstar$$\bigstar$}

\textbf{带皮亚诺余项的泰勒公式}（局部展开，用于求极限）：设 $f$ 在 $x_0$ 处有 $n$ 阶导数，则

\[f(x)=\sum_{k=0}^{n}\frac{f^{(k)}(x_0)}{k!}(x-x_0)^k+o\big((x-x_0)^n\big)\qquad(x\to x_0).\]

\textbf{带拉格朗日余项的泰勒公式}（用于证明与估计）：设 $f$ 在含 $x_0$ 的区间 $I$ 上有 $n+1$ 阶导数，则对任意 $x\in I$，存在 $\xi$ 介于 $x_0$ 与 $x$ 之间，使

\[f(x)=\sum_{k=0}^{n}\frac{f^{(k)}(x_0)}{k!}(x-x_0)^k+\frac{f^{(n+1)}(\xi)}{(n+1)!}(x-x_0)^{n+1}.\]

取 $x_0=0$ 即得麦克劳林公式。

\textbf{常用麦克劳林展开（必须背到指定阶）}：

\[\begin{aligned}
e^x&=1+x+\frac{x^2}{2!}+\cdots+\frac{x^n}{n!}+o(x^n),\\
\sin x&=x-\frac{x^3}{3!}+\frac{x^5}{5!}-\cdots+(-1)^{n}\frac{x^{2n+1}}{(2n+1)!}+o(x^{2n+1}),\\
\cos x&=1-\frac{x^2}{2!}+\frac{x^4}{4!}-\cdots+(-1)^{n}\frac{x^{2n}}{(2n)!}+o(x^{2n}),\\
\ln(1+x)&=x-\frac{x^2}{2}+\frac{x^3}{3}-\cdots+(-1)^{n-1}\frac{x^n}{n}+o(x^n),\\
(1+x)^\alpha&=1+\alpha x+\frac{\alpha(\alpha-1)}{2!}x^2+\cdots+\frac{\alpha(\alpha-1)\cdots(\alpha-n+1)}{n!}x^n+o(x^n),\\
\frac{1}{1-x}&=1+x+x^2+\cdots+x^n+o(x^n),\qquad \frac{1}{1+x}=1-x+x^2-\cdots+(-1)^nx^n+o(x^n).
\end{aligned}\]

\textbf{用法一：求极限}。原则是"展开到分母同阶"：分母是 $x^k$ 就把分子展到 $x^k$ 项；若分子中两部分的低阶项恰好抵消，要继续展到第一个不抵消的阶。

\begin{callkey}{必背}
极限与证明中用到频率最高的几个展开，必须背到指定阶并记清余项类型：

\[e^x=1+x+\frac{x^2}{2}+\frac{x^3}{6}+o(x^3),\qquad \ln(1+x)=x-\frac{x^2}{2}+\frac{x^3}{3}+o(x^3)\]

\[\sin x=x-\frac{x^3}{6}+o(x^3),\qquad \cos x=1-\frac{x^2}{2}+\frac{x^4}{24}+o(x^4)\]

\[(1+x)^{\alpha}=1+\alpha x+\frac{\alpha(\alpha-1)}{2}x^2+o(x^2),\qquad \frac{1}{1-x}=1+x+x^2+o(x^2)\]

求极限用皮亚诺余项；证明不等式或存在 $\xi$ 的等式用拉格朗日余项。

\end{callkey}

\textbf{用法二：证明不等式}。在合适的点展开，利用余项的符号推出不等号方向。

\textbf{用法三：证明含高阶导数的中值问题}。\textbf{展开点的选择是关键}：

\begin{itemize}
\item 结论中同时出现 $f(a)$、$f(b)$ 与 $f\left(\frac{a+b}{2}\right)$，取 $x_0=\frac{a+b}{2}$ 展开；
\item 结论只含 $f(a)$、$f(b)$，分别在这两个端点展开，再相减（消去 $f'$ 项）或相加；
\item 已知某点的导数值，就选该点为展开点。
\end{itemize}
\textbf{例（求极限）}：求 $\lim\limits_{x\to 0}\frac{\cos x-e^{-x^2/2}}{x^4}$。 展到 $x^4$ 项：

\[\cos x=1-\frac{x^2}{2}+\frac{x^4}{24}+o(x^4),\qquad
e^{-x^2/2}=1-\frac{x^2}{2}+\frac{x^4}{8}+o(x^4).\]

相减得 $\left(\frac{1}{24}-\frac18\right)x^4+o(x^4)=-\frac{x^4}{12}+o(x^4)$，故原极限为 $-\frac{1}{12}$。

\textbf{例（证明不等式）}：证明 $x>0$ 时 $e^x>1+x+\frac{x^2}{2}$。 在 $x_0=0$ 处展到二阶带拉格朗日余项：

\[e^x=1+x+\frac{x^2}{2}+\frac{e^{\xi}}{6}x^3,\qquad \xi\in(0,x).\]

因为 $\frac{e^{\xi}}{6}x^3>0$，故 $e^x>1+x+\frac{x^2}{2}$。

\textbf{例（泰勒法证中值问题）}：设 $f$ 在 $[0,1]$ 上有二阶导数，$f(0)=f(1)=0$，$f\left(\frac12\right)=1$，证明存在 $\xi\in(0,1)$ 使 $f''(\xi)\le-8$。 在 $x_0=\frac12$ 处分别对 $x=0$ 和 $x=1$ 展开（带拉格朗日余项）：

\[0=f(0)=1-\frac12f'\!\left(\frac12\right)+\frac{1}{8}f''(\xi_1),\qquad
0=f(1)=1+\frac12f'\!\left(\frac12\right)+\frac{1}{8}f''(\xi_2).\]

于是 $f''(\xi_1)=4f'\!\left(\frac12\right)-8$，$f''(\xi_2)=-4f'\!\left(\frac12\right)-8$。若 $f'\!\left(\frac12\right)\le0$，取 $\xi=\xi_1$ 得 $f''(\xi_1)\le-8$；若 $f'\!\left(\frac12\right)>0$，取 $\xi=\xi_2$ 得 $f''(\xi_2)<-8$。两种情况都得证。

\textbf{例（由极限条件反求导数）}：设 $f$ 在 $x=0$ 处三阶可导，且 $\lim\limits_{x\to0}\frac{f(x)-x}{x^3}=1$，求 $f'(0),f''(0),f'''(0)$。 把 $f$ 在 0 处展到三阶：$f(x)=f(0)+f'(0)x+\frac{f''(0)}{2}x^2+\frac{f'''(0)}{6}x^3+o(x^3)$。代入极限并与 $x^3$ 比较系数：$f(0)=0$，$f'(0)=1$，$f''(0)=0$，$\frac{f'''(0)}{6}=1$，即 $f'''(0)=6$。

\textbf{例（泰勒法定阶）}：设 $f$ 在 $[0,1]$ 上二阶可导，$f(0)=f(1)=0$，$f'(0)=1$，证明存在 $\xi\in(0,1)$ 使 $f''(\xi)=-2$。 结论中只出现端点 $0,1$，且已知 $0$ 处的导数值，故在 $x_0=0$ 处对 $x=1$ 展开：

\[f(1)=f(0)+f'(0)(1-0)+\frac{f''(\xi)}{2}(1-0)^2,\]

即 $0=0+1+\frac{f''(\xi)}{2}$，故 $f''(\xi)=-2$，$\xi\in(0,1)$。

\textbf{泰勒展开的定阶与定点口诀}：

\begin{itemize}
\item \textbf{定阶}：结论里出现的最高阶导数就是需要展开到的阶数（要证 $f''(\xi)$ 的等式，就展到二阶并带拉格朗日余项）；
\item \textbf{定点}：结论里出现哪些点，就在这些点（或其中信息最多的点）展开；出现 $\frac{a+b}{2}$ 就在中点展开；只出现端点 $a,b$ 时，可以在同一点对 $x=a$、$x=b$ 各展一次，再相减或相加。
\end{itemize}
\begin{callwarn}{易错}
\begin{itemize}
\item 展开阶数不够是极限算错的头号原因；要展到"分母的阶数"，并检查低阶项是否抵消。
\item 皮亚诺余项只适用于 $x\to x_0$ 的局部极限；涉及区间上函数值比较（证明不等式）必须用拉格朗日余项。
\item 用泰勒证明中值时，结论中的 $f''(\xi)$ 来自某个展开点的余项，两个余项中要选符号符合需要的那一个。
\item 展开时 $(x-x_0)$ 的符号不能丢，例如在 $x_0=\frac12$ 处代 $x=0$ 时 $x-x_0=-\frac12$。
\item 由极限条件反求 $f^{(n)}(0)$ 时，务必把低阶项全部对消后再比较系数。
\end{itemize}
\end{callwarn}

\begin{callexam}{真题考法}
泰勒公式是数二证明题的核心工具，主要两种考法：一是"求极限"，用皮亚诺余项展开到分母同阶；二是"证明存在 $\xi$ 使某个含二阶导数的等式（或不等式）成立"，用拉格朗日余项并按结论中出现的点选展开点。也曾考"给出某分式极限的值，反求 $f^{(n)}(0)$"。

\end{callexam}

\subsection{证明题的套路}

\subsubsection{中值定理证明题的辅助函数构造 $\bigstar$$\bigstar$$\bigstar$$\bigstar$$\bigstar$}

总原则只有一句话：\textbf{把要证的结论"还原"成某个函数的导数等于 0}。因为一旦化成 $F'(\xi)=0$，就可以用罗尔定理或费马引理。

\textbf{第一类：还原法（原函数法）} 把结论中的 $\xi$ 换成 $x$，移项使一端为 0，得到形如 $G(x)=0$ 的式子；若 $G(x)$ 恰好是某个 $F(x)$ 的导数，则 $F$ 就是辅助函数。等价的操作是"对 $G(x)$ 积分"（或直接观察出原函数）。

\textbf{第二类：微分方程法} 把结论（同时含 $f'(\xi)$ 与 $f(\xi)$）看成关于 $y=y(x)$ 的微分方程，解出通解，再把通解整理成 $F(x,y)=C$ 的形式，则 $F(x,f(x))$ 就是辅助函数。例如结论 $f'(\xi)+g(\xi)f(\xi)=0$ 对应的一阶线性齐次方程 $y'+g(x)y=0$ 有通解 $ye^{\int g\,dx}=C$。

\textbf{第三类：万能构造（一阶线性型，最好用）} 若结论形如

\[f'(\xi)+g(\xi)f(\xi)=0,\]

则取

\[F(x)=f(x)e^{\int g(x)\,dx},\]

因为 $F'(x)=e^{\int g}\left[f'(x)+g(x)f(x)\right]$，括号内的符号与结论一致，对 $F$ 用罗尔定理即得。

常用原型（务必背熟）：

\begingroup\small\setlength{\tabcolsep}{4.0pt}
\begin{longtable}{@{}L{0.4728\linewidth}L{0.4728\linewidth}@{}}
\toprule
\textbf{结论形式} & \textbf{辅助函数 $F(x)$} \\
\midrule
\endhead
$f'(\xi)=0$ & $f(x)$ \\
$f'(\xi)+f(\xi)=0$ & $f(x)e^{x}$ \\
$f'(\xi)-f(\xi)=0$ & $f(x)e^{-x}$ \\
$\xi f'(\xi)+f(\xi)=0$ & $xf(x)$ \\
$\xi f'(\xi)-f(\xi)=0$ & $\frac{f(x)}{x}$ \\
$\xi f'(\xi)+nf(\xi)=0$ & $x^nf(x)$ \\
$f'(\xi)+2\xi f(\xi)=0$ & $f(x)e^{x^2}$ \\
$f'(\xi)\cos\xi+f(\xi)\sin\xi=0$ & $f(x)\cos x$ \\
$f'(\xi)\sin\xi-f(\xi)\cos\xi=0$ & $\frac{f(x)}{\sin x}$ \\
\bottomrule
\end{longtable}
\endgroup

\begin{callkey}{万能构造口诀}
看到"$f'(\xi)+g(\xi)f(\xi)=0$"就写

\[F(x)=f(x)e^{\int g(x)\,dx},\]

再回去验证 $F'(\xi)=e^{\int g}\left[f'(\xi)+g(\xi)f(\xi)\right]=0$。

口诀：\textbf{"导数相加配指数，导数相减配负指数；$\xi$ 当系数配幂函数（乘 $x^n$），除 $\xi$ 就配分母（除以 $x$）。"}

\end{callkey}

\textbf{第四类：常数 $K$ 值法} 适用于"结论中区间端点与中值可以分离"的情形。步骤：

\begin{enumerate}
\item 把结论整理成"一端只含端点、一端只含中值"的等式，令中值那端的值为 $K$；
\item 把 $K$ 中的 $\xi$ 换成 $x$，构造 $F(x)$，使 $F$ 在端点（有时还需再插一个已知点）取相同的值；
\item 对 $F$ 用罗尔定理，并把 $K$ 代回。
\end{enumerate}
\textbf{第五类：零点定理 / 介值定理法} 当结论只要求"存在 $\xi$ 使某个数值等式成立"而不含导数时，构造连续函数 $F(x)$，验证 $F(a)F(b)<0$（零点定理）或利用介值定理，即可直接得结论。

\textbf{第六类：双中值问题（$\xi\ne\eta$）} 三种处理方式：

\begin{enumerate}
\item \textbf{几何法}：把 $f'(\xi)$、$f'(\eta)$ 看作首尾相连两条弦的斜率，插入分点 $c$，使两条弦的区间不重叠；
\item \textbf{插分点法}：在 $[a,b]$ 内插入分点 $c$，把结论中的 $\xi$ 换成 $c$、$\eta$ 换成 $c$，解出 $c$，再在 $[a,c]$ 与 $[c,b]$ 上分别用拉格朗日（或柯西）中值定理；
\item \textbf{降为单中值}：先把其中一个中值固定（例如先产生 $\eta$），把结论化成一个单中值问题，用还原法/万能构造求出辅助函数产生另一个中值。
\end{enumerate}
\textbf{第七类：反复罗尔} 结论含 $f^{(n)}(\xi)$ 时，通常先构造 $F$，使 $F$ 在 $n+1$ 个点处取值相同，然后连续使用 $n$ 次罗尔定理。

\textbf{若构造出的辅助函数用不上条件}，有三条出路：$\circled{1}$ 对已知条件作处理（如化为极限、用推广罗尔定理）；$\circled{2}$ 二次构造（对 $F'$ 再构造原函数）；$\circled{3}$ 换一个辅助函数重来。

\textbf{例（二次构造 / 反复罗尔）}：设 $f$ 在 $[0,1]$ 上二阶可导，$f(0)=f(1)=0$，证明存在 $\xi\in(0,1)$ 使

\[2f'(\xi)+\xi f''(\xi)=0 .\]

思路：先观察结论，它恰是 $F(x)=x^2f'(x)$ 的导数除以 $x$ 后的形式，因为

\[F'(x)=2xf'(x)+x^2f''(x)=x\big[2f'(x)+xf''(x)\big].\]

所以只要找到一点使 $F'=0$ 且该点不为 $0$ 即可。由 $f(0)=f(1)=0$ 与罗尔定理，存在 $c\in(0,1)$ 使 $f'(c)=0$。于是

\[F(0)=0^2f'(0)=0,\qquad F(c)=c^2f'(c)=0 .\]

对 $F$ 在 $[0,c]$ 上用罗尔定理，得 $\xi\in(0,c)\subset(0,1)$ 使 $F'(\xi)=0$，即 $\xi[2f'(\xi)+\xi f''(\xi)]=0$。由于 $\xi>0$，故 $2f'(\xi)+\xi f''(\xi)=0$，证毕。

\textbf{要点}：结论含二阶导数时，往往先构造一个 $F$，使得要证的式子等于 $F'$；再想办法让 $F$ 在两个点取相同的值（常用罗尔定理或已知条件制造第二个等值点），最后对 $F$ 用一次罗尔。

\textbf{例（还原法）}：设 $f$ 在 $[0,1]$ 上连续、在 $(0,1)$ 内可导，$f(0)=0$，$f(1)=1$，证明存在 $\xi\in(0,1)$ 使 $f'(\xi)=2\xi$。 把结论写成 $f'(\xi)-2\xi=0$，积分得原函数 $F(x)=f(x)-x^2$。$F(0)=0$，$F(1)=f(1)-1=0$，对 $F$ 用罗尔定理得 $\xi$ 使 $F'(\xi)=f'(\xi)-2\xi=0$。

\textbf{例（万能构造）}：设 $f$ 在 $[0,1]$ 上连续、在 $(0,1)$ 内可导，$f(0)=f(1)=0$，证明存在 $\xi\in(0,1)$ 使 $f'(\xi)+2\xi f(\xi)=0$。 结论属于 $f'+g f=0$ 型，$g(x)=2x$，取 $F(x)=f(x)e^{x^2}$。$F(0)=0,F(1)=0$，罗尔定理给出 $\xi$ 使

\[F'(\xi)=e^{\xi^2}\left[f'(\xi)+2\xi f(\xi)\right]=0,\]

即得 $f'(\xi)+2\xi f(\xi)=0$。

\textbf{例（常数 $K$ 值法：用 K 值法再证一次拉格朗日中值定理）}：证明存在 $\xi\in(a,b)$ 使 $f'(\xi)=\frac{f(b)-f(a)}{b-a}$。 结论右端只含端点值、与 $\xi$ 无关，把它冻结为常数

\[K=\frac{f(b)-f(a)}{b-a}.\]

于是要证的是"存在 $\xi$ 使 $f'(\xi)-K=0$"，自然构造

\[F(x)=f(x)-Kx .\]

由 $K$ 的定义，$F(b)-F(a)=f(b)-f(a)-K(b-a)=0$，即 $F(a)=F(b)$。对 $F$ 用罗尔定理，得 $\xi\in(a,b)$ 使 $F'(\xi)=f'(\xi)-K=0$，证毕。

\textbf{使用要点}：$K$ 值法适用的前提是"结论中的区间端点部分与中值部分可以分离"。核心动作就是把只含端点的那个式子记作 $K$，然后再构造含 $K$ 的辅助函数；若结论中还出现第二个已知点（如中点），则让 $F$ 在该点与端点处取相同的值，最后反复使用罗尔定理。

\textbf{例（双中值，插分点法）}：设 $f$ 在 $[0,1]$ 上连续、在 $(0,1)$ 内可导，$f(0)=0$，$f(1)=1$，证明存在互不相同的 $\xi,\eta\in(0,1)$ 使

\[\frac{1}{f'(\xi)}+\frac{1}{f'(\eta)}=2 .\]

思路：先把分点 $c$ 定出来。由 $f$ 在 $[0,1]$ 上连续、$f(0)=0<f(1)=1$，介值定理给出 $c\in(0,1)$ 使 $f(c)=\frac12$。在 $[0,c]$、$[c,1]$ 上分别用拉格朗日中值定理：

\[f(c)-f(0)=f'(\xi)\,c,\qquad f(1)-f(c)=f'(\eta)\,(1-c).\]

代入 $f(c)=\frac12$ 得 $f'(\xi)=\frac{1}{2c}$，$f'(\eta)=\frac{1}{2(1-c)}$，两式均不为零（分子为 $\frac12>0$）。于是

\[\frac{1}{f'(\xi)}+\frac{1}{f'(\eta)}=2c+2(1-c)=2 .\]

又 $\xi\in(0,c)$、$\eta\in(c,1)$，两个子区间不相交，故 $\xi\ne\eta$。证毕。

\begin{callwarn}{易错}
\begin{itemize}
\item 构造出的 $F$ 必须在\textbf{闭区间连续、开区间可导}且\textbf{端点值相等}，答题时要逐句交代这三个条件。
\item 结论要求 $\xi\in(a,b)$，不能写 $[a,b]$；双中值题必须保证 $\xi\ne\eta$，靠"分点把区间分成两个不相交的子区间"实现。
\item 万能构造中 $e^{\int g\,dx}$ 的积分常数可任取（常取 0），但要保证 $F$ 有意义。
\item 常数 $K$ 值法里 $K$ 是待定常数，最后必须代回并说明它与中值的关系。
\item 辅助函数构造出来后要检查能不能用上全部条件；用不上时不要硬算，改走二次构造。
\end{itemize}
\end{callwarn}

\begin{callexam}{真题考法}
证明题（10 分上下）的固定套路：第一问用零点定理或罗尔定理证明"存在某点使某等式成立"，第二问用万能构造或泰勒证明"存在 $\xi$ 使含导数（或二阶导数）的等式成立"。近年偏好三种形式：含 $f'(\xi)+g(\xi)f(\xi)=0$ 的万能构造、含两个不同中值 $\xi\ne\eta$ 的双中值问题、含二阶导数需二次构造的问题。

\end{callexam}

\subsubsection{不等式证明与放缩技巧 $\bigstar$$\bigstar$$\bigstar$$\bigstar$}

\textbf{四条主路（按题目给的信息选工具）}：

\begin{enumerate}
\item \textbf{单调性法}：要证 $f(x)>g(x)$（或 $f(x)>0$），令 $\varphi(x)=f(x)-g(x)$，证明 $\varphi'(x)>0$ 且端点值 $\varphi(a)\ge0$。给的是\textbf{一阶导}信息时首选。
\item \textbf{最值法}：证明 $f(x)\ge m$，转化为求 $f$ 的最小值；适合区间内部取到极值的情形。
\item \textbf{中值定理法}：结论中出现\textbf{同一函数的差} $f(b)-f(a)\Rightarrow$ 用拉格朗日；出现\textbf{两个不同函数}的差之比 $\Rightarrow$ 用柯西；条件给的是\textbf{二阶或更高阶导数}信息 $\Rightarrow$ 用泰勒公式。
\item \textbf{凹凸性法}：$f''>0$（凹）时曲线在切线上方、在割线下方，
\end{enumerate}
\[f(x)\ge f(x_0)+f'(x_0)(x-x_0),\qquad
f(\lambda x_1+(1-\lambda)x_2)\le\lambda f(x_1)+(1-\lambda)f(x_2);\]

$f''<0$（凸）时不等号反向。

\textbf{选工具的对照表（考场上直接对号入座）}：

\begingroup\small\setlength{\tabcolsep}{4.0pt}
\begin{longtable}{@{}L{0.5058\linewidth}L{0.4399\linewidth}@{}}
\toprule
\textbf{题目给出的信息} & \textbf{首选方法} \\
\midrule
\endhead
只给 $f$ 的单调性或一阶导符号 & 单调性法（作差构造函数） \\
给 $f''$ 的符号 & 凹凸性法或泰勒展开 \\
给 $\vert f'\vert$ 的上界 & 拉格朗日中值定理放缩 \\
给 $f(a)$、$f(b)$ 的具体值且结论含 $f'(\xi)$ & 拉格朗日（或柯西）中值定理 \\
结论含 $f''(\xi)$ 或更高阶 & 泰勒公式（拉格朗日余项） \\
两端为常数、形式为"$A\le B$"且 $A,B$ 都是函数值 & 最值法 \\
\bottomrule
\end{longtable}
\endgroup

\textbf{常用放缩不等式}：

\begin{itemize}
\item $e^x\ge 1+x$（$x\in\mathbb R$），$e^x>1+x+\frac{x^2}{2}$（$x>0$）；
\item $\frac{x}{1+x}<\ln(1+x)<x$（$x>0$），$\ln x\le x-1$（$x>0$）；
\item $\sin x<x<\tan x$（$0<x<\frac\pi2$）；
\item 均值不等式 $\frac{2}{\frac1a+\frac1b}\le\sqrt{ab}\le\frac{a+b}{2}\le\sqrt{\frac{a^2+b^2}{2}}$（$a,b>0$）；
\item 分区间、分类讨论（例如 $0<x<1$ 与 $x\ge1$ 分别处理）。
\end{itemize}
\textbf{例}：证明 $x>0$ 时 $\frac{x}{1+x}<\ln(1+x)<x$。 右侧：令 $f(x)=x-\ln(1+x)$，$f'(x)=1-\frac{1}{1+x}=\frac{x}{1+x}>0$，$f(0)=0$，故 $f(x)>0$。 左侧：令 $g(x)=\ln(1+x)-\frac{x}{1+x}$，$g'(x)=\frac{1}{1+x}-\frac{1}{(1+x)^2}=\frac{x}{(1+x)^2}>0$，$g(0)=0$，故 $g(x)>0$。合起来即得。

\textbf{例（凹凸性）}：证明 $x>0$ 时 $e^x\ge ex$，当且仅当 $x=1$ 时取等号。 令 $f(x)=e^x-ex$，$f'(x)=e^x-e$，$f''(x)=e^x>0$。$f'(1)=0$，且 $f'$ 由负变正，故 $x=1$ 是极小值点；$f(1)=0$ 是最小值，于是 $f(x)\ge 0$。

\textbf{例（用凹凸性/Jensen 型不等式）}：证明对任意实数 $a,b$ 有 $\frac{e^a+e^b}{2}\ge e^{\frac{a+b}{2}}$。 $f(x)=e^x$ 满足 $f''(x)=e^x>0$，故 $f$ 是凹函数，由凹凸性的定义（取 $\lambda=\frac12$）直接得

\[f\!\left(\frac{a+b}{2}\right)\le\frac{f(a)+f(b)}{2},\]

即 $e^{\frac{a+b}{2}}\le\frac{e^a+e^b}{2}$。

\textbf{例（用最值法）}：证明 $x>0$ 时 $\ln x\le x-1$。 令 $f(x)=x-1-\ln x\ (x>0)$，$f'(x)=1-\frac1x$。$0<x<1$ 时 $f'<0$，$x>1$ 时 $f'>0$，故 $x=1$ 是极小值点，$f(1)=0$ 也是最小值，于是 $f(x)\ge0$，即 $\ln x\le x-1$，等号仅在 $x=1$ 成立。

\textbf{例（用拉格朗日放缩）}：设 $f$ 在 $[a,b]$ 上可导且 $|f'(x)|\le M$，证明 $|f(b)-f(a)|\le M(b-a)$。 由拉格朗日中值定理，$|f(b)-f(a)|=|f'(\xi)|(b-a)\le M(b-a)$。这类题的关键是\textbf{先写出中值等式，再用条件给的上界放缩}。

\begin{callwarn}{易错}
\begin{itemize}
\item 用单调性证不等式必须\textbf{先写清区间}，再在该区间上论证 $\varphi'$ 的符号，最后代端点值。
\item 用中值定理时 $\xi$ 落在区间内部，放缩方向由 $f'$ 的单调性或上下界决定，不能凭感觉。
\item 用泰勒证明不等式必须用\textbf{拉格朗日余项}并判断余项的符号，用皮亚诺余项无法确定符号。
\item 双向不等式要拆成两个不等式分别证明。
\item 忘记讨论等号成立的条件（取等点）。
\end{itemize}
\end{callwarn}

\begin{callexam}{真题考法}
不等式证明常作为证明题的最后一问，或在选择/填空里以"比较大小"的形式出现。选工具的口诀：给一阶导信息用单调性或拉格朗日；给二阶导信息用凹凸性或泰勒；给定常数上下界（如 $|f'|\le M$）用拉格朗日放缩。陷阱是"用一阶导信息硬套泰勒"或"忘记讨论端点值"。

\end{callexam}

\subsubsection{方程根的个数与零点存在性讨论 $\bigstar$$\bigstar$$\bigstar$$\bigstar$}

\textbf{存在性}：用零点定理。若 $f$ 在 $[a,b]$ 上连续且 $f(a)f(b)<0$，则 $(a,b)$ 内至少有一个根。推广形式：若 $\lim_{x\to a^+}f(x)$ 与 $\lim_{x\to b^-}f(x)$ 异号（允许为无穷），则在 $(a,b)$ 内至少有一个根。另一个基本事实：\textbf{实系数奇次方程至少有一个实根}（因为 $x\to\pm\infty$ 时两端异号）。

\textbf{唯一性}：单调性 + 存在性。若 $f'(x)>0$（严格单调递增），则 $f$ 至多一个零点，结合零点定理即得"有且仅有一个根"。

\textbf{根的个数}：用反复罗尔的推广——"导数零点定理"：

\begin{quote}
若 $f(x)=0$ 至多有 $k$ 个根，则 $f'(x)=0$ 至多有 $k-1$ 个根；反之，若 $f^{(n)}(x)=0$ 至多有 $k$ 个根，则 $f(x)=0$ 至多有 $k+n$ 个根。

\end{quote}
常用特例：$f^{(n)}(x)\ne0$ 恒成立 $\Rightarrow f(x)=0$ 至多有 $n$ 个实根。

\textbf{标准解题流程}：

\begin{enumerate}
\item 判断有没有根（零点定理，或利用奇次方程结论）；
\item 求 $f'$，定出 $f$ 的单调区间，得到 $f$ 的走势（必要时求极值）；
\item 在每个单调区间上判断是否有根：结合端点值、极值的符号；
\item 若含参数，按参数与极值的大小关系分类讨论。
\end{enumerate}
\textbf{例}：证明方程 $x^5+x+1=0$ 有且仅有一个实根。 存在性：$f(0)=1>0$，$f(-1)=-1<0$，由零点定理在 $(-1,0)$ 内至少有一个根。 唯一性：$f'(x)=5x^4+1>0$ 恒成立，$f$ 严格单调递增，故至多一个根。综上有且仅有一个实根。

\textbf{例}：讨论方程 $2^x-x-1=0$ 的实根个数。 令 $f(x)=2^x-x-1$，则 $f(0)=0$，$f(1)=0$，已找到两个根。又

\[f'(x)=2^x\ln 2-1,\qquad f''(x)=2^x(\ln 2)^2>0 .\]

由 $f''>0$ 知 $f'$ 严格单调递增，故 $f'$ 至多有一个零点，从而 $f$ 至多有两个零点。因此该方程恰有两个实根 $x=0$ 与 $x=1$。（顺带可看出 $f$ 在 $x=\log_2\frac{1}{\ln 2}\approx0.53$ 处取极小值，该极小值小于 $0$，与"两个根"一致。）

\textbf{例（含参数讨论根的个数）}：讨论方程 $x^3-3x+a=0$ 的实根个数。 令 $f(x)=x^3-3x$，问题等价于求曲线 $y=f(x)$ 与水平线 $y=-a$ 的交点个数。由 $f'(x)=3x^2-3$ 得驻点 $x=\pm1$，且 $f(-1)=2$ 为极大值，$f(1)=-2$ 为极小值。结合 $f$ 的单调性（在 $(-\infty,-1)$、$(1,+\infty)$ 上递增，在 $(-1,1)$ 上递减）与 $x\to\pm\infty$ 时 $f\to\pm\infty$，得

\begingroup\small\setlength{\tabcolsep}{4.0pt}
\begin{longtable}{@{}L{0.5109\linewidth}L{0.4348\linewidth}@{}}
\toprule
\textbf{参数范围} & \textbf{交点个数（实根个数）} \\
\midrule
\endhead
$-2<-a<2$，即 $-2<a<2$ & $3$ \\
$-a=2$ 或 $-a=-2$，即 $a=\pm2$ & $2$（含一个二重根） \\
$-a>2$ 或 $-a<-2$，即 $a<-2$ 或 $a>2$ & $1$ \\
\bottomrule
\end{longtable}
\endgroup

\textbf{要点}：含参的根个数问题，标准做法就是"把方程整理成 $\varphi(x)=a$ 的形式 $\to$ 研究 $\varphi$ 的单调性与极值 $\to$ 用水平线与 $\varphi$ 的图像相交情况分类讨论"。分类的临界点正是 $\varphi$ 的极值（以及 $\varphi$ 的间断点）。

\textbf{含参讨论}：把方程整理成 $\varphi(x)=a$ 的形式，研究 $\varphi$ 的单调性与极值，按水平线 $y=a$ 与 $\varphi$ 的极值的相对位置分类，得到根的个数。

\begin{callwarn}{易错}
\begin{itemize}
\item 只证"有根"不够，个数问题必须再用单调性或"导数零点定理"说明"至多几个根"。
\item $f'(x)\ge0$ 只能得单调不减；要严格单调需要 $f'(x)>0$（或 $f'\ge0$ 且等号仅孤立成立）。
\item 零点定理的区间端点必须落在定义域内且函数在该闭区间连续。
\item 含参数讨论时容易漏掉相切（重根）的临界情形，导致分类不完整。
\item 用"至多几个根"的推广时要数清求导的阶数与已知零点的个数。
\end{itemize}
\end{callwarn}

\begin{callexam}{真题考法}
解答题常见的两种形式：$\circled{1}$ 证明方程"有且仅有 $n$ 个实根"；$\circled{2}$ 讨论根的个数随参数的变化。固定骨架是"求导 $\to$ 定单调区间 $\to$ 看极值（最值）符号 $\to$ 分类讨论"。选择题则常问"下列方程有几个实根"，可用奇次方程有根 + 单调性快速判断。

\end{callexam}

\section{导数的应用}

\subsection{洛必达法则}

\subsubsection{洛必达法则及其适用条件与失效情形 $\bigstar$$\bigstar$$\bigstar$$\bigstar$}

\textbf{基本形式}（$\frac00$ 型）：设函数 $f,g$ 满足

\begin{enumerate}
\item $\lim\limits_{x\to x_0}f(x)=\lim\limits_{x\to x_0}g(x)=0$（或两者都趋于 $\infty$）；
\item 在 $x_0$ 的某去心邻域内 $f,g$ 可导，且 $g'(x)\ne0$；
\item $\lim\limits_{x\to x_0}\frac{f'(x)}{g'(x)}$ 存在（或为 $\infty$），
\end{enumerate}
则

\[\lim_{x\to x_0}\frac{f(x)}{g(x)}=\lim_{x\to x_0}\frac{f'(x)}{g'(x)}.\]

关键：结论成立的前提是\textbf{右侧极限存在或为无穷}。若右侧振荡不存在，不能反过来断言原极限不存在。

\textbf{七种未定式与转化方法}：

\begingroup\small\setlength{\tabcolsep}{4.0pt}
\begin{longtable}{@{}L{0.3565\linewidth}L{0.5891\linewidth}@{}}
\toprule
\textbf{类型} & \textbf{转化方法} \\
\midrule
\endhead
$\frac00$ 型 & 直接洛必达 \\
$\frac{\infty}{\infty}$ 型 & 直接洛必达 \\
$0\cdot\infty$ 型 & 化为 $\frac{0}{1/\infty}$ 或 $\frac{\infty}{1/0}$ 后用洛必达 \\
$\infty-\infty$ 型 & 通分、有理化、提公因式化为商的形式 \\
$0^0,\allowbreak\ 1^\infty,\allowbreak\ \infty^0$ 型 & 取对数：$\lim f^{g}=e^{\lim g\ln f}$，化为 $0\cdot\infty$ \\
\bottomrule
\end{longtable}
\endgroup

\textbf{使用流程}：$\circled{1}$ 判断类型（不是未定式直接算，不能洛必达）；$\circled{2}$ 化简（等价无穷小替换、提取非零因子、变量代换）；$\circled{3}$ 分子分母分别求导；$\circled{4}$ 求极限；$\circled{5}$ 若仍是未定式且条件满足，可以再用一次。

\textbf{优先级}：先用等价无穷小与泰勒化简，再用洛必达；洛必达与泰勒可以交替使用。\textbf{数列极限不能直接用洛必达}（自变量是离散的），要么用海涅定理把它化成函数极限，要么用 Stolz 定理。

\textbf{四种失效情形}：

\begin{enumerate}
\item \textbf{循环求导}：如 $\lim\limits_{x\to+\infty}\frac{e^x-e^{-x}}{e^x+e^{-x}}$，洛必达后回到自身，应分子分母同除 $e^x$，得极限 $1$；
\item \textbf{振荡不存在}：如 $\lim\limits_{x\to0}\frac{x^2\sin\frac1x}{\sin x}$，洛必达后出现 $\cos\frac1x$ 振荡，应改用"无穷小乘有界量"：原式 $=\frac{x}{\sin x}\cdot x\sin\frac1x\to 0$；
\item \textbf{条件不满足}：如 $\lim\limits_{x\to0}\frac{|x|}{x}$，分子在 $0$ 处不可导，不能洛必达，应讨论左右极限；
\item \textbf{不是未定式}：如 $\frac{0}{\infty}$、$\frac{\infty}{0}$、$\frac{1}{0}$ 等，直接判断结果，洛必达在此处是错的。
\end{enumerate}
\textbf{例（$1^\infty$ 型）}：求 $\lim\limits_{x\to0^+}x^x$。 $x^x=e^{x\ln x}$，而

\[\lim_{x\to0^+}x\ln x=\lim_{x\to0^+}\frac{\ln x}{\frac1x}
=\lim_{x\to0^+}\frac{\frac1x}{-\frac{1}{x^2}}
=\lim_{x\to0^+}(-x)=0,\]

故 $\lim\limits_{x\to0^+}x^x=e^0=1$。

\textbf{例（$\infty-\infty$ 型）}：求 $\lim\limits_{x\to1}\left(\frac{1}{\ln x}-\frac{1}{x-1}\right)$。 通分得 $\frac{x-1-\ln x}{(x-1)\ln x}$，为 $\frac00$ 型。令 $t=x-1\to0$，用泰勒：分子 $=t-\ln(1+t)=\frac{t^2}{2}+o(t^2)$，分母 $=t\ln(1+t)=t^2+o(t^2)$，故极限为 $\frac12$。

\textbf{例（$1^\infty$ 型与泰勒配合）}：求 $\lim\limits_{x\to0}\left(\frac{\sin x}{x}\right)^{\frac{1}{x^2}}$。 记 $y=\left(\frac{\sin x}{x}\right)^{1/x^2}$，则

\[\ln y=\frac{1}{x^2}\ln\frac{\sin x}{x}=\frac{1}{x^2}\ln\!\left(1+\frac{\sin x-x}{x}\right).\]

由 $\sin x=x-\frac{x^3}{6}+o(x^3)$ 得 $\frac{\sin x-x}{x}=-\frac{x^2}{6}+o(x^2)$，再用 $\ln(1+u)\sim u$：

\[\ln y\sim\frac{1}{x^2}\cdot\left(-\frac{x^2}{6}\right)=-\frac16\ \Longrightarrow\ \lim_{x\to0}y=e^{-1/6}.\]

\textbf{要点}：$1^\infty$ 型的标准三步是"取对数 $\to$ 化 $\frac00$ 或 $0\cdot\infty$ $\to$ 求指数上的极限"，其中指数上的极限多半要用等价无穷小或泰勒，而不是硬洛必达。也可以直接用结论 $\lim f^{g}=e^{\lim g(f-1)}$（当 $f\to1,g\to\infty$ 时）。

\textbf{洛必达法则的答题模板（保证不被扣条件分）}：

\begin{enumerate}
\item 先写"当 $x\to x_0$ 时，分子趋于 $0$（或 $\infty$），分母趋于 $0$（或 $\infty$），这是 $\frac00$（或 $\frac{\infty}{\infty}$）型未定式"；
\item 再写"$f,g$ 在 $x_0$ 的去心邻域内可导，且 $g'(x)\ne0$"；
\item 然后写"由洛必达法则"，分子分母分别求导；
\item 每次使用都重复第 1、2 步；只有当最后一步得到的极限存在或为 $\infty$ 时，才能回代得结论。
\end{enumerate}
\textbf{洛必达法则的适用性与失效情形的统一理解}：洛必达法则的本质是"用导数之比代替函数之比"，它成立需要两个前提——\textbf{分子分母同时趋于 $0$ 或同时趋于 $\infty$}（这样才能用柯西中值定理），以及\textbf{导数之比的极限存在}。因此凡是"导数之比振荡或循环"、"分子分母不可导"、"根本不是未定式"的情形都不能用，而要回到等价无穷小、泰勒展开或"有界量乘无穷小"这些更基本的工具。

\begin{callwarn}{易错}
\begin{itemize}
\item 洛必达前必须验类型与条件（可导、分母导数非零）；不是未定式时使用一定错。
\item 每次洛必达后都要重新判断是否仍是未定式，不能连续无脑求导。
\item 右侧极限不存在时不能下"原极限不存在"的结论。
\item 数列极限、含绝对值的点、分段点处慎用洛必达，优先用导数定义或等价无穷小。
\item $0\cdot\infty$ 型必须先化形再洛必达，直接对乘积求导不成立。
\end{itemize}
\end{callwarn}

\begin{callexam}{真题考法}
极限题的主力工具。常见形式是给一个 $\frac00$ 或 $\frac{\infty}{\infty}$ 型的复杂分式，需要先用等价无穷小与泰勒化简再洛必达；也有概念型选择题，问"下列使用洛必达法则的过程是否正确"，考查条件检验与失效情形。

\end{callexam}

\subsection{单调性、极值与最值}

\subsubsection{函数单调性的判别 $\bigstar$$\bigstar$$\bigstar$$\bigstar$}

\textbf{判定定理}：设 $f$ 在 $[a,b]$ 上连续，在 $(a,b)$ 内可导。

\begin{itemize}
\item 若 $(a,b)$ 内 $f'(x)>0$，则 $f$ 在 $[a,b]$ 上严格单调递增；
\item 若 $(a,b)$ 内 $f'(x)<0$，则 $f$ 在 $[a,b]$ 上严格单调递减；
\item 若 $f'(x)\ge0$ 且等号仅在有限个点处成立，同样可判定严格单调递增（递减同理）。
\end{itemize}
\textbf{用单调性证明不等式的固定三步}：

\begin{enumerate}
\item 把要证的不等式移项，令 $\varphi(x)=$ 左端 $-$ 右端（使结论变成 $\varphi(x)>0$ 或 $\ge0$）；
\item 求 $\varphi'(x)$，判断它在所给区间上的符号（常常需要再求一次导，即"二次求导定号"）；
\item 由 $\varphi$ 单调及端点值 $\varphi(a)\ge0$ 得出结论，并说明等号成立的条件。
\end{enumerate}
\textbf{含参数单调性讨论的流程}：$\circled{1}$ 求定义域；$\circled{2}$ 求 $f'$ 并\textbf{通分、因式分解}成便于定号的形式；$\circled{3}$ 找出 $f'$ 的零点（含参数，需按参数分类）；$\circled{4}$ 用零点分割定义域，逐段定号；$\circled{5}$ 分类写出结论（分类要完备，做到"不重不漏"）。

\textbf{求单调区间的流程}：

\begin{enumerate}
\item 写出定义域（含分母、对数、根号、反正弦等限制）；
\item 求 $f'(x)$，解 $f'(x)=0$ 得驻点，并找出 $f'$ \textbf{不存在}的点；
\item 用这些点把定义域分成若干区间，在每个区间内取一个测试值判断 $f'$ 的符号；
\item 写出单调递增、单调递减区间，\textbf{不能}用并集符号把两段合并。
\end{enumerate}
\textbf{例}：求 $f(x)=x^3-3x$ 的单调区间与极值。 $f'(x)=3x^2-3=3(x-1)(x+1)$，驻点 $x=\pm1$。 $x<-1$ 时 $f'>0$；$-1<x<1$ 时 $f'<0$；$x>1$ 时 $f'>0$。 故单调递增区间为 $(-\infty,-1)$ 与 $(1,+\infty)$，单调递减区间为 $(-1,1)$；极大值 $f(-1)=2$，极小值 $f(1)=-2$。

\textbf{例（含参数）}：讨论 $f(x)=x^3-ax\ (a>0)$ 的单调性。 $f'(x)=3x^2-a$，零点 $x=\pm\sqrt{\frac a3}$。故 $f$ 在 $\left(-\infty,-\sqrt{\frac a3}\right)$ 与 $\left(\sqrt{\frac a3},+\infty\right)$ 上单调递增，在 $\left(-\sqrt{\frac a3},\sqrt{\frac a3}\right)$ 上单调递减。

\begin{callwarn}{易错}
\begin{itemize}
\item 单调区间不能用并集符号连接（如写成 $(-\infty,-1)\cup(1,+\infty)$ 表示"单调递增区间"是错的）。
\item 导数为零的点未必是单调性改变点（如 $x^3$ 在 $0$ 处）。
\item 求单调区间前必须先写定义域，否则会漏掉间断点造成的分段。
\item 含参数时要对参数分类讨论（如 $a=0$、$a>0$、$a<0$）。
\end{itemize}
\end{callwarn}

\begin{callexam}{真题考法}
基础题型，常与极值、最值、不等式证明合并考查；解答题的第一步通常是"求 $f$ 的单调区间"。陷阱是定义域、不可导点与并集符号。

\end{callexam}

\subsubsection{函数的极值（含不可导点、驻点判定） $\bigstar$$\bigstar$$\bigstar$$\bigstar$}

\textbf{定义}：若存在 $x_0$ 的某去心邻域，使其中任意 $x$ 都有 $f(x)<f(x_0)$（或 $f(x)>f(x_0)$），则称 $f(x_0)$ 为 $f$ 的极大值（极小值），$x_0$ 为极大值点（极小值点），统称极值与极值点。

\textbf{必要条件}：若 $f$ 在 $x_0$ 处可导且 $x_0$ 是极值点，则 $f'(x_0)=0$。 由此有两条重要结论：

\begin{itemize}
\item 可导的极值点一定是\textbf{驻点}；但驻点不一定是极值点（如 $x^3$ 在 $0$ 处）；
\item $f$ 的极值点只能是\textbf{驻点或不可导点}。
\end{itemize}
\textbf{第一充分条件}（看 $f'$ 是否变号）：设 $f$ 在 $x_0$ 处连续，在 $x_0$ 的某去心邻域内可导。

\begin{itemize}
\item $f'$ 由正变负 $\Rightarrow x_0$ 为极大值点；
\item $f'$ 由负变正 $\Rightarrow x_0$ 为极小值点；
\item $f'$ 不变号 $\Rightarrow x_0$ 不是极值点。
\end{itemize}
\textbf{第二充分条件}（看二阶导）：设 $f'(x_0)=0$，$f''(x_0)\ne0$。

\begin{itemize}
\item $f''(x_0)<0\Rightarrow f(x_0)$ 为极大值；
\item $f''(x_0)>0\Rightarrow f(x_0)$ 为极小值。
\end{itemize}
\textbf{第三充分条件}（高阶导）：设 $f'(x_0)=f''(x_0)=\cdots=f^{(n-1)}(x_0)=0$，$f^{(n)}(x_0)\ne0$。

\begin{itemize}
\item $n$ 为奇数：$x_0$ 不是极值点；
\item $n$ 为偶数：$x_0$ 是极值点；$f^{(n)}(x_0)<0$ 为极大值，$f^{(n)}(x_0)>0$ 为极小值。
\end{itemize}
\textbf{求极值的流程}：$\circled{1}$ 求定义域与 $f'$；$\circled{2}$ 求驻点与 $f'$ 不存在的点；$\circled{3}$ 用第一充分条件列表判断（或对驻点用第二充分条件）；$\circled{4}$ 写出极值（写函数值）与极值点（写横坐标）。

\textbf{例}：求 $f(x)=x^3-3x^2+1$ 的极值。 $f'(x)=3x^2-6x=3x(x-2)$，驻点 $x=0,2$。$f''(x)=6x-6$，$f''(0)=-6<0$，故 $x=0$ 取极大值 $f(0)=1$；$f''(2)=6>0$，故 $x=2$ 取极小值 $f(2)=-3$。

\textbf{例（不可导点）}：求 $f(x)=|x^2-4x+3|$ 的极值点。 $f(x)=|(x-1)(x-3)|$。$x=1,3$ 是不可导点，$f$ 在这两点取得极小值 $0$。在 $x=2$ 处 $f(2)=1$，且在 $x=2$ 附近 $f(x)\le1$，故 $x=2$ 是极大值点，极大值 $1$。

\textbf{四组必须分清的概念（选择题重灾区）}：

\begin{enumerate}
\item \textbf{极值 vs 最值}：最值不一定是极值（端点处的最值不是极值），极值也不一定是最值；但\textbf{在区间内部取得的最值一定是极值}；若连续函数在某区间上只有一个极值点，则该极值就是最值。
\item \textbf{极值点 vs 驻点}：极值点不一定是驻点（不可导点也可能是极值点，如 $|x|$ 的 $x=0$），驻点不一定是极值点（如 $x^3$ 的 $x=0$）；但\textbf{可导的极值点一定是驻点}，且\textbf{函数的极值点只能是驻点或不可导点}。
\item \textbf{极值点 vs 拐点}：极值点不一定是拐点，拐点也不一定是极值点；可以证明\textbf{可导的极值点一定不是拐点}（因为若在 $x_0$ 处 $f'$ 由正变负且凹凸性也改变，$f$ 在 $x_0$ 两侧都会单调上升，与极值矛盾）。
\item \textbf{零点 vs 驻点 vs 极值点}：三者都是"横坐标 $x=a$"的形式，而\textbf{拐点是点，要写成 $(x_0,f(x_0))$}。
\end{enumerate}
\textbf{例（含参极值）}：已知 $f(x)=x^3+ax^2+bx$ 在 $x=1$ 处取极小值 $-1$，求 $a,b$。 由 $f(1)=1+a+b=-1$ 得 $a+b=-2$；由 $f'(1)=3+2a+b=0$，把 $a+b=-2$ 代入得 $3+a-2=0$，故 $a=-1$，$b=-1$。检验：$f'(x)=3x^2-2x-1$，$f''(1)=4>0$，确为极小值点，符合题意。

\begin{callwarn}{易错}
\begin{itemize}
\item 极值点写横坐标 $x=x_0$，极值写函数值 $f(x_0)$，两者不能混。
\item $f''(x_0)=0$ 时第二充分条件失效，必须回到第一充分条件或高阶导判别。
\item 不可导点也可能是极值点，只解 $f'(x)=0$ 会漏解。
\item 极值是局部概念：极值不一定是最值，最值也不一定是极值（可能在端点取得）。
\end{itemize}
\end{callwarn}

\begin{callexam}{真题考法}
解答题的基础问法：求单调区间、极值、凹凸区间与拐点。选择题常给"$f$ 在 $x_0$ 取极大值"并配一个关于 $f'(x_0)$ 或 $f''(x_0)$ 的极限条件，让考生判断参数符号，方向是 $f'(x_0)=0$ 且极值类型与 $f''(x_0)$ 的符号相反。

\end{callexam}

\subsubsection{最大值、最小值及其应用 $\bigstar$$\bigstar$$\bigstar$}

\textbf{闭区间上连续函数的最值}：$f$ 在 $[a,b]$ 上连续必取到最大值与最小值。求法：

\begin{enumerate}
\item 求 $f'(x)=0$ 的根（驻点）；
\item 求 $f'$ 不存在但 $f$ 有定义的点；
\item 计算上述所有点以及\textbf{两个端点} $a,b$ 处的函数值；
\item 比较大小，最大者为最大值，最小者为最小值。
\end{enumerate}
\textbf{开区间或无穷区间}：若在区间内有唯一极值点，则该极值即为最值；否则要比较端点的极限值 $\lim_{x\to a^+}f(x)$ 与 $\lim_{x\to b^-}f(x)$，若为 $\infty$ 则该侧无最值。

\textbf{应用问题的建模步骤}：

\begin{enumerate}
\item 设变量并写清取值范围；
\item 建立目标函数（面积、体积、周长、距离、成本等）；
\item 用约束条件消元化为单变量函数；
\item 求导找驻点；实际问题中唯一驻点通常就是最值点（可用"实际问题必有最值"直接断定）；
\item 回到实际问题，写清答案与单位。
\end{enumerate}
\textbf{例}：面积为 $100\ \mathrm{m}^2$ 的长方形，求周长的最小值。 设长为 $x$，则宽为 $\frac{100}{x}$，$x>0$，周长

\[L(x)=2\left(x+\frac{100}{x}\right),\qquad L'(x)=2\left(1-\frac{100}{x^2}\right).\]

$L'(x)=0$ 得 $x=10$（舍去 $x=-10$）。$L''(10)=\frac{400}{1000}>0$，故 $x=10$ 为极小值点；由于是唯一驻点，最小值 $L(10)=40\ \mathrm{m}$。

\textbf{例}：求 $f(x)=x^3-3x$ 在 $[0,3]$ 上的最大值与最小值。 驻点 $x=1$（在区间内），比较 $f(0)=0$，$f(1)=-2$，$f(3)=18$，故最大值为 $18$，最小值为 $-2$。

\textbf{例（应用：最短距离）}：求点 $(0,3)$ 到曲线 $y=x^2$ 的最短距离。 设曲线上点为 $(t,t^2)$，则距离的平方

\[D(t)=t^2+(t^2-3)^2=t^4-5t^2+9 .\]

$D'(t)=4t^3-10t=2t(2t^2-5)$，令 $D'=0$ 得 $t=0$ 或 $t=\pm\sqrt{\frac52}$。比较：$D(0)=9$，$D\!\left(\pm\sqrt{\frac52}\right)=\frac{25}{4}-\frac{25}{2}+9=\frac{11}{4}$。故最短距离为 $\sqrt{\frac{11}{4}}=\frac{\sqrt{11}}{2}$。（也可直接用"曲线上点到定点距离最短处，其连线与该点处切线垂直"来验证。）

\textbf{例（应用：体积最大）}：用长为 $L$ 的铁丝围成一个矩形，问长宽各为多少时面积最大。 设长为 $x$，则宽为 $\frac L2-x\ (0<x<\frac L2)$，面积

\[S(x)=x\left(\frac L2-x\right)=\frac L2x-x^2,\qquad S'(x)=\frac L2-2x=0\Rightarrow x=\frac L4 .\]

此时宽也是 $\frac L4$，即围成正方形时面积最大，最大值 $\frac{L^2}{16}$（$S''=-\frac12<0$，且是唯一驻点，故为最大值）。

\textbf{应用题的答题规范}：写清"设变量$\to$写约束$\to$建目标函数$\to$求导$\to$判断极值类型$\to$回答实际问题（带单位）"，最后一句要落到实际问题上（如"当长为 $\frac L4$ 时面积最大"），只写数值不写结论会失分。

\begin{callwarn}{易错}
\begin{itemize}
\item 闭区间求最值必须把\textbf{端点值}与驻点值、不可导点值一起比较；漏端点是最常见失分点。
\item 开区间与无穷区间要先判断最值是否存在，再求极值。
\item 实际问题不要漏写自变量的取值范围，最后要回答实际问题（带单位）。
\item 若驻点唯一且实际问题必有最值，可直接断定，不必再算二阶导。
\end{itemize}
\end{callwarn}

\begin{callexam}{真题考法}
数二的常见应用背景：几何（面积、体积、最短距离）、物理（最值速度/功率）、简单经济问题。也会直接考"求 $f$ 在 $[a,b]$ 上的最大值与最小值"。

\end{callexam}

\subsection{图形、渐近线与曲率}

\subsubsection{凹凸性、拐点、渐近线 $\bigstar$$\bigstar$$\bigstar$}

\textbf{凹凸的定义}（统考约定）：设 $f$ 在区间 $I$ 上连续，若对 $I$ 上任意两点 $x_1,x_2$ 与任意 $\lambda\in(0,1)$ 有

\[f\big(\lambda x_1+(1-\lambda)x_2\big)\le \lambda f(x_1)+(1-\lambda)f(x_2),\]

则称 $f$ 在 $I$ 上是\textbf{凹}的（曲线的弦在曲线上方，形如碗）；不等号反向则称\textbf{凸}的（弦在曲线下方，形如拱）。考纲明确："在区间 $(a,b)$ 内，当 $f''(x)>0$ 时图形是凹的，当 $f''(x)<0$ 时图形是凸的。"

\textbf{判别法}（$f$ 二阶可导）：$f''(x)>0\Rightarrow$ 凹；$f''(x)<0\Rightarrow$ 凸。

\textbf{凹凸性的等价刻画}（便于判断与证明，均以"凹"为例）：

\begin{itemize}
\item $f''(x)>0$ 在区间上恒成立；
\item $f'(x)$ 在区间上单调递增；
\item 曲线位于其任意一点切线的\textbf{上方}：$f(x)\ge f(x_0)+f'(x_0)(x-x_0)$；
\item 曲线位于其任意两点割线的\textbf{下方}：$f(\lambda x_1+(1-\lambda)x_2)\le\lambda f(x_1)+(1-\lambda)f(x_2)$。
\end{itemize}
证明不等式时，前两条给出"构函数求导"的思路，后两条给出"用切线/割线放缩"的思路。

\textbf{拐点}：连续曲线上凹凸性发生改变的点 $(x_0,f(x_0))$。判别：

\begin{itemize}
\item \textbf{必要条件}：$f$ 在 $x_0$ 二阶可导且 $(x_0,f(x_0))$ 为拐点 $\Rightarrow f''(x_0)=0$；
\item \textbf{第一充分条件}：$f''$ 在 $x_0$ 两侧异号；
\item \textbf{第二充分条件}：$f''(x_0)=0$ 且 $f'''(x_0)\ne0$；
\item \textbf{高阶导判别}：$f''(x_0)=\cdots=f^{(n-1)}(x_0)=0$，$f^{(n)}(x_0)\ne0$，则 $n$ 为奇数时是拐点，$n$ 为偶数时不是；
\item $f''$ 不存在的点也可能是拐点（如 $y=\sqrt[3]{x}$ 在 $x=0$）。
\end{itemize}
拐点是\textbf{点}，必须写成坐标 $(x_0,f(x_0))$。

\textbf{三条渐近线}：

\begin{itemize}
\item \textbf{水平渐近线}：若 $\lim\limits_{x\to+\infty}f(x)=b$ 或 $\lim\limits_{x\to-\infty}f(x)=b$，则 $y=b$ 是水平渐近线。左右两侧要分别计算，可能得到不同的水平渐近线。
\item \textbf{铅直渐近线}：若 $\lim\limits_{x\to a^+}f(x)=\infty$ 或 $\lim\limits_{x\to a^-}f(x)=\infty$，则 $x=a$ 是铅直渐近线。常出现在无定义点、$\ln$ 的定义域端点、分母零点。
\item \textbf{斜渐近线}：若
\end{itemize}
\[a=\lim_{x\to+\infty}\frac{f(x)}{x},\qquad b=\lim_{x\to+\infty}\big[f(x)-ax\big]\]

都存在且 $a\ne0$，则 $y=ax+b$ 是 $x\to+\infty$ 方向的斜渐近线；$x\to-\infty$ 要另算一遍（可能得到另一条）。

\textbf{斜渐近线的速算法}：把 $f(x)$ 拆成"一次式 $ax+b$ + 趋于 0 的部分"，剩下的常数就是 $b$，拆的过程中可以用泰勒展开或有理化。例如

\[\frac{x^2}{x+1}=x-1+\frac{1}{x+1},\]

故 $y=x-1$ 是斜渐近线。

\textbf{有理函数的速判}：设 $f(x)=\frac{P_m(x)}{Q_n(x)}$，则 $m<n$ 时 $y=0$ 是水平渐近线；$m=n$ 时 $y=\frac{a_m}{b_n}$ 是水平渐近线；$m=n+1$ 时用长除法得斜渐近线。

\textbf{同一侧，水平渐近线与斜渐近线不能同时存在}；但可以一侧是水平、另一侧是斜。

\textbf{例}：求 $y=\frac{x^2}{x-1}$ 的所有渐近线。 $x\to1$ 时 $y\to\infty$，故 $x=1$ 是铅直渐近线。又

\[\frac{x^2}{x-1}=x+1+\frac{1}{x-1},\]

$x\to\pm\infty$ 时 $\frac{1}{x-1}\to0$，故 $y=x+1$ 是两侧共有的斜渐近线。

\textbf{例}：求 $y=\frac1x+\ln(1+e^x)$ 的渐近线条数。 $x\to0$ 时 $\frac1x\to\infty$，故 $x=0$ 是铅直渐近线（$x\to0$ 时 $\ln(1+e^x)$ 有界）。 $x\to-\infty$ 时 $\frac1x\to0$，$\ln(1+e^x)\to0$，故 $y=0$ 是 $x\to-\infty$ 方向的水平渐近线。 $x\to+\infty$ 时 $\ln(1+e^x)=x+\ln(1+e^{-x})$，故 $y=x$ 是 $x\to+\infty$ 方向的斜渐近线。 共有三条渐近线。

\textbf{高次幂乘积型函数的拐点速判}：若 $f(x)=(x-a)^ng(x)$，其中 $g(a)\ne0$ 且 $g$ 在 $a$ 附近充分光滑，则

\begin{itemize}
\item $n=1$ 或 $n=2$ 时：$f''(a)\ne0$，$a$ 不是拐点（$n=2$ 时 $a$ 是极值点）；
\item $n\ge3$ 且 $n$ 为\textbf{奇数}时：$a$ 是拐点的横坐标；
\item $n\ge3$ 且 $n$ 为\textbf{偶数}时：$a$ 不是拐点。
\end{itemize}
依据是 $f''(x)$ 在 $x=a$ 附近的主项为 $n(n-1)(x-a)^{n-2}g(a)$：$n-2$ 为奇数（即 $n$ 为奇数）时 $f''$ 变号，故为拐点；$n-2$ 为偶数时 $f''$ 不变号。

\textbf{例}：判断 $f(x)=(x-1)^3(x-2)^4(x-3)$ 的拐点个数（只看 $x=1,2,3$ 附近的性态）。 对因子 $(x-1)^3$：$n=3$ 为奇数，$x=1$ 是拐点；对 $(x-2)^4$：$n=4$ 为偶数，$x=2$ 不是拐点；对 $(x-3)^1$：$n=1$，$x=3$ 不是拐点（该处 $f''(3)\ne0$）。因此 $x=1$ 是拐点横坐标。这个方法在"数一数有几个拐点"的题目里可以秒杀。

\begin{callwarn}{易错}
\begin{itemize}
\item 求拐点容易漏掉 $f''$ \textbf{不存在}的点。
\item 拐点必须写坐标，只写横坐标会被扣分。
\item 斜渐近线的 $a,b$ 要分 $x\to+\infty$ 与 $x\to-\infty$ 分别计算，两侧结果可能不同。
\item 只算出 $a$ 存在就断定有斜渐近线是错的，必须 $b$ 也存在。
\item 同一侧不能同时有水平与斜渐近线；判断时应先在每一侧确定是哪一种。
\end{itemize}
\end{callwarn}

\begin{callexam}{真题考法}
选择题与填空题的高频点：求渐近线的条数、由凹凸性判断 $f''$ 的符号、判断拐点个数、由导函数图形反推原函数图形。陷阱集中在"分左右两侧讨论"与"漏掉 $f''$ 不存在的点"。

\end{callexam}

\subsubsection{函数图形的描绘 $\bigstar$$\bigstar$}

\textbf{作图五步}：

\begin{enumerate}
\item 确定定义域，考察奇偶性与周期性（可以只研究一半或一个周期，减少工作量）；
\item 求 $f'$ 与 $f''$，求出 $f'=0$、$f''=0$ 的点以及 $f'$、$f''$ 不存在的点；
\item 列表：用上述点把定义域分段，逐段标出 $f'$、$f''$ 的符号，得到单调性、凹凸性、极值与拐点；
\item 求渐近线（水平、铅直、斜），并算出关键点（与坐标轴的交点、极值点、拐点）；
\item 描点连线：用渐近线定"骨架"，用单调性与凹凸性定"形状"，画出图形。
\end{enumerate}
\textbf{列表模板}：一行为 $x$ 的分段点；一行为 $f'$ 的符号与 $f$ 的升降（极值）；一行为 $f''$ 的符号与 $f$ 的凹凸（拐点）。

\textbf{例}：作 $y=x^3-3x$ 的图形。 定义域为 $\mathbb R$，是奇函数，只需研究 $x\ge0$。$y'=3(x^2-1)$，$y''=6x$。

\begin{itemize}
\item $x>1$：$y'>0$，$y''>0$，递增且凹；
\item $0<x<1$：$y'<0$，$y''>0$，递减且凹；
\item $x=1$ 处取极小值 $-2$；$x=0$ 处 $y''=0$ 且左右变号，故 $(0,0)$ 是拐点；
\item 由奇对称性得 $x=-1$ 处取极大值 $2$，$x<-1$ 递增且凸，$-1<x<0$ 递减且凸。
\end{itemize}
无渐近线。连线即得图形。

\textbf{例（含渐近线的作图）}：作 $y=x+\frac1x$ 的图形。 定义域为 $\{x\ne0\}$，是奇函数。$y'=1-\frac{1}{x^2}=\frac{x^2-1}{x^2}$，$y''=\frac{2}{x^3}$。

\begin{itemize}
\item $x\to0$ 时 $y\to\infty$，故 $x=0$（$y$ 轴）是铅直渐近线；
\item $x\to\pm\infty$ 时 $y-x=\frac1x\to0$，故 $y=x$ 是斜渐近线；
\item $x>1$ 时 $y'>0$、$y''>0$：递增且凹；$0<x<1$ 时 $y'<0$、$y''>0$：递减且凹；
\item $x=1$ 处取极小值 $2$（由对称性 $x=-1$ 处取极大值 $-2$）；
\item $x<0$ 时 $y''<0$，无拐点。
\end{itemize}
作图时先画两条渐近线（$y$ 轴与 $y=x$），再在两个分支上按单调性、凹凸性描点连线。

\textbf{作图题的答题规范}：

\begin{enumerate}
\item 必须写出定义域，并考察奇偶性、周期性；
\item 必须列出 $f'$、$f''$ 的零点与不存在的点，并用表格或分段说明符号；
\item 必须写出渐近线方程、极值点、拐点坐标、与坐标轴的交点；
\item 图形只需"大致正确"：单调性、凹凸性、渐近线三者的走向不能画错。
\end{enumerate}
\begin{callwarn}{易错}
\begin{itemize}
\item 列表时符号区间必须与定义域的分段一致，不能漏段，也不能把分界点写错。
\item 作图题要标出与坐标轴的交点、极值点、拐点的坐标。
\item 若有渐近线，先画渐近线再描曲线。
\item 利用奇偶性可以只画一半，但要注意对称的是"关于原点"还是"关于 $y$ 轴"。
\end{itemize}
\end{callwarn}

\begin{callexam}{真题考法}
数二直接考"描绘函数图形"的整题较少，多与单调性、极值、拐点、渐近线合并在解答题中，或要求"画出函数的大致图形"（往往只需给出单调性、凹凸性与渐近线的表格）。

\end{callexam}

\subsubsection{曲率、曲率圆与曲率半径的概念与计算 $\bigstar$$\bigstar$}

\textbf{弧微分}（三种坐标形式）：

\[ds=\sqrt{1+y'^2}\,dx,\qquad
ds=\sqrt{\varphi'^2(t)+\psi'^2(t)}\,dt,\qquad
ds=\sqrt{r^2(\theta)+r'^2(\theta)}\,d\theta,\]

分别对应直角坐标 $y=f(x)$、参数方程 $x=\varphi(t),y=\psi(t)$、极坐标 $r=r(\theta)$。

\textbf{曲率}：设曲线二阶可导，则曲率

\[K=\frac{|y''|}{\left(1+y'^2\right)^{3/2}}.\]

参数方程形式：

\[K=\frac{|\varphi'(t)\psi''(t)-\varphi''(t)\psi'(t)|}{\left[\varphi'^2(t)+\psi'^2(t)\right]^{3/2}}.\]

\textbf{曲率半径}：当 $K\ne0$ 时

\[R=\frac1K=\frac{\left(1+y'^2\right)^{3/2}}{|y''|}.\]

\begin{callkey}{曲率两公式}
\[K=\frac{|y''|}{\left(1+y'^2\right)^{3/2}},\qquad R=\frac1K\]

参数方程下：$K=\frac{|x'y''-x''y'|}{\left(x'^2+y'^2\right)^{3/2}}$（这里 $x'$ 表示对参数 $t$ 求导）。

\end{callkey}

\textbf{曲率圆（密切圆）}：在曲线凹的一侧的法线上取点 $D$，使 $|MD|=R$，以 $D$ 为圆心、$R$ 为半径的圆称为曲线在该点的曲率圆。曲率圆与曲线在该点有相同的切线（一阶导数）和相同的凹凸性（二阶导数）。曲率中心（曲率圆的圆心）的坐标为

\[\alpha=x-\frac{y'(1+y'^2)}{y''},\qquad \beta=y+\frac{1+y'^2}{y''}.\]

例如单位圆上任意点（$y''>0$ 取上半圆）算出的曲率中心都是原点，可用来检验公式用得对不对。

\textbf{计算流程}：$\circled{1}$ 求 $y',y''$（参数方程则先算 $\varphi',\psi',\varphi'',\psi''$）；$\circled{2}$ 代入公式求 $K$；$\circled{3}$ 需要时取倒数得 $R$。

\textbf{例}：求 $y=\ln x$ 在点 $(1,0)$ 处的曲率与曲率半径。 $y'=\frac1x$，$y''=-\frac{1}{x^2}$，在 $x=1$ 处 $y'=1$，$y''=-1$，故

\[K=\frac{|-1|}{(1+1)^{3/2}}=\frac{1}{2\sqrt2},\qquad R=2\sqrt2 .\]

\textbf{例（参数方程）}：求摆线 $x=a(t-\sin t)$，$y=a(1-\cos t)$ 在 $t=\frac{\pi}{2}$ 处的曲率。 $\varphi'=a(1-\cos t)$，$\psi'=a\sin t$，$\varphi''=a\sin t$，$\psi''=a\cos t$。于是

\[\varphi'\psi''-\varphi''\psi'=a^2\big[(1-\cos t)\cos t-\sin^2t\big]=a^2(\cos t-1).\]

当 $t=\frac\pi2$ 时，$\varphi'=\psi'=a$，分子为 $|a^2(0-1)|=a^2$，分母为 $(a^2+a^2)^{3/2}=2\sqrt2\,a^3$，故

\[K=\frac{a^2}{2\sqrt2\,a^3}=\frac{1}{2\sqrt2\,a}.\]

\textbf{圆弧的曲率}：半径为 $R$ 的圆上各点曲率恒为 $\frac1R$，曲率半径为 $R$；直线曲率为 $0$。

\textbf{例（求弧微分）}：求曲线 $y=x^{3/2}$ 的弧微分，并写出它在 $x\in[0,1]$ 上的表达式。 $y'=\frac32x^{1/2}$，故

\[ds=\sqrt{1+\frac94x}\,dx .\]

\textbf{例（隐函数曲线的曲率）}：求 $x^2+y^2=1$ 在点 $\left(\frac{\sqrt2}{2},\frac{\sqrt2}{2}\right)$ 处的曲率。 对 $x^2+y^2=1$ 两边求导：$2x+2yy'=0$，故 $y'=-\frac xy$；再求导得

\[y''=-\frac{y-xy'}{y^2}=-\frac{y+x\cdot\frac xy}{y^2}=-\frac{y^2+x^2}{y^3}=-\frac{1}{y^3}.\]

代入点得 $y'=-1$，$y''=-\frac{1}{\left(\frac{\sqrt2}{2}\right)^3}=-2\sqrt2$。于是

\[K=\frac{2\sqrt2}{\left(1+1\right)^{3/2}}=\frac{2\sqrt2}{2\sqrt2}=1,\qquad R=1 .\]

这与"单位圆的曲率处处为 $1$"一致，可以用来检验计算是否正确。

\textbf{注意}：数二的曲率题一般只要求"会计算"，因此复习时只要把直角坐标与参数方程两个公式用熟即可；极坐标曲线按"先化为参数方程"处理（这一步在数二属于了解层次，不必深究）。

\begin{callwarn}{易错}
\begin{itemize}
\item 曲率公式的分子是 $|y''|$，分母是 $(1+y'^2)^{3/2}$，不要把分母写成 $1+y''$ 或 $1+y'^2$。
\item 曲率半径是曲率的倒数，不是相反数；曲率半径总为正。
\item 参数方程求曲率时，$x'y''-x''y'$ 的顺序不能颠倒，计算后要取绝对值。
\item 直线曲率为 0，此时曲率半径不存在（不是 $0$）。
\end{itemize}
\end{callwarn}

\begin{callexam}{真题考法}
填空题常直接给一个函数求某点的曲率或曲率半径，有时结合参数方程或极坐标曲线；也可能问"曲率圆的圆心坐标"。数二对曲率只要求"会计算"，一般不与其他考点综合命题。

\end{callexam}

\subsection*{本章易错清单}
\addcontentsline{toc}{subsection}{本章易错清单}

\begin{enumerate}
\item 把 $f'(x_0)$ 与 $[f(x_0)]'$ 混淆：后者是常数的导数，恒为 $0$。
\item 记反"可导必连续、连续不一定可导"的方向，或错误地认为"一元函数可导与可微是充分不必要关系"。
\item $|f(x)|$ 在 $f(x_0)=0$ 处的可导性必须用定义判断（等价于 $f'(x_0)=0$），不能只凭"图像有尖角"下结论。
\item 分段函数在分段点直接对表达式求导后代入，而不用定义求左右导数。
\item "过点 $P$ 的切线"忘记设切点 $(t,f(t))$，直接当成"在 $P$ 处的切线"。
\item 参数方程求二阶导漏除 $\frac{dx}{dt}$；隐函数求二阶导漏掉 $y$ 是 $x$ 的函数而产生的项。
\item 幂指函数 $u(x)^{v(x)}$ 直接套 $(x^\alpha)'$ 或 $(a^x)'$ 的公式。
\item 罗尔定理漏验"闭区间连续、开区间可导、端点等值"三个条件，或把 $\xi\in(a,b)$ 写成闭区间。
\item 泰勒展开阶数不足（未展到分母同阶）导致极限算错；证明不等式时误用皮亚诺余项。
\item 用柯西中值定理时忘记验证 $g'(x)\ne0$，或分子分母用了不同的 $\xi$。
\item 求极值只解 $f'(x)=0$ 而漏掉不可导点；$f''(x_0)=0$ 时仍用第二充分条件下结论。
\item 闭区间求最值漏掉端点函数值；单调区间用并集符号连接。
\item 斜渐近线只算 $a$ 不算 $b$，不分 $x\to+\infty$ 与 $x\to-\infty$ 讨论。
\item 拐点只写横坐标不写坐标 $(x_0,f(x_0))$，或漏掉 $f''$ 不存在的点。
\item 洛必达法则用在非未定式上，或在循环求导、振荡、不可导点等条件下硬用。
\item 曲率公式分母写成 $1+y''$，或把曲率半径写成 $-\frac1K$。
\end{enumerate}
\chapter{一元函数积分学}
\markboth{一元函数积分学}{一元函数积分学}
\subsection*{本章考纲要求}
\addcontentsline{toc}{subsection}{本章考纲要求}

\begin{quote}
\textbf{考试内容}

\begin{enumerate}
\item 原函数和不定积分的概念
\item 不定积分的基本性质
\item 基本积分公式
\item 定积分的概念和基本性质
\item 定积分中值定理
\item 积分上限的函数及其导数
\item 牛顿-莱布尼茨（Newton-Leibniz）公式
\item 不定积分和定积分的换元积分法与分部积分法
\item 有理函数、三角函数有理式和简单无理函数的积分
\item 反常（广义）积分
\item 定积分的应用
\end{enumerate}
\textbf{考试要求}

\begin{enumerate}
\item 理解原函数的概念，理解不定积分和定积分的概念。
\item 掌握不定积分的基本公式，掌握不定积分和定积分的性质及定积分中值定理，掌握换元积分法与分部积分法。
\item 会求有理函数、三角函数有理式和简单无理函数的积分。
\item 理解积分上限的函数，会求它的导数，掌握牛顿-莱布尼茨公式。
\item 理解反常积分的概念，了解反常积分收敛的比较判别法，会计算反常积分。
\item 掌握用定积分表达和计算一些几何量与物理量（平面图形的面积、平面曲线的弧长、旋转体的体积及侧面积、平行截面面积为已知的立体体积、功、引力、压力、质心、形心等）及函数平均值。
\end{enumerate}
\end{quote}
\subsection*{本章星级总表}
\addcontentsline{toc}{subsection}{本章星级总表}

\begingroup\small\setlength{\tabcolsep}{4.0pt}
\begin{longtable}{@{}L{0.4434\linewidth}L{0.1127\linewidth}L{0.1127\linewidth}L{0.1127\linewidth}L{0.1127\linewidth}@{}}
\toprule
\textbf{考点} & \textbf{星级} & \textbf{考纲要求} & \textbf{难度} & \textbf{常见题型} \\
\midrule
\endhead
原函数与不定积分的概念、基本积分公式表 & $\bigstar$$\bigstar$$\bigstar$ & 理解/掌握 & 基础 & 填空 \\
不定积分的性质与换元积分法（第一类、第二类） & $\bigstar$$\bigstar$$\bigstar$$\bigstar$$\bigstar$ & 掌握 & 中等 & 填空、解答 \\
分部积分法及其推广 & $\bigstar$$\bigstar$$\bigstar$$\bigstar$$\bigstar$ & 掌握 & 中等 & 填空、解答 \\
有理函数的积分（部分分式） & $\bigstar$$\bigstar$$\bigstar$$\bigstar$ & 会 & 中等 & 解答 \\
三角函数有理式与简单无理函数的积分 & $\bigstar$$\bigstar$$\bigstar$$\bigstar$ & 会 & 难 & 解答 \\
定积分的概念（黎曼和定义）与几何意义 & $\bigstar$$\bigstar$$\bigstar$ & 理解 & 中等 & 选择 \\
定积分的基本性质、积分中值定理与估值 & $\bigstar$$\bigstar$$\bigstar$$\bigstar$ & 掌握 & 中等 & 选择、证明 \\
变限积分函数及其导数 & $\bigstar$$\bigstar$$\bigstar$$\bigstar$$\bigstar$ & 理解/会 & 难 & 选择、解答 \\
牛顿-莱布尼茨公式 & $\bigstar$$\bigstar$$\bigstar$$\bigstar$$\bigstar$ & 掌握 & 基础 & 填空、解答 \\
定积分的换元法与分部积分法 & $\bigstar$$\bigstar$$\bigstar$$\bigstar$$\bigstar$ & 掌握 & 中等 & 解答 \\
对称区间奇偶性、周期性、区间再现等化简技巧 & $\bigstar$$\bigstar$$\bigstar$$\bigstar$ & 会 & 中等 & 解答 \\
分段函数与含绝对值、含根号的积分 & $\bigstar$$\bigstar$$\bigstar$$\bigstar$ & 会 & 中等 & 解答 \\
定积分计算综合（递推公式、Wallis 公式类） & $\bigstar$$\bigstar$$\bigstar$$\bigstar$ & 会 & 难 & 解答 \\
反常（广义）积分的概念与计算 & $\bigstar$$\bigstar$$\bigstar$$\bigstar$ & 理解/会 & 中等 & 填空、解答 \\
反常积分敛散性的判别（比较判别法、p 积分） & $\bigstar$$\bigstar$$\bigstar$$\bigstar$ & 了解/会 & 难 & 选择、解答 \\
定积分的几何应用：平面图形面积 & $\bigstar$$\bigstar$$\bigstar$$\bigstar$ & 掌握 & 中等 & 解答 \\
定积分的几何应用：旋转体体积与平行截面面积已知的立体体积 & $\bigstar$$\bigstar$$\bigstar$$\bigstar$ & 掌握 & 中等 & 解答 \\
定积分的几何应用：平面曲线弧长、旋转体侧面积 & $\bigstar$$\bigstar$$\bigstar$ & 掌握 & 中等 & 解答 \\
定积分的物理应用：功、引力、压力 & $\bigstar$$\bigstar$$\bigstar$ & 掌握 & 中等 & 解答 \\
质心、形心与函数平均值 & $\bigstar$$\bigstar$$\bigstar$ & 掌握 & 中等 & 解答 \\
积分中值问题与积分等式证明 & $\bigstar$$\bigstar$$\bigstar$$\bigstar$ & 会用 & 难 & 证明 \\
积分不等式与积分放缩 & $\bigstar$$\bigstar$$\bigstar$$\bigstar$ & 会用 & 难 & 证明 \\
\bottomrule
\end{longtable}
\endgroup

\section{不定积分}

一元函数积分学有两条主线：不定积分（求原函数，是求导的逆运算）与定积分（求累积量，是黎曼和的极限）。本章先讲不定积分，再讲定积分，最后讲反常积分与定积分的应用。全章的核心工具是三大积分法——凑微分、换元、分部积分，以及微积分基本定理。学习本章的正确顺序是：先把基本积分表背到条件反射，再掌握三大积分法的选择逻辑，然后理解变限积分与牛顿-莱布尼茨公式，最后把定积分当成工具去解决面积、体积、弧长、功、压力等问题。做题时最忌讳“见题就换元”，一定要先观察被积函数的结构。

\subsection{原函数与基本积分表}

\subsubsection{原函数与不定积分的概念、基本积分公式表 $\bigstar$$\bigstar$$\bigstar$}

\textbf{原函数的定义}：设函数 $f(x)$ 在区间 $I$ 上有定义，若存在区间 $I$ 上的可导函数 $F(x)$，使得对一切 $x\in I$ 都有

\[F'(x)=f(x),\]

则称 $F(x)$ 是 $f(x)$ 在区间 $I$ 上的一个\textbf{原函数}。

\textbf{原函数存在定理}：若 $f(x)$ 在区间 $I$ 上连续，则 $f(x)$ 在 $I$ 上必有原函数（且它的一个原函数就是变限积分 $\int_a^x f(t)\,dt$）。反之，若 $f(x)$ 在 $I$ 上有\textbf{第一类间断点}（可去间断点或跳跃间断点），则 $f(x)$ 在 $I$ 上\textbf{不存在}原函数（导函数不能有第一类间断点）；若 $f$ 只有\textbf{第二类间断点}，则原函数可能存在，也可能不存在，需具体分析。例如

\[f(x)=\begin{cases}2x\sin\dfrac{1}{x^2}-\dfrac{2}{x}\cos\dfrac{1}{x^2}, & x\neq 0,\\[2mm] 0, & x=0\end{cases}\]

在 $\mathbb{R}$ 上有第二类（振荡）间断点，但它有原函数 $F(x)=x^2\sin\frac{1}{x^2}\ (x\neq 0),\ F(0)=0$，可自行求导验证。

\textbf{理解要点}：原函数存在定理的价值在于它把“求原函数”从“凭经验猜”变成“用积分构造”——只要 $f$ 连续，$\int_a^{x}f(t)\,dt$ 就是它的一个原函数。注意它的逆否形式：若 $f$ 在某区间上存在第一类间断点，则 $f$ 在该区间上不存在原函数，所以符号函数 $\operatorname{sgn}x$、取整函数 $\lfloor x\rfloor$ 在整个实数轴上都没有原函数（但在不含间断点的子区间上各段都有原函数）。这是概念选择题的常见考点。另外要养成一个自检习惯：不定积分算完，对结果求导看是否等于被积函数。这一步只需几秒钟，却能拦住绝大多数低级错误。

\textbf{不定积分}：$f(x)$ 在区间 $I$ 上的全体原函数称为 $f(x)$ 的\textbf{不定积分}，记作

\[\int f(x)\,dx=F(x)+C,\]

其中 $C$ 为任意常数（称为\textbf{积分常数}）。若 $F,G$ 都是 $f$ 的原函数，则 $F(x)-G(x)\equiv C$，即原函数在"相差一个常数"的意义下唯一。

\textbf{互逆关系}（做题时最常用的两条）：

\[\frac{d}{dx}\int f(x)\,dx=f(x),\qquad \int F'(x)\,dx=F(x)+C.\]

\textbf{几何意义}：$y=F(x)+C$ 是\textbf{一族平行曲线}，它们在同一个横坐标 $x$ 处的切线斜率都等于 $f(x)$，常数 $C$ 只决定这条曲线沿 $y$ 轴上下的位置。因此“求不定积分”相当于“由斜率反推曲线族的形状”；要确定具体是哪一条曲线，还需要一个初值条件，这正是微分方程初值问题的由来。此外，不定积分的运算性质还有两条：$\frac{d}{dx}\int f(x)dx=f(x)$ 与 $\int F'(x)dx=F(x)+C$。

\textbf{基本积分公式表}（必须默写，共 16 组）：

\begingroup\small\setlength{\tabcolsep}{4.0pt}
\begin{longtable}{@{}L{0.3095\linewidth}L{0.3095\linewidth}L{0.3095\linewidth}@{}}
\toprule
\textbf{类型} & \textbf{公式} & \textbf{条件} \\
\midrule
\endhead
幂函数 & $\displaystyle\int x^{a}\,dx=\frac{x^{a+1}}{a+1}+C$ & $a\neq -1$ \\
幂函数 & $\displaystyle\int \frac{dx}{x}=\ln\lvert x\rvert+C$ & $x\neq 0$ \\
指数函数 & $\displaystyle\int e^{x}\,dx=e^{x}+C$，$\displaystyle\int a^{x}\,dx=\frac{a^{x}}{\ln a}+C$ & $a>0,\allowbreak\ a\neq 1$ \\
三角 & $\displaystyle\int \sin x\,dx=-\cos x+C$，$\displaystyle\int \cos x\,dx=\sin x+C$ & — \\
三角 & $\displaystyle\int \sec^2x\,dx=\tan x+C$，$\displaystyle\int \csc^2x\,dx=-\cot x+C$ & — \\
三角 & $\displaystyle\int \sec x\tan x\,dx=\sec x+C$，$\displaystyle\int \csc x\cot x\,dx=-\csc x+C$ & — \\
三角 & $\displaystyle\int \tan x\,dx=-\ln\lvert \cos x\rvert+C$，$\displaystyle\int \cot x\,dx=\ln\lvert \sin x\rvert+C$ & — \\
三角 & $\displaystyle\int \sec x\,dx=\ln\lvert \sec x+\tan x\rvert+C$，$\displaystyle\int \csc x\,dx=\ln\lvert \csc x-\cot x\rvert+C$ & — \\
反三角 & $\displaystyle\int \frac{dx}{1+x^2}=\arctan x+C$，$\displaystyle\int \frac{dx}{\sqrt{1-x^2}}=\arcsin x+C$ & — \\
推广 & $\displaystyle\int \frac{dx}{a^2+x^2}=\frac{1}{a}\arctan\frac{x}{a}+C$，$\displaystyle\int \frac{dx}{\sqrt{a^2-x^2}}=\arcsin\frac{x}{a}+C$ & $a>0$ \\
分式 & $\displaystyle\int \frac{dx}{x^2-a^2}=\frac{1}{2a}\ln\left\lvert\frac{x-a}{x+a}\right\rvert+C$ & $a>0$ \\
根式 & $\displaystyle\int \frac{dx}{\sqrt{x^2+a^2}}=\ln\left(x+\sqrt{x^2+a^2}\right)+C$ & — \\
根式 & $\displaystyle\int \frac{dx}{\sqrt{x^2-a^2}}=\ln\left\lvert x+\sqrt{x^2-a^2}\right\rvert+C$ & $\lvert x\rvert>a>0$ \\
根式 & $\displaystyle\int \sqrt{a^2-x^2}\,dx=\frac{x}{2}\sqrt{a^2-x^2}+\frac{a^2}{2}\arcsin\frac{x}{a}+C$ & $a>0$ \\
根式 & $\displaystyle\int \sqrt{x^2+a^2}\,dx=\frac{x}{2}\sqrt{x^2+a^2}+\frac{a^2}{2}\ln\left(x+\sqrt{x^2+a^2}\right)+C$ & — \\
双曲型 & $\displaystyle\int \frac{dx}{\sqrt{x^2-a^2}}$ 与幂型 $\displaystyle\int x^{a}\,dx$ 的适用区间要按定义域讨论 & 分段处理 \\
\bottomrule
\end{longtable}
\endgroup

\begin{callkey}{基本积分表使用要诀}
\begin{itemize}
\item 幂函数、指数、三角、反三角四类公式必须默写；其中 $\int\frac{dx}{x}=\ln\lvert x\rvert+C$、$\int\frac{dx}{1+x^2}=\arctan x+C$、$\int\frac{dx}{\sqrt{1-x^2}}=\arcsin x+C$ 出现频率最高。
\item 所有公式都可以整体代换使用：若 $\int f(x)dx=F(x)+C$，则 $\int f(ax+b)dx=\frac{1}{a}F(ax+b)+C\ (a\neq0)$。
\item 结果必须带任意常数 $C$；验证办法是对结果求导。
\end{itemize}
\end{callkey}

\textbf{解题套路}（拿到一个不定积分先想三件事）：

\begin{enumerate}
\item \textbf{能不能直接查表}：把被积函数拆成基本公式的线性组合（分部积分前也可先拆项）。
\item \textbf{能不能凑微分}：看是否存在"内层函数 $\times$ 内层函数的导数"的结构。
\item \textbf{能不能换元或分部}：含根号用代换，两类函数相乘用分部积分。
\end{enumerate}
\textbf{典型例子}

\[\int\left(3x^2-\frac{2}{x}+5\right)dx=x^3-2\ln\lvert x\rvert+5x+C.\]

自己求导验证：$(x^3-2\ln\lvert x\rvert+5x)'=3x^2-\frac{2}{x}+5$，正确。

再补充两种最基本的化简动作：

\begin{itemize}
\item \textbf{拆项}（分母有什么，分子就拆什么）：$\displaystyle\int\frac{x^2}{1+x^2}dx=\int\left(1-\frac{1}{1+x^2}\right)dx=x-\arctan x+C$。
\item \textbf{线性代换}：若 $\int f(x)dx=F(x)+C$，则 $\int f(ax+b)dx=\frac{1}{a}F(ax+b)+C\ (a\neq0)$。例如 $\int e^{3x-1}dx=\frac{1}{3}e^{3x-1}+C$、$\int\cos 2x\,dx=\frac{1}{2}\sin 2x+C$。这是“一步到位”的提速技巧，不必每次换元。
\end{itemize}
\textbf{一道概念例题}：设 $f$ 在 $[a,b]$ 上连续，$F$ 是 $f$ 的一个原函数，问 $\int_a^xf(t)\,dt$ 与 $F(x)$ 是什么关系？答：二者只相差一个常数，$\int_a^xf(t)dt=F(x)-F(a)$。若再给出 $F(a)=0$，则 $\int_a^xf(t)dt$ 就是 $f$ 在 $[a,x]$ 上取零初值的那个原函数。这一句是变限积分题目的思想起点。

\begin{callwarn}{易错点}
\begin{itemize}
\item \textbf{不定积分一定要写 $C$}，漏写 $C$ 是填空题最常见的扣分点。
\item $\int\frac{dx}{x}=\ln\lvert x\rvert+C$ 中的绝对值不能丢。
\item \textbf{原函数与不定积分不是一个东西}：原函数是一个具体函数，不定积分是一族函数（含 $C$）。
\item 若 $f$ 在含第一类间断点的区间上没有原函数，就不要硬求；这是选择题常考的概念陷阱。
\item 分段函数的原函数在分段点处必须\textbf{连续}（原函数可导必连续），各段常数 $C_1,C_2$ 不独立。
\end{itemize}
\end{callwarn}

\begin{callexam}{真题考法}
本考点以\textbf{填空题}为主：直接算基本公式的线性组合、判定"$F(x)$ 是否为 $f(x)$ 的原函数"、由 $\int f(x)\,dx=F(x)+C$ 反求 $f$ 或其导数（对两边求导即可）。概念题常把"连续/可积/有原函数"三者混在一起设陷阱。

\end{callexam}

\subsection{换元积分法}

\subsubsection{不定积分的性质与换元积分法（第一类、第二类） $\bigstar$$\bigstar$$\bigstar$$\bigstar$$\bigstar$}

\textbf{基本性质}（线性性）：

\[\int\left[\alpha f(x)+\beta g(x)\right]dx=\alpha\int f(x)\,dx+\beta\int g(x)\,dx.\]

\textbf{换元法的本质}：换元法就是把复合函数求导法则倒过来用。由 $\frac{d}{dx}F[\varphi(x)]=F'[\varphi(x)]\varphi'(x)$ 两边积分，立即得到第一类换元公式；第二类换元则是把 $x=\varphi(t)$ 代入，用 $\varphi'(t)dt$ 替换 $dx$，把“关于 $x$ 不好积的式子”翻译成“关于 $t$ 好积的式子”。判断用哪一类，只需看被积函数中“谁像是谁的导数”：出现 $\varphi(x)\varphi'(x)$ 结构就是第一类；出现根号、需要主动把 $x$ 换成 $t$ 就是第二类。

\textbf{第一类换元法（凑微分法）}：若 $f$ 有原函数 $F$，$u=\varphi(x)$ 可导，则

\[\int f\!\left[\varphi(x)\right]\varphi'(x)\,dx=\int f(u)\,du\Big\vert_{u=\varphi(x)}=F\!\left[\varphi(x)\right]+C.\]

关键在于把 $\varphi'(x)\,dx$ 凑成 $d\varphi(x)$。必须背下来的\textbf{凑微分表}：

\begingroup\small\setlength{\tabcolsep}{4.0pt}
\begin{longtable}{@{}L{0.4728\linewidth}L{0.4728\linewidth}@{}}
\toprule
\textbf{微分形式} & \textbf{等价变换} \\
\midrule
\endhead
$x\,dx$ & $\frac{1}{2}d(x^2)$ \\
$\frac{dx}{x}$ & $d(\ln\lvert x\rvert)$ \\
$\frac{dx}{\sqrt{x}}$ & $2\,d(\sqrt{x})$ \\
$e^{x}dx$ & $d(e^{x})$ \\
$a^{x}dx$ & $\frac{1}{\ln a}d(a^{x})$ \\
$\cos x\,dx$ & $d(\sin x)$ \\
$\sin x\,dx$ & $-d(\cos x)$ \\
$\sec^2x\,dx$ & $d(\tan x)$ \\
$\frac{dx}{1+x^2}$ & $d(\arctan x)$ \\
$\frac{dx}{\sqrt{1-x^2}}$ & $d(\arcsin x)$ \\
$\frac{dx}{x^2}$ & $-d\left(\frac{1}{x}\right)$ \\
$\sec x\tan x\,dx$ & $d(\sec x)$ \\
\bottomrule
\end{longtable}
\endgroup

这张表要形成条件反射：看到 $x\,dx$ 就想 $d(x^2)$，看到 $\frac{dx}{x}$ 就想 $d(\ln\lvert x\rvert)$，看到 $\frac{dx}{1+x^2}$ 就想 $d(\arctan x)$。考试中绝大多数凑微分题都是这张表里某一行“变形 $+$ 常数配平”的结果。

\textbf{第二类换元法（变量代换法）}：设 $x=\varphi(t)$ 单调、可导且 $\varphi'(t)\neq 0$，则

\[\int f(x)\,dx=\int f\!\left[\varphi(t)\right]\varphi'(t)\,dt,\]

求出右端关于 $t$ 的原函数后，再用 $t=\varphi^{-1}(x)$ \textbf{回代}。

常用代换：

\begin{itemize}
\item \textbf{三角代换}：$\sqrt{a^2-x^2}$ 令 $x=a\sin t$；$\sqrt{a^2+x^2}$ 令 $x=a\tan t$；$\sqrt{x^2-a^2}$ 令 $x=a\sec t$。
\item \textbf{根式代换}：$\sqrt[n]{ax+b}$ 令 $t=\sqrt[n]{ax+b}$；$\sqrt{\frac{ax+b}{cx+d}}$ 令 $t=\sqrt{\frac{ax+b}{cx+d}}$。
\item \textbf{倒代换}：分母次数较高（比分子高 2 次以上）时令 $x=\frac{1}{t}$。
\end{itemize}
\textbf{解第二类换元的四步}：$\circled{1}$ 选代换并\textbf{书面写出 $t$ 的取值范围}；$\circled{2}$ 求 $dx=\varphi'(t)dt$；$\circled{3}$ 化为关于 $t$ 的积分并算出；$\circled{4}$ 用直角三角形或恒等式回代。回代时画直角三角形最稳妥：以 $x=a\sin t$ 为例，画一个斜边为 $a$、对边为 $x$ 的直角三角形，邻边就是 $\sqrt{a^2-x^2}$，于是 $\cos t=\frac{\sqrt{a^2-x^2}}{a}$、$t=\arcsin\frac{x}{a}$，所有回代关系一次写清，不必死记。

\textbf{典型例子}

（1）凑微分：$\displaystyle\int xe^{x^2}dx=\frac{1}{2}\int e^{x^2}d(x^2)=\frac{1}{2}e^{x^2}+C$。

（2）三角代换：令 $x=a\sin t\ \left(t\in\left[-\frac{\pi}{2},\frac{\pi}{2}\right]\right)$，

\[\int\sqrt{a^2-x^2}\,dx=a^2\int\cos^2t\,dt=\frac{a^2}{2}\left(t+\frac{\sin 2t}{2}\right)+C=\frac{x}{2}\sqrt{a^2-x^2}+\frac{a^2}{2}\arcsin\frac{x}{a}+C.\]

（3）倒代换：对 $\displaystyle\int\frac{dx}{x^2\sqrt{x^2+1}}\ (x>0)$，令 $x=\frac{1}{t}$，则

\[\int\frac{dx}{x^2\sqrt{x^2+1}}=-\int\frac{t\,dt}{\sqrt{1+t^2}}=-\sqrt{1+t^2}+C=-\frac{\sqrt{x^2+1}}{x}+C.\]

（4）整体根式代换：$\displaystyle\int\frac{dx}{1+\sqrt[3]{x+1}}$，令 $t=\sqrt[3]{x+1}$（即 $x=t^3-1$，$dx=3t^2dt$），

\[\int\frac{3t^2}{1+t}dt=3\int\left(t-1+\frac{1}{1+t}\right)dt=\frac{3}{2}t^2-3t+3\ln\lvert 1+t\rvert+C,\qquad t=\sqrt[3]{x+1}.\]

根式代换的要领是“整个根号一起换”，换完之后 $x$ 与 $dx$ 都要用 $t$ 表示，不能只换一半。

\textbf{换元失败的情形}：$\int\frac{dx}{x\sqrt{1+x^2}}$ 若直接令 $x=\tan t$，会得到 $\int\frac{dt}{\sin t}$，仍要再凑一次微分；更好的做法是令 $x=\frac{1}{t}$（倒代换），或直接凑微分 $\frac{dx}{x\sqrt{x^2+1}}=-\frac{d(1/x)}{\sqrt{(1/x)^2+1}}$。这说明换元的优劣取决于“是否一步到位”，多练才能形成条件反射。

\begin{callwarn}{易错点}
\begin{itemize}
\item \textbf{换元必换限（定积分）／回代（不定积分）}：不定积分的结果最后必须换回 $x$，定积分则不回代但必须把上下限一起换掉。
\item 三角代换漏写 $dx$（如 $x=a\sin t$ 时 $dx=a\cos t\,dt$）是最常见的错误。
\item 三角代换后 $t$ 的范围必须保证 $x=\varphi(t)$ \textbf{单调}，否则回代时符号会错。
\item 凑微分时注意常数配平：$\int e^{2x}dx=\frac{1}{2}e^{2x}+C$，不要漏 $\frac{1}{2}$。
\end{itemize}
\end{callwarn}

\begin{callexam}{真题考法}
这是本章分值最高的考点之一。小题直接考凑微分（如 $\int xe^{-x^2}dx$、$\int\frac{\ln x}{x}dx$），大题常把换元作为第一步，后面再接分部积分或定积分计算。含 $\sqrt{a^2-x^2}$、$\sqrt{x^2+a^2}$ 的积分几乎一定用三角代换，并要求把回代过程写清楚（画直角三角形）。

\end{callexam}

\subsection{分部积分法}

\subsubsection{分部积分法及其推广 $\bigstar$$\bigstar$$\bigstar$$\bigstar$$\bigstar$}

\textbf{核心公式}（由乘积求导公式逆用得到）：

\[\int u\,dv=uv-\int v\,du,\qquad\text{即}\qquad \int u(x)v'(x)\,dx=u(x)v(x)-\int u'(x)v(x)\,dx.\]

适用条件：$u,v$ 在区间上连续可导，且 $\int v\,du$ 比原积分更容易积出。

\textbf{公式的来历}：由乘积求导公式 $(uv)'=u'v+uv'$ 两边积分得 $uv=\int v\,du+\int u\,dv$，移项即得分部积分公式。它的本质是\textbf{把求导的负担从一个因子转移到另一个因子上}：$u$ 求导后要变简单，$dv$ 积分后不能变复杂。判断“能不能用分部积分”，就看 $\int v\,du$ 是否比原积分更容易算。

\textbf{选取 $u$ 的口诀——“反对幂指三”}：反三角函数、对数函数、幂函数、指数函数、三角函数，\textbf{排在前面的取作 $u$}（指数函数与三角函数同时出现时用循环法，顺序无所谓）。背后的道理是：\textbf{当 $u$ 的那个函数求导后要变简单}（对数、反三角求导后变成有理式，幂函数求导后降次）。例如 $\int x\arctan x\,dx$ 应取 $u=\arctan x$：$\arctan x$ 求导后变成有理式 $\frac{1}{1+x^2}$；若取 $u=x$，则要先算出 $\int\arctan x\,dx$ 才能继续，反而更慢。

\textbf{三类必会题型}：

\begin{enumerate}
\item \textbf{降阶型}（幂函数 $\times$ 指/三角）：取幂函数为 $u$，反复分部直到幂函数变成常数。例如
\end{enumerate}
\[\int x^2e^{x}dx=e^{x}(x^2-2x+2)+C.\]

列表法的写法：左边一列写 $u=x^2$ 并逐次求导（$2x,\ 2,\ 0$），右边一列写 $dv=e^{x}dx$ 并逐次积分（$e^{x},\ e^{x},\ e^{x}$），对角相乘、符号正负相间，得 $x^2e^{x}-2xe^{x}+2e^{x}$。符号一旦记错结果必然不符，算完可用求导验证。

多次分部可用\textbf{列表法}加速：上方写多项式并逐次求导（$x^2,\ 2x,\ 2,\ 0$），下方写 $e^x$ 并逐次积分，对角相乘、符号正负相间，最后相加即得。

\begin{enumerate}
\item \textbf{循环型}（指数 $\times$ 三角）：连续两次分部后出现与原积分相同的项，移项解出：
\end{enumerate}
\[\int e^{ax}\sin bx\,dx=\frac{e^{ax}\left(a\sin bx-b\cos bx\right)}{a^2+b^2}+C,\qquad \int e^{ax}\cos bx\,dx=\frac{e^{ax}\left(a\cos bx+b\sin bx\right)}{a^2+b^2}+C.\]

例如 $\int e^{x}\cos x\,dx=\frac{e^{x}\left(\sin x+\cos x\right)}{2}+C$。

\begin{enumerate}
\item \textbf{构造 $dv=dx$ 型}（对数、反三角单独出现）：把被积函数看成 "$\ln x\cdot 1$"，取 $dv=dx$：
\end{enumerate}
\[\int\ln x\,dx=x\ln x-x+C,\qquad \int\arctan x\,dx=x\arctan x-\frac{1}{2}\ln(1+x^2)+C.\]

再看一个“边界项恰好干净”的定积分例子：

\[\int_0^1x\ln(1+x)\,dx=\left[\frac{x^2}{2}\ln(1+x)\right]_0^1-\frac{1}{2}\int_0^1\frac{x^2}{1+x}dx=\frac{1}{2}\ln 2-\frac{1}{2}\left(\ln 2-\frac{1}{2}\right)=\frac{1}{4}.\]

\textbf{递推公式的建立}：把 $\int \sin^n x\,dx$ 拆成 $\int \sin^{n-1}x\cdot\sin x\,dx$ 分部，得

\[I_n=\int\sin^n x\,dx=-\frac{1}{n}\sin^{n-1}x\cos x+\frac{n-1}{n}I_{n-2}\quad(n\geq 2).\]

这个公式的推导只用了两个动作：先拆出一个 $\sin x$ 与 $dx$ 凑成 $-d(\cos x)$，再分部积分一次，把 $\sin^{n-2}x\cos^2x$ 写成 $\sin^{n-2}x(1-\sin^2x)$，移项即得。凡是“降次型”递推，都是这一套动作的重复。

同理 $I_n=\int\frac{dx}{\left(x^2+a^2\right)^n}$、$I_n=\int_0^1x^ne^{-x}dx$ 等也能建立递推，见"定积分计算综合"一节。

\begin{callwarn}{易错点}
\begin{itemize}
\item $u$ 与 $dv$ 选反了会越积越复杂（如 $\int x\ln x\,dx$ 若取 $u=x$ 就做不下去），必须按口诀取 $u$。
\item 循环型移项时不要漏掉 $+C$；移项前两边的积分是"同一个不定积分"，可当作未知数解。
\item 分部积分\textbf{可以随时与换元法合用}：如 $\int\sin\sqrt{x}\,dx$ 先令 $t=\sqrt{x}$，再分部。
\item 定积分的分部积分公式是 $\int_a^bu\,dv=\left[uv\right]_a^b-\int_a^bv\,du$，边界项不能漏。
\end{itemize}
\end{callwarn}

\begin{callexam}{真题考法}
分部积分是解答题的主力工具，常以三种面貌出现：$\circled{1}$ 与换元法结合的综合积分；$\circled{2}$ 建立并证明递推公式（如 $I_n=\int_0^{\pi/2}\sin^nx\,dx$ 的递推）；$\circled{3}$ 在定积分中与对称性、区间再现配合。小题中 $\int x^ne^{x}dx$、$\int x\arctan x\,dx$、$\int e^x\sin x\,dx$ 属于必得分题。

\end{callexam}

\subsection{特殊类型函数的积分}

\subsubsection{有理函数的积分（部分分式） $\bigstar$$\bigstar$$\bigstar$$\bigstar$}

\textbf{有理函数}：$R(x)=\frac{P(x)}{Q(x)}$，其中 $P,Q$ 为多项式。若 $\deg P<\deg Q$ 称为\textbf{真分式}，否则称为\textbf{假分式}。假分式先做\textbf{多项式除法}，拆成"多项式 $+$ 真分式"。

\textbf{为什么可以这样做}：有理函数的原函数一定能用初等函数表示——因为分解之后的每一块只可能是 $\frac{A}{(x-a)^k}$ 或 $\frac{Bx+C}{(x^2+px+q)^m}$，前者的原函数是对数或幂函数，后者配方后化为对数与反正切。所以有理函数积分是一条\textbf{机械化}的路线：只要分解正确，剩下的只是“查表”。这也正是考纲专门要求“会求有理函数的积分”的原因。

\textbf{标准流程}：

\begin{enumerate}
\item \textbf{化真分式}：$\frac{P}{Q}=S(x)+\frac{R(x)}{Q(x)}$（$\deg R<\deg Q$）。
\item \textbf{分母因式分解}（在实数范围内分解到底）：$Q(x)=\prod_i(x-a_i)^{k_i}\prod_j\left(x^2+p_jx+q_j\right)^{m_j}$。
\item \textbf{部分分式分解}：分母中的 $(x-a)^k$ 必须写成 $k$ 项之和，$k$ 次幂因式\textbf{从低次到高次全部写出}：
\end{enumerate}
\[\frac{R(x)}{Q(x)}=\sum_i\sum_{r=1}^{k_i}\frac{A_{ir}}{(x-a_i)^{r}}+\sum_j\sum_{s=1}^{m_j}\frac{B_{js}x+C_{js}}{\left(x^2+p_jx+q_j\right)^{s}}.\]

\begin{enumerate}
\item \textbf{待定系数}：通分后与原分子比较系数，或用\textbf{特值法}（取 $x=a_i$ 及 $0,1,-1$ 等好算的点）求系数；单根处的系数有速算公式
\end{enumerate}
\[A_i=\lim_{x\to a_i}(x-a_i)\frac{R(x)}{Q(x)}.\]

\begin{enumerate}
\item \textbf{逐项积分}：只用到三类基本积分
\end{enumerate}
\[\int\frac{dx}{x-a}=\ln\lvert x-a\rvert+C,\qquad \int\frac{dx}{(x-a)^k}=\frac{-1}{(k-1)(x-a)^{k-1}}+C\ (k\geq2),\]

\[\int\frac{Bx+C}{x^2+px+q}\,dx\ \xrightarrow{\ \text{配方、拆分母导数}\ }\ \text{对数项}+\text{反正切项}.\]

\textbf{分子一次的分式要"拆出分母导数"}：把分子写成 $\frac{B}{2}\left(2x+p\right)+\left(C-\frac{Bp}{2}\right)$，第一部分凑成 $d(\ln)$，第二部分配方后凑成 $\arctan$。若分母是 $(x^2+a^2)^n\ (n\geq2)$，就用前面给出的递推公式逐次降次；数二中 $n=2$ 的情形已足够，\textbf{配方 $+$ 三角代换}是通用出路。记住一条：有理函数积分永远是“四步走”，不会无路可走。

\textbf{典型例子}

（1）重根必须列全：$\displaystyle\int\frac{dx}{x(x-1)^2}$。设

\[\frac{1}{x(x-1)^2}=\frac{A}{x}+\frac{B}{x-1}+\frac{C}{(x-1)^2},\]

特值法 $x=0$ 得 $A=1$，$x=1$ 得 $C=1$，比较 $x^2$ 系数得 $1+B=0$，故 $B=-1$。于是

\[\int\frac{dx}{x(x-1)^2}=\ln\lvert x\rvert-\ln\lvert x-1\rvert-\frac{1}{x-1}+C.\]

（2）分子一次拆分母导数：

\[\int\frac{2x+5}{x^2+3x+5}\,dx=\ln\left(x^2+3x+5\right)+\frac{4}{\sqrt{11}}\arctan\frac{2x+3}{\sqrt{11}}+C.\]

（3）分子分母同除 $x^2$ 型的拆项：$\displaystyle\int\frac{dx}{x^2(x^2+1)}=\int\left(\frac{1}{x^2}-\frac{1}{x^2+1}\right)dx=-\frac{1}{x}-\arctan x+C$。

（4）假分式先做除法：$\displaystyle\int\frac{x^3}{x^2+1}dx=\int\left(x-\frac{x}{x^2+1}\right)dx=\frac{x^2}{2}-\frac{1}{2}\ln\left(x^2+1\right)+C$。

（5）组合积分法（分子恰是分母导数的倍数）：$\displaystyle\int\frac{2x+1}{x^2+x+3}dx=\ln\left(x^2+x+3\right)+C$，因为 $(x^2+x+3)'=2x+1$。凡发现分子是分母的导数，立刻写对数，不必做部分分式。

\textbf{做题顺序小结}：先看分子与分母的次数（假分式先做除法），再分解分母，再定系数（特值法优先），最后逐项积分。四步中任何一步偷懒，后面的计算都会崩掉。

\begin{callwarn}{易错点}
\begin{itemize}
\item 分解式\textbf{漏项}是最致命的错误：分母 $(x-a)^k$ 对应 $k$ 项，二次因式的 $m$ 次幂对应 $m$ 项。
\item 每个部分分式的分子次数必须比对应因式\textbf{低一次}：一次因式配常数，二次因式配一次式 $Bx+C$。
\item 假分式不先做除法就分解，一定做不出来。
\item 分母 $x^4+1$、$x^4+4$ 一类要先在实数域上分解（$x^4+4=(x^2-2x+2)(x^2+2x+2)$），或用"同除 $x^2$、凑 $d\left(x\pm\frac{1}{x}\right)$"的技巧。
\end{itemize}
\end{callwarn}

\begin{callexam}{真题考法}
有理函数积分以\textbf{解答题}为主，常作为综合题的第一步（先通分分解，再逐项积分）。命中率最高的两类是"分母为二次质因式、分子为一次式"与"分母含重根"。不要忘记：分子是分母导数的整数倍时，直接得对数（\textbf{组合积分法}）。

\end{callexam}

\subsubsection{三角函数有理式与简单无理函数的积分 $\bigstar$$\bigstar$$\bigstar$$\bigstar$}

\textbf{三角函数有理式}：$R(\sin x,\cos x)=\frac{P(\sin x,\cos x)}{Q(\sin x,\cos x)}$。总原则是\textbf{化成有理函数}再积分。

\textbf{万能代换}：令 $t=\tan\frac{x}{2}$，则

\[\sin x=\frac{2t}{1+t^2},\qquad \cos x=\frac{1-t^2}{1+t^2},\qquad dx=\frac{2\,dt}{1+t^2}.\]

万能代换是"保底方案"，任何三角有理式都能用它化成有理函数，但计算量大，\textbf{能用特殊代换就别用万能代换}。

\textbf{特殊代换的判定（看奇偶性，能省一半计算）}：

\begingroup\small\setlength{\tabcolsep}{4.0pt}
\begin{longtable}{@{}L{0.4728\linewidth}L{0.4728\linewidth}@{}}
\toprule
\textbf{被积函数满足的性质} & \textbf{代换} \\
\midrule
\endhead
$R(-\sin x,\cos x)=-R(\sin x,\cos x)$ & $t=\cos x$ \\
$R(\sin x,-\cos x)=-R(\sin x,\cos x)$ & $t=\sin x$ \\
$R(-\sin x,-\cos x)=R(\sin x,\cos x)$ & $t=\tan x$ \\
只含 $\sin^2x,\cos^2x,\sin x\cos x$ & $t=\tan x$（更高效） \\
\bottomrule
\end{longtable}
\endgroup

\textbf{八条常用手法}：$\circled{1}$ 万能代换；$\circled{2}$ 分子分母同乘共轭式（平方差）；$\circled{3}$ 用 $\sin^2x+\cos^2x=1$ 制造可凑微分的项；$\circled{4}$ 用二倍角/降幂公式 $\cos^2x=\frac{1+\cos 2x}{2}$、$\sin^2x=\frac{1-\cos 2x}{2}$；$\circled{5}$ 积化和差公式处理 $\sin mx\cos nx$；$\circled{6}$ 辅助角公式 $a\sin x+b\cos x=\sqrt{a^2+b^2}\sin(x+\varphi)$；$\circled{7}$ 分母齐次时同除 $\cos^n x$，全部化为 $\tan x$；$\circled{8}$ \textbf{组合积分法}——把分子拆成 "$A\times$ 分母 $+\ B\times$ 分母的导数"。

\textbf{奇偶次幂的处理}：$\sin^mx\cos^nx$ 中若 $m$ 为奇数，拆出一个 $\sin x$ 与 $dx$ 凑成 $-d(\cos x)$，其余用 $\sin^2x=1-\cos^2x$ 换元；若 $n$ 为奇数同理；若 $m,n$ 都是偶数，用降幂公式降次。积化和差公式也常在这里派上用场：$\sin mx\cos nx=\frac{1}{2}\left[\sin(m+n)x+\sin(m-n)x\right]$，$\sin mx\sin nx=\frac{1}{2}\left[\cos(m-n)x-\cos(m+n)x\right]$，$\cos mx\cos nx=\frac{1}{2}\left[\cos(m-n)x+\cos(m+n)x\right]$。不同频率的正弦余弦相乘时，先化积为和，再逐项积分。

\textbf{三角有理式的解题顺序}（很重要）：$\circled{1}$ 先看被积函数对 $\sin x,\cos x$ 的奇偶性，能用 $t=\cos x$、$t=\sin x$、$t=\tan x$ 就绝不用万能代换；$\circled{2}$ 看分母是否齐次（同除 $\cos^nx$）或能否组合积分（分子拆成 $A\cdot$ 分母 $+\ B\cdot$ 分母的导数）；$\circled{3}$ 看是否含高次幂（奇次凑微分、偶次降幂）；$\circled{4}$ 以上都不行才用万能代换。万能代换一定做得出来，但计算量最大，属于“保底”而非常规武器。

\textbf{简单无理函数}：含 $\sqrt[n]{ax+b}$ 的令 $t=\sqrt[n]{ax+b}$；含 $\sqrt{\frac{ax+b}{cx+d}}$ 的令 $t=\sqrt{\frac{ax+b}{cx+d}}$；含 $\sqrt{ax^2+bx+c}$ 的\textbf{先配方}再三角代换，也可用倒代换 $x=\frac{1}{t}$。

\textbf{几个反复出现的形式}：$\int\frac{dx}{x^2\sqrt{x^2+1}}$、$\int\frac{dx}{x\sqrt{x^2+1}}$、$\int\sqrt{\frac{1-x}{1+x}}dx$ 的统一处理原则是“先有理化或倒代换，再化为有理函数或作三角代换”，不要一上来就硬套万能代换。

\textbf{典型例子}

（1）组合积分法：$\displaystyle\int\frac{\sin x+3\cos x}{\sin x+\cos x}\,dx$。设分子 $=A(\sin x+\cos x)+B(\cos x-\sin x)$，解得 $A=2,B=1$，于是

\[\int\frac{\sin x+3\cos x}{\sin x+\cos x}dx=2x+\ln\lvert\sin x+\cos x\rvert+C.\]

（2）奇次幂凑微分：$\displaystyle\int\sin^3x\,dx=\int\left(\cos^2x-1\right)d(\cos x)=\frac{\cos^3x}{3}-\cos x+C$。

（3）万能代换：$t=\tan\frac{x}{2}$，

\[\int\frac{dx}{2+\sin x}=\int\frac{dt}{t^2+t+1}=\frac{2}{\sqrt{3}}\arctan\frac{2t+1}{\sqrt{3}}+C,\quad t=\tan\frac{x}{2}.\]

（4）根式代换：令 $t=\sqrt{x}$（$x\geq0$），

\[\int\frac{dx}{1+\sqrt{x}}=\int\frac{2t}{1+t}dt=2\sqrt{x}-2\ln\left(1+\sqrt{x}\right)+C.\]

\begin{callwarn}{易错点}
\begin{itemize}
\item 万能代换写完 $\sin x,\cos x,dx$ 三件套后一定要检查 $dx$ 是否换成了 $\frac{2dt}{1+t^2}$。
\item 三角代换后必须写出 $t$ 的范围（如 $t=\tan\frac{x}{2}$ 时 $x\neq(2k+1)\pi$），保证回代合法。
\item 组合积分法拆分子时，$B$ 配的是\textbf{分母的导数}，不是 $\cos x-\sin x$ 的某种变形，要当场求导核对。
\item $\sqrt{x^2}=\lvert x\rvert$，去掉根号时别把符号弄丢。
\end{itemize}
\end{callwarn}

\begin{callexam}{真题考法}
本考点以\textbf{解答题}出现，难度偏高，常与换元、分部积分串联。判定技巧是：先看分母是否齐次（同除 $\cos^nx$ 走 $\tan$ 路线），再看能否组合积分（分子拆分母及其导数），最后才考虑万能代换。含 $\sqrt{\frac{x}{1-x}}$、$\sqrt{x^2+a^2}$ 的积分要会配方与代换。

\end{callexam}

\section{定积分的概念与性质}

不定积分解决"求原函数"，定积分解决"求累积量"。二者的桥梁是微积分基本定理（变限积分求导 $+$ 牛顿-莱布尼茨公式）。这一部分的重点不是计算，而是\textbf{概念与条件的准确}：定积分要求“有限区间、有界函数”；变限积分的可导性取决于被积函数在该点是否连续；牛顿-莱布尼茨公式要求原函数在闭区间上连续。真题中相当多的失分来自“条件没看清”，而不是算不出来。

\subsection{定积分的定义与性质}

\subsubsection{定积分的概念（黎曼和定义）与几何意义 $\bigstar$$\bigstar$$\bigstar$}

\textbf{定义}：设 $f(x)$ 在 $[a,b]$ 上有界。任取分点

\[a=x_0<x_1<\cdots<x_n=b,\]

记 $\Delta x_i=x_i-x_{i-1}$，$\lambda=\max\limits_{1\leq i\leq n}\Delta x_i$，在每个小区间上任取 $\xi_i\in[x_{i-1},x_i]$。若极限

\[\lim_{\lambda\to 0}\sum_{i=1}^{n}f(\xi_i)\Delta x_i\]

存在，且\textbf{与分割方式、与 $\xi_i$ 的取法都无关}，则称 $f$ 在 $[a,b]$ 上\textbf{黎曼可积}，该极限值称为 $f$ 在 $[a,b]$ 上的\textbf{定积分}，记作

\[\int_a^bf(x)\,dx=\lim_{\lambda\to0}\sum_{i=1}^nf(\xi_i)\Delta x_i.\]

\textbf{可积性}：

\begin{itemize}
\item 必要条件：可积必有界（有界不一定可积）。
\item 充分条件：$f$ 在 $[a,b]$ 上连续；或 $f$ 有界且只有有限个间断点；或 $f$ 在 $[a,b]$ 上单调有界。
\item 记忆链条：可导 $\Rightarrow$ 连续 $\Rightarrow$ 可积 $\Rightarrow$ 有界。注意链条的方向是单向的：有界不一定可积，可积不一定连续（有有限个跳跃间断点的有界函数仍可积），连续不一定可导。选择题常在这四个概念的箭头上做文章。
\end{itemize}
\textbf{理解定义的三层意思}：$\circled{1}$ \textbf{分割任意}——把 $[a,b]$ 任意分成 $n$ 份，不要求等分；$\circled{2}$ \textbf{取点任意}——在每个小区间上任取 $\xi_i$，不要求取端点或中点；$\circled{3}$ \textbf{极限与取法无关}——只有极限存在且不依赖前两步的取法，这个极限才叫定积分，$f$ 才叫可积。正因为有这样的“任意性”，定积分才能代表一个与分割方式无关的客观量（面积、质量、路程、功）。计算时为了方便\textbf{可以}取等分、取右端点，但那只是手段，不是定义。

\textbf{几何意义}：当 $f(x)\geq0$ 时，$\int_a^bf\,dx$ 等于曲边梯形的\textbf{面积}；当 $f$ 变号时，$\int_a^bf\,dx$ 表示 $x$ 轴上方图形的面积减去下方图形的面积，即\textbf{代数面积}。

\textbf{用定义求和式极限}：把 $n$ 项和改成 $\frac{1}{n}\sum f\left(\frac{i}{n}\right)$ 的形式，

\[\lim_{n\to\infty}\frac{1}{n}\sum_{i=1}^{n}f\left(\frac{i}{n}\right)=\int_0^1f(x)\,dx.\]

对应关系：$\lim\sum\to\int_0^1$，$\frac{1}{n}\to dx$，$\frac{i}{n}\to x$。

\textbf{三个必须满足的条件}：区间是 $[0,1]$、$n$ 等分、取端点（或取任意点 $\xi_i$，此时 $\xi_i\to x$）。若和式的分母不是齐次的 $n$（如含 $n^2+k$），通常要先放缩再用夹逼准则，不能硬套定义。

\textbf{典型例子}

\[\lim_{n\to\infty}\frac{1}{n}\sum_{k=1}^{n}\frac{1}{1+\frac{k}{n}}=\int_0^1\frac{dx}{1+x}=\ln 2.\]

\[\lim_{n\to\infty}\frac{1}{n}\sum_{k=1}^{n}\sin\frac{k\pi}{n}=\int_0^1\sin\pi x\,dx=\frac{2}{\pi}.\]

\textbf{几何意义的直接用法}：$\int_{-2}^{2}\sqrt{4-x^2}\,dx$ 不必求原函数，它是上半圆 $x^2+y^2=4\ (y\geq0)$ 的面积，等于 $2\pi$；同理 $\int_{-1}^{1}\sqrt{1-x^2}\,dx=\frac{\pi}{2}$。看到被积函数形如 $\sqrt{R^2-x^2}$ 且区间对称，先想面积。

\begin{callwarn}{易错点}
\begin{itemize}
\item 定积分是\textbf{代数面积}，不是几何面积；求几何面积时要加绝对值（见"几何应用"）。
\item 定积分的值只与 $f$、$a$、$b$ 有关，与积分变量用什么字母无关：$\int_a^bf(x)dx=\int_a^bf(t)dt$。
\item 用定义求极限必须严格化为 $\frac{1}{n}\sum f\left(\frac{i}{n}\right)$ 的\textbf{标准形状}（区间 $[0,1]$、$n$ 等分、取端点）；分母不是齐次 $n$ 的时候要先放缩再用夹逼。
\item 不要用"取右端点"代替定义的任意性，只能作为计算手段。
\end{itemize}
\end{callwarn}

\begin{callexam}{真题考法}
本考点以\textbf{选择题}为主，两种考法：$\circled{1}$ 用定义求 $n$ 项和的极限（常与夹逼准则配合）；$\circled{2}$ 判断某和式的极限对应哪个定积分。记住：任意取点 $\xi_i$ 也可写成 $\xi_i\to x$，但区间必须是 $[0,1]$、必须 $n$ 等分。

\end{callexam}

\subsubsection{定积分的基本性质、积分中值定理与估值 $\bigstar$$\bigstar$$\bigstar$$\bigstar$}

设 $f,g$ 在 $[a,b]$ 上可积。

\textbf{八条核心性质}：

\begin{enumerate}
\item \textbf{线性性}：$\int_a^b\left[\alpha f(x)+\beta g(x)\right]dx=\alpha\int_a^bf(x)dx+\beta\int_a^bg(x)dx$。
\item \textbf{区间可加性}：$\int_a^bf+\int_b^cf=\int_a^cf$（对任意 $a,b,c$ 成立，只要积分存在）。约定 $\int_a^af=0$，$\int_a^bf=-\int_b^af$。
\item \textbf{保号性}：若 $f(x)\geq0$，则 $\int_a^bf(x)dx\geq0$（$a<b$）。
\item \textbf{单调性}：若 $f(x)\leq g(x)$，则 $\int_a^bf\,dx\leq\int_a^bg\,dx$。
\item \textbf{绝对值不等式}：$\left\lvert\int_a^bf(x)dx\right\rvert\leq\int_a^b\lvert f(x)\rvert dx$。
\item \textbf{估值定理}：设 $m\leq f(x)\leq M$，则 $m(b-a)\leq\int_a^bf(x)dx\leq M(b-a)$。
\item \textbf{常数积分}：$\int_a^bk\,dx=k(b-a)$。
\item \textbf{奇偶性}：$f$ 为奇函数时 $\int_{-a}^af(x)dx=0$；$f$ 为偶函数时 $\int_{-a}^af(x)dx=2\int_0^af(x)dx$。
\end{enumerate}
\textbf{这些性质怎么用}：性质 1～8 是“计算与估计工具”，其中最常考的是保号性、单调性、绝对值不等式、估值定理与中值定理。真题里最常见的用法有三类：$\circled{1}$ 用保号性、单调性比较两个积分的大小；$\circled{2}$ 用估值定理给积分定出一个范围；$\circled{3}$ 用中值定理把积分号去掉，变成函数值的问题。注意所有性质的前提都是 $f$ \textbf{可积}，而不是连续。

\textbf{积分中值定理}：若 $f(x)$ 在 $[a,b]$ 上\textbf{连续}，则至少存在一点 $\xi\in[a,b]$，使

\[\int_a^bf(x)\,dx=f(\xi)(b-a).\]

\textbf{推广形式（第二积分中值定理的常用形式）}：若 $f$ 连续，$g$ 可积且\textbf{在 $[a,b]$ 上不变号}，则存在 $\xi\in[a,b]$，使

\[\int_a^bf(x)g(x)\,dx=f(\xi)\int_a^bg(x)\,dx.\]

它的几何意义是：曲边梯形的面积等于以 $f(\xi)$ 为高、同底矩形的面积，故 $f(\xi)=\frac{1}{b-a}\int_a^bf\,dx$ 就是 $f$ 在 $[a,b]$ 上的\textbf{平均值}。记忆口诀是“一点代表整体”：一个点的函数值可以代表整段的平均水平，前提是 $f$ 连续、权函数不变号。理解了这一点，后面“用积分中值定理求极限”“用平均值证明等式”的题目就都有了下手处。

\textbf{辅助工具（柯西-施瓦茨积分不等式，证题时常用）}：$f,g$ 在 $[a,b]$ 上连续，则

\[\left(\int_a^bf(x)g(x)\,dx\right)^2\leq\int_a^bf^2(x)\,dx\cdot\int_a^bg^2(x)\,dx,\]

当且仅当 $f,g$ 线性相关时取等号。

例如取 $f(x)=x$、$g(x)=1$，区间 $[0,1]$：左边 $=\left(\int_0^1x\,dx\right)^2=\frac{1}{4}$，右边 $=\int_0^1x^2dx\cdot\int_0^11\,dx=\frac{1}{3}$，确有 $\frac{1}{4}<\frac{1}{3}$。

\textbf{典型例子}

（1）比较大小：$x\in[0,1]$ 时 $x\geq x^2$，故 $\int_0^1x\,dx\geq\int_0^1x^2\,dx$。

（2）估值：$x\in[0,1]$ 时 $e^{-1}\leq e^{-x^2}\leq 1$，故 $\frac{1}{e}\leq\int_0^1e^{-x^2}dx\leq 1$。

（3）推广中值定理：$f$ 在 $[0,1]$ 上连续，取 $g(x)=x\geq0$（不变号），则存在 $\xi\in[0,1]$ 使

\[\int_0^1xf(x)\,dx=f(\xi)\int_0^1x\,dx=\frac{1}{2}f(\xi).\]

\begin{callwarn}{易错点}
\begin{itemize}
\item 估值定理要求 $m,M$ 是 $f$ 在 $[a,b]$ 上的\textbf{最小值与最大值}，用端点值代入往往不等号失效。
\item 保号性不可逆：$\int_a^bf\geq0$ 推不出 $f\geq0$。
\item 中值定理要求 $f$ \textbf{连续}；只有可积时定理不一定成立。若题目要 $\xi\in(a,b)$，需要额外论证（如先说明 $f$ 不恒为常数）。
\item 推广形式必须验证 $g$ \textbf{不变号}；$g$ 变号时结论可能不成立。
\end{itemize}
\end{callwarn}

\begin{callexam}{真题考法}
以\textbf{选择}与\textbf{证明}为主：比较两个积分的大小、由积分不等式反推函数性质、用中值定理证明存在性。命题人常把"$f$ 连续"改成"$f$ 可积"来设陷阱；也常考"由 $\int_a^bf\geq0$ 判断 $f$ 的符号"这类概念题。

\end{callexam}

\subsection{微积分基本定理}

\subsubsection{变限积分函数及其导数 $\bigstar$$\bigstar$$\bigstar$$\bigstar$$\bigstar$}

\textbf{定义}：设 $f(x)$ 在 $[a,b]$ 上可积，称

\[\Phi(x)=\int_a^xf(t)\,dt\qquad (x\in[a,b])\]

为\textbf{变上限积分函数}（积分上限的函数）。

\textbf{这一节为什么是重中之重}：变限积分把“积分”变成了“函数”，使我们可以对它求导，从而把积分问题（求累积量）转化为微分问题。它是连接微分学与积分学的唯一桥梁；历年真题都会在这一节设置考点，且常与极限、单调性、极值、微分方程综合命题。下面三条结论（连续、可导、原函数存在）的条件必须能一字不差地复述。

\textbf{两条基本定理}：

\begin{enumerate}
\item \textbf{连续性}：若 $f$ 在 $[a,b]$ 上可积且有界（$\lvert f\rvert\leq M$），则 $\Phi$ 在 $[a,b]$ 上连续，且满足
\end{enumerate}
\[\left\lvert\Phi(x+h)-\Phi(x)\right\rvert=\left\lvert\int_x^{x+h}f(t)dt\right\rvert\leq M\lvert h\rvert.\]

\begin{enumerate}
\item \textbf{可导性（微积分基本定理）}：若 $f$ 在点 $x_0$ 连续（或 $f$ 在 $[a,b]$ 上连续），则 $\Phi$ 在 $x_0$ 可导，且
\end{enumerate}
\[\Phi'(x_0)=f(x_0),\qquad\text{即}\qquad \frac{d}{dx}\int_a^xf(t)\,dt=f(x).\]

由此得\textbf{原函数存在定理}：连续函数必有原函数，且 $\int_a^xf(t)dt$ 就是它的一个原函数。

\textbf{变限积分可导性的细则（选择题高频）}：设 $x_0\in(a,b)$ 是 $f$ 的间断点，则

\begingroup\small\setlength{\tabcolsep}{4.0pt}
\begin{longtable}{@{}L{0.4199\linewidth}L{0.5258\linewidth}@{}}
\toprule
\textbf{$x_0$ 处 $f$ 的间断类型} & \textbf{变限积分 $\Phi$ 在 $x_0$ 的性质} \\
\midrule
\endhead
连续 & $\Phi$ 可导，且 $\Phi'(x_0)=f(x_0)$ \\
可去间断点（$\lim\limits_{x\to x_0}f(x)=A\neq f(x_0)$） & $\Phi$ 可导，但 $\Phi'(x_0)=A\neq f(x_0)$ \\
跳跃间断点（左右极限不相等） & $\Phi$ 连续但\textbf{不可导}（左右导数分别为 $f(x_0^+)$、$f(x_0^-)$） \\
第二类间断点 & 需具体讨论，可能可导 \\
\bottomrule
\end{longtable}
\endgroup

原因是 $\Phi'(x_0)=\lim\limits_{h\to0}\frac{1}{h}\int_{x_0}^{x_0+h}f(t)dt$，它只"看到" $f$ 在 $x_0$ 附近的极限值，而不是 $f(x_0)$ 本身。

\textbf{一般求导公式}（最常用）：

\[\frac{d}{dx}\int_{\varphi(x)}^{\psi(x)}f(t)\,dt=f\!\left[\psi(x)\right]\psi'(x)-f\!\left[\varphi(x)\right]\varphi'(x).\]

\begin{callkey}{变限积分求导公式（必背）}
\[\frac{d}{dx}\int_{\varphi(x)}^{\psi(x)}f(t)\,dt=f\!\left[\psi(x)\right]\psi'(x)-f\!\left[\varphi(x)\right]\varphi'(x).\]

\begin{itemize}
\item $f$ 连续 $\Rightarrow$ $\Phi(x)=\int_a^xf(t)dt$ 可导且 $\Phi'(x)=f(x)$，$\Phi$ 就是 $f$ 的一个原函数。
\item $f$ 可积 $\Rightarrow$ $\Phi$ 连续；$f$ 有跳跃间断点 $\Rightarrow$ $\Phi$ 连续但不可导。
\item 被积函数含 $x$ 时先提取或换元，再求导。
\end{itemize}
\end{callkey}

\textbf{被积函数含 $x$ 时的两种处理}：

\begin{itemize}
\item \textbf{能提就提}：$\displaystyle\int_0^x(x-t)f(t)\,dt=x\int_0^xf(t)dt-\int_0^xtf(t)dt$，再求导：
\end{itemize}
\[F'(x)=\int_0^xf(t)dt+xf(x)-xf(x)=\int_0^xf(t)dt,\qquad F''(x)=f(x).\]

\begin{itemize}
\item \textbf{不能提就换元}：把 $x$ 从被积函数中换出去（如令 $t=xu$），再求导。判断口诀是“\textbf{积分时 $t$ 是变量，求导时 $x$ 是变量}”：只要被积表达式里明明写着 $x$，就不能直接对 $x$ 求导，必须先把 $x$ 移出积分号或通过换元把它变成积分变量的函数。这是本章出错率最高的一处。
\end{itemize}
\textbf{二重变限积分求极限}：两条路。$\circled{1}$ \textbf{交换积分次序}，把二重变限积分化为单重变限积分（在 $x$ 出现处求导）；$\circled{2}$ 若被积函数在该点为\textbf{非零因子}，可用二重积分中值定理。例如

\[\lim_{x\to0}\frac{1}{x^2}\int_0^x\left(\int_0^te^{-u^2}du\right)dt\ \xrightarrow{\ \text{洛必达两次}\ }\ \lim_{x\to0}\frac{e^{-x^2}}{2}=\frac{1}{2}.\]

交换次序的做法是：把 $\int_0^xdt\int_0^te^{-u^2}du$ 看成区域 $0\leq u\leq t\leq x$ 上的二重积分，交换后为 $\int_0^xdu\int_u^xe^{-u^2}dt=\int_0^x(x-u)e^{-u^2}du$，再对 $x$ 用两次洛必达法则，同样得 $\frac{1}{2}$。两条路殊途同归，考场上挑计算量小的一条。

\textbf{典型例子}

（1）$\displaystyle\frac{d}{dx}\int_{x^2}^{x^3}\sin t\,dt=3x^2\sin x^3-2x\sin x^2$。

（2）$F(x)=\int_0^xe^{-t^2}dt$，则 $F'(x)=e^{-x^2}$，$F''(x)=-2xe^{-x^2}$。

（3）$\displaystyle\lim_{x\to0}\frac{\int_0^x\arctan t\,dt}{x^2}=\lim_{x\to0}\frac{\arctan x}{2x}=\frac{1}{2}$。

（4）被积函数含 $x$ 的处理：$G(x)=\displaystyle\int_0^x\sin(x-t)\,dt$，令 $u=x-t$ 得 $G(x)=\int_0^x\sin u\,du$，故 $G'(x)=\sin x$；也可直接算出 $G(x)=1-\cos x$，同样得 $G'(x)=\sin x$。

\begin{callwarn}{易错点}
\begin{itemize}
\item \textbf{被积函数含 $x$ 时不能直接套求导公式}（积分变量是 $t$，求导变量是 $x$，二者不同），必须先提取或换元。
\item 上限是复合函数时别漏 $\psi'(x)$：$\frac{d}{dx}\int_0^{\sin x}f(t)dt=f(\sin x)\cos x$。
\item $\Phi(x)$ 可导 $\neq$ $f$ 连续：$f$ 有跳跃间断点时 $\Phi$ 连续但不可导；$f$ 有可去间断点时 $\Phi$ 可导但 $\Phi'\neq f$。
\item 变限积分的下限是 $x$ 时别漏负号：$\frac{d}{dx}\int_x^bf(t)dt=-f(x)$。
\item 用洛必达法则前先确认是 $\frac{0}{0}$ 或 $\frac{\infty}{\infty}$ 型。
\end{itemize}
\end{callwarn}

\begin{callexam}{真题考法}
这是本章\textbf{最高频}的考点，选择、填空、解答通吃。典型考法：$\circled{1}$ 变限积分与极限结合（洛必达法则）；$\circled{2}$ 由 $\Phi(x)=\int_a^xf(t)dt$ 研究单调性、极值、凹凸性；$\circled{3}$ 含参变限积分（被积函数含 $x$）求导；$\circled{4}$ 已知 $\int_0^x f(t)dt=g(x)$ 反求 $f$；$\circled{5}$ 变限积分定义的隐函数求导；$\circled{6}$ 二重变限积分求极限。做这类题的第一句话往往是"先求导"。

\end{callexam}

\subsubsection{牛顿-莱布尼茨公式 $\bigstar$$\bigstar$$\bigstar$$\bigstar$$\bigstar$}

\textbf{公式}：若 $f(x)$ 在 $[a,b]$ 上连续，$F(x)$ 是 $f(x)$ 在 $[a,b]$ 上的任意一个原函数，则

\[\int_a^bf(x)\,dx=F(b)-F(a)=\left[F(x)\right]_a^b.\]

\textbf{公式的意义}：它把“求和取极限”这件苦差事换成“求原函数后代两端点”这一件轻松事，是定积分能够被广泛计算的根本原因。\textbf{证明思路}：令 $\Phi(x)=\int_a^xf(t)dt$，由变限积分求导得 $\Phi'(x)=f(x)$，于是 $\Phi(x)$ 与 $F(x)$ 只相差一个常数，代入 $x=a$ 定出这个常数后令 $x=b$，即得 $\int_a^bf=F(b)-F(a)$。可见牛顿-莱布尼茨公式是变限积分求导的直接推论。\textbf{使用公式的三步}：$\circled{1}$ 看清被积函数在 $[a,b]$ 上是否连续（有无瑕点、跳跃点）；$\circled{2}$ 求出一个原函数 $F$（分段时哪一段用哪个原函数要写清）；$\circled{3}$ 代入端点作差 $F(b)-F(a)$。凡是含绝对值、分段、$\min/\max$ 的被积函数，第$\circled{1}$步就要先分段。

\textbf{三个前提}（缺一不可）：$\circled{1}$ $f$ 在 $[a,b]$ 上连续（可减弱为可积且存在原函数）；$\circled{2}$ $F'(x)=f(x)$ 在 $[a,b]$ 上成立；$\circled{3}$ $F$ 在区间端点单侧连续，即 $F$ 在 $[a,b]$ 上连续。

\textbf{推广形式}：

\begin{itemize}
\item 分段函数：$\int_a^bf(x)dx=\sum\limits_{k=1}^n\int_{c_{k-1}}^{c_k}f(x)dx$（$a=c_0<c_1<\cdots<c_n=b$）。
\item 与变限积分配合：$\int_a^bf(x)dx=F(x)\big\vert_a^b$，其中 $F(x)=\int_a^xf(t)dt$。
\end{itemize}
\textbf{典型例子}

\[\int_1^e\ln x\,dx=\left[x\ln x-x\right]_1^e=(e-e)-(0-1)=1.\]

分段函数用公式的例子：设 $f(x)=\begin{cases}2x, & 0\leq x\leq1,\\ 1, & 1<x\leq2,\end{cases}$ 则

\[\int_0^2f(x)dx=\int_0^12x\,dx+\int_1^21\,dx=1+1=2.\]

要点是分段点只算一次，每一段用该段自己的原函数。

\begin{callwarn}{易错点}
\begin{itemize}
\item \textbf{含瑕点或跳跃间断点时不能直接用公式}。反例：$\int_{-1}^{1}\frac{dx}{x^2}$，若误用 $F(x)=-\frac{1}{x}$ 得 $-2$，实际该积分发散；$F$ 在 $x=0$ 不连续，公式的前提被破坏。
\item 端点处 $F$ 必须连续。如 $\int_0^1\left\lfloor x\right\rfloor$ 型、含绝对值的积分，都要先分段。
\item 写答案时注意符号：$F(b)-F(a)$ 不要写成 $F(a)-F(b)$。
\end{itemize}
\end{callwarn}

\begin{callexam}{真题考法}
牛顿-莱布尼茨公式是\textbf{一切定积分计算题的地基}，本身极少单独命题，但错用它（在有瑕点或分段点时）是丢分的重灾区。选择题常给出一个"看起来能用公式、其实被积函数在区间内有瑕点"的积分来考概念。

\end{callexam}

\section{定积分的计算}

定积分的计算方法与不定积分一一对应，但有三点不同：换元要换限、分部积分有边界项、可利用奇偶性周期性与区间再现简化。本章的计算题几乎都在这三点上做文章。做题前先读一遍被积函数与积分区间携带的“对称性信息”（奇偶性、周期性、是否关于区间中点对称、能否用区间再现），往往能省掉一半工作量。另外要养成一个习惯：算之前先估计量级。例如被积函数在区间上恒小于 $1$，积分值一定小于区间长度；被积函数恒正，积分值一定为正。这些粗判能帮你及时发现自己算错了。

\subsection{定积分的换元与分部}

\subsubsection{定积分的换元法与分部积分法 $\bigstar$$\bigstar$$\bigstar$$\bigstar$$\bigstar$}

\textbf{定积分换元与不定积分换元的两点不同}：$\circled{1}$ 定积分换元后\textbf{不必回代}，因为最后得到一个数，但\textbf{必须把上下限一起换掉}；$\circled{2}$ 定积分对被积函数的要求更宽——只要 $f$ 连续、$\varphi$ 有连续导数并把 $[\alpha,\beta]$ 映满 $[a,b]$，即使 $\varphi$ 不单调，等式仍成立（数二题目一般都用单调代换，以保证换限清晰）。

\textbf{换元积分法}：设 $f(x)$ 在 $[a,b]$ 上连续，$x=\varphi(t)$ 在 $[\alpha,\beta]$ 上单调、有连续导数，且 $\varphi(\alpha)=a$，$\varphi(\beta)=b$，则

\[\int_a^bf(x)\,dx=\int_{\alpha}^{\beta}f\!\left[\varphi(t)\right]\varphi'(t)\,dt.\]

\textbf{五步操作}：$\circled{1}$ 选代换；$\circled{2}$ 由 $a=\varphi(\alpha)$、$b=\varphi(\beta)$ 定出新限；$\circled{3}$ 求 $dx=\varphi'(t)dt$；$\circled{4}$ 换元后计算；$\circled{5}$ 定积分的结果是一个数，\textbf{不必回代}。

常见代换与不定积分类似：$\sqrt{a^2-x^2}\to x=a\sin t$，$\sqrt{a^2+x^2}\to x=a\tan t$，$\sqrt{x^2-a^2}\to x=a\sec t$，$\sqrt[n]{ax+b}\to t=\sqrt[n]{ax+b}$，分母次数高用 $x=\frac{1}{t}$。

\textbf{分部积分法}：设 $u(x),v(x)$ 在 $[a,b]$ 上有连续导数，则

\[\int_a^bu\,dv=\left[uv\right]_a^b-\int_a^bv\,du,\qquad\text{即}\qquad\int_a^bu(x)v'(x)dx=\left[u(x)v(x)\right]_a^b-\int_a^bu'(x)v(x)dx.\]

选 $u$ 仍按"反对幂指三"；\textbf{边界项 $\left[uv\right]_a^b$ 一定要写出来}，常常正好等于 $0$，能把问题简化。

\textbf{典型例子}

（1）换元：令 $x=\sin t\ (t\in[0,\frac{\pi}{2}])$，

\[\int_0^1\sqrt{1-x^2}\,dx=\int_0^{\pi/2}\cos^2t\,dt=\frac{1}{2}\cdot\frac{\pi}{2}=\frac{\pi}{4}.\]

（2）换元再分部：$\displaystyle\int_0^1xe^{-x}dx=\left[-xe^{-x}\right]_0^1+\int_0^1e^{-x}dx=-\frac{1}{e}+\left(1-\frac{1}{e}\right)=1-\frac{2}{e}$。

（3）分部：$\displaystyle\int_0^{\pi/2}x\sin x\,dx=\left[-x\cos x\right]_0^{\pi/2}+\int_0^{\pi/2}\cos x\,dx=0+1=1$。

\textbf{含参数的换元技巧}：$\displaystyle\int_0^{\pi/2}\frac{dx}{a^2\cos^2x+b^2\sin^2x}$，令 $t=\tan x$，化为 $\displaystyle\int_0^{+\infty}\frac{dt}{a^2+b^2t^2}=\frac{\pi}{2ab}$（$a,b>0$）——这是数二中常见的化归手法。

换元的另一个常用场景是\textbf{把区间搬到题设条件给定的区间上}：已知 $f$ 在 $[0,1]$ 上连续而要算 $\int_1^2f(x-1)dx$，作 $t=x-1$ 即可回到 $\int_0^1f(t)dt$。区间不合适时先想换元，这是定积分计算的一条通用经验。

\begin{callwarn}{易错点}
\begin{itemize}
\item \textbf{换元必换限}，且新限要与新变量对应；只换被积函数不换限是最常见的错误。
\item 定积分换元要求 $\varphi$ 单调（或至少把 $[\alpha,\beta]$ 映满 $[a,b]$），否则会多算或漏算。
\item 分部积分的\textbf{边界项}不能省；两个边界都在 $0$ 处时要明确指出。
\item 定积分中被积函数含绝对值、$\min$、$\max$、$\left\lfloor\cdot\right\rfloor$ 时，必须先分段再套公式。
\end{itemize}
\end{callwarn}

\begin{callexam}{真题考法}
本考点是\textbf{解答题的常规配置}，一般与换元、分部、对称性组合出现。命题人喜欢把换元设计得"顺手"：看到 $\sqrt{1-x^2}$、$\sqrt{e^x-1}$、$\ln x$、$\arcsin x$ 就分别想到三角代换、整体代换、分部积分。分区间的题目务必把每段的上下限写清楚。

\end{callexam}

\subsection{定积分的计算技巧}

\subsubsection{对称区间奇偶性、周期性、区间再现等化简技巧 $\bigstar$$\bigstar$$\bigstar$$\bigstar$}

\textbf{1. 奇偶性（对称区间"偶倍奇零"）}

\[\int_{-a}^{a}f(x)dx=\begin{cases}0, & f\ \text{为奇函数},\\[1mm]2\displaystyle\int_0^{a}f(x)dx, & f\ \text{为偶函数}.\end{cases}\]

注意"奇函数 $\times$ 偶函数 $=$ 奇函数"，即使被积函数含有 $e^x$、$x^3\cos x$ 等，也常可通过\textbf{拆分}得到奇偶项：如 $f(x)=\frac{1}{1+e^x}$ 不奇不偶，但 $f(x)+f(-x)=1$，于是 $\int_{-a}^a\frac{dx}{1+e^x}=a$。

\textbf{2. 周期性}

\[\int_a^{a+T}f(x)dx=\int_0^{T}f(x)dx,\qquad \int_a^{a+nT}f(x)dx=n\int_0^{T}f(x)dx\quad (n\in\mathbb{N}).\]

\textbf{3. 区间再现公式}

\[\int_a^bf(x)\,dx=\int_a^bf(a+b-x)\,dx\qquad(\text{令 }x=a+b-t).\]

常用推论：

\begin{itemize}
\item $\displaystyle\int_0^{\pi}xf(\sin x)dx=\frac{\pi}{2}\int_0^{\pi}f(\sin x)dx$；
\item $\displaystyle\int_0^{\pi/2}f(\sin x)dx=\int_0^{\pi/2}f(\cos x)dx$；
\item 若 $f(x)+f(a+b-x)=1$，则 $\displaystyle\int_a^bf(x)dx=\frac{b-a}{2}$。
\end{itemize}
\textbf{4. 平移与构造对称区间（"区间减函数加、区间加函数减"）}：当题设条件给的是 $[0,1]$ 上的性质，而所求积分区间是 $[1,2]$ 时，作平移 $x=t+1$ 把区间搬回去；当需要凑出对称区间 $[-a,a]$ 时，把被积函数平移构造。

\textbf{典型例子}

（1）$\displaystyle\int_0^{\pi}\frac{x\sin x}{1+\cos^2x}dx$。令 $x=\pi-t$，

\[I=\int_0^{\pi}\frac{(\pi-t)\sin t}{1+\cos^2t}dt=\pi\int_0^{\pi}\frac{\sin t}{1+\cos^2t}dt-I,\]

故 $2I=\pi\int_{-1}^{1}\frac{du}{1+u^2}=\pi\cdot\frac{\pi}{2}$（$u=\cos t$），得 $I=\frac{\pi^2}{4}$。

（2）$\displaystyle\int_0^{\pi/2}\frac{dx}{1+\tan^{2025}x}=\frac{\pi}{4}$。因 $\tan\left(\frac{\pi}{2}-x\right)=\cot x$，两式相加被积函数恒为 $1$。

（3）$\displaystyle\int_{-1}^{1}\left(x^5+x^3\cos x+x^2\right)dx=0+\int_{-1}^{1}x^2dx=\frac{2}{3}$。

（4）构造对称区间：已知 $f$ 在 $[0,1]$ 上有某性质，要算 $\int_1^2f(x-1)dx$，作平移 $t=x-1$ 即化为 $\int_0^1f(t)dt$；反之，若被积函数含 $\frac{1}{1+e^x}$ 这类不对称因子，可在 $[-a,a]$ 上用 $f(x)+f(-x)=1$ 的配对技巧，两式相加即得结果。

\textbf{三个必须记住的结论}：$\circled{1}$ $\int_{-a}^{a}f=0$（$f$ 为奇函数）；$\circled{2}$ $\int_0^{\pi}xf(\sin x)dx=\frac{\pi}{2}\int_0^{\pi}f(\sin x)dx$；$\circled{3}$ $\int_0^{\pi/2}f(\sin x)dx=\int_0^{\pi/2}f(\cos x)dx$。它们覆盖了数二中绝大多数“看着麻烦、实际一步解决”的对称性题目。

\begin{callwarn}{易错点}
\begin{itemize}
\item 奇偶性只适用于\textbf{关于原点对称}的区间；区间是 $[0,2]$ 时要先平移。
\item 周期函数只有"积满一个整周期"才与起点无关，积半个周期不能这样算。
\item 区间再现后要与原式\textbf{相加或相减}，不能直接换完就算完。
\item 对称区间中被积函数含参数时（如 $\frac{1}{1+a^x}$），要用 $f(x)+f(-x)$ 的技巧，不要硬算。
\end{itemize}
\end{callwarn}

\begin{callexam}{真题考法}
本考点以\textbf{解答题中的关键一步}出现，也单独出现在小题中。最典型的两种设问：$\circled{1}$ 求 $\int_0^{\pi}xf(\sin x)dx$ 型；$\circled{2}$ 求 $\int_0^{\pi/2}\frac{dx}{1+\tan^n x}$ 型。化简顺序建议：先看奇偶性，再看周期性，再用区间再现。

\end{callexam}

\subsubsection{分段函数与含绝对值、含根号的积分 $\bigstar$$\bigstar$$\bigstar$$\bigstar$}

\textbf{为什么要分段}：定积分的定义与计算公式都建立在“被积函数在区间上有一个统一表达式”这一前提上。遇到绝对值、分段函数、$\min/\max$、去根号时，被积函数在不同部分规则不同，必须先按“规则改变的位置”拆分区间，使每一段上被积函数有一个干净的表达式，才能套牛顿-莱布尼茨公式。分界点通常是\textbf{零点、交点、分段点本身}。

\textbf{处理原则：分段是万能的，先找"分界点"}。

\begin{itemize}
\item \textbf{含绝对值的积分}：令绝对值内的式子为零，解出分界点，把积分区间拆开，在每一段上去掉绝对值符号。
\item \textbf{分段函数的积分}：按分段点拆分区间；若积分限本身含参数，需按参数与分段点的位置\textbf{分类讨论}。
\item \textbf{含 $\min\{f,g\}$、$\max\{f,g\}$ 的积分}：先解 $f=g$ 求出分界点，再在每一段上判断大小从而确定被积函数。
\item \textbf{含根号的积分}：$\sqrt{x^2}=\lvert x\rvert$；$\sqrt{\sin^2x}=\lvert\sin x\rvert$；$\sqrt{(x-a)^2}=\lvert x-a\rvert$，去根号就是去绝对值。
\item \textbf{分段函数求原函数}：分段求不定积分，各段常数分别记为 $C_1,C_2,\dots$，再用"原函数在分段点处连续"把常数串起来（有时还要验证分段点处可导）。
\end{itemize}
\textbf{典型例子}

（1）$\displaystyle\int_0^4\lvert x-1\rvert dx=\int_0^1(1-x)dx+\int_1^4(x-1)dx=\frac{1}{2}+\frac{9}{2}=5$。

（2）$\displaystyle\int_0^2\max\{1,x\}\,dx=\int_0^11\,dx+\int_1^2x\,dx=1+\frac{3}{2}=\frac{5}{2}$。

（3）去根号：$\displaystyle\int_0^{2\pi}\sqrt{1-\cos 2x}\,dx=\int_0^{2\pi}\sqrt{2\sin^2x}\,dx=\sqrt{2}\int_0^{2\pi}\lvert\sin x\rvert dx=4\sqrt{2}$。

（4）分段函数的原函数：设 $f(x)=\begin{cases}x+1, & x<0,\\ e^{x}, & x\geq0,\end{cases}$ 求原函数 $F$。

$x<0$ 时 $F(x)=\frac{x^2}{2}+x+C_1$；$x\geq0$ 时 $F(x)=e^x+C_2$。由 $F$ 在 $x=0$ 连续：$C_1=1+C_2$。取 $C_1=0$，则 $C_2=-1$，故

\[F(x)=\begin{cases}\dfrac{x^2}{2}+x, & x<0,\\[2mm] e^{x}-1, & x\geq0.\end{cases}\]

验证 $F'(0)=f(0)=1$（左导数 $=1$，右导数 $=e^0=1$），说明拼接正确。

若题目还给出初值条件（如 $F(1)=0$），就用它把剩下的那个常数定下来；若题目问“$f$ 的所有原函数”，答案中要保留一个任意常数，并说明它在各段之间的关系。这类题评分很细，分段点处的连续性说明必须写出来。

\begin{callwarn}{易错点}
\begin{itemize}
\item $\sqrt{x^2}=\lvert x\rvert$ 而不是 $x$；这是"含根号积分"最大的陷阱。
\item 分段时区间\textbf{不重不漏}：分界点只能属于其中一段；写 $\int_0^4=\int_0^1+\int_1^4$，不要把 $x=1$ 算两次。
\item 分段函数的原函数：\textbf{不能用一个 $C$ 覆盖所有段}；只保证连续不一定保证可导，题目要求"原函数"时必须验证分段点可导。
\item 积分限含参数时，"分段点是否落在积分区间内"要分类讨论，不能默认落在里面。
\end{itemize}
\end{callwarn}

\begin{callexam}{真题考法}
本考点是\textbf{解答题的常客}：分段函数的原函数（含初值条件）、含绝对值的定积分、含 $\min/\max$ 的定积分都是高频。命题人常把分界点设在区间的"不好看"位置（如 $x=\frac{1}{2}$、$x=\ln 2$），要求把每段算准。

\end{callexam}

\subsubsection{定积分计算综合（递推公式、Wallis 公式类） $\bigstar$$\bigstar$$\bigstar$$\bigstar$}

\textbf{为什么要建立递推}：像 $\int_0^{\pi/2}\sin^nx\,dx$ 这样的积分，对每个具体的 $n$ 都能硬算，但要把 $n$ 当作一般参数就必须建立递推关系。做法固定：把 $\sin^nx$ 拆成 $\sin^{n-1}x\cdot\sin x$，用分部积分把次数降下来，得到 $I_n$ 与 $I_{n-2}$ 的关系，再给出初值 $I_0,I_1$，递推到底即可。这种“降次—递推”的思路在真题证明题中反复出现。

\textbf{Wallis（华里士）公式}：

\[\int_0^{\pi/2}\sin^nx\,dx=\int_0^{\pi/2}\cos^nx\,dx=\begin{cases}\dfrac{(n-1)!!}{n!!}\cdot\dfrac{\pi}{2}, & n\ \text{为正偶数},\\[3mm]\dfrac{(n-1)!!}{n!!}, & n\ \text{为正奇数}.\end{cases}\]

其中 $(2k)!!=2\cdot4\cdots2k$，$(2k-1)!!=1\cdot3\cdots(2k-1)$，并规定 $0!!=1$。常用数值备忘：

\[\int_0^{\pi/2}\sin^2x\,dx=\frac{\pi}{4},\quad \int_0^{\pi/2}\sin^3x\,dx=\frac{2}{3},\quad \int_0^{\pi/2}\sin^4x\,dx=\frac{3\pi}{16},\quad \int_0^{\pi/2}\sin^5x\,dx=\frac{8}{15}.\]

\begin{callkey}{Wallis 公式（必背）}
\[\int_0^{\pi/2}\sin^nx\,dx=\int_0^{\pi/2}\cos^nx\,dx=\begin{cases}\dfrac{(n-1)!!}{n!!}\cdot\dfrac{\pi}{2}, & n\ \text{为偶数},\\[2mm]\dfrac{(n-1)!!}{n!!}, & n\ \text{为奇数}.\end{cases}\]

\begin{itemize}
\item 区间必须先化为 $\left[0,\frac{\pi}{2}\right]$；$\int_0^{\pi}$ 乘 $2$，$\int_0^{2\pi}$ 乘 $4$。
\item 递推式 $I_n=\frac{n-1}{n}I_{n-2}$，初值 $I_0=\frac{\pi}{2}$，$I_1=1$。
\end{itemize}
\end{callkey}

\textbf{区间推广}：

\[\int_0^{\pi}\sin^nx\,dx=2\int_0^{\pi/2}\sin^nx\,dx,\qquad \int_0^{2\pi}\sin^nx\,dx=4\int_0^{\pi/2}\sin^nx\,dx\ (n\geq2).\]

\textbf{区间推广的用法}：凡是被积函数含 $\sin^nx$ 或 $\cos^nx$ 而区间是 $[0,\pi]$、$[0,2\pi]$、$[-\pi,\pi]$ 的，都先归到 $\int_0^{\pi/2}$ 上再用 Wallis 公式；若区间有平移（如 $\left[\frac{\pi}{4},\frac{3\pi}{4}\right]$），先用周期性或换元把它搬回 $[0,\pi]$。

\textbf{递推公式}：由分部积分得

\[I_n=\int_0^{\pi/2}\sin^nx\,dx=\frac{n-1}{n}I_{n-2}\quad(n\geq2),\qquad I_0=\frac{\pi}{2},\ I_1=1.\]

同类递推（都要会推）：

\[J_n=\int_0^{\pi/4}\tan^nx\,dx=\frac{1}{n-1}-J_{n-2}\quad(n\geq2),\]

\[K_n=\int\frac{dx}{\left(x^2+a^2\right)^n}=\frac{x}{2a^2(n-1)\left(x^2+a^2\right)^{n-1}}+\frac{2n-3}{2a^2(n-1)}K_{n-1}.\]

\textbf{典型例子}

\[\int_0^{2\pi}\sin^4x\,dx=4\cdot\frac{3\pi}{16}=\frac{3\pi}{4},\]

\[\int_0^{\pi/2}\sin^6x\,dx=\frac{5\cdot3\cdot1}{6\cdot4\cdot2}\cdot\frac{\pi}{2}=\frac{5\pi}{32}.\]

验证前者与直接降幂一致：$\sin^4x=\frac{3-4\cos2x+\cos4x}{8}$，在一个整周期上积分为 $\frac{3}{8}\cdot2\pi=\frac{3\pi}{4}$，吻合。

\textbf{一道递推证明的模板}：设 $I_n=\int_0^{\pi/2}\sin^nx\,dx\ (n\geq2)$，则

\[\begin{aligned}
I_n&=\int_0^{\pi/2}\sin^{n-1}x\,d(-\cos x)
=\left[-\sin^{n-1}x\cos x\right]_0^{\pi/2}+(n-1)\int_0^{\pi/2}\sin^{n-2}x\cos^2x\,dx\\
&=(n-1)\left(I_{n-2}-I_n\right),
\end{aligned}\]

移项即得 $I_n=\frac{n-1}{n}I_{n-2}$。这个推导要能闭卷写出来。

\begin{callwarn}{易错点}
\begin{itemize}
\item Wallis 公式的区间必须先把积分\textbf{化为 $[0,\frac{\pi}{2}]$}；$\int_0^{\pi}$、$\int_0^{2\pi}$ 上要乘 $2$ 或 $4$。
\item 偶数次与奇数次的公式容易记反：\textbf{偶数带 $\frac{\pi}{2}$，奇数不带}。
\item 递推公式要写清 $n$ 的起点与初值 $I_0,I_1$，否则递推不到底。
\item 递推题不要忘记先在答案中给出 $I_0,I_1$，再写一般项。
\end{itemize}
\end{callwarn}

\begin{callexam}{真题考法}
本考点在\textbf{解答题}中常以"证明递推公式并求值"的形式出现，也可能作为大题中一步（如计算 $\int_0^{\pi/2}\sin^nx\,dx$ 型积分或面积/体积中含高次三角函数的积分）。命题角度固定：分部积分 + 移项 + 递推。

\end{callexam}

\section{反常积分}

反常积分是把定积分推广到"无穷区间"与"无界函数"两种情形，是数二的高频考点（尤其是敛散性判别）。这一节的题量与分值都不小，而且“会算”与“会判”是两种不同的能力：计算题按定义取极限即可，判敛题则要在奇点附近做等价化简。两者都要求先\textbf{找齐奇点}，这一步绝对不能省。

\subsection{反常积分的概念与计算}

\subsubsection{反常（广义）积分的概念与计算 $\bigstar$$\bigstar$$\bigstar$$\bigstar$}

\textbf{为什么需要“广义”积分}：定积分要求积分区间是\textbf{有限闭区间}且被积函数\textbf{有界}。当区间无限（如 $\int_1^{+\infty}$）或被积函数在某点附近无界（如 $\int_0^1\frac{dx}{\sqrt{x}}$）时，这两个前提被破坏，原来的定义失效。处理办法是用\textbf{极限}重新定义：先在上限为 $b$ 的有限区间上积分，再令 $b\to+\infty$；或先在去掉瑕点的小区间上积分，再令其长度 $\varepsilon\to0^+$。极限存在且有限就称收敛，否则称发散。

\textbf{两类反常积分}：

\textbf{（1）无穷限反常积分}

\[\int_a^{+\infty}f(x)\,dx=\lim_{b\to+\infty}\int_a^bf(x)\,dx,\]

\[\int_{-\infty}^{b}f(x)\,dx=\lim_{a\to-\infty}\int_a^bf(x)\,dx,\qquad \int_{-\infty}^{+\infty}f(x)\,dx=\int_{-\infty}^{c}f+\int_{c}^{+\infty}f.\]

右端\textbf{每个}极限都存在且有限时，才称反常积分\textbf{收敛}，否则称\textbf{发散}。

\textbf{（2）无界函数反常积分（瑕积分）}：若 $f$ 在 $a$ 的右邻域内无界（$a$ 称为\textbf{瑕点}），则

\[\int_a^bf(x)\,dx=\lim_{\varepsilon\to0^+}\int_{a+\varepsilon}^{b}f(x)\,dx.\]

瑕点在上限时类似；\textbf{瑕点在区间内部时，必须以瑕点为界拆成两个瑕积分}，两边都收敛才收敛。

\textbf{p 积分（基准，必须背）}：

\begingroup\small\setlength{\tabcolsep}{4.0pt}
\begin{longtable}{@{}L{0.3095\linewidth}L{0.3095\linewidth}L{0.3095\linewidth}@{}}
\toprule
\textbf{积分} & \textbf{收敛} & \textbf{发散} \\
\midrule
\endhead
$\displaystyle\int_1^{+\infty}\frac{dx}{x^p}$ & $p>1$ & $p\leq1$ \\
$\displaystyle\int_0^{1}\frac{dx}{x^p}$ & $p<1$ & $p\geq1$ \\
$\displaystyle\int_2^{+\infty}\frac{dx}{x\ln^qx}$ & $q>1$ & $q\leq1$ \\
\bottomrule
\end{longtable}
\endgroup

直观记忆：\textbf{无穷限时函数趋于零越快越容易收敛（$p$ 越大越好）；瑕积分时函数趋于无穷越慢越容易收敛（$p$ 越小越好）}，两类判据方向相反。

\begin{callkey}{反常积分 p 判据（必背，两类方向相反）}
\begin{itemize}
\item 无穷限：$\displaystyle\int_1^{+\infty}\frac{dx}{x^p}$ 当 $p>1$ 收敛，当 $p\leq1$ 发散。
\item 瑕积分（瑕点 $x=a$）：$\displaystyle\int_a^b\frac{dx}{(x-a)^p}$ 当 $p<1$ 收敛，当 $p\geq1$ 发散。
\item 对数型：$\displaystyle\int_2^{+\infty}\frac{dx}{x\ln^qx}$ 当 $q>1$ 收敛。
\item 操作顺序：先找齐所有奇点（含 $+\infty,-\infty$ 与全部瑕点），拆成单奇点积分，再在奇点附近等价化简后与基准比较。
\end{itemize}
\end{callkey}

\textbf{计算步骤}：$\circled{1}$ 找出所有奇点（$+\infty,-\infty$ 与所有瑕点）；$\circled{2}$ 用奇点把积分拆成若干"只含一个奇点"的积分；$\circled{3}$ 对每一段先求原函数、再取极限；$\circled{4}$ 只要有一段发散，原积分发散。写解答时建议显式写出极限过程（用 $\lim\limits_{b\to+\infty}$、$\lim\limits_{\varepsilon\to0^+}$），并注明“该极限存在且有限，故收敛”。直接对含瑕点的积分套原函数代上下限，即使答案对，过程也会被扣分。

\textbf{典型例子}

\[\int_0^{+\infty}e^{-x}dx=\lim_{b\to+\infty}\left[-e^{-x}\right]_0^b=1,\qquad \int_0^{+\infty}xe^{-x}dx=1,\]

\[\int_{-\infty}^{+\infty}\frac{dx}{1+x^2}=\left[\arctan x\right]_{-\infty}^{+\infty}=\pi,\qquad \int_0^1\ln x\,dx=\left[x\ln x-x\right]_{0^+}^{1}=-1,\]

\[\int_1^{+\infty}\frac{\ln x}{x^2}dx=\left[-\frac{\ln x}{x}-\frac{1}{x}\right]_1^{+\infty}=0-(-1)=1.\]

\textbf{瑕积分的典型计算}：$\displaystyle\int_0^1\frac{dx}{\sqrt{x}}=2\lim_{\varepsilon\to0^+}\left[\sqrt{x}\right]_\varepsilon^1=2$（收敛）；而 $\displaystyle\int_0^1\frac{dx}{x}=+\infty$（发散）。二者只差在 $p$ 是小于 $1$ 还是等于 $1$，可见 p 判据必须记准。

\begin{callwarn}{易错点}
\begin{itemize}
\item 计算前\textbf{先找瑕点}：分母为零的点、$\ln x$ 在 $x\to0^+$ 处、$\tan x$ 的无穷间断点等都可能是瑕点。
\item 瑕点在内部必须\textbf{拆开}，绝对不能在含瑕点的区间上直接用牛顿-莱布尼茨公式。
\item $\int_{-\infty}^{+\infty}f$ 与 $\int_1^{+\infty}\frac{dx}{x^p}$ 这类"多奇点"积分，只要一段发散，整体就发散。
\item 用极限求值时注意 $x\ln x\to0$（$x\to0^+$）、$\frac{\ln x}{x}\to0$（$x\to+\infty$）这类端点极限。
\end{itemize}
\end{callwarn}

\begin{callexam}{真题考法}
\textbf{填空/解答}常见：直接按定义计算反常积分（含 $\int_0^{+\infty}\frac{dx}{x^2+a^2}$、$\int_0^1\frac{dx}{\sqrt{x(1-x)}}$ 型），或讨论含参数的反常积分何时收敛。容易被忽略的是"瑕点在内部"与"无穷限 $+$ 瑕点同时出现"的情形。

\end{callexam}

\subsection{敛散性判别}

\subsubsection{反常积分敛散性的判别（比较判别法、p 积分） $\bigstar$$\bigstar$$\bigstar$$\bigstar$}

考纲对比较判别法的要求是"了解（会判断即可）"，但考试中它比"算出来"更高效。

\textbf{比较判别法}（要求被积函数\textbf{非负}）：

\begin{itemize}
\item 若 $0\leq f(x)\leq g(x)$ 且 $\int_a^{+\infty}g$ 收敛，则 $\int_a^{+\infty}f$ 收敛；
\item 若 $0\leq g(x)\leq f(x)$ 且 $\int_a^{+\infty}g$ 发散，则 $\int_a^{+\infty}f$ 发散。
\end{itemize}
口诀：\textbf{大收则小收，小发则大发}。瑕积分的比较判别法形式完全相同（只需把 $x\to+\infty$ 换成 $x\to a^+$）。

\textbf{极限比较判别法}：设 $f,g\geq0$，$\lim\limits_{x\to+\infty}\frac{f(x)}{g(x)}=L$（瑕点时 $x\to a^+$），则

\begin{itemize}
\item $0<L<+\infty$：$\int f$ 与 $\int g$ \textbf{同敛散}；
\item $L=0$：$\int g$ 收敛 $\Rightarrow$ $\int f$ 收敛（$f$ 是比 $g$ 高阶的无穷小）；
\item $L=+\infty$：$\int g$ 发散 $\Rightarrow$ $\int f$ 发散（$f$ 是比 $g$ 低阶的无穷小）。
\end{itemize}
\textbf{比较判别法的实质}：把复杂的被积函数在奇点附近“换成形状相似的基准函数” $\frac{1}{x^p}$ 或 $\frac{1}{(x-a)^p}$，再用 p 积分的结论。所以判敛题的核心动作不是“算积分”，而是\textbf{在奇点附近做等价替换、抓大头、提非零因子}，把被积函数化简到基准形式。考纲只要求“了解”比较判别法，即会判断即可，不必背定理的完整证明。

\textbf{判敛四步法}：

\begin{enumerate}
\item \textbf{找齐所有奇点}（$+\infty,-\infty$ 与所有瑕点），拆成只含单个奇点的积分；
\item 在奇点附近把 $f$ 用\textbf{等价无穷小/等价无穷大/抓大头/提非零因子}化简；
\item 化为 $\frac{C}{x^p}$ 或 $\frac{C}{(x-a)^p}$ 的基准形式；
\item 用 p 积分判据下结论。
\end{enumerate}
\textbf{幂对数基准（进阶，考试常用）}：

\begingroup\small\setlength{\tabcolsep}{4.0pt}
\begin{longtable}{@{}L{0.3479\linewidth}L{0.5978\linewidth}@{}}
\toprule
\textbf{形式} & \textbf{敛散性} \\
\midrule
\endhead
$\displaystyle\int_2^{+\infty}\frac{dx}{x\ln^qx}$ & $q>1$ 收敛，$q\leq1$ 发散 \\
$\displaystyle\int_0^{1/2}\frac{dx}{x\lvert\ln x\rvert^q}$ & $q>1$ 收敛，$q\leq1$ 发散 \\
$\displaystyle\int_2^{+\infty}\frac{\ln^ax}{x^p}dx$ & 与 $\displaystyle\int_2^{+\infty}\frac{dx}{x^p}$ 同敛散（对数幂次 $a$ 不影响） \\
$\displaystyle\int_2^{+\infty}\frac{dx}{\ln^ax}$ & 发散（对数衰减太慢） \\
$\displaystyle\int_0^{1/2}\frac{dx}{\lvert\ln x\rvert^a}$ & 收敛（对数增长太慢） \\
\bottomrule
\end{longtable}
\endgroup

\textbf{绝对收敛定理}：$\int_a^{+\infty}\lvert f(x)\rvert dx$ 收敛 $\Rightarrow$ $\int_a^{+\infty}f(x)dx$ 收敛。被积函数变号时，先考察 $\lvert f\rvert$；若 $\int f$ 收敛而 $\int\lvert f\rvert$ 发散，称为\textbf{条件收敛}。

\textbf{典型例子}

（1）两端都不好判断的积分：$\displaystyle\int_0^{+\infty}\frac{dx}{\sqrt{x+x^3}}$。$x\to0^+$ 时 $\sqrt{x+x^3}\sim\sqrt{x}$，$\int_0^1x^{-1/2}dx$ 收敛（$p=\frac{1}{2}<1$）；$x\to+\infty$ 时 $\sqrt{x+x^3}\sim x^{3/2}$，$\int_1^{+\infty}x^{-3/2}dx$ 收敛（$p=\frac{3}{2}>1$）。故原积分收敛。

（2）$\displaystyle\int_0^{+\infty}\frac{dx}{e^x-1}$。$x\to0^+$ 时 $\frac{1}{e^x-1}\sim\frac{1}{x}$，$\int_0^1\frac{dx}{x}$ 发散，故原积分\textbf{发散}。

（3）已知收敛求参数：$\displaystyle\int_0^{+\infty}\frac{\arctan x\cdot\ln(1+x)}{x^p}dx$ 收敛，求 $p$。

\begin{itemize}
\item $x\to0^+$：$\arctan x\sim x$，$\ln(1+x)\sim x$，被积函数 $\sim x^{2-p}$，$\int_0^1x^{2-p}dx$ 收敛 $\iff 2-p>-1 \iff p<3$。
\item $x\to+\infty$：$\arctan x\to\frac{\pi}{2}$，$\ln(1+x)\sim\ln x$，被积函数 $\sim\frac{\pi}{2}\cdot\frac{\ln x}{x^p}$，由幂对数基准知收敛 $\iff p>1$。
\end{itemize}
综合得 $1<p<3$。

\textbf{两个易混的等价替换}：$\ln(1+x)\sim x$ 只在 $x\to0$ 时成立，$x\to+\infty$ 时应写成 $\ln(1+x)\sim\ln x$；同理 $e^{x}-1\sim x$ 只在 $x\to0$ 用，$x\to+\infty$ 时 $e^{x}-1\sim e^{x}$。判敛题里“等价替换的方向取决于奇点在哪一端”，这是最容易失分的地方。

\begin{callwarn}{易错点}
\begin{itemize}
\item 比较判别法要求被积函数\textbf{不变号}；$f$ 变号时不能直接比较，要先看 $\lvert f\rvert$（绝对收敛）。
\item 等价替换只在\textbf{奇点附近}使用，不能在整个积分区间上替换。
\item 多奇点不拆区间直接比较，结论一定不可靠。
\item 无穷限与瑕积分的 p 判据方向相反，是最容易记混的一点。
\end{itemize}
\end{callwarn}

\begin{callexam}{真题考法}
\textbf{选择 / 解答}高频。选择题问"下列积分收敛的是"，解答题常问"已知收敛，求参数取值范围"。标准做法是"拆成单奇点 $\to$ 在奇点附近等价化简 $\to$ 用 p 判据"，答题时要把两端（$x\to0^+$ 与 $x\to+\infty$，或瑕点与无穷远）分开讨论。

\end{callexam}

\section{定积分的应用与积分证明}

定积分的应用统一用\textbf{微元法}：在区间上任取一点 $x$，让它"波动" $dx$，取一小块近似（写出微元），再把微元沿区间积分。几何量包括面积、体积、弧长、侧面积，物理量包括功、压力、引力，此外还有质心形心与平均值。应用题的难点不在积分本身，而在“把几何量或物理量写成微元”这一步。写微元有三条经验：$\circled{1}$ 先画图、定坐标系，让边界表达式最简单；$\circled{2}$ 用“以直代曲、以不变代变”构造近似；$\circled{3}$ 用对称性减少计算量。积分只是最后一步。

\subsection{几何应用}

\subsubsection{定积分的几何应用：平面图形面积 $\bigstar$$\bigstar$$\bigstar$$\bigstar$}

\textbf{直角坐标}：

\begin{itemize}
\item $X$ 型区域（上下边界）：$S=\displaystyle\int_a^b\left\lvert f_{\text{上}}(x)-f_{\text{下}}(x)\right\rvert dx$；
\item $Y$ 型区域（左右边界）：$S=\displaystyle\int_c^d\left\lvert\varphi_{\text{右}}(y)-\varphi_{\text{左}}(y)\right\rvert dy$。
\end{itemize}
\textbf{极坐标}（曲边扇形）：$S=\displaystyle\frac{1}{2}\int_{\alpha}^{\beta}r^2(\theta)\,d\theta$。

\textbf{参数方程}（曲线逆时针围成）：$x=x(t),\ y=y(t)$，则

\[S=\int_{t_1}^{t_2}y(t)x'(t)\,dt.\]

\textbf{用定积分求面积的前提}：面积是几何量，非负，所以公式里必须带绝对值，或严格保证“上减下”的顺序正确。平面区域要么“竖着切”（$X$ 型，对 $x$ 积分），要么“横着切”（$Y$ 型，对 $y$ 积分）；切法选对，计算量可以相差很多。若边界表示成 $x=g(y)$ 更简单，就改用 $Y$ 型；若区域是圆、圆弧、花瓣，一律优先极坐标。

\textbf{解题套路}：$\circled{1}$ 画图，求出所有交点并定出积分区间；$\circled{2}$ 判断用 $X$ 型还是 $Y$ 型（选被积函数更简单的那种，必要时反解 $x=g(y)$）；$\circled{3}$ 含 $x^2+y^2$、圆弧、花瓣的一律优先极坐标；$\circled{4}$ 边界是参数方程的直接用参数公式；$\circled{5}$ 分段求交点后分段积分。

\textbf{典型例子}

（1）$y=x^2$ 与 $y^2=x$ 围成区域：交点 $(0,0),(1,1)$，

\[S=\int_0^1\left(\sqrt{x}-x^2\right)dx=\frac{2}{3}-\frac{1}{3}=\frac{1}{3}.\]

（2）心形线 $r=a(1+\cos\theta)$ 所围面积：

\[S=\frac{1}{2}\int_0^{2\pi}a^2(1+\cos\theta)^2d\theta=\frac{a^2}{2}\left(2\pi+0+\pi\right)=\frac{3\pi a^2}{2}.\]

（3）双纽线 $r^2=a^2\cos 2\theta$ 所围面积：由对称性 $S=4\cdot\frac{1}{2}\int_0^{\pi/4}a^2\cos2\theta\,d\theta=a^2$。

（4）椭圆 $x=a\cos t,\ y=b\sin t$：$S=\pi ab$。第（4）问的推导值得记住：$S=\int_0^{2\pi}b\sin t\cdot(-a\sin t)\,dt=ab\int_0^{2\pi}\sin^2t\,dt=\pi ab$，其中用到 $\int_0^{2\pi}\sin^2t\,dt=\pi$。这是参数方程求面积的典型演算，也是椭圆面积公式的来源。

（5）摆线一拱 $\begin{cases}x=a(t-\sin t),\\ y=a(1-\cos t)\end{cases}(0\leq t\leq2\pi)$ 与 $x$ 轴围成的面积：

\[S=\int_0^{2\pi}a(1-\cos t)\cdot a(1-\cos t)\,dt=a^2\cdot3\pi=3\pi a^2.\]

\textbf{定积分限的方法}：把两条曲线联立解交点，交点的横坐标就是分段点；极坐标曲线则令 $r=0$ 解出角度，再用对称性只算一部分并乘以倍数。凡区域关于 $x$ 轴或 $y$ 轴对称，务必先用对称性缩小积分区间，往往能省掉一半以上的计算。

\begin{callwarn}{易错点}
\begin{itemize}
\item 极坐标面积公式的系数是 $\frac{1}{2}$，与直角坐标面积公式不同，最容易漏。
\item 上下边界交叉时必须\textbf{分段}，不能一路积到底；被积函数交叉处取绝对值。
\item 参数方程面积公式要求曲线\textbf{逆时针}方向，方向反了要加负号。
\item 极坐标中求"公共部分面积"要先解 $r_1=r_2$ 定出角度分界点。
\end{itemize}
\end{callwarn}

\begin{callexam}{真题考法}
\textbf{必考解答题之一}。常见设问：由两条曲线围成的面积、极坐标曲线围成的面积、由参数方程给出的曲线围成的面积、求区域面积再求其旋转体体积。陷阱多半在"交点没找全"与"上下边界判断反了"。

\end{callexam}

\subsubsection{定积分的几何应用：旋转体体积与平行截面面积已知的立体体积 $\bigstar$$\bigstar$$\bigstar$$\bigstar$}

\textbf{（1）旋转体体积}

\begin{itemize}
\item \textbf{绕 $x$ 轴（圆盘法）}：$V=\displaystyle\pi\int_a^by^2\,dx$；
\item \textbf{绕 $x$ 轴、上下两条曲线（垫圈法）}：$V=\displaystyle\pi\int_a^b\left[y_2^2(x)-y_1^2(x)\right]dx$；
\item \textbf{绕 $y$ 轴（柱壳法，保留 $y=f(x)$ 形式）}：$V=\displaystyle2\pi\int_a^b\lvert x\rvert\,\lvert f(x)\rvert\,dx$；
\item \textbf{绕 $y$ 轴（圆盘法，已解出 $x=g(y)$）}：$V=\displaystyle\pi\int_c^dx^2\,dy$；
\item \textbf{绕直线 $x=k$}：$V=\displaystyle2\pi\int_a^b\lvert x-k\rvert\lvert f(x)\rvert dx$；
\item \textbf{绕直线 $y=k$}：$V=\displaystyle\pi\int_a^b\left[f(x)-k\right]^2dx$；
\item \textbf{参数方程} $x=x(t),y=y(t)$ 绕 $x$ 轴：$V=\displaystyle\pi\int_{t_1}^{t_2}y^2(t)x'(t)\,dt$。
\end{itemize}
\textbf{（2）平行截面面积已知的立体体积}

\[V=\int_a^bA(x)\,dx,\]

\textbf{平行截面法的关键是“选对截面方向”}：截面垂直于 $x$ 轴时用 $A(x)$，垂直于 $y$ 轴时用 $B(y)$；方向选得好，截面面积就简单。已知截面是三角形、半圆、正方形时，先写出边长与 $x$ 的关系，再套对应的面积公式。

其中 $A(x)$ 是过点 $x$ 且垂直于 $x$ 轴的截面面积。

\textbf{典型例子}

（1）$y=x^2\ (0\leq x\leq1)$ 绕 $x$ 轴：$V=\pi\int_0^1x^4dx=\frac{\pi}{5}$。

（2）同曲线绕 $y$ 轴（柱壳法）：$V=2\pi\int_0^1x\cdot x^2dx=\frac{\pi}{2}$；用圆盘法验证：$x=\sqrt{y}$，$V=\pi\int_0^1y\,dy=\frac{\pi}{2}$，一致。

（3）球体：$y=\sqrt{R^2-x^2}$ 绕 $x$ 轴，$V=\pi\int_{-R}^{R}(R^2-x^2)dx=\frac{4}{3}\pi R^3$。

（4）平行截面：以圆 $x^2+y^2\leq R^2$ 为底，垂直于 $x$ 轴的截面都是等边三角形。截面边长 $=\ 2\sqrt{R^2-x^2}$，故 $A(x)=\frac{\sqrt{3}}{4}\left(2\sqrt{R^2-x^2}\right)^2=\sqrt{3}\left(R^2-x^2\right)$，

\[V=\sqrt{3}\int_{-R}^{R}\left(R^2-x^2\right)dx=\frac{4\sqrt{3}}{3}R^3.\]

\textbf{两种方法的选择原则}：绕坐标轴旋转时，若曲线以 $y=f(x)$ 给出而反解 $x=g(y)$ 困难，用柱壳法；若能轻松反解且区域是 $Y$ 型，用圆盘法。绕非坐标轴旋转时先把“到轴的距离”写清楚：绕 $x=k$ 时半径是 $\lvert x-k\rvert$，绕 $y=k$ 时半径是 $\lvert y-k\rvert$。

\begin{callwarn}{易错点}
\begin{itemize}
\item 垫圈法里"外半径的平方减内半径的平方"，顺序反了会得到负体积。
\item 绕\textbf{非坐标轴}旋转时，微元半径是\textbf{点到旋转轴的距离}，不是到原点的距离。
\item 曲线与旋转轴相交（如 $y=\sin x$ 与 $x$ 轴）要在交点处\textbf{分段}积分。
\item 平行截面法不要漏了截面面积函数中的常数（如等边三角形面积系数 $\frac{\sqrt{3}}{4}$）。
\end{itemize}
\end{callwarn}

\begin{callexam}{真题考法}
\textbf{必考解答题之一}，且常与面积计算连在一起命题：先求交点与面积，再求绕某直线旋转所得体积。命题人偏爱"绕 $y$ 轴用柱壳法更快"的情形，以及"绕 $y=k$ 或 $x=k$ 这类非坐标轴"的平移情形。

\end{callexam}

\subsubsection{定积分的几何应用：平面曲线弧长、旋转体侧面积 $\bigstar$$\bigstar$$\bigstar$}

\textbf{弧微分}：$ds=\sqrt{dx^2+dy^2}$。弧微分是弧长与侧面积共同的基础，记牢三个版本：直角坐标 $ds=\sqrt{1+y'^2}\,dx$，参数方程 $ds=\sqrt{x'^2+y'^2}\,dt$，极坐标 $ds=\sqrt{r^2+r'^2}\,d\theta$。三者可以互相推导，不必分别死记。

\textbf{弧长}：

\[L=\int_a^b\sqrt{1+\left[y'(x)\right]^2}\,dx\quad(\text{直角坐标}),\]

\[L=\int_{\alpha}^{\beta}\sqrt{\left[x'(t)\right]^2+\left[y'(t)\right]^2}\,dt\quad(\text{参数方程}),\]

\[L=\int_{\alpha}^{\beta}\sqrt{r^2(\theta)+\left[r'(\theta)\right]^2}\,d\theta\quad(\text{极坐标}).\]

\textbf{旋转体侧面积}（绕 $x$ 轴）：

\[S=2\pi\int_a^b\lvert y\rvert\sqrt{1+\left[y'(x)\right]^2}\,dx=2\pi\int_a^b\lvert y\rvert\,ds,\]

参数方程形式：$S=2\pi\displaystyle\int_{\alpha}^{\beta}\lvert y(t)\rvert\sqrt{\left[x'(t)\right]^2+\left[y'(t)\right]^2}\,dt$；

极坐标形式（绕极轴）：$S=2\pi\displaystyle\int_{\alpha}^{\beta}\lvert r\sin\theta\rvert\sqrt{r^2+r'^2}\,d\theta$。

绕 $y$ 轴旋转时：$S=2\pi\displaystyle\int_a^b\lvert x\rvert\sqrt{1+\left[y'(x)\right]^2}\,dx$。

\textbf{常用特殊值}（可作为验算）：

\begin{itemize}
\item 圆 $x^2+y^2=R^2$ 的周长 $2\pi R$；球面面积 $4\pi R^2$。
\item 摆线一拱的弧长 $8a$；星形线 $x^{2/3}+y^{2/3}=a^{2/3}$ 的全长 $6a$；心形线 $r=a(1+\cos\theta)$ 的全长 $8a$。
\end{itemize}
\textbf{典型例子}

（1）星形线全长：$x=a\cos^3t,\ y=a\sin^3t$，

\[ds=\sqrt{9a^2\cos^4t\sin^2t+9a^2\sin^4t\cos^2t}\,dt=3a\lvert\sin t\cos t\rvert dt,\]

\[L=4\int_0^{\pi/2}3a\sin t\cos t\,dt=12a\cdot\frac{1}{2}=6a.\]

（2）球面面积：$S=2\pi\displaystyle\int_{-R}^{R}\sqrt{R^2-x^2}\cdot\sqrt{1+\frac{x^2}{R^2-x^2}}\,dx=2\pi\int_{-R}^{R}R\,dx=4\pi R^2$。

（3）心形线全长：$r^2+r'^2=a^2\left[(1+\cos\theta)^2+\sin^2\theta\right]=4a^2\cos^2\frac{\theta}{2}$，

\[L=\int_0^{2\pi}2a\left\lvert\cos\frac{\theta}{2}\right\rvert d\theta=4a\int_0^{\pi}\lvert\cos u\rvert du=8a.\]

\textbf{弧长与侧面积的配合}：算旋转体侧面积时先把 $ds$ 算出来（这一步与弧长共用），再乘 $2\pi\lvert y\rvert$ 积分。这样一道题可以同时得到弧长与侧面积两个结果，也是命题人常用的“一题两问”设计。

\begin{callwarn}{易错点}
\begin{itemize}
\item \textbf{弧长公式里的根号不能丢}，侧面积公式里\textbf{既要 $\lvert y\rvert$ 又要弧微分}，两者缺一不可。
\item 极坐标侧面积公式是 $2\pi\lvert r\sin\theta\rvert\sqrt{r^2+r'^2}$，不要漏 $\lvert r\sin\theta\rvert$。
\item 侧面积与"旋转一周后所得曲面的面积"是同一件事；不要与旋转体体积的公式混用（一个含 $ds$，一个含 $dx$ 与平方）。
\item 曲线有尖点、拐折（如 $\lvert x\rvert$ 型）时分段求弧长。
\end{itemize}
\end{callwarn}

\begin{callexam}{真题考法}
多为\textbf{填空或解答中的一步}：求摆线一拱弧长、星形线周长、球面面积、圆锥侧面面积等。命题人常把弧长与侧面积放进同一题（先算 $ds$，再算 $S$），也会与参数方程联考。

\end{callexam}

\subsection{物理应用与平均值}

\subsubsection{定积分的物理应用：功、引力、压力 $\bigstar$$\bigstar$$\bigstar$}

\textbf{物理应用的通法}：物理量的计算都遵循“微元法”——在区间上任取一点，让变量“波动” $dx$，把这一小段上的物理量近似看作常量（这就是\textbf{以直代曲、以不变代变}），写出微元表达式 $dW$、$dF$，再沿区间积分。判断微元写得对不对，可以用量纲检查：功的量纲是“力 $\times$ 长度”，压力的量纲是“力”。

\begin{callkey}{微元法四步（所有应用题的通用流程）}
\begin{enumerate}
\item 选变量与坐标系（常取深度或位移为变量），让边界表达式最简单。
\item 取微元：让变量波动 $dx$，把这一小段上的量近似看作常量。
\item 写出微元表达式：$dW=F(x)dx$、$dF=\rho gh\,L(h)dh$、$dV=A(x)dx$ 等。
\item 确定积分限并积分；最后用对称性或几何意义检验结果是否合理。
\end{enumerate}
\end{callkey}

\textbf{微元法四步}：$\circled{1}$ 建立坐标系、选取积分变量（常取深度或位移为变量）；$\circled{2}$ 取任意微元并写出微元表达式；$\circled{3}$ 统一量纲核对；$\circled{4}$ 在相应区间上积分。

\textbf{（1）变力做功}

\[W=\int_a^bF(x)\,dx.\]

抽水（或提升物体）做功：设水深 $x$ 处厚度为 $dx$ 的一层水的截面积为 $A(x)$，则

\[dW=\rho g\,x\,A(x)\,dx,\qquad W=\rho g\int_{x_1}^{x_2}xA(x)\,dx,\]

其中 $x$ 是\textbf{该微元到出口（或提升目标高度）的距离}，不是"水的深度"本身。

弹簧做功：$F=kx$，$W=\displaystyle\int_{x_1}^{x_2}kx\,dx$。

\textbf{（2）液体静压力}

压强 $P=\rho gh$（$h$ 为液面下深度），压力微元 $dF=P\,dA=\rho gh\,L(h)\,dh$，故

\[F=\rho g\int_{h_1}^{h_2}h\,L(h)\,dh,\]

其中 $L(h)$ 是深度 $h$ 处平板的\textbf{宽度}。

\textbf{压力的两种常见题型}：$\circled{1}$ 竖直平板（闸门）——宽度 $L(h)$ 随深度变化，直接积分；$\circled{2}$ 曲面（容器侧壁）——先写出侧面积微元，再对压力分量积分。数二以竖直平板为主，注意“水面距顶端多少米”这句话决定了积分的下限。

\textbf{（3）引力}

连续质量分布对质点的引力是\textbf{向量}，必须分解分量后积分。以细杆为例：杆长 $L$、质量 $M$、线密度 $\rho=\frac{M}{L}$，质点 $m$ 位于杆一端垂线上距离 $a$ 处，取杆上 $[x,x+dx]$ 一小段，其质量为 $\rho\,dx$，到质点的距离为 $\sqrt{a^2+x^2}$，由对称性只需沿垂直方向的分量：

\[dF_{\perp}=\frac{Gm\rho\,dx}{a^2+x^2}\cdot\frac{a}{\sqrt{a^2+x^2}},\qquad F=\int_0^L\frac{Gm\rho a}{\left(a^2+x^2\right)^{3/2}}dx=\frac{Gm\rho L}{a\sqrt{a^2+L^2}}=\frac{GMm}{a\sqrt{a^2+L^2}}.\]

\textbf{典型例子}

（1）弹簧劲度系数 $k=200$ N/m，由原长 $0.5$ m 拉长到 $0.8$ m：

\[W=\int_0^{0.3}200x\,dx=100\times0.09=9\ \text{J}.\]

（2）变力 $F(x)=x^2+1$（N）沿 $x$ 轴把物体从 $x=1$ 移到 $x=3$：

\[W=\int_1^3\left(x^2+1\right)dx=\frac{26}{3}+2=\frac{32}{3}\ \text{J}.\]

（3）抽水：半径 $R=2$ m 的半球形容器（开口向上）盛满水，抽到容器口。以容器口平面为原点、向下为 $x$ 轴，深度 $x$ 处截面圆半径为 $\sqrt{R^2-x^2}$，

\[W=\rho g\pi\int_0^{R}x\left(R^2-x^2\right)dx=\frac{\rho g\pi R^4}{4}=1000\times9.8\times\pi\times4\approx1.23\times10^{5}\ \text{J}.\]

（4）压力：矩形闸门宽 $3$ m、高 $4$ m，闸门顶端在水面下 $1$ m（即深度从 $1$ m 到 $5$ m），

\[F=\rho g\int_1^5h\times3\,dh=3\rho g\cdot\frac{25-1}{2}=36\rho g\approx3.53\times10^{5}\ \text{N}.\]

（5）引力：细杆长 $L$、质量 $M$，质点 $m$ 在杆一端延长线上距该端 $a$ 处，

\[F=\int_a^{a+L}\frac{GmM/L}{x^2}dx=\frac{GMm}{L}\left(\frac{1}{a}-\frac{1}{a+L}\right)=\frac{GMm}{a(a+L)}.\]

\textbf{常见容器的截面半径}（拿来就用）：半球形容器（口朝上）深度 $x$ 处截面半径为 $\sqrt{R^2-x^2}$；圆锥形容器（顶点朝下、底半径 $R$、高 $H$）深度 $x$ 处截面半径为 $R\left(1-\frac{x}{H}\right)$；圆柱形容器的截面半径恒为 $R$。写微元时直接套用，把时间留给积分本身。

\begin{callwarn}{易错点}
\begin{itemize}
\item 抽水做功中被提升的\textbf{高度}是"微元到出口的距离"，若容器不满，出口在水面处，高度要从水面量起。
\item 压力题必须先写出\textbf{宽度函数} $L(h)$，宽度随深度变化（梯形、圆形闸门）是常见陷阱。
\item 引力必须\textbf{分解分量}，直接用 $F=G\frac{m_1m_2}{r^2}$ 会算错；对称性可以消去某些分量。
\item 单位统一（N、J、m、kg），$\rho_{\text{水}}=1000\ \text{kg}/\text{m}^3$，$g=9.8\ \text{m/s}^2$；题目给 $g=10$ 就用 $10$。
\end{itemize}
\end{callwarn}

\begin{callexam}{真题考法}
\textbf{解答题常客}，尤其"变力做功"与"水压力"几乎每隔几年就会出现一种。命题形式：给出容器形状（半球、圆锥、梯形闸门）求抽水做功或求侧壁总压力；也可能给变力函数求功。技巧是把坐标系原点选在"液面"或"出口"上，让微元表达式最简。

\end{callexam}

\subsubsection{质心、形心与函数平均值 $\bigstar$$\bigstar$$\bigstar$}

\textbf{质心与形心的区别与联系}：质心的定义是“质量分布的中心”，需要对密度加权；形心的定义是“几何形状的中心”，只与区域的形状有关。密度均匀时二者重合，所以均匀薄片的质心可以直接用面积代替质量来计算。求质心的一条重要捷径是\textbf{对称性}：只要图形（或密度分布）关于某条直线对称，质心就在这条直线上，对应坐标可以直接写出，只需再算另一个坐标。

\textbf{质心与形心}：密度均匀的物体的质心就是\textbf{形心}（几何中心）。若密度不均匀，则用密度加权。

\textbf{平面薄片（面密度 $\rho(x,y)$）}：

\[\bar{x}=\frac{\displaystyle\iint_Dx\rho(x,y)\,d\sigma}{\displaystyle\iint_D\rho(x,y)\,d\sigma},\qquad \bar{y}=\frac{\displaystyle\iint_Dy\rho(x,y)\,d\sigma}{\displaystyle\iint_D\rho(x,y)\,d\sigma}.\]

密度均匀（$\rho\equiv$ 常数）时用面积 $A$ 代替质量：

\[\bar{x}=\frac{1}{A}\int_a^bx\left[y_2(x)-y_1(x)\right]dx,\qquad \bar{y}=\frac{1}{2A}\int_a^b\left[y_2^2(x)-y_1^2(x)\right]dx,\qquad A=\int_a^b\left[y_2-y_1\right]dx.\]

\textbf{平面曲线（均匀，弧长为 $L$）}：

\[\bar{x}=\frac{1}{L}\int_Lx\,ds,\qquad \bar{y}=\frac{1}{L}\int_Ly\,ds,\qquad L=\int_Lds.\]

\textbf{对称性}：若区域关于直线 $x=a$ 对称，则 $\bar{x}=a$；关于 $y=b$ 对称，则 $\bar{y}=b$；有两条对称轴时交点即为质心。

\textbf{组合体}：$\bar{x}=\frac{\sum A_ix_i}{\sum A_i}$，$\bar{y}=\frac{\sum A_iy_i}{\sum A_i}$（把复杂图形拆成矩形、三角形等标准图形）。

\textbf{古尔丁定理（旋转体体积的另一种算法）}：平面图形绕不与它相交的直线旋转一周，

\[V=2\pi\bar{r}A,\]

其中 $\bar{r}$ 是\textbf{形心到旋转轴的距离}，$A$ 为区域面积。同理，平面曲线绕轴旋转所得曲面面积 $S=2\pi\bar{r}L$。古尔丁定理有两个用处：一是能快速检验用积分算出的体积是否正确；二是在求形心时可以反过来用——已知体积与面积去求形心的位置。考场上常见“求形心”与“求旋转体体积”互推的设问。

\textbf{函数的平均值}：$f$ 在 $[a,b]$ 上的平均值为

\[\bar{f}=\frac{1}{b-a}\int_a^bf(x)\,dx,\]

它正是积分中值定理中的 $f(\xi)$。加权平均值：$\bar{f}=\frac{\int_a^bf(x)w(x)dx}{\int_a^bw(x)dx}$（$w\geq0$ 为权函数）。

\textbf{典型例子}

（1）半圆域 $x^2+y^2\leq R^2,\ y\geq0$ 的形心：由对称性 $\bar{x}=0$，$A=\frac{\pi R^2}{2}$，

\[\bar{y}=\frac{1}{2A}\int_{-R}^{R}\left(R^2-x^2\right)dx=\frac{1}{\pi R^2}\cdot\frac{4R^3}{3}=\frac{4R}{3\pi}.\]

（2）用古尔丁定理验证球体积：半圆域绕 $x$ 轴旋转，

\[V=2\pi\cdot\frac{4R}{3\pi}\cdot\frac{\pi R^2}{2}=\frac{4}{3}\pi R^3,\qquad\text{与直接积分一致}.\]

（3）用古尔丁定理验证摆线旋转体体积：摆线一拱与 $x$ 轴围成区域 $A=3\pi a^2$，其形心 $\bar{y}=\frac{5a}{6}$，绕 $x$ 轴旋转

\[V=2\pi\cdot\frac{5a}{6}\cdot3\pi a^2=5\pi^2a^3,\]

与 $V=\pi\int_0^{2\pi}y^2dx=\pi a^3\int_0^{2\pi}(1-\cos t)^3dt=5\pi^2a^3$ 完全一致。

（4）平均值：$f(x)=x^2$ 在 $[0,3]$ 上 $\bar{f}=\frac{1}{3}\int_0^3x^2dx=3$；$f(x)=\lvert x-1\rvert$ 在 $[0,2]$ 上 $\bar{f}=\frac{1}{2}\left(\frac{1}{2}+\frac{1}{2}\right)=\frac{1}{2}$。

\textbf{物理量的平均值}：平均速度 $\bar{v}=\frac{1}{t_2-t_1}\int_{t_1}^{t_2}v(t)dt$，平均功率 $\bar{P}=\frac{1}{T}\int_0^TP(t)dt$，平均电流 $\bar{I}=\frac{1}{T}\int_0^TI(t)dt$。它们的共同点是“总量除以总长度”，与积分中值定理是同一件事。

\begin{callwarn}{易错点}
\begin{itemize}
\item \textbf{质心与形心不是一回事}：密度不均匀时质心要用 $\rho$ 加权，不能用面积 $A$ 代替质量。
\item 求曲线质心用 $ds$（弧微分），求平面区域质心用 $dx$（面积微元），二者不能混。
\item 形心的 $\bar{y}$ 公式里有系数 $\frac{1}{2}$（来自 $\frac{y_2+y_1}{2}\cdot(y_2-y_1)$），漏掉必错。
\item 平均值是"积分除以区间长度"，别忘除以 $b-a$；对称区间上奇函数的平均值为 $0$。
\end{itemize}
\end{callwarn}

\begin{callexam}{真题考法}
以\textbf{填空/解答}为主：求平面图形（半圆、抛物线弓形、摆线一拱）的形心，求函数平均值，或先用古尔丁定理求旋转体体积（比积分更快）。命题人偏爱"利用对称性直接写出 $\bar{x}$，只算 $\bar{y}$"的设计。

\end{callexam}

\subsection{积分等式与不等式}

\subsubsection{积分中值问题与积分等式证明 $\bigstar$$\bigstar$$\bigstar$$\bigstar$}

\textbf{积分等式证明的总思路}：等式里出现 $\int_a^bf(x)dx$ 时，相当于出现了一个“平均值”，因此第一步通常是想“能不能用中值定理把它变成 $f(\xi)$”；等式里若出现 $f$ 与 $f'$、$f''$ 的混合，第一步则是“构造变限积分函数 $F(x)=\int_a^xf(t)dt$”，把关于 $f$ 的问题转化为关于 $F$ 的导数问题，再用微分中值定理。凡是结论中出现“存在 $\xi$”，都要想清楚用哪个定理（罗尔、拉格朗日、柯西、介值、零点）能保证存在。

\textbf{五种常用方法}：

\begin{enumerate}
\item \textbf{直接使用（推广的）积分中值定理}：要证的等式中出现"$f(\xi)\times$ 某积分"结构时首选。例如 $f$ 连续、$g$ 不变号时 $\int_a^bfg=f(\xi)\int_a^bg$。
\item \textbf{构造变限积分函数}：设 $F(x)=\int_a^xf(t)dt$，则 $F$ 可导，用罗尔定理、拉格朗日中值定理或单调性处理关于 $f$ 的等式。
\item \textbf{牛顿-莱布尼茨公式 $+$ 介值/零点定理}：若 $\int_a^bf=0$，则 $f$ 在 $[a,b]$ 上的平均值为 $0$，故 $f$ 必取到 $0$（连续时）。
\item \textbf{拆分区间用两次中值定理}（双中值问题）：把 $[a,b]$ 拆成 $[a,c]$、$[c,b]$，分别用中值定理，再比较。
\item \textbf{构造辅助函数}：把结论改写成 $H(\xi)=0$ 的形式，验证 $H$ 在某两点取值相同（罗尔定理）或异号（零点定理）。
\end{enumerate}
\textbf{典型例子}

（1）设 $f$ 在 $[0,1]$ 上连续，证明存在 $\xi\in[0,1]$ 使

\[\int_0^1xf(x)\,dx=\frac{1}{2}f(\xi).\]

证明：取 $g(x)=x$ 在 $[0,1]$ 上非负（不变号），$f$ 连续，由推广的积分中值定理，

\[\int_0^1xf(x)dx=f(\xi)\int_0^1x\,dx=\frac{1}{2}f(\xi).\]

（2）设 $f$ 在 $[0,1]$ 上连续且 $\int_0^1f(x)dx=0$，证明存在 $\xi\in(0,1)$ 使 $f(\xi)=0$。

证明：若 $f\equiv0$ 结论显然。否则由积分中值定理存在 $c\in[0,1]$ 使 $f(c)=0$；若 $c$ 不在内部，说明 $f$ 在端点非零。若 $f$ 恒不小于 $0$ 且不恒为 $0$，则 $\int_0^1f>0$，矛盾；同理不能恒不大于 $0$。故 $f$ 在 $(0,1)$ 内必取正值也取负值，由零点定理存在 $\xi\in(0,1)$ 使 $f(\xi)=0$。

（3）设 $f$ 在 $[0,1]$ 上连续，$F(x)=\int_0^xf(t)dt$，证明存在 $\xi\in(0,1)$ 使 $f(\xi)=\frac{F(1)}{1-0}$（即 $f$ 取到平均值）：直接对 $F$ 用拉格朗日中值定理。

（4）设 $f$ 在 $[0,1]$ 上可导，$f(0)=0$，证明存在 $\xi\in(0,1)$ 使 $f(\xi)=(1-\xi)f'(\xi)$。

证明：把结论改写成 $f(\xi)+(\xi-1)f'(\xi)=0$，构造 $F(x)=(x-1)f(x)$，则 $F'(x)=f(x)+(x-1)f'(x)$，且 $F(0)=F(1)=0$。由罗尔定理，存在 $\xi\in(0,1)$ 使 $F'(\xi)=0$，即 $f(\xi)=(1-\xi)f'(\xi)$。这里“把结论凑成某个函数的导数等于零”是构造辅助函数的通用思路。

\textbf{双中值问题的处理}：若结论中出现两个不同的点 $\xi,\eta$，通常把区间拆成两段（如 $[a,c]$、$[c,b]$，其中 $c$ 常取中点或由题设确定），分别用中值定理后再设法比较；若两个点分别出现在不同的函数上，就对每个函数各用一次中值定理。

\begin{callwarn}{易错点}
\begin{itemize}
\item 中值定理给出的 $\xi$ 只保证\textbf{存在}，无法确定具体位置；求"中值点"的具体值需要额外条件（如 $f$ 单调）。
\item 推广形式必须验证 $g$ \textbf{不变号}；$g=x$ 在 $[0,1]$ 上、$g=\sin x$ 在 $[0,\pi]$ 上都不变号，可以直接用。
\item 要说明 $\xi$ 在\textbf{开区间}内，必须再排除端点的可能性，不能只写闭区间。
\item 等式两边都含积分时，先化为同一个变限积分函数再比较，不要两边各用一次中值定理。
\end{itemize}
\end{callwarn}

\begin{callexam}{真题考法}
\textbf{证明题}高频。常见设问："证明存在 $\xi$ 使 $\int_a^bfg=f(\xi)\int_a^bg$"、"已知 $\int_a^bf=0$ 证 $f$ 有零点"、"$\int_0^1f(x)g(x)dx=f(\xi)\int_0^1g(x)dx$"、"双中值 $\xi,\eta$ 的问题"。答题的第一步几乎总是写出变限积分 $F(x)=\int_a^xf(t)dt$ 或判断 $g$ 是否不变号。

\end{callexam}

\subsubsection{积分不等式与积分放缩 $\bigstar$$\bigstar$$\bigstar$$\bigstar$}

\textbf{积分不等式的总思路}：先看能不能把被积函数\textbf{逐点比较}（保序性），不行就\textbf{估值}（找上下界），再不行就\textbf{用中值定理把积分化为函数值}；涉及平方结构就用\textbf{柯西-施瓦茨不等式}；涉及参数或端点变化就\textbf{构造变限积分函数用单调性}。无论用哪一种，放缩方向必须一路保持一致。

\textbf{六种常用手法}：

\begin{enumerate}
\item \textbf{保序性直接比较被积函数}：在 $[a,b]$ 上若 $f\leq g$，则 $\int_a^bf\leq\int_a^bg$。这是最基本的一招，常常配合"证明 $g-f\geq0$"。
\item \textbf{估值定理}：用 $m\leq f(x)\leq M$ 得 $m(b-a)\leq\int_a^bf\leq M(b-a)$；$m,M$ 通过求导找最值或利用单调性确定。
\item \textbf{中值定理转化为函数值}：$\int_a^bf=f(\xi)(b-a)$，再对 $f(\xi)$ 放缩。
\item \textbf{柯西-施瓦茨不等式}：$\left(\int_a^bfg\right)^2\leq\int_a^bf^2\int_a^bg^2$。凡是要证"积分的平方 $\leq$ 两个积分之积"的题，第一步就是它。
\item \textbf{构造变限积分函数用单调性}：令 $F(x)=\int_a^xf(t)dt$（或 $F(x)=$ 不等式两端的差），证 $F'\geq0$ 且 $F(a)\geq0$（或 $F$ 的某个初值），从而 $F(x)\geq0$。
\item \textbf{利用基本初等函数的放缩}：$x\in(0,+\infty)$ 时 $\sin x\leq x$、$\ln(1+x)\leq x$、$e^x\geq1+x$、$\frac{x}{1+x}\leq\ln(1+x)\leq x$；$x\in[a,b]$ 时 $\frac{1}{b}\leq\frac{1}{x}\leq\frac{1}{a}$。
\end{enumerate}
\textbf{典型例子}

（1）$\displaystyle\int_0^1e^{-x^2}dx\in\left[\frac{1}{e},1\right]$：由 $e^{-1}\leq e^{-x^2}\leq1$ 估值。

（2）$\displaystyle\int_0^{\pi/2}\frac{\sin x}{x}dx\leq\frac{\pi}{2}$：因 $\sin x\leq x$ 得 $\frac{\sin x}{x}\leq1$（并注意补充定义 $x=0$ 处的值）。

（3）$\displaystyle\int_0^1\frac{x}{1+x^2}dx\leq\frac{1}{2}\ln 2\leq\frac{1}{2}$：由 $\frac{1}{1+x^2}\leq1$ 得第一个不等号，直接算出 $\frac{1}{2}\ln2$ 得第二个。

（4）柯西-施瓦茨：对连续函数 $f$，

\[\left(\int_0^1xf(x)\,dx\right)^2\leq\int_0^1f^2(x)\,dx\cdot\int_0^1x^2\,dx=\frac{1}{3}\int_0^1f^2(x)\,dx.\]

（5）构造变限积分：设 $f$ 在 $[0,1]$ 上连续且 $f(x)\geq0$，则 $F(x)=\int_0^xf(t)dt$ 单调不减，从而 $F(1)\geq F(0)=0$，即 $\int_0^1f\geq0$（保号性的再证明）。

\textbf{估值与中值的联合使用}：要证 $\int_a^bf\leq M$ 型结论时，先用 $f$ 的最值作第一层放缩（$m(b-a)\leq\int_a^bf\leq M(b-a)$）；若不够精细，再用 $f$ 的单调性与中值定理把 $m,M$ 换成更紧的界。每一步放缩都要写清依据（保序性、估值定理还是基本初等函数不等式）。

（6）\textbf{用单调性放缩}：设 $f$ 在 $[0,1]$ 上可导且 $f'(x)>0$，则 $f$ 单调增加，故 $f(0)\leq f(x)\leq f(1)$，由估值定理得 $f(0)\leq\int_0^1f(x)dx\leq f(1)$，即函数的平均值介于两端函数值之间——“平均值夹在两端之间”正是单调性的直接推论。

\begin{callwarn}{易错点}
\begin{itemize}
\item 放缩方向必须一致：想证"$\leq$"就要一路往大处放，中途反向放缩会毁掉整个证明。
\item 柯西-施瓦茨不等式是"平方后比较"，不能用来证明非平方形式的不等式。
\item 用 $\sin x\leq x$ 时注意 $x>0$；$x<0$ 时要变号。
\item 使用构造法证明时，别忘验证\textbf{初始条件}（如 $F(a)\geq0$），只证 $F'\geq0$ 是不够的。
\item 中值定理做不等式时 $\xi$ 的位置不确定，放缩要写成只与 $a,b$ 有关的形式。
\end{itemize}
\end{callwarn}

\begin{callexam}{真题考法}
\textbf{证明题}高频，常与中值定理题交替出现。设问形式：证明某积分不超过某个常数、证明积分关于参数的单调性、证明 $\int_0^1f^2\geq\left(\int_0^1f\right)^2$ 型。答题时的通用起手式是"先在区间上确定被积函数的界，再用估值定理"或"构造变限积分函数并求导"。

\end{callexam}

\subsection*{本章易错清单}
\addcontentsline{toc}{subsection}{本章易错清单}

\begin{enumerate}
\item \textbf{不定积分漏写 $C$}；分段函数的原函数各段常数不同，必须用"原函数在分段点连续"把 $C_1,C_2$ 串起来，并验证分段点可导。
\item $\int\frac{dx}{x}=\ln\lvert x\rvert+C$ 的\textbf{绝对值}不能丢；$\sqrt{x^2}=\lvert x\rvert$，去根号就是去绝对值。
\item 定积分\textbf{换元必换限}，只换被积函数不换上下限；不定积分换元后必须回代。
\item 分部积分只按"反对幂指三"取 $u$，取反了会越算越复杂；循环型移项时不要漏 $+C$；定积分的边界项不能省。
\item 被积函数含 $x$ 的变限积分\textbf{不能直接求导}，要"能提就提、不能提就换元"；上限为复合函数时别漏 $\psi'(x)$。
\item 变限积分的可导性与被积函数的连续性不是一回事：跳跃间断点上变限积分\textbf{连续但不可导}，可去间断点上变限积分可导但 $\Phi'\neq f$。
\item 定积分是\textbf{代数面积}，求几何面积要加绝对值并分段；上下边界交叉时必须分段。
\item 含瑕点（分母为零的点、$\ln x$ 在 $0$ 处等）或有跳跃间断点的区间\textbf{不能直接用牛顿-莱布尼茨公式}；多奇点必须拆开，一段发散即整体发散。
\item 两类反常积分的 p 判据\textbf{方向相反}：无穷限要 $p>1$ 收敛，瑕积分要 $p<1$ 收敛，极易记混。
\item 比较判别法要求被积函数\textbf{非负}，变号时要先看绝对收敛；等价无穷小替换只在\textbf{奇点附近}使用。
\item 旋转体体积：垫圈法"外$^2-$内$^2$"顺序别反；绕非坐标轴时半径是到\textbf{旋转轴}的距离；曲线与轴相交要分段。
\item 侧面积公式 $S=2\pi\int\lvert y\rvert\,ds$ 中 \textbf{$\lvert y\rvert$ 与弧微分缺一不可}；极坐标侧面积不要漏 $\lvert r\sin\theta\rvert$。
\item Wallis 公式要求区间先化为 $[0,\frac{\pi}{2}]$，偶次带 $\frac{\pi}{2}$、奇次不带，别记反。
\item 物理应用：抽水做功中"提升高度"是微元到出口的距离；压力题必须先写出宽度函数 $L(h)$；引力必须分解分量。
\item 形心公式中 $\bar{y}$ 的系数 $\frac{1}{2}$ 与平均值公式中的 $\frac{1}{b-a}$ 都不能漏；质心只有在密度均匀时才等于形心。
\item 算完务必\textbf{验证}：不定积分求导回代验证，定积分用几何意义（面积、体积非负）与量纲检查。本章计算量大，逐步核对比整题重算更省时间。
\end{enumerate}
\chapter{多元函数微积分学}
\markboth{多元函数微积分学}{多元函数微积分学}
\subsection*{本章考纲要求}
\addcontentsline{toc}{subsection}{本章考纲要求}

\begin{quote}
\textbf{考试内容}

\begin{enumerate}
\item 多元函数的概念
\item 二元函数的几何意义
\item 二元函数的极限与连续的概念
\item 有界闭区域上二元连续函数的性质
\item 多元函数的偏导数和全微分
\item 多元复合函数、隐函数的求导法
\item 二阶偏导数
\item 多元函数的极值和条件极值、最大值和最小值
\item 二重积分的概念、基本性质和计算
\end{enumerate}
\textbf{考试要求}

\begin{enumerate}
\item 了解多元函数的概念，了解二元函数的几何意义。
\item 了解多元函数的极限与连续的概念，了解有界闭区域上二元连续函数的性质。
\item 了解多元函数偏导数与全微分的概念，会求多元复合函数一阶、二阶偏导数，会求全微分，了解隐函数存在定理，会求多元隐函数的偏导数。
\item 了解多元函数极值和条件极值的概念，掌握多元函数极值存在的必要条件，了解二元函数极值存在的充分条件，会求二元函数的极值，会用拉格朗日乘数法求条件极值，会求简单多元函数的最大值和最小值，并会解决一些简单的应用问题。
\item 理解二重积分的概念，了解二重积分的基本性质，了解二重积分的中值定理，掌握二重积分的计算方法(直角坐标、极坐标)。
\end{enumerate}
\end{quote}
\subsection*{本章星级总表}
\addcontentsline{toc}{subsection}{本章星级总表}

\begingroup\small\setlength{\tabcolsep}{4.0pt}
\begin{longtable}{@{}L{0.4597\linewidth}L{0.1086\linewidth}L{0.1086\linewidth}L{0.1086\linewidth}L{0.1086\linewidth}@{}}
\toprule
\textbf{考点} & \textbf{星级} & \textbf{考纲要求} & \textbf{难度} & \textbf{常见题型} \\
\midrule
\endhead
多元函数的概念、二元函数的几何意义 & $\bigstar$ & 了解 & 基础 & 选择 \\
二元函数的极限与连续、有界闭区域上连续函数性质 & $\bigstar$$\bigstar$ & 了解 & 中等 & 选择 \\
偏导数与全微分的概念、可微的判定 & $\bigstar$$\bigstar$$\bigstar$ & 了解/会 & 难 & 选择 \\
多元复合函数的求导法（链式法则、抽象函数） & $\bigstar$$\bigstar$$\bigstar$$\bigstar$ & 会 & 难 & 解答 \\
多元隐函数的求导法（含隐函数存在定理） & $\bigstar$$\bigstar$$\bigstar$ & 会 & 中等 & 解答 \\
二阶偏导数与混合偏导相等的条件 & $\bigstar$$\bigstar$$\bigstar$ & 会 & 中等 & 填空、解答 \\
多元函数的极值、条件极值与拉格朗日乘数法 & $\bigstar$$\bigstar$$\bigstar$$\bigstar$ & 会 & 难 & 解答 \\
多元函数的最大值、最小值及简单应用 & $\bigstar$$\bigstar$$\bigstar$ & 会 & 中等 & 解答、应用 \\
二重积分的概念、基本性质与中值定理 & $\bigstar$$\bigstar$$\bigstar$ & 理解/了解 & 中等 & 选择 \\
二重积分的计算：直角坐标、交换积分次序 & $\bigstar$$\bigstar$$\bigstar$$\bigstar$$\bigstar$ & 掌握 & 中等 & 解答 \\
二重积分的计算：极坐标 & $\bigstar$$\bigstar$$\bigstar$$\bigstar$$\bigstar$ & 掌握 & 中等 & 解答 \\
二重积分的对称性、奇偶性与轮换对称性化简 & $\bigstar$$\bigstar$$\bigstar$$\bigstar$ & 会 & 中等 & 解答 \\
分区域积分、含绝对值与最值函数的二重积分 & $\bigstar$$\bigstar$$\bigstar$$\bigstar$ & 会 & 难 & 解答 \\
\bottomrule
\end{longtable}
\endgroup

\section{多元函数的基本概念}

本章前半部分（概念、极限、连续、偏导与全微分）在数学二中大多只要求"了解"，但它们决定了后面复合求导、极值与二重积分的所有前提。复习策略是：\textbf{概念题的判定逻辑必须清楚，计算题不能出错}。尤其是"可微"这一条，是选择题反复设卡的地方。

\subsection{多元函数的概念与几何意义}

\subsubsection{多元函数的概念、二元函数的几何意义 $\bigstar$}

设 $D$ 是平面上的一个点集。若对 $D$ 中每一点 $(x,y)$，按照某个对应法则 $f$ 都有唯一确定的实数 $z$ 与之对应，则称 $f$ 是定义在 $D$ 上的\textbf{二元函数}，记作

\[z=f(x,y),\qquad (x,y)\in D\]

其中 $x,y$ 称为\textbf{自变量}，$z$ 称为\textbf{因变量}，$D$ 称为函数的\textbf{定义域}，全体函数值构成的集合称为\textbf{值域}。三元及以上的函数统称多元函数，考纲只要求掌握到二元。

\textbf{定义域的求法}与一元函数完全一致，只是要写成平面区域（并画出草图）：

\begin{itemize}
\item 分式：分母不为零；
\item 偶次根式：被开方式 $\ge 0$；
\item 对数：真数 $>0$；
\item $\arcsin u,\ \arccos u$：要求 $\vert u\vert \le 1$；
\item 正切 $\tan u$：$u\ne k\pi+\frac{\pi}{2}$；
\item 多个限制条件同时成立时取\textbf{交集}。
\end{itemize}
\textbf{几何意义}：二元函数 $z=f(x,y)$ 的图形是空间直角坐标系中的一张\textbf{曲面}；定义域 $D$ 就是这张曲面在 $xOy$ 平面上的\textbf{投影}。把 $f(x,y)$ 取常数得到的曲线 $f(x,y)=c$ 称为\textbf{等高线}，它是曲面被平面 $z=c$ 截得的截痕在 $xOy$ 面上的投影。

例：求 $f(x,y)=\sqrt{4-x^{2}-y^{2}}+\ln y$ 的定义域。

由偶次根式得 $x^{2}+y^{2}\le 4$，由对数得 $y>0$。故定义域是上半圆盘

\[D=\{(x,y)\mid x^{2}+y^{2}\le 4,\ y>0\}\]

即半径为 $2$ 的上半圆（不含直径所在直线 $y=0$）。这里 $z\ge 0$ 的部分是上半球面，投影正是这个上半圆盘。 若函数的对应法则以解析式给出而未说明定义域，则约定定义域为使该式有意义的一切点构成的集合，称为\textbf{自然定义域}（统考按此理解）。例如 $g(x,y,z)=\frac{z}{\sqrt{9-x^{2}-y^{2}-z^{2}}}$ 的自然定义域是球面 $x^{2}+y^{2}+z^{2}=9$ 的内部（不含球面），它是一个开区域。

\textbf{表示法与常见曲面}：二元函数的表示法与一元类似，可用解析式、表格或图像给出。画图时注意 $D$ 是平面上的区域、画在 $xOy$ 面内，而曲面 $z=f(x,y)$ 画在空间中，两者不要混。对常见曲面要形成条件反射：$z=x^{2}+y^{2}$ 是开口向上的旋转抛物面；$z=\sqrt{R^{2}-x^{2}-y^{2}}$ 是上半球面；$z=xy$ 是马鞍面（双曲抛物面）；$z^{2}=x^{2}+y^{2}$ 要限定 $z\ge 0$（或 $z\le 0$）才是函数图像，它是圆锥面的一半。这些图形在极值题（判断鞍点）与二重积分题（判断区域形状）中反复用到。 多元函数与一元函数的根本差别在于\textbf{趋近一个点的方式有无穷多种}，这一点决定了后面所有概念（极限、连续、可微）都必须"对一切方向同时成立"。理解这一句，本章概念的判定题就有了统一的出发点：凡是只验证了部分方向就下结论的做法，几乎都是错的。

\begin{callkey}{三个"图形—区域"对应}
\begin{itemize}
\item 球的内部 $x^{2}+y^{2}+z^{2}<R^{2}$，在二重积分中压成圆盘 $x^{2}+y^{2}\le R^{2}$；
\item 圆锥面 $z=\sqrt{x^{2}+y^{2}}$ 的投影仍是圆盘，被积函数要按锥面上的表达式代入；
\item 抛物面 $z=x^{2}+y^{2}$ 在 $xOy$ 面内的投影是圆盘，它上方立体的体积正好用二重积分表示。
\end{itemize}
\end{callkey}

\begin{callwarn}{易错点}
定义域是平面上的一块\textbf{区域}，不是空间中的曲面，答"定义域是上半球面"是典型错误。另外，多元函数的定义域必须写成集合或不等式组，只写"$y>0$"会丢分。

\end{callwarn}

\begin{callexam}{真题考法}
以选择、填空形式考查：给出二元（偶尔三元）函数求定义域并画图，或反过来由图形写不等式。常见陷阱是漏掉对数真数、分母、根号中的隐性限制，以及忘记若干个条件取交集。 复习时建议把定义域题与后面的"区域画法"一起练：二元函数的定义域本身就是二重积分的积分区域，把这道小题练熟，后面画图定限会顺很多。

\end{callexam}

\subsection{二元函数的极限与连续}

\subsubsection{二元函数的极限与连续、有界闭区域上连续函数性质 $\bigstar$$\bigstar$}

\textbf{二重极限的定义}：设 $f(x,y)$ 在点 $P_0(x_0,y_0)$ 的某去心邻域内有定义，若存在常数 $A$，使对任意 $\varepsilon>0$，总存在 $\delta>0$，当

\[0<\sqrt{(x-x_0)^{2}+(y-y_0)^{2}}<\delta\]

时恒有 $\vert f(x,y)-A\vert <\varepsilon$，则称 $A$ 是 $f(x,y)$ 当 $(x,y)\to(x_0,y_0)$ 时的极限，记作

\[\lim_{(x,y)\to(x_0,y_0)}f(x,y)=A\]

核心要素有三条：\textbf{去心邻域}（不要求 $P_0$ 处有定义）、\textbf{路径无关}（以任何方式趋近 $P_0$ 极限都相同）、\textbf{二维距离}代替一元的绝对值。数学二对严格的 $\varepsilon$-$\delta$ 语言只要求"了解"，重点是会用下列方法判定极限的存在与不存在。

\textbf{求二重极限的常用方法}

\begin{enumerate}
\item \textbf{夹逼准则（最常用）}：把 $\vert f(x,y)-A\vert$ 放大成只含 $\sqrt{x^{2}+y^{2}}$（或某个趋于 $0$ 的量）的式子，若放大后的式子趋于 $0$，则极限为 $A$。常用配凑：
\end{enumerate}
\[0\le \frac{x^{2}}{x^{2}+y^{2}}\le 1,\qquad \vert xy\vert\le \frac{x^{2}+y^{2}}{2}\]

\begin{enumerate}
\item \textbf{极坐标法}：令 $x=x_0+r\cos\theta,\ y=y_0+r\sin\theta$，若 $f$ 化为 $g(r,\theta)$ 且 $\vert g(r,\theta)-A\vert\le \varphi(r)$ 对一切 $\theta$ 成立、$\varphi(r)\to 0\ (r\to 0^{+})$，则极限为 $A$。\textbf{关键是放大后的界必须与 $\theta$ 无关}。
\item \textbf{等价无穷小替换}：整体结构中的 $\sin u\sim u$、$1-\cos u\sim \frac{u^{2}}{2}$、$\ln(1+u)\sim u$、$e^{u}-1\sim u$ 都可以照用，其中 $u$ 是趋于 $0$ 的二元表达式（如 $u=xy$）。
\end{enumerate}
\textbf{证明二重极限不存在的方法}

\begin{itemize}
\item \textbf{路径法}：取两条不同的路径（如 $y=kx$、$y=kx^{2}$、$y=x^{2}$），算出两个不同的极限，即证不存在；
\item \textbf{极坐标法}：化出 $g(r,\theta)$ 后若结果与 $\theta$ 有关，则不存在。
\end{itemize}
\textbf{求二重极限的四步流程}

\begin{enumerate}
\item \textbf{先试路径}：取 $y=kx$、$y=kx^{2}$、$y=x^{2}$ 等，若两条路径得到不同极限，立即断定极限不存在，这是最省时的做法；
\item \textbf{再试夹逼}：若路径法没有发现矛盾，就把表达式放大成只含 $\sqrt{x^{2}+y^{2}}$ 的整体，或放成"有界量乘无穷小"的形式；
\item \textbf{改用极坐标}：当分子分母都含平方和时，化到 $g(r,\theta)$；若 $\theta$ 被消去、余下只含 $r$，则极限就是 $r\to 0$ 时的值；
\item \textbf{先化简再判定}：用等价无穷小或重要极限结构（$\frac{\sin u}{u}\to 1$、$(1+u)^{1/u}\to e$）简化后再回到第 2、3 步。
\end{enumerate}
\textbf{二重极限与累次极限的区别}（了解）：$\lim_{x\to x_0}\lim_{y\to y_0}f(x,y)$ 与 $\lim_{y\to y_0}\lim_{x\to x_0}f(x,y)$ 是两次一元极限，称为\textbf{累次极限}，与二重极限不是一回事。三者关系是：二重极限存在且两个累次极限都存在时三者相等；但累次极限存在推不出二重极限存在。考纲范围内的极限题一律以二重极限的定义为准。 还要强调：\textbf{路径法只能否定存在性，永远不能肯定存在性}；夹逼与极坐标既可以肯定存在（给出与方向无关的估计），也可以在估计失败时提示"可能不存在"，此时再回头用路径法找两条不同路径的反例。两类问法的写法完全不同：问"是否存在"要给出路径并比较；问"求极限"要给出夹逼或极坐标的估计过程，不能只写一句"极限为 $0$"。

\textbf{连续性}：若 $\displaystyle\lim_{(x,y)\to(x_0,y_0)}f(x,y)=f(x_0,y_0)$，则称 $f$ 在 $P_0$ 处连续。三要素为：$P_0$ 处有定义、该点极限存在、极限值等于函数值。和、差、积、商（分母非零）与复合仍保持连续，初等函数在其定义域内连续——这是判断连续区域的通用工具。

\textbf{有界闭区域上二元连续函数的性质}（了解）：设 $D$ 为有界闭区域，$f$ 在 $D$ 上连续，则

\begin{itemize}
\item \textbf{有界性}：存在 $M>0$，使 $\vert f(P)\vert\le M$ 对一切 $P\in D$ 成立；
\item \textbf{最值定理}：$f$ 在 $D$ 上必取得最大值与最小值；
\item \textbf{介值定理}：若 $D$ 连通，$c$ 介于两个函数值之间，则存在 $P\in D$ 使 $f(P)=c$。
\end{itemize}
这三条是后面"闭区域上求最值"的理论依据。

例 1：求 $\displaystyle\lim_{(x,y)\to(0,0)}\frac{x^{2}y}{x^{2}+y^{2}}$。

因为

\[\left\vert \frac{x^{2}y}{x^{2}+y^{2}}\right\vert =\vert y\vert\cdot\frac{x^{2}}{x^{2}+y^{2}}\le \vert y\vert\to 0\]

由夹逼准则，极限为 $0$。

例 2：证明 $\displaystyle\lim_{(x,y)\to(0,0)}\frac{xy}{x^{2}+y^{2}}$ 不存在。

沿 $y=x$：$\frac{x\cdot x}{x^{2}+x^{2}}=\frac12$；沿 $y=0$：$\frac{0}{x^{2}}=0$。两条路径极限不同，故极限不存在。

例 3：求 $\displaystyle\lim_{(x,y)\to(0,0)}\frac{1-\cos(xy)}{x^{2}y^{2}}$。

由 $1-\cos u\sim \frac{u^{2}}{2}$，令 $u=xy$，得原式 $=\frac{(xy)^{2}/2}{x^{2}y^{2}}=\frac12$。

\begin{callwarn}{易错点}
三个高频陷阱。其一，沿\textbf{所有直线} $y=kx$ 极限都相同，仍不能断言二重极限存在：例如 $\frac{x^{2}y}{x^{4}+y^{2}}$ 沿 $y=kx$ 时化为 $\frac{kx}{x^{2}+k^{2}}\to 0$，但沿 $y=x^{2}$ 时化为 $\frac{x^{4}}{2x^{4}}=\frac12$，故极限不存在。必须用夹逼或极坐标给出与方向无关的统一估计。其二，极坐标法中若化简结果仍含 $\theta$，只能说明"沿该路径族极限依赖方向"，据此可判定不存在，但不能用它证明存在。其三，函数在某点不连续并不能推出极限不存在（可能只是"极限存在但与函数值不等"）。

\end{callwarn}

\begin{callexam}{真题考法}
选择题常考两件事：判断给定二重极限是否存在（给一个含平方和分母的分式）；用极限存在的条件反求参数。陷阱集中在"只验证了直线路径就下结论"以及忘记讨论 $\theta$ 是否被消去。二重极限与累次极限的关系只需知道：二重极限存在且累次极限也存在时三者相等，但累次极限存在推不出二重极限存在。

\end{callexam}

\subsection{偏导数与全微分}

\subsubsection{偏导数与全微分的概念、可微的判定 $\bigstar$$\bigstar$$\bigstar$}

\textbf{偏导数的定义}：设 $z=f(x,y)$ 在 $P_0(x_0,y_0)$ 的某邻域内有定义，则

\[f_x(x_0,y_0)=\lim_{\Delta x\to 0}\frac{f(x_0+\Delta x,y_0)-f(x_0,y_0)}{\Delta x}\]

\[f_y(x_0,y_0)=\lim_{\Delta y\to 0}\frac{f(x_0,y_0+\Delta y)-f(x_0,y_0)}{\Delta y}\]

\textbf{求法}：对 $x$ 求偏导时，把 $y$ 当作常数，用一元求导法则即可；对 $y$ 同理。但\textbf{分段函数在分界点、以及偏导数不连续的点，必须用定义式（上面两个极限）计算}。

\textbf{几何意义}：$f_x(x_0,y_0)$ 是曲面 $z=f(x,y)$ 与平面 $y=y_0$ 的交线在 $x=x_0$ 处的切线斜率；$f_y(x_0,y_0)$ 是曲面与 $x=x_0$ 的交线在 $y=y_0$ 处的切线斜率。

\textbf{全微分的定义}：若

\[\Delta z=f(x_0+\Delta x,y_0+\Delta y)-f(x_0,y_0)=A\Delta x+B\Delta y+o(\rho)\]

其中 $\rho=\sqrt{(\Delta x)^{2}+(\Delta y)^{2}}$，$A,B$ 与 $\Delta x,\Delta y$ 无关，则称 $f$ 在 $P_0$ \textbf{可微}，并称 $A\Delta x+B\Delta y$ 为全微分，记作

\[dz=A\Delta x+B\Delta y=f_x(x_0,y_0)\,dx+f_y(x_0,y_0)\,dy\]

\textbf{三个概念的关系（必须背下来）}

\begin{itemize}
\item 可微 $\Rightarrow$ 偏导数存在，且 $A=f_x,\ B=f_y$；
\item 可微 $\Rightarrow$ 连续；
\item 偏导数存在 $\not\Rightarrow$ 连续，偏导数存在 $\not\Rightarrow$ 可微；
\item 偏导数连续 $\Rightarrow$ 可微（充分条件，最实用）。
\end{itemize}
关系链可记为：

\[\text{偏导数连续}\Rightarrow \text{可微}\Rightarrow \begin{cases}\text{函数连续}\\[2pt] \text{偏导数存在}\end{cases}\]

\textbf{判定可微的三步流程}

\begin{enumerate}
\item 先求 $f_x(x_0,y_0),f_y(x_0,y_0)$：若有一个不存在，立即判定\textbf{不可微}；
\item 若都存在，计算极限
\end{enumerate}
\[\lim_{\rho\to 0}\frac{\Delta z-\left[f_x(x_0,y_0)\Delta x+f_y(x_0,y_0)\Delta y\right]}{\rho}\]

若该极限为 $0$，则可微；若不为 $0$ 或不存在，则不可微；

\begin{enumerate}
\item 若 $f_x,f_y$ 在 $P_0$ 的某邻域内存在且连续，则直接由充分条件判定可微。
\end{enumerate}
\textbf{全微分的近似应用}：当 $\vert \Delta x\vert,\vert \Delta y\vert$ 很小时，

\[f(x_0+\Delta x,y_0+\Delta y)\approx f(x_0,y_0)+f_x(x_0,y_0)\Delta x+f_y(x_0,y_0)\Delta y\]

例 1：设 $z=x^{y}\ (x>0)$，求 $dz$。

\[z_x=yx^{y-1},\qquad z_y=x^{y}\ln x,\qquad dz=yx^{y-1}dx+x^{y}\ln x\,dy\]

\textbf{用定义求偏导的典型情形}：分段函数在分界点、偏导表达式在分界点两侧不同的点，都必须用定义式。例如 $f(x,y)=\sqrt{x^{2}+y^{2}}$ 在原点处

\[\frac{f(\Delta x,0)-f(0,0)}{\Delta x}=\frac{\vert \Delta x\vert}{\Delta x}\]

该式的极限不存在，故 $f_x(0,0)$ 不存在（同理 $f_y(0,0)$ 不存在）。这说明函数在原点连续并不能保证偏导数存在。

可微的定义还可等价地写成：$f$ 在 $P_0$ 可微当且仅当存在常数 $A,B$，使 $\Delta z-A\Delta x-B\Delta y=o(\rho)$。由于可微时必有 $A=f_x,\ B=f_y$，所以只需验证上面第 2 步的极限。若题目已知 $f_x,f_y$ 在某邻域内连续（初等函数在其定义域内即如此），可直接判定可微并用公式求 $dz$，不必再做极限运算。

例 2：判断 $f(x,y)=\dfrac{xy}{\sqrt{x^{2}+y^{2}}}\ ((x,y)\ne(0,0))$，$f(0,0)=0$ 在原点是否可微。

用定义：$f(\Delta x,0)=\frac{\Delta x\cdot 0}{\vert \Delta x\vert}=0$，故 $f_x(0,0)=0$；同理 $f_y(0,0)=0$。此时 $\Delta z=\frac{\Delta x\Delta y}{\rho}$，于是

\[\frac{\Delta z-f_x\Delta x-f_y\Delta y}{\rho}=\frac{\Delta x\Delta y}{\rho^{2}}\]

沿 $\Delta y=\Delta x$ 时该式为 $\frac{\Delta x^{2}}{2\Delta x^{2}}=\frac12\ne 0$，故极限不为 $0$，$f$ 在原点\textbf{不可微}。可见两个偏导数都存在也保证不了可微。

例 3：用全微分计算 $(1.04)^{2.02}$ 的近似值。

取 $f(x,y)=x^{y}$，$(x_0,y_0)=(1,2)$，$\Delta x=0.04,\ \Delta y=0.02$。此时 $f(1,2)=1$，$f_x=yx^{y-1}=2$，$f_y=x^{y}\ln x=0$。故

\[(1.04)^{2.02}\approx 1+2\times 0.04+0\times 0.02=1.08\]

例 4：证明 $f(x,y)=(x^{2}+y^{2})\sin\frac{1}{x^{2}+y^{2}}\ ((x,y)\ne(0,0))$，$f(0,0)=0$ 在原点可微，但偏导数在原点不连续。

由 $f(\Delta x,0)=(\Delta x)^{2}\sin\frac{1}{(\Delta x)^{2}}$ 得 $f_x(0,0)=\lim_{\Delta x\to 0}\Delta x\sin\frac{1}{(\Delta x)^{2}}=0$，同理 $f_y(0,0)=0$。而

\[\frac{\Delta z-0}{\rho}=\rho\sin\frac{1}{\rho^{2}}\to 0\]

故可微。但当 $(x,y)\ne(0,0)$ 时

\[f_x=2x\sin\frac{1}{x^{2}+y^{2}}-\frac{2x}{x^{2}+y^{2}}\cos\frac{1}{x^{2}+y^{2}}\]

沿 $y=0$ 看，第二项为 $\frac{2}{x}\cos\frac{1}{x^{2}}$，$x\to 0$ 时无极限，故 $f_x$ 在原点不连续。这恰好说明"偏导数连续"只是可微的\textbf{充分}条件，不是必要条件。

\begin{callkey}{可微判定必背三句}
\begin{enumerate}
\item 偏导存在且 $\lim_{\rho\to 0}\frac{\Delta z-f_x\Delta x-f_y\Delta y}{\rho}=0$ $\iff$ 可微；
\item 偏导连续 $\Rightarrow$ 可微（充分不必要）；
\item 可微 $\Rightarrow$ 连续，可微 $\Rightarrow$ 偏导存在（必要不充分）。
\end{enumerate}
\end{callkey}

\begin{callwarn}{易错点}
四个易错点。第一，把"偏导存在"当成"可微"或"连续"，这是一元函数"可导即连续"的思维惯性所致，二元函数完全失效。第二，判断可微时只算了两个偏导就下结论，漏掉第 2 步的极限验证。第三，分段函数在分界点的偏导必须用定义式算，直接对表达式求导再代点几乎必错。第四，$dz=f_x\,dx+f_y\,dy$ 只在可微时成立，不可微点写成这个式子是无意义的。

\end{callwarn}

\begin{callexam}{真题考法}
选择题的固定套路：给出一个分段定义的二元函数（分子分母同次，或含 $\vert xy\vert$），问它在原点的连续性、偏导存在性、可微性中哪些成立。答题时严格按"偏导是否存在 $\to$ 极限是否为零 $\to$ 偏导是否连续"的顺序逐条验证，不要跳步。另有一类小题是求某点处的全微分或估计绝对误差。 全微分的近似计算（如估算 $1.02^{2.99}$、估计测量误差）也是小题来源，写法固定：选定 $(x_0,y_0)$、求出 $f_x,f_y$、代入 $f+\Delta f$。答题时要说明"$\Delta x,\Delta y$ 很小，故可用全微分近似"，并保留足够位数。

\end{callexam}

\section{多元函数的微分法}

这一板块是本章的第一个计算主战场。链式法则"路径相乘、分叉相加"是所有求导公式的来源；隐函数求导则把公式法、两边求导法、全微分法三种手段统一起来。数学二的解答题常以"抽象函数 + 二阶偏导 + 验证恒等式"的形式出现。

\subsection{复合函数与隐函数的求导}

\subsubsection{多元复合函数的求导法（链式法则、抽象函数） $\bigstar$$\bigstar$$\bigstar$$\bigstar$}

\textbf{基本链式法则}：设 $z=f(u,v)$ 可微，$u=u(x,y)$、$v=v(x,y)$ 具有偏导数，则

\[\frac{\partial z}{\partial x}=\frac{\partial z}{\partial u}\frac{\partial u}{\partial x}+\frac{\partial z}{\partial v}\frac{\partial v}{\partial x},\qquad \frac{\partial z}{\partial y}=\frac{\partial z}{\partial u}\frac{\partial u}{\partial y}+\frac{\partial z}{\partial v}\frac{\partial v}{\partial y}\]

\textbf{全导数}：若 $z=f(x,y)$，而 $y=y(x)$，则

\[\frac{dz}{dx}=\frac{\partial f}{\partial x}+\frac{\partial f}{\partial y}\frac{dy}{dx}\]

注意此处的 $\frac{\partial f}{\partial x}$ 与 $\frac{dz}{dx}$ 不是一回事：前者是 $f$ 对第一个位置的偏导，后者是沿曲线 $y=y(x)$ 的全导数。若 $u=u(t),v=v(t)$，则 $\frac{dz}{dt}=\frac{\partial z}{\partial u}\frac{du}{dt}+\frac{\partial z}{\partial v}\frac{dv}{dt}$。

\textbf{树状图法（万能做法）}\textbf{复合函数的三种常见结构}（必须分清，因为记号与公式都不同）

\begin{enumerate}
\item \textbf{一元中间变量}：$z=f(u,v)$，$u=u(t)$，$v=v(t)$，最终 $z$ 只是 $t$ 的一元函数，用全导数公式，不能写偏导符号；
\item \textbf{多元中间变量}：$z=f(u,v)$，$u=u(x,y)$，$v=v(x,y)$，$z$ 是 $x,y$ 的二元函数，用上面的链式法则；
\item \textbf{混合结构}：$z=f(x,y,u)$，$u=u(x,y)$，此时 $x$ 既通过 $u$ 影响 $z$、又直接是 $f$ 的自变量，两条路径都要算：
\end{enumerate}
\[\frac{\partial z}{\partial x}=f_1'+f_3'\frac{\partial u}{\partial x},\qquad \frac{\partial z}{\partial y}=f_2'+f_3'\frac{\partial u}{\partial y}\]

第 3 种结构下，$f_1'$ 表示 $f$ 对\textbf{第一个位置（即 $x$ 本身）}的偏导，与 $\frac{\partial z}{\partial x}$ 不是一回事——后者还要加上经由 $u$ 传来的那一项。这是"位置记号"最容易出错的地方。

：先画变量关系图——根是 $z$，中间是 $u,v$，叶是 $x,y$。规则是

\begin{itemize}
\item \textbf{路径相乘}：从 $z$ 出发沿一条路径走到 $x$，每一步的偏导相乘；
\item \textbf{分叉相加}：所有通向 $x$ 的路径的结果相加。
\end{itemize}
\textbf{抽象函数的两个关键动作}

\begin{enumerate}
\item 把 $\frac{\partial z}{\partial u}$ 记作 $f_u'$（或 $f_1'$），它仍然是一个以 $u,v$ 为中间变量的二元函数；
\item \textbf{对一阶偏导再求导时，必须继续对 $u,v$ 用链式法则}，这是所有错误的来源：
\end{enumerate}
\[\frac{\partial^{2}z}{\partial x^{2}}=f_{uu}''\left(\frac{\partial u}{\partial x}\right)^{2}+2f_{uv}''\frac{\partial u}{\partial x}\frac{\partial v}{\partial x}+f_{vv}''\left(\frac{\partial v}{\partial x}\right)^{2}+f_u'\frac{\partial^{2}u}{\partial x^{2}}+f_v'\frac{\partial^{2}v}{\partial x^{2}}\]

其中混合项系数为 $2$，因为 $f_{uv}''=f_{vu}''$（在连续的前提下）。若 $u,v$ 都是 $x,y$ 的\textbf{线性}函数，则 $\frac{\partial^{2}u}{\partial x^{2}}=\frac{\partial^{2}v}{\partial x^{2}}=0$，最后两项消失。

例 1：设 $z=f(2x-3y,\ e^{xy})$，$f$ 具有二阶连续偏导，求 $\frac{\partial z}{\partial x}$ 与 $\frac{\partial z}{\partial y}$。

令 $u=2x-3y,\ v=e^{xy}$，则 $u_x=2,\ u_y=-3,\ v_x=ye^{xy},\ v_y=xe^{xy}$，于是

\[\frac{\partial z}{\partial x}=2f_u'+ye^{xy}f_v',\qquad \frac{\partial z}{\partial y}=-3f_u'+xe^{xy}f_v'\]

例 2：设 $z=f(x+y,\ xy)$，$f$ 二阶连续可导，求 $\frac{\partial^{2}z}{\partial x\partial y}$。

令 $u=x+y,\ v=xy$，则 $u_x=u_y=1,\ v_x=y,\ v_y=x$，故

\[z_x=f_u'+yf_v'\]

再对 $y$ 求导，注意 $f_u',f_v'$ 仍是 $u,v$ 的函数：

\[z_{xy}=f_{uu}''\cdot 1+f_{uv}''\cdot x+f_v'+y\left(f_{vu}''\cdot 1+f_{vv}''\cdot x\right)=f_{uu}''+(x+y)f_{uv}''+xyf_{vv}''+f_v'\]

例 3：设 $z=f\left(x,\ \frac yx\right)$，$f$ 具有二阶连续偏导，求 $\frac{\partial z}{\partial x}$。

令 $u=x,\ v=\frac yx$，则 $u_x=1,\ v_x=-\frac{y}{x^{2}}$，故

\[\frac{\partial z}{\partial x}=f_u'-\frac{y}{x^{2}}f_v'\]

这类"第二个中间变量是分式"的结构常被改装成 $f\left(\frac yx\right)$（齐次结构）或 $f\left(xy,\ \frac xy\right)$，做法完全一样：先认准中间变量，再算它们对自变量的偏导。

\begin{callkey}{链式法则口诀}
"根到叶，路径相乘；同层分叉，结果相加；抽象函数再求导，$f_u'$ 还要走一遍。"

\end{callkey}

另外，二阶偏导的交叉项在连续时 $f_{uv}''=f_{vu}''$，展开后可直接合并写成 $2f_{uv}''$，既省书写又不易出错。

\begin{callwarn}{易错点}
第一，$f_u'$ 不是常数，对 $x$ 或 $y$ 再求导时它本身要按链式法则展开，漏掉这一层是失分最多的地方。第二，记号 $f_1',f_2'$ 表示对第 1、2 个\textbf{中间变量}求导，与 $\frac{\partial z}{\partial x}$ 完全不同，切勿混用。第三，$z=f(x+y,xy)$ 这类"和与积"结构的二阶导要写全含 $f_v'$ 的那一项，不要只写二次项。第四，使用链式法则前要先说明 $f$ 可微（题目一般已给"具有二阶连续偏导"）。 最后强调一句：本章所有"求偏导"的题，第一步都是\textbf{先分清自变量、中间变量与复合层次}，再动笔。树状图画出来，公式几乎不会写错；不画图直接写，二阶偏导十有八九要漏项。

\end{callwarn}

\begin{callexam}{真题考法}
解答题的标准形态：给出含抽象函数 $f$ 的复合表达式，要求一阶、二阶偏导；或者给出一个偏微分方程要你验证 $\frac{\partial z}{\partial x}$ 满足它。另一类是"由 $z=f(u,v)$ 满足某方程，求 $f$ 的偏导的某个关系"，思路仍是先写出一阶偏导再代入。陷阱集中在缺少二阶项与漏乘中间变量的导数。

\end{callexam}

\subsubsection{多元隐函数的求导法（含隐函数存在定理） $\bigstar$$\bigstar$$\bigstar$}

\textbf{一个方程确定一个隐函数（核心）}

设 $F(x,y,z)=0$ 确定了可微函数 $z=z(x,y)$，且 $F_z\ne 0$，则

\[\frac{\partial z}{\partial x}=-\frac{F_x}{F_z},\qquad \frac{\partial z}{\partial y}=-\frac{F_y}{F_z}\]

特别地，由 $F(x,y)=0$ 确定 $y=y(x)$ 时，$\frac{dy}{dx}=-\frac{F_x}{F_y}$。

\textbf{隐函数存在定理（了解）}：若 $F(x,y,z)$ 在点 $P_0(x_0,y_0,z_0)$ 的某邻域内有连续偏导数，$F(P_0)=0$，且 $F_z(P_0)\ne 0$，则方程 $F=0$ 在 $P_0$ 附近唯一确定一个连续可微函数 $z=z(x,y)$，使 $z_0=z(x_0,y_0)$，并有上面的偏导公式。\textbf{关键词是"$F$ 对该变量的偏导数不为零"}——分母不为零才是能"解出"该变量的条件。

\textbf{三种求法（都要会用）}

\begin{itemize}
\item \textbf{公式法}：把 $x,y,z$ 都当独立变量，分别求 $F_x,F_y,F_z$ 后代入。注意此时求 $F_x$ 时 $z$ 不随 $x$ 变，绝不能把 $z$ 当成 $x$ 的函数；
\item \textbf{两边求导法}：把 $z$ 视为 $x,y$ 的函数，方程两边对 $x$（或 $y$）求导，再把 $z$ 的偏导当作未知量解方程；
\item \textbf{全微分法}：对 $F(x,y,z)=0$ 取全微分得 $F_x\,dx+F_y\,dy+F_z\,dz=0$，整理出 $dz=z_x\,dx+z_y\,dy$，其系数即为所求。
\end{itemize}
数学二要求的是"一个方程确定的隐函数"。涉及两个方程确定的隐函数组时，用两边求导法把它化为线性方程组解出即可，思路相同，不必专门记忆行列式公式。

例 1：设 $x^{2}+y^{2}+z^{2}=4z$ 确定 $z=z(x,y)$，求 $\frac{\partial z}{\partial x}$。

令 $F=x^{2}+y^{2}+z^{2}-4z$，则 $F_x=2x,\ F_z=2z-4$，故

\[\frac{\partial z}{\partial x}=-\frac{2x}{2z-4}=\frac{x}{2-z}\]

验证：两边直接对 $x$ 求导得 $2x+2zz_x=4z_x$，即 $z_x=\frac{2x}{4-2z}=\frac{x}{2-z}$，一致。

例 2：设 $e^{z}-xyz=0$ 确定 $z=z(x,y)$，求 $\frac{\partial z}{\partial x},\frac{\partial z}{\partial y}$。

$F=e^{z}-xyz$，$F_x=-yz$，$F_y=-xz$，$F_z=e^{z}-xy$，故

\[\frac{\partial z}{\partial x}=\frac{yz}{e^{z}-xy},\qquad \frac{\partial z}{\partial y}=\frac{xz}{e^{z}-xy}\]

例 3：对上例继续求 $\frac{\partial^{2}z}{\partial x\partial y}$ 的基本思路。

把 $z_x=\frac{yz}{e^{z}-xy}$ 再对 $y$ 求导时，其中的 $z$ 与 $z_x$ 都是 $x,y$ 的函数，必须用商的求导法则并代入 $z_y$，最后结果中一般仍含 $z$，用原方程可继续化简。 例 4：设 $z=z(x,y)$ 由 $z^{3}-3xyz=1$ 确定，求 $z_x,z_y$。

$F=z^{3}-3xyz-1$，$F_x=-3yz$，$F_y=-3xz$，$F_z=3z^{2}-3xy$，故

\[z_x=\frac{yz}{z^{2}-xy},\qquad z_y=\frac{xz}{z^{2}-xy}\]

若题目要求某点 $(x_0,y_0,z_0)$ 处的偏导值，应先确认 $z_0$ 满足原方程再代入；\textbf{尽量在求完一阶偏导后一次性代点}，不要中途把 $z$ 代成具体数值而丢掉后续求导所需的函数关系。

\begin{callkey}{隐函数求导三条铁律}
$\circled{1}$ 公式法把 $x,y,z$ 都当独立变量；$\circled{2}$ 两边求导法把 $z$ 当 $x,y$ 的函数；$\circled{3}$ 只有 $F_z\ne 0$ 才能用公式，否则改用两边求导或说明该点不能确定隐函数。

\end{callkey}

\begin{callwarn}{易错点}
第一，公式法必须"三个变量都独立"，求 $F_x$ 时把 $z$ 当常数；两边求导法则相反，要把 $z$ 当成 $x,y$ 的函数。两种方法在一道题里混着用是高频错误。第二，$F_z=0$ 时公式失效，此时该点处不能确定隐函数，不能强行代入。第三，求二阶偏导时容易漏掉 $\frac{\partial z}{\partial x}$ 中 $z$ 对 $y$ 的依赖项。第四，结果中允许保留 $z$（只要能由原方程消去），不必强行解出 $z$。

\end{callwarn}

\begin{callexam}{真题考法}
解答题常见两种问法：一是求 $\frac{\partial z}{\partial x},\frac{\partial z}{\partial y}$ 在某点的值（先求偏导再代点）；二是由隐函数确定 $z$ 后验证某个恒等式（如 $x\frac{\partial z}{\partial x}+y\frac{\partial z}{\partial y}=z$）。陷阱是在等式两边求导时把 $z$ 当常数，或者用公式法时对 $z$ 也求了偏导。 还有一种问法把隐函数与极值结合："由方程 $F(x,y,z)=0$ 确定 $z=z(x,y)$，求 $z$ 的极值"。做法是先用驻点条件 $z_x=z_y=0$ 与 $F=0$ 联立求出候选点，再作判别。判断隐函数能否确定时，只需检查 $F_z$ 在该点是否非零。

\end{callexam}

\subsection{高阶偏导数}

\subsubsection{二阶偏导数与混合偏导相等的条件 $\bigstar$$\bigstar$$\bigstar$}

\textbf{定义}：对偏导数再求偏导即得二阶偏导数，共有四个：

\[f_{xx}=\frac{\partial}{\partial x}\left(\frac{\partial f}{\partial x}\right),\quad f_{xy}=\frac{\partial}{\partial y}\left(\frac{\partial f}{\partial x}\right),\quad f_{yx}=\frac{\partial}{\partial x}\left(\frac{\partial f}{\partial y}\right),\quad f_{yy}=\frac{\partial}{\partial y}\left(\frac{\partial f}{\partial y}\right)\]

\textbf{混合偏导相等的条件}：若 $f_{xy}$ 与 $f_{yx}$ 在区域 $D$ 内\textbf{连续}，则在 $D$ 内 $f_{xy}=f_{yx}$。初等函数在其定义域内满足这一条件，所以通常求二阶混合偏导时可以任选求导次序；但\textbf{分段函数在分界点必须分别用定义计算两个混合偏导}，二者可能不同。经典反例是

\[f(x,y)=\frac{xy(x^{2}-y^{2})}{x^{2}+y^{2}}\quad((x,y)\ne(0,0)),\qquad f(0,0)=0\]

在原点处 $f_{xy}(0,0)=-1$，而 $f_{yx}(0,0)=1$。

\textbf{计算策略}：先求较容易化简的那个一阶偏导，再求另一个；求二阶时优先选择让表达式变简单的次序。

例 1：设 $z=x^{3}y^{2}+e^{xy}$，求 $z_{xy}$。

\[z_x=3x^{2}y^{2}+ye^{xy}\ \Longrightarrow\ z_{xy}=6x^{2}y+e^{xy}+xye^{xy}=6x^{2}y+(1+xy)e^{xy}\]

例 2：验证 $z=\arctan\frac{y}{x}$ 满足 $z_{xx}+z_{yy}=0$。

\[z_x=-\frac{y}{x^{2}+y^{2}},\quad z_{xx}=\frac{2xy}{(x^{2}+y^{2})^{2}};\qquad z_y=\frac{x}{x^{2}+y^{2}},\quad z_{yy}=-\frac{2xy}{(x^{2}+y^{2})^{2}}\]

两式相加得 $0$，验证成立。 例 3：设 $u(x,y)=f(x+y)+g(x-y)$，其中 $f,g$ 二阶可导，验证 $u_{xx}-u_{yy}=0$。

\[u_x=f'+g',\qquad u_{xx}=f''+g'';\qquad u_y=f'-g',\qquad u_{yy}=f''+g''\]

两式相减得 $u_{xx}-u_{yy}=0$。

\textbf{分段函数在分界点求混合偏导的标准做法}：把 $f_{xy}(x_0,y_0)$ 写成 $\lim_{h\to 0}\frac{f_x(x_0,y_0+h)-f_x(x_0,y_0)}{h}$，其中 $f_x(x_0,y_0)$ 必须由定义求得，而 $h\ne 0$ 时的 $f_x(x_0,y_0+h)$ 用表达式求。两项来源不同，正是错误的来源。 补充一条常用结论：若 $f(x,y)$ 是关于 $x,y$ 的次数不超过 $n$ 的多项式，则 $n+1$ 阶及其以上的偏导恒为零。求高阶偏导前先判断"后面还会不会有非零项"，可以省下大量计算。 记号方面要读准：$\frac{\partial^{2}z}{\partial x\partial y}$ 读作"先对 $x$ 求偏导、再对 $y$ 求偏导"，也常写作 $z_{xy}$ 或 $f_{xy}$；而 $\frac{\partial^{2}z}{\partial y\partial x}=z_{yx}$，两者在连续时相等，但书写时不建议随意对调，以免与自己前面的推导脱节。

\begin{callwarn}{易错点}
第一，$f_{xy}$ 表示"先对 $x$ 后对 $y$"，与 $f_{yx}$ 的次序相反，读题要看清。第二，混合偏导相等是有条件的，遇到分段函数、分界点必须回到定义式分别计算。第三，复合函数的二阶偏导要展开成二次项加一次项（见上一考点公式），漏项最常见。第四，求 $z_{xx}$ 时若 $z_x$ 中含隐函数 $z$，该部分对 $x$ 求导不是常数，仍需继续求导。

\end{callwarn}

\begin{callexam}{真题考法}
填空题常直接要求某个二阶偏导（尤其混合偏导）的值；解答题则常以"验证 $u(x,y)$ 满足 $u_{xx}+u_{yy}=0$（调和方程）"或"化简某个二阶偏微分方程"的形式出现。陷阱集中在复合函数二阶导漏掉 $f_u'$ 的再求导项。 验证恒等式的题要把二阶偏导\textbf{算到底}再代回原式；中途用原方程化简易掩盖计算错误。若最终结果与待证式子差一个符号，多半是某一步漏了负号，应回到一阶偏导逐步检查，而不是反复调整末几步。

\end{callexam}

\section{多元函数的极值与最值}

极值问题分两类：\textbf{无条件极值}（在区域内部找驻点，用 $AC-B^{2}$ 判别）与\textbf{条件极值}（在约束下找最值，用拉格朗日乘数法）。闭区域上的最值问题则是两者的合成：内部驻点、边界条件极值、角点全部算出后比较大小。

\subsection{极值与条件极值}

\subsubsection{多元函数的极值、条件极值与拉格朗日乘数法 $\bigstar$$\bigstar$$\bigstar$$\bigstar$}

\paragraph{无条件极值}\mbox{}{\par}

\textbf{定义}：若存在 $P_0$ 的某邻域，使该邻域内一切点都有 $f(P)\le f(P_0)$（或 $\ge$），则称 $f$ 在 $P_0$ 取\textbf{极大值}（或极小值）。

\textbf{必要条件}：若 $f$ 在 $P_0$ 可微且取极值，则

\[f_x(x_0,y_0)=0,\qquad f_y(x_0,y_0)=0\]

满足该条件的点称为\textbf{驻点}。注意：偏导数不存在的点也可能是极值点，所以"驻点"与"极值点"是两个不同的概念。

\textbf{充分条件（二元函数的 $AC-B^{2}$ 判别法）}：设 $f_x(x_0,y_0)=f_y(x_0,y_0)=0$，记

\[A=f_{xx}(x_0,y_0),\qquad B=f_{xy}(x_0,y_0),\qquad C=f_{yy}(x_0,y_0)\]

则

\begin{itemize}
\item $AC-B^{2}>0$ 时，$P_0$ 是极值点：$A<0$ 为极大值，$A>0$ 为极小值；
\item $AC-B^{2}<0$ 时，$P_0$ 不是极值点；
\item $AC-B^{2}=0$ 时，判别法\textbf{失效}，需用定义、配方或其他方法另作判断。
\end{itemize}
\textbf{$AC-B^{2}=0$ 时怎么办}：判别法失效，只能另想办法，常用手段有三种。$\circled{1}$ 回到定义，直接比较 $f$ 与 $f(P_0)$ 的大小，例如 $f=x^{2}y^{2}$ 在原点，$f\ge 0=f(0,0)$，故为极小值；$\circled{2}$ 沿不同路径考察符号，若存在两条路径分别使 $f-f(P_0)$ 取正、取负，则不是极值点；$\circled{3}$ 配方成完全平方之和（或差）的形式再作判断。

\textbf{求解步骤}：$\circled{1}$ 解方程组 $f_x=0,\ f_y=0$ 求出全部驻点；$\circled{2}$ 计算 $A,B,C$；$\circled{3}$ 逐个判别并求出极值。

例：求 $f(x,y)=x^{3}+y^{3}-3xy$ 的极值。

\[f_x=3x^{2}-3y=0,\qquad f_y=3y^{2}-3x=0\]

由第一式 $y=x^{2}$ 代入第二式得 $x^{4}-x=0$，即 $x(x^{3}-1)=0$，驻点为 $(0,0)$ 与 $(1,1)$。又

\[A=f_{xx}=6x,\quad B=f_{xy}=-3,\quad C=f_{yy}=6y\]

在 $(0,0)$：$AC-B^{2}=0-9=-9<0$，非极值点。在 $(1,1)$：$AC-B^{2}=36-9=27>0$ 且 $A=6>0$，故取极小值

\[f(1,1)=1+1-3=-1\]

\paragraph{条件极值与拉格朗日乘数法}\mbox{}{\par}

\textbf{问题}：求目标函数 $z=f(x,y)$ 在约束条件 $\varphi(x,y)=0$ 下的极值。

\textbf{步骤}：

\begin{enumerate}
\item 构造拉格朗日函数 $L(x,y,\lambda)=f(x,y)+\lambda\varphi(x,y)$；
\item 解方程组
\end{enumerate}
\[L_x=f_x+\lambda\varphi_x=0,\qquad L_y=f_y+\lambda\varphi_y=0,\qquad L_\lambda=\varphi(x,y)=0\]

\begin{enumerate}
\item 所得 $(x,y)$ 都是\textbf{可能的}极值点，再结合实际问题或与端点值比较确定极值与最值。
\end{enumerate}
若有两个约束 $\varphi=0,\ \psi=0$，则构造 $L=f+\lambda\varphi+\mu\psi$，解 $L_x=L_y=L_\lambda=L_\mu=0$。

\textbf{实用技巧}

\begin{itemize}
\item \textbf{目标函数的等价变形}：把根号、指数、对数、绝对值换成更简单的单调形式，极值点不变：$\sqrt{u}\to u$，$e^{u}\to u$，$\ln u\to u$，$\vert u\vert\to u^{2}$；
\item \textbf{消 $\lambda$ 用除法时，必须单独讨论除式为零的情形}；
\item \textbf{边界不封闭}（如约束是线段、半圆）时，必须另外比较端点的函数值；
\item \textbf{能消元就消元}：若约束容易解出 $y=y(x)$，代入目标函数转化为一元函数极值，比乘数法更快；
\item \textbf{对称性}：若目标函数与约束关于 $x,y$ 对称，最值点常满足 $x=y$，可先用它简化方程。
\end{itemize}
例 1：求 $z=xy$ 在 $x+y=1$ 下的最大值。

\[L=xy+\lambda(x+y-1),\qquad L_x=y+\lambda=0,\quad L_y=x+\lambda=0,\quad x+y=1\]

由前两式得 $x=y$，代入约束得 $x=y=\frac12$，最大值 $z=\frac14$。

例 2：求 $z=xy$ 在 $\frac{x^{2}}{4}+y^{2}=1$ 上的最大值与最小值。

\[L=xy+\lambda\left(\frac{x^{2}}{4}+y^{2}-1\right)\]

\[L_x=y+\frac{\lambda x}{2}=0,\qquad L_y=x+2\lambda y=0\]

当 $x\ne 0,y\ne 0$ 时，由第一式得 $\lambda=-\frac{2y}{x}$，代入第二式得 $x-\frac{4y^{2}}{x}=0$，即 $x^{2}=4y^{2}$。与约束联立得 $y^{2}=\frac12$，$x^{2}=2$，且 $x,y$ 同号。故

\[z_{\max}=2\times\frac12=1,\qquad z_{\min}=-1\]

从例 2 可以总结出条件极值题的完整套路：\textbf{先化简目标函数、再写拉格朗日函数、然后解方程组、最后比较端点}。解方程组时多用"两式相比消去 $\lambda$"或"两式相减提出公因子"这两个技巧，能把二元二次方程组化为一元方程。

\begin{callkey}{极值必须默写的两条结论}
\begin{itemize}
\item \textbf{无条件极值}：可微且取极值 $\Rightarrow$ 偏导全为零（驻点）；判别用 $AC-B^{2}$：$>0$ 为极值（$A>0$ 极小、$A<0$ 极大），$<0$ 不是极值，$=0$ 判别法失效。
\item \textbf{条件极值}：构造 $L=f+\lambda\varphi$，解 $L_x=L_y=L_\lambda=0$；所得只是候选点，必须与端点、角点比较才能定最值。
\end{itemize}
\end{callkey}

\begin{callwarn}{易错点}
第一，拉格朗日乘数法给出的只是\textbf{必要条件}，解出的点必须与端点、边界点比较才能断言最值。第二，消 $\lambda$ 时两边同除某个因子（如 $x$ 或 $y$）之前要讨论该因子为零的情形，否则可能丢解。第三，约束是线段或半圆弧时，端点（角点）的函数值必须单独算。第四，$AC-B^{2}=0$ 时判别法失效，不能凭"大于等于零就是极值"下结论，要回到定义或配方，例如 $f=x^{2}y^{2}$ 在原点取极小值但 $AC-B^{2}=0$。

\end{callwarn}

\begin{callexam}{真题考法}
解答题的高频形态：求某函数在给定约束下的最值，或应用题（体积、表面积、距离最大最小）建模后用乘数法。得分关键在目标函数是否简化、方程组是否解全、端点是否补算。选择题则常考"某点是极大值点、极小值点还是非极值点"的判别，以及"驻点但非极值点"的构造。 做题顺序建议：求出一个驻点就立刻算出 $A,B,C$ 并判别，再处理下一个驻点。把多个驻点的计算堆到最后，容易张冠李戴地把某个点的 $A,B,C$ 用到另一个点上。

\end{callexam}

\subsection{最大值与最小值}

\subsubsection{多元函数的最大值、最小值及简单应用 $\bigstar$$\bigstar$$\bigstar$}

\textbf{存在性}：由有界闭区域上连续函数的性质，若 $f$ 在有界闭区域 $D$ 上连续，则 $f$ 在 $D$ 上必取得最大值与最小值。

\textbf{闭区域上求最值的三步法}

\begin{enumerate}
\item \textbf{内部驻点}：解 $f_x=f_y=0$，求出落在 $D$ 内部的驻点，算出函数值（不必判别是极大还是极小）；
\item \textbf{边界上的最值}：把边界曲线参数化代入 $f$ 化为一元函数，或用拉格朗日乘数法在边界约束下求候选点；\textbf{分段光滑的边界要逐段处理}；
\item \textbf{比较所有候选值}：把内部驻点、边界候选点以及角点（边界交点）的函数值放在一起，最大者为最大值，最小者为最小值。
\end{enumerate}
\textbf{应用题建模步骤}：$\circled{1}$ 设变量并写出目标函数（如面积、表面积、体积、距离）；$\circled{2}$ 找出约束等式；$\circled{3}$ 用拉格朗日乘数法或消元法求驻点；$\circled{4}$ 由实际意义判断是最大还是最小。距离最值常转化为距离平方的最值，以避开根号。

例 1：求 $z=xy$ 在闭区域 $D=\{(x,y)\mid x^{2}+y^{2}\le 1,\ x\ge 0,\ y\ge 0\}$ 上的最大值与最小值。

内部：$z_x=y=0,\ z_y=x=0$ 得驻点 $(0,0)$，$z=0$。

边界：在圆弧上令 $x=\cos t,\ y=\sin t,\ t\in[0,\frac{\pi}{2}]$，则 $z=\sin t\cos t=\frac12\sin 2t\le \frac12$，在 $t=\frac{\pi}{4}$ 处取到；在两条直径 $x=0$ 与 $y=0$ 上 $z=0$。

比较得 $z_{\max}=\frac12$，$z_{\min}=0$。

例 2：设长方体三边长 $x,y,z$，体积 $V=xyz$ 固定，求表面积 $S=2(xy+yz+zx)$ 的最小值。

构造 $L=2(xy+yz+zx)+\lambda(xyz-V)$，令 $L_x=2(y+z)+\lambda yz=0$ 并写出轮换的两式。两两相减可得 $(y-z)(2-\lambda x)=0$ 之类的等式，结合实际问题排除 $2-\lambda x=0$（否则与 $L_x=0$ 矛盾），得 $x=y=z=V^{1/3}$。故当长方体为正方体时表面积最小，$S_{\min}=6V^{2/3}$。 例 3：求 $f(x,y)=x^{2}+2y^{2}-x$ 在闭圆盘 $D:\ x^{2}+y^{2}\le 1$ 上的最值。

内部：$f_x=2x-1=0,\ f_y=4y=0$，得驻点 $\left(\frac12,0\right)$，$f=-\frac14$。

边界 $x^{2}+y^{2}=1$：令 $x=\cos t,\ y=\sin t$，则

\[f=\cos^{2}t+2\sin^{2}t-\cos t=2-\cos^{2}t-\cos t\]

记 $u=\cos t\in[-1,1]$，则 $f=2-u^{2}-u$，这是开口向下的二次函数，在 $u=-\frac12$ 处取最大值 $\frac94$，在 $u=1$ 处取最小值 $0$。

比较三个候选值 $-\frac14,\ \frac94,\ 0$，得 $f_{\max}=\frac94$，$f_{\min}=-\frac14$。 从例 3 可以总结出闭区域求最值的完整写法：\textbf{内部驻点、边界最值、端点（角点）三部分各自求完再比较}，缺任何一部分都可能答错。边界是圆或椭圆时，参数化后化为三角函数的最值往往比乘数法更快；边界是直线段时，直接解出其中一个变量代入即可。

\begin{callwarn}{易错点}
第一，只求内部驻点就作答是最常见的错误——最值往往出现在边界上。第二，闭区域的分段边界要逐段求候选点，边界交角点不能漏。第三，参数化边界时参数范围要与实际边界一致（如上半圆对应 $t\in[0,\pi]$）。第四，无界区域上最值可能不存在，需考察函数在无穷远处的行为。第五，目标函数取平方或取对数后，结论必须回到原函数（如距离平方的最小值点即距离最小点，但最小值本身要开方）。

\end{callwarn}

\begin{callexam}{真题考法}
应用题集中在几何与极值：过定点的最大体积、用料最省的容器、点到曲线（曲面）的最短距离、矩形绕边旋转体积最大等。解题套路固定为"建模—约束—乘数法—判断"，答题时要写出"由实际问题，该驻点即为最值点"这一句。纯计算题则常把区域设为三角形、圆盘、椭圆的截段。 边界是最值题的主要失分点：圆盘用参数化，矩形用四条边分别讨论，三角形用三条边加三个顶点。牢记"边界上的最值只可能在端点或边界驻点处取到"，先把候选点列全，再比较函数值。

\end{callexam}

\section{二重积分}

二重积分是本章的绝对重点，数学二几乎每年必考，且常常是解答题中占比最高的一题。计算的通用流程只有四步：\textbf{画区域 $\to$ 判对称 $\to$ 选坐标系 $\to$ 定限计算}。其中"画区域"和"定限"决定成败，务必练到看到边界方程就能立刻画出草图。

\begin{callkey}{二重积分计算总流程（四步）}
$\circled{1}$ \textbf{画区域}：求交点、画边界、说明属于 $X$ 型、$Y$ 型还是需要分块； $\circled{2}$ \textbf{判对称}：区域关于坐标轴或 $y=x$ 对称时，先用偶倍奇零与轮换对称把能划掉的项划掉； $\circled{3}$ \textbf{选坐标系}：区域或函数含 $x^{2}+y^{2}$、$\frac yx$ 就上极坐标，其余用直角坐标； $\circled{4}$ \textbf{定限计算}：外层限为常数、内层限为外层变量的函数，下限小于上限，再逐次积分。 四步的顺序不能颠倒：先判对称再定限，往往能把计算量减半。

\end{callkey}

\subsection{二重积分的概念与性质}

\subsubsection{二重积分的概念、基本性质与中值定理 $\bigstar$$\bigstar$$\bigstar$}

\textbf{定义}：设 $f(x,y)$ 在有界闭区域 $D$ 上有界，把 $D$ 任意分割成 $n$ 个小区域 $\Delta\sigma_i$，在每个小区域上任取一点 $(\xi_i,\eta_i)$，作和式 $\sum_{i=1}^{n}f(\xi_i,\eta_i)\Delta\sigma_i$。若当各小区域直径的最大值 $\lambda\to 0$ 时该和式的极限存在，且与分割方式、取点方式无关，则称此极限为 $f$ 在 $D$ 上的二重积分：

\[\iint_{D}f(x,y)\,d\sigma=\iint_{D}f(x,y)\,dxdy=\lim_{\lambda\to 0}\sum_{i=1}^{n}f(\xi_i,\eta_i)\Delta\sigma_i\]

\textbf{可积性}：有界闭区域上的连续函数必可积。

\textbf{几何意义}：当 $f(x,y)\ge 0$ 时，$\iint_{D}f\,d\sigma$ 表示以 $D$ 为底、以曲面 $z=f(x,y)$ 为顶的\textbf{曲顶柱体的体积}；当 $f\equiv 1$ 时，$\iint_{D}1\,d\sigma$ 等于区域 $D$ 的\textbf{面积}。这两条是几何应用题与结果验算的依据。

\textbf{基本性质}（设 $f,g$ 在 $D$ 上可积）

\begin{itemize}
\item \textbf{线性性}：$\iint_{D}\left[\alpha f+\beta g\right]d\sigma=\alpha\iint_{D}f\,d\sigma+\beta\iint_{D}g\,d\sigma$；
\item \textbf{区域可加性}：若 $D=D_1\cup D_2$ 且 $D_1,D_2$ 至多沿边界相交，则 $\iint_{D}=\iint_{D_1}+\iint_{D_2}$；
\item \textbf{保序性}：若 $f\le g$，则 $\iint_{D}f\,d\sigma\le \iint_{D}g\,d\sigma$；
\item \textbf{绝对值不等式}：$\left\vert \iint_{D}f\,d\sigma\right\vert \le \iint_{D}\vert f\vert\,d\sigma$；
\item \textbf{面积公式}：$\iint_{D}1\,d\sigma=S_D$。
\end{itemize}
\textbf{中值定理}：若 $f(x,y)$ 在有界闭区域 $D$ 上连续，则存在 $(\xi,\eta)\in D$，使

\[\iint_{D}f(x,y)\,d\sigma=f(\xi,\eta)\cdot S_D\]

其中 $S_D$ 是 $D$ 的面积。它常用于求含二重积分的极限：把积分整体替换为 $f(\xi,\eta)S_D$，再让区域收缩到一个点，$f(\xi,\eta)$ 就被迫趋于该点的函数值。

例 1：$\iint_{D}1\,d\sigma$，其中 $D:\ x^{2}+y^{2}\le 1$，其值为 $\pi$（单位圆面积）。

例 2：设 $f$ 连续，求 $\displaystyle\lim_{t\to 0^{+}}\frac{1}{\pi t^{2}}\iint_{x^{2}+y^{2}\le t^{2}}(x+y+2)\,d\sigma$。

由中值定理，积分为 $\left.(x+y+2)\right\vert_{(\xi,\eta)}\cdot\pi t^{2}$，故原式 $=\xi+\eta+2$。当 $t\to 0^{+}$ 时 $(\xi,\eta)\to(0,0)$，极限为 $2$。 例 3（比较大小）：设 $D:\ x^{2}+y^{2}\le 1$，比较 $I_1=\iint_{D}(x^{2}+y^{2})^{2}d\sigma$ 与 $I_2=\iint_{D}(x^{2}+y^{2})^{3}d\sigma$ 的大小。因为在 $D$ 上恒有 $0\le x^{2}+y^{2}\le 1$，所以 $(x^{2}+y^{2})^{3}\le(x^{2}+y^{2})^{2}$，由保序性得 $I_2\le I_1$。

例 4（用定义求 $n$ 项和的极限）：求 $\displaystyle\lim_{n\to\infty}\frac{1}{n^{2}}\sum_{i=1}^{n}\sum_{j=1}^{n}\sin\frac{i+j}{n}$。按对应关系 $\frac in\to x$、$\frac jn\to y$、$\frac{1}{n^{2}}\to dxdy$，原式 $=\iint_{[0,1]\times[0,1]}\sin(x+y)\,dxdy$，再化为累次积分计算即可。这类题的关键是先把和式整理成"$\frac{1}{n^{2}}\sum\sum$"的标准形态。

\begin{callkey}{四条基本性质的对应题型}
比较大小用\textbf{保序性}；拆分区域用\textbf{可加性}；验算结果用"$f\equiv 1$ 时为面积"；求极限用\textbf{中值定理}。四条性质分别对应四类小题，考场上先判断题目属于哪一类再动笔。

\end{callkey}

\begin{callwarn}{易错点}
中值定理要求 $f$ 在 $D$ 上\textbf{连续}，且 $(\xi,\eta)$ 只是 $D$ 内某点，位置不能确定——所以在极限题中要靠"区域收缩迫使该点趋于目标点"来收尾，不能把 $(\xi,\eta)$ 直接当作 $(0,0)$ 代入而不说明理由。此外，用中值定理求极限时若被积函数的极限值为零，中值定理可能失效，此时应改用交换积分次序。

\end{callwarn}

\begin{callexam}{真题考法}
选择题常考性质比较大小：在同一区域上比较两个含参数二重积分的大小（用保序性），或给一个积分值与区域面积，要求推断函数的取值范围。另有一类是二重积分定义求 $n$ 项和的极限——把 $\frac{1}{n^{2}}\sum\sum$ 改写成 $\iint_{[0,1]\times[0,1]}$，对应关系是 $\frac in\to x,\ \frac jn\to y,\ \frac{1}{n^{2}}\to dxdy$。 用中值定理求这类极限时，必须写清"由 $f$ 连续及二重积分中值定理，存在 $(\xi,\eta)\in D$ 使……"，再说明 $(\xi,\eta)\to(0,0)$ 的理由；过程完整才能得满分，只写结论不给依据会被扣分。

\end{callexam}

\subsection{二重积分的计算}

\subsubsection{二重积分的计算：直角坐标、交换积分次序 $\bigstar$$\bigstar$$\bigstar$$\bigstar$$\bigstar$}

\textbf{两类区域与两个公式}

\begin{itemize}
\item \textbf{$X$ 型区域}（上下由曲线夹、左右由直线夹）：$D=\{(x,y)\mid a\le x\le b,\ \varphi_1(x)\le y\le \varphi_2(x)\}$，则
\end{itemize}
\[\iint_{D}f(x,y)\,dxdy=\int_{a}^{b}dx\int_{\varphi_1(x)}^{\varphi_2(x)}f(x,y)\,dy\]

\begin{itemize}
\item \textbf{$Y$ 型区域}（左右由曲线夹、上下由直线夹）：$D=\{(x,y)\mid c\le y\le d,\ \psi_1(y)\le x\le \psi_2(y)\}$，则
\end{itemize}
\[\iint_{D}f(x,y)\,dxdy=\int_{c}^{d}dy\int_{\psi_1(y)}^{\psi_2(y)}f(x,y)\,dx\]

\textbf{画图定限四步（穿针定限法）}

\begin{enumerate}
\item 画出 $D$ 的所有边界曲线，求出它们的交点，明确区域的位置与形状；
\item 判断 $D$ 是否同时是 $X$ 型与 $Y$ 型，还是需要分块；
\item 定限：\textbf{外层积分的上下限必须是常数，内层积分的上下限是外层变量的函数，且必须满足"下限小、上限大"}（从左到右、从下到上）；
\item 先算内层（对另一个变量视为常数），再算外层。
\end{enumerate}
\textbf{画区域的操作顺序}：先分别令 $x=0$、$y=0$ 求边界曲线与坐标轴的交点，再联立相邻两条边界方程求交点；遇到含有 $x^{2}+y^{2}$、$\vert x\vert+\vert y\vert$、$\frac{x^{2}}{a^{2}}+\frac{y^{2}}{b^{2}}$ 的边界，就直接按圆、菱形、椭圆来画。区域画不准时用"描特殊点 + 判对称"两步确认。

\begin{callwarn}{易错点}
画图时最易犯的错误是把"函数图像"与"积分区域"混为一谈：$D$ 是 $xOy$ 面上的区域，而 $z=f(x,y)$ 是空间的曲面；累次积分的上下限只由 $D$ 的边界决定，与曲面形状无关。

\end{callwarn}

\textbf{选择积分次序的三条原则}

\begin{itemize}
\item \textbf{内层好积}：形如 $\int e^{x^{2}}dx$、$\int\frac{\sin y}{y}dy$ 这种"积不出来"的表达式，要设法让它作为\textbf{外层}积分变量被处理，即交换次序；
\item \textbf{分块要少}：宁可选能一次表示区域的次序，避免把 $D$ 拆成两块；
\item \textbf{积分限简单}：优先选上下限表达式更简单的次序。
\end{itemize}
\textbf{交换积分次序的步骤}：$\circled{1}$ 由原累次积分的上下限写出区域 $D$ 的不等式表示；$\circled{2}$ 画出 $D$ 的草图；$\circled{3}$ 按新的次序重新确定内外限（即把 $D$ 改写成另一种类型）；$\circled{4}$ 写出新的累次积分并计算。\textbf{交换次序后必须重新画图定限，绝不能把原来的限互换位置照抄。}

例 1：计算 $\iint_{D}xy\,dxdy$，其中 $D$ 由 $y=x^{2}$ 与 $y=x$ 围成。

两条曲线交于 $(0,0)$ 与 $(1,1)$，且 $D$ 为 $X$ 型：$0\le x\le 1,\ x^{2}\le y\le x$。

\[\iint_{D}xy\,dxdy=\int_{0}^{1}x\,dx\int_{x^{2}}^{x}y\,dy=\int_{0}^{1}x\cdot\frac{x^{2}-x^{4}}{2}dx=\frac12\left(\frac14-\frac16\right)=\frac{1}{24}\]

例 2（换序经典）：计算 $\displaystyle\int_{0}^{1}dx\int_{x}^{1}e^{-y^{2}}dy$。

内层 $\int e^{-y^{2}}dy$ 积不出来。区域 $D:\ 0\le x\le 1,\ x\le y\le 1$，改写为 $Y$ 型：$0\le y\le 1,\ 0\le x\le y$。于是

\[\int_{0}^{1}dx\int_{x}^{1}e^{-y^{2}}dy=\int_{0}^{1}dy\int_{0}^{y}e^{-y^{2}}dx=\int_{0}^{1}ye^{-y^{2}}dy=\left[-\frac12e^{-y^{2}}\right]_{0}^{1}=\frac{1-e^{-1}}{2}\]

例 3：用面积验算定限。计算 $\displaystyle\int_{0}^{1}dx\int_{x^{2}}^{x}dy$，得

\[\int_{0}^{1}\left(x-x^{2}\right)dx=\frac12-\frac13=\frac16\]

其几何意义是 $y=x^{2}$ 与 $y=x$ 所围区域的面积，可与例 1 的区域互相印证。 例 4：计算 $\displaystyle\iint_{D}y\,dxdy$，其中 $D$ 由 $y=x,\ y=2x,\ x=2,\ x=4$ 围成。

$D$ 为 $X$ 型：$2\le x\le 4,\ x\le y\le 2x$，故

\[\iint_{D}y\,dxdy=\int_{2}^{4}dx\int_{x}^{2x}y\,dy=\int_{2}^{4}\frac{4x^{2}-x^{2}}{2}dx=\frac32\int_{2}^{4}x^{2}dx=\frac32\cdot\frac{64-8}{3}=28\]

例 5：计算 $\displaystyle\iint_{D}x\,dxdy$，其中 $D$ 由 $y=x^{2}$ 与 $x=y^{2}$ 围成。

按 $Y$ 型定限：$0\le y\le 1,\ y^{2}\le x\le \sqrt{y}$，故

\[\iint_{D}x\,dxdy=\int_{0}^{1}dy\int_{y^{2}}^{\sqrt{y}}x\,dx=\frac12\int_{0}^{1}\left(y-y^{4}\right)dy=\frac12\left(\frac12-\frac15\right)=\frac{3}{20}\]

同一区域按 $X$ 型定限为 $0\le x\le 1,\ x^{2}\le y\le\sqrt{x}$，两种次序都能算且结果相同，可作为定限是否写对的验算手段。

\begin{callkey}{换序题的识别}
看到内层是 $\int e^{x^{2}}dx$、$\int\frac{\sin x}{x}dx$、$\int\frac{1}{\ln x}dx$、$\int\sqrt{1+y^{3}}dy$ 这类"原函数不是初等函数"的积分，就必须换序。标准写法三步：由原限写出 $D$ 的不等式 $\to$ 画出 $D$ $\to$ 按新次序重新定限。

\end{callkey}

\begin{callwarn}{易错点}
第一，内外层上下限必须"下限小、上限大"，写成 $\int_{x}^{x^{2}}$ 会差一个负号；外层限必须是常数，写成含 $y$ 的表达式直接出错。第二，交换次序时不重新画图、直接把上下限对调是最常见的失分点。第三，被积函数可分离时（如 $f(x)g(y)$ 在矩形域上）可直接拆成两个定积分之积，但非矩形域不能这样拆。第四，分块积分时各块要"不重不漏"，边界上的曲线归属任一块即可（边界面积为零不影响积分值）。第五，先积内层时把外层变量视为常数，不要顺手"消掉"含外层变量的因子。

\end{callwarn}

\begin{callexam}{真题考法}
解答题必考，两种题型：直接计算给定区域的二重积分；给出一个内层积不出的累次积分，要求交换次序后计算。第二类几乎是在明示"必须换序"，看到 $\int e^{x^{2}}dx$、$\int\frac{\sin x}{x}dx$、$\int\frac{1}{\ln x}dx$ 这类内层就要立刻想到换序。此外还有"交换积分次序并求值"的填空题。 计算二重积分的完整答题模板是：$\circled{1}$ 画区域草图并说明它属于哪一型；$\circled{2}$ 写出累次积分；$\circled{3}$ 逐步计算；$\circled{4}$ 给出结果。若是换序题，新老两个区域都要画出来，这是阅卷的给分点。

\end{callexam}

\subsubsection{二重积分的计算：极坐标 $\bigstar$$\bigstar$$\bigstar$$\bigstar$$\bigstar$}

\textbf{变换公式}：令 $x=r\cos\theta,\ y=r\sin\theta$，则面积元素变为

\[d\sigma=dxdy=r\,dr\,d\theta\]

于是

\[\iint_{D}f(x,y)\,dxdy=\iint_{D'}f(r\cos\theta,r\sin\theta)\,r\,dr\,d\theta\]

面积元素的来源：极坐标下的小区域近似为边长分别是 $dr$ 与 $r\,d\theta$ 的小矩形（弧形区域在无限小时趋于矩形），故面积为 $r\,dr\,d\theta$。\textbf{$r$ 这个因子是极坐标法最易漏的地方}。

\textbf{常规定限（先 $r$ 后 $\theta$）}：若 $D'=\{(r,\theta)\mid \alpha\le\theta\le\beta,\ r_1(\theta)\le r\le r_2(\theta)\}$，则

\[\iint_{D}f\,d\sigma=\int_{\alpha}^{\beta}d\theta\int_{r_1(\theta)}^{r_2(\theta)}f(r\cos\theta,r\sin\theta)\,r\,dr\]

定限方法：从原点出发引一条射线，看它穿过区域时 $\theta$ 的取值范围（外层），再看沿这条射线 $r$ 从哪条曲线进入、从哪条曲线离开（内层）。区域包含原点时内层下限为 $0$。

\textbf{二次定限（先 $\theta$ 后 $r$，即极坐标下的换序）}：任取半径 $r$ 作圆弧与区域相交，看交点对应的 $\theta$ 范围。这属于极坐标中的"换序"情形，出现频率低于先 $r$ 后 $\theta$，但题目若给出先 $\theta$ 后 $r$ 的累次积分，就必须按此法重新定限。

\textbf{适用极坐标的四种信号}

\begin{enumerate}
\item 积分区域是圆、圆环、扇形，或圆心不在原点的圆；
\item 被积函数中含 $x^{2}+y^{2}$，或形如 $f\left(\frac yx\right)$ 的齐次结构（齐次结构在极坐标下 $r$ 会被消去）；
\item 区域边界以极坐标方程给出：心形线 $r=a(1+\cos\theta)$、双纽线 $r^{2}=a^{2}\cos 2\theta$、玫瑰线 $r=a\cos n\theta$（$n$ 为奇数时 $n$ 叶，偶数时 $2n$ 叶）；
\item 区域关于原点有旋转对称性。
\end{enumerate}
\textbf{偏心圆口诀（务必记牢）}：圆心在 $x$ 轴上的圆 $(x-a)^{2}+y^{2}=a^{2}$ 化为 $r=2a\cos\theta$；圆心在 $y$ 轴上的圆 $x^{2}+(y-a)^{2}=a^{2}$ 化为 $r=2a\sin\theta$。规律是"$r=$ 符号 $\times$ 直径 $\times$（$\cos$ 配 $x$ 轴、$\sin$ 配 $y$ 轴）"：系数 $2a$ 是\textbf{直径}，$a$ 是半径；$\cos$ 对应圆心在 $x$ 轴、$\sin$ 对应圆心在 $y$ 轴；圆心在正半轴取正号，负半轴取负号。使用 $r=2a\cos\theta$ 时 $\theta\in\left[-\frac{\pi}{2},\frac{\pi}{2}\right]$，否则 $r$ 会取负值。椭圆域 $\frac{x^{2}}{a^{2}}+\frac{y^{2}}{b^{2}}\le 1$ 可用广义极坐标 $x=ar\cos\theta,\ y=br\sin\theta$（此时 $d\sigma=abr\,dr\,d\theta$），作为极坐标的延伸备用。

例 1：计算 $\iint_{D}(x^{2}+y^{2})\,d\sigma$，其中 $D:\ x^{2}+y^{2}\le 2x$。

区域是圆心 $(1,0)$、半径 $1$ 的圆，即 $r\le 2\cos\theta$，$\theta\in\left[-\frac{\pi}{2},\frac{\pi}{2}\right]$。于是

\[\iint_{D}(x^{2}+y^{2})\,d\sigma=\int_{-\pi/2}^{\pi/2}d\theta\int_{0}^{2\cos\theta}r^{2}\cdot r\,dr=\int_{-\pi/2}^{\pi/2}\frac{16\cos^{4}\theta}{4}d\theta=4\int_{-\pi/2}^{\pi/2}\cos^{4}\theta\,d\theta=4\cdot\frac{3\pi}{8}=\frac{3\pi}{2}\]

例 2：计算 $\iint_{D}e^{-(x^{2}+y^{2})}d\sigma$，其中 $D:\ x^{2}+y^{2}\le a^{2}$。

\[\int_{0}^{2\pi}d\theta\int_{0}^{a}e^{-r^{2}}r\,dr=2\pi\cdot\left[-\frac12e^{-r^{2}}\right]_{0}^{a}=\pi\left(1-e^{-a^{2}}\right)\]

此例说明极坐标能把"积不出来的 $e^{-x^{2}}$"化为可直接积分的 $e^{-r^{2}}\cdot r$，这是极坐标不可替代的价值。

例 3：计算 $\displaystyle\int_{0}^{\pi/2}d\theta\int_{0}^{2\sin\theta}r\,dr$。

先算内层：$\int_{0}^{2\sin\theta}r\,dr=2\sin^{2}\theta$；再算外层：$\int_{0}^{\pi/2}2\sin^{2}\theta\,d\theta=2\cdot\frac{\pi}{4}=\frac{\pi}{2}$。该积分的几何意义是圆 $x^{2}+y^{2}\le 2y$ 在第一象限部分的面积，等于单位圆面积 $\pi$ 的一半，结果吻合，可用于验算。 例 4：计算 $\displaystyle\iint_{D}\sqrt{x^{2}+y^{2}}\,d\sigma$，其中 $D$ 为心形线 $r=1+\cos\theta$ 围成的区域。

\[\int_{0}^{2\pi}d\theta\int_{0}^{1+\cos\theta}r\cdot r\,dr=\frac13\int_{0}^{2\pi}(1+\cos\theta)^{3}d\theta\]

展开 $(1+\cos\theta)^{3}=1+3\cos\theta+3\cos^{2}\theta+\cos^{3}\theta$，在 $[0,2\pi]$ 上 $\int_{0}^{2\pi}\cos\theta\,d\theta=0$、$\int_{0}^{2\pi}\cos^{3}\theta\,d\theta=0$、$\int_{0}^{2\pi}\cos^{2}\theta\,d\theta=\pi$，故原积分 $=\frac13\left(2\pi+3\pi\right)=\frac{5\pi}{3}$。

例 5（极坐标下的换序）：把 $\displaystyle\int_{0}^{\pi/2}d\theta\int_{0}^{2\cos\theta}f(r\cos\theta,r\sin\theta)\,r\,dr$ 化成先 $r$ 后 $\theta$ 的累次积分。

该区域是圆 $x^{2}+y^{2}\le 2x$ 在第一象限的部分。改为固定 $r$：当 $r\in[0,2]$ 时，满足 $r\le 2\cos\theta$ 即 $\cos\theta\ge\frac r2$ 的 $\theta$ 范围是 $0\le\theta\le\arccos\frac r2$，故

\[\int_{0}^{2}dr\int_{0}^{\arccos(r/2)}f(r\cos\theta,r\sin\theta)\,r\,d\theta\]

用 $f\equiv 1$ 验算：两种次序算出的面积都等于圆面积的一半 $\frac{\pi}{2}$。

\begin{callkey}{极坐标要背熟的两个"半径"}
圆 $x^{2}+y^{2}=2ax$ 给 $r=2a\cos\theta$，$\theta\in\left[-\frac{\pi}{2},\frac{\pi}{2}\right]$；圆 $x^{2}+y^{2}=2ay$ 给 $r=2a\sin\theta$，$\theta\in[0,\pi]$。这两个式子覆盖绝大多数偏心圆考题，记牢可省去大量推导。

\end{callkey}

\begin{callwarn}{易错点}
第一，忘记乘 $r$ 是极坐标法第一大错，$d\sigma=r\,dr\,d\theta$ 中的 $r$ 不能省。第二，$r=2\cos\theta$ 之类的偏心圆，$\theta$ 的范围是 $\left[-\frac{\pi}{2},\frac{\pi}{2}\right]$，若写成 $[0,2\pi]$ 会出现 $r<0$ 的无意义区间。第三，内层限是 $r$ 的函数（表示从哪条曲线进、哪条曲线出），外层限必须是常数。第四，只有当区域与函数都适合极坐标时才换，矩形域上被积函数不含 $x^{2}+y^{2}$ 时用极坐标反而更繁。第五，$\cos^{4}\theta$ 这类函数在对称区间上要善用偶函数性质：$\int_{-\pi/2}^{\pi/2}\cos^{4}\theta\,d\theta=2\int_{0}^{\pi/2}\cos^{4}\theta\,d\theta$。

\end{callwarn}

\begin{callexam}{真题考法}
解答题的高频设置：圆域、扇环域、偏心圆域上的二重积分；被积函数含 $x^{2}+y^{2}$、$\frac{y}{x}$、$\arctan\frac yx$。另一类是"给直角坐标的累次积分要求改用极坐标计算"，答题时必须把区域画出来重新定限。极坐标与对称性联用是压轴常见组合。 改用极坐标时要把变换的三件事都写出来：$x=r\cos\theta$、$y=r\sin\theta$、$d\sigma=r\,dr\,d\theta$，缺任何一件都可能被判跳步。定限之后，若内层是关于 $r$ 的幂函数，先算 $r$ 的积分通常能大幅简化后面的三角运算。

\end{callexam}

\subsection{二重积分的化简与分块}

\subsubsection{二重积分的对称性、奇偶性与轮换对称性化简 $\bigstar$$\bigstar$$\bigstar$$\bigstar$}

对称性是二重积分中"少算一半甚至直接得零"的技巧，用之前必须\textbf{同时}核对区域与函数两个条件。

\textbf{普通对称性（偶倍奇零）}

若 $D$ 关于 $y$ 轴对称（即 $(x,y)\in D\Leftrightarrow(-x,y)\in D$），则按被积函数\textbf{关于 $x$ 的奇偶性}处理：

\[\iint_{D}f(x,y)\,d\sigma=\begin{cases}0, & f(-x,y)=-f(x,y)\\[4pt] 2\displaystyle\iint_{D_1}f(x,y)\,d\sigma, & f(-x,y)=f(x,y)\end{cases}\]

其中 $D_1$ 是 $D$ 位于 $x\ge 0$ 的部分。若 $D$ 关于 $x$ 轴对称，则看关于 $y$ 的奇偶性，结论同理。

记忆要点：\textbf{区域关于哪条轴对称，就看函数关于哪个变量的奇偶性}；奇函数得零，偶函数折半翻倍。

\textbf{轮换对称性}：若把 $D$ 中的 $x$ 与 $y$ 互换后区域不变（即 $D$ 关于直线 $y=x$ 对称），则

\[\iint_{D}f(x,y)\,d\sigma=\iint_{D}f(y,x)\,d\sigma\]

特别地 $\iint_{D}g(x)\,d\sigma=\iint_{D}g(y)\,d\sigma$。这常用来"配对求和"：例如在关于 $y=x$ 对称的区域上，

\[\iint_{D}(x+y)\,d\sigma=2\iint_{D}x\,d\sigma\]

\textbf{判对称的操作方法}：把区域的不等式组与边界方程中的 $x$ 换成 $-x$，若整体不变则关于 $y$ 轴对称；把 $y$ 换成 $-y$ 不变则关于 $x$ 轴对称；把 $x,y$ 互换不变则关于 $y=x$ 对称；把 $(x,y)$ 换成 $(-x,-y)$ 不变则关于原点中心对称。

\textbf{形心公式（线性被积函数的速算工具）}：设 $D$ 的面积为 $S$，形心坐标为 $(\bar x,\bar y)$，则

\[\iint_{D}x\,d\sigma=\bar x\,S,\qquad \iint_{D}y\,d\sigma=\bar y\,S\]

于是对线性函数 $f(x,y)=ax+by+c$ 有

\[\iint_{D}(ax+by+c)\,d\sigma=(a\bar x+b\bar y+c)\,S\]

规则图形（圆、矩形、三角形、正多边形）的形心是对称轴的交点，无需另算。

例 1：计算 $\displaystyle\iint_{D}\left(2x+3y^{2}+5\right)d\sigma$，其中 $D:\ x^{2}+y^{2}\le 1$。

$D$ 关于 $y$ 轴对称，而 $2x$ 关于 $x$ 为奇函数，故 $\iint_{D}2x\,d\sigma=0$。又 $D$ 关于 $y=x$ 对称，故 $\iint_{D}x^{2}d\sigma=\iint_{D}y^{2}d\sigma$，而

\[\iint_{D}\left(x^{2}+y^{2}\right)d\sigma=\int_{0}^{2\pi}d\theta\int_{0}^{1}r^{2}\cdot r\,dr=2\pi\cdot\frac14=\frac{\pi}{2}\]

所以 $\iint_{D}y^{2}d\sigma=\frac{\pi}{4}$。故原式 $=3\cdot\frac{\pi}{4}+5\pi=\frac{23\pi}{4}$。

例 2：设 $D$ 为 $x+y\le 1,\ x\ge 0,\ y\ge 0$ 围成的三角形，计算 $\displaystyle\iint_{D}(6x+3y)\,d\sigma$。

$D$ 是直角三角形，形心为 $\left(\frac13,\frac13\right)$，面积 $S=\frac12$。由形心公式

\[\iint_{D}x\,d\sigma=\frac13\cdot\frac12=\frac16,\qquad \iint_{D}y\,d\sigma=\frac16\]

故原式 $=6\cdot\frac16+3\cdot\frac16=\frac32$。直接按累次积分计算也得同样结果，可作验算。 例 3：计算 $\displaystyle\iint_{D}\left(x^{3}y+\vert x\vert y^{2}\right)d\sigma$，其中 $D:\ x^{2}+y^{2}\le 1$。

$D$ 关于 $y$ 轴对称，而 $x^{3}y$ 关于 $x$ 为奇函数，故 $\iint_{D}x^{3}y\,d\sigma=0$。又 $D$ 关于 $x$ 轴与 $y$ 轴都对称，$\vert x\vert y^{2}$ 关于 $x$、$y$ 均为偶函数，于是可以四等分折到第一象限再乘 $4$：

\[\iint_{D}\vert x\vert y^{2}d\sigma=4\int_{0}^{\pi/2}d\theta\int_{0}^{1}\left(r\cos\theta\right)\left(r^{2}\sin^{2}\theta\right)r\,dr=\frac45\int_{0}^{\pi/2}\cos\theta\sin^{2}\theta\,d\theta=\frac45\cdot\frac13=\frac{4}{15}\]

\begin{callkey}{对称性的使用顺序}
先判区域对称类型（关于 $x$ 轴、$y$ 轴、原点、直线 $y=x$），再判被积函数相应变量的奇偶性；能得零的直接划掉，能折半的折半，最后才动手积分。若区域关于 $y=x$ 对称，还可把 $\iint_{D}x\,d\sigma$ 与 $\iint_{D}y\,d\sigma$ 合并成 $\frac12\iint_{D}(x+y)\,d\sigma$。

\end{callkey}

\begin{callwarn}{易错点}
第一，奇偶性必须区域对称与函数奇偶\textbf{同时}成立，区域不对称时一律不能用。第二，区域关于 $y$ 轴对称要看 $x$ 的奇偶性，这是最容易记反的一条。第三，轮换对称不能对函数的任意部分乱用，必须是整个被积函数或明确的对称配对。第四，形心公式只对 $x$、$y$ 的一次式有效（$x^{2}$、$xy$ 等不能用），且必须先用对称性求形心或确认图形规则。第五，折叠成一半时 $D_1$ 要取区域的"一半"，不能随手取一个子块。

\end{callwarn}

\begin{callexam}{真题考法}
解答题中的化简工具，通常与直角坐标或极坐标联合：先用对称性去掉奇函数项，再用轮换对称配对求 $x^{2}$ 与 $y^{2}$ 的积分，最后用形心公式处理线性项。选择题则常直接问"某二重积分的值等于多少"，正确做法是先判对称再看是否为零。 用对称性一定要写出理由，例如"因为 $D$ 关于 $y$ 轴对称，且被积函数关于 $x$ 为奇函数，故积分为零"。解答题中只给结论不给理由要扣分；而选择题里，先判对称往往就是最快的解法。

\end{callexam}

\subsubsection{分区域积分、含绝对值与最值函数的二重积分 $\bigstar$$\bigstar$$\bigstar$$\bigstar$}

当被积函数含绝对值、$\max$、$\min$、取整符号或本身是分段函数时，不能直接套公式，必须先"去掉非线性符号"，即\textbf{按分界线把区域拆开}。

\textbf{通用步骤}

\begin{enumerate}
\item 令绝对值内部（或 $\max$ 的两项之差）等于零，求出\textbf{分界线}；
\item 用分界线把 $D$ 分成若干子区域 $D_1,D_2,\dots$，要求 $D_1\cup D_2\cup\cdots=D$ 且两两至多沿边界相交；
\item 在每个子区域上把被积函数写成\textbf{无绝对值、无最值符号}的普通表达式；
\item 由区域可加性分别计算后相加。
\end{enumerate}
\textbf{理论化简式}（可用来验算分区是否正确）

\[\max\{u,v\}=\frac{u+v+\vert u-v\vert}{2},\qquad \min\{u,v\}=\frac{u+v-\vert u-v\vert}{2}\]

\textbf{取整函数} $[x+y]$ 的情形：按 $x+y=1,2,3,\dots$ 等整数值作分界线把区域切成若干条带，逐条计算。

例 1：计算 $\displaystyle\iint_{D}\left\vert y-x^{2}\right\vert dxdy$，其中 $D=[-1,1]\times[0,1]$。

分界线为 $y=x^{2}$。在 $D_1=\{y\ge x^{2}\}$ 上被积函数为 $y-x^{2}$，在 $D_2=\{y<x^{2}\}$ 上为 $x^{2}-y$。按 $x$ 型定限：

\[\iint_{D_1}\left(y-x^{2}\right)dxdy=\int_{-1}^{1}\int_{x^{2}}^{1}\left(y-x^{2}\right)dy\,dx=\frac12\int_{-1}^{1}\left(1-x^{2}\right)^{2}dx=\frac12\left(2-\frac43+\frac25\right)=\frac{8}{15}\]

\[\iint_{D_2}\left(x^{2}-y\right)dxdy=\int_{-1}^{1}\int_{0}^{x^{2}}\left(x^{2}-y\right)dy\,dx=\int_{-1}^{1}\frac{x^{4}}{2}dx=\frac15\]

两式相加得原积分 $=\frac{8}{15}+\frac15=\frac{11}{15}$。

例 2：计算 $\displaystyle\iint_{D}\max\{x,y\}\,d\sigma$，其中 $D=[0,1]\times[0,1]$。

用直线 $y=x$ 把正方形分成两块。在 $y\ge x$ 的三角形上 $\max\{x,y\}=y$，故

\[\int_{0}^{1}dy\int_{0}^{y}y\,dx=\int_{0}^{1}y^{2}dy=\frac13\]

由区域关于 $y=x$ 的对称性，另一块上的积分也是 $\frac13$，故原积分 $=\frac23$。

例 3：分段函数情形。若 $f(x,y)=x+y$（当 $x\ge 0$）、$f(x,y)=x-y$（当 $x<0$），在 $D=[-1,1]\times[0,1]$ 上积分时，只需用直线 $x=0$ 把 $D$ 分成左右两个矩形，分别代入对应的表达式再相加即可。 例 4：计算 $\displaystyle\iint_{D}\vert x+y\vert\,d\sigma$，其中 $D=[-1,1]\times[-1,1]$。

分界线 $x+y=0$ 把正方形分成两个三角形。利用区域关于原点的中心对称性、以及 $\vert x+y\vert$ 在 $(x,y)\to(-x,-y)$ 下不变，可把积分化为第一象限部分的 $4$ 倍：

\[\iint_{D}\vert x+y\vert\,d\sigma=4\int_{0}^{1}\int_{0}^{1}(x+y)\,dxdy=4\left(\frac12+\frac12\right)=4\]

若区域改成 $D=[0,1]\times[0,1]$，则 $x+y\ge 0$ 恒成立，绝对值可直接去掉，积分值只有 $1$。可见\textbf{先看清区域形状、再决定是否分区}，是这类题的关键。

\begin{callkey}{分区化解题模板}
$\circled{1}$ 求分界线（令绝对值内部或两函数之差为零）；$\circled{2}$ 画分界线并切分区域；$\circled{3}$ 逐块写出无符号的表达式；$\circled{4}$ 逐块定限积分；$\circled{5}$ 相加。含 $\max$、$\min$、取整符号、分段定义的情形全部套用这五步。

\end{callkey}

\begin{callwarn}{易错点}
第一，分界线求错会导致整题符号错，务必解方程而不是凭图像目测。第二，分区后要保证\textbf{不重不漏}，遗漏子区域或多算一次都会错。第三，把 $\max$ 或绝对值直接按一种情况计算是最严重的错误。第四，分区后每一块都要\textbf{重新定限}，不能照抄原区域的内外层界限。第五，含参数的分界线（如 $y=ax$）要按参数的范围分类讨论。

\end{callwarn}

\begin{callexam}{真题考法}
这是二重积分中难度较高的一类，常在解答题中作为较难的一问出现，也常与对称性结合（例如先判对称把奇函数项去掉，再对剩余的绝对值部分分区）。典型素材：$\iint\vert x+y\vert$、$\iint\vert x^{2}+y^{2}-1\vert$、$\iint\max\{xy,1\}$、$\iint\min\{x,y\}$。答题模板是"分界线—分区—去符号—分别积分—相加"，五步写清楚就能拿下大部分分数。 还有一个常被忽略的细节：含绝对值或 $\max$ 的题，如果区域本身对称、被积函数也有奇偶性，应先用对称性把不贡献的部分去掉，再对剩余部分分区；这样常常能省掉一半计算量。 本板块的复习优先级很明确：\textbf{直角坐标与极坐标定限（$\bigstar$$\bigstar$$\bigstar$$\bigstar$$\bigstar$）必须练到条件反射}，对称性化简（$\bigstar$$\bigstar$$\bigstar$$\bigstar$）要能在读题时立刻判断可用与否，而概念、性质与中值定理（$\bigstar$$\bigstar$$\bigstar$）以理解和选择填空为主。把历年解答题中的二重积分题集中做一遍，把每道题的区域都亲手画出来，是这一章最有效的复习方式。

\end{callexam}

\subsection*{本章易错清单}
\addcontentsline{toc}{subsection}{本章易错清单}

\begin{enumerate}
\item \textbf{定义域写成曲面}：二元函数的定义域是 $xOy$ 平面上的区域，多个限制条件要取交集并画出草图。
\item \textbf{用直线路径验证极限就下结论}：沿所有直线极限相同仍不能说明二重极限存在，必须用夹逼或极坐标给出与方向无关的估计；只有沿两条路径极限不同才能断定不存在。
\item \textbf{混淆偏导、连续与可微，并漏掉可微的极限验证}：偏导数存在既推不出连续，也推不出可微（只有偏导连续才是可微的充分条件）；判定可微必须验证 $\lim_{\rho\to 0}\frac{\Delta z-f_x\Delta x-f_y\Delta y}{\rho}=0$，分段函数在分界点还要用定义求偏导。
\item \textbf{复合函数求二阶偏导漏项}：$f_u'$ 仍是 $u,v$ 的函数，再求导要继续用链式法则；二阶交叉项前有系数 $2$。
\item \textbf{隐函数求导时变量角色混乱}：公式法要把 $x,y,z$ 当独立变量（求 $F_x$ 时 $z$ 不变）；两边求导法要把 $z$ 当 $x,y$ 的函数；两种方法不能混用，且 $F_z=0$ 时公式失效。
\item \textbf{混合偏导相等不加条件}：$f_{xy}=f_{yx}$ 需要二者连续；分段函数在分界点必须分别用定义计算。
\item \textbf{求极值时只算驻点}：$AC-B^{2}=0$ 时判别法失效，要另作判断；偏导不存在的点也可能是极值点。
\item \textbf{拉格朗日乘数法解完就下结论}：乘数法只给必要条件，必须补算端点、角点并比较才能确定最值；消 $\lambda$ 时同除因子要先讨论该因子为零。
\item \textbf{闭区域求最值只查内部}：最值常出现在边界上，必须逐段处理边界并比较角点。
\item \textbf{极坐标忘记乘 $r$}：$d\sigma=r\,dr\,d\theta$；偏心圆 $r=2a\cos\theta$ 的 $\theta$ 范围是 $\left[-\frac{\pi}{2},\frac{\pi}{2}\right]$。
\item \textbf{交换积分次序不重新画图}：换序必须由原限写出区域再重新定限；内外层限都要保证"下限小、上限大"，外层限必须是常数。
\item \textbf{含绝对值或 $\max$ 的积分不分区域}：必须先求分界线，按区域可加性拆开后逐块去掉符号再积分，且每块都要重新定限。
\end{enumerate}
\chapter{常微分方程}
\markboth{常微分方程}{常微分方程}
\subsection*{本章考纲要求}
\addcontentsline{toc}{subsection}{本章考纲要求}

\begin{quote}
\textbf{考试内容}

\begin{enumerate}
\item 常微分方程的基本概念
\item 变量可分离的微分方程
\item 齐次微分方程
\item 一阶线性微分方程
\item 可降阶的高阶微分方程
\item 线性微分方程解的性质及解的结构定理
\item 二阶常系数齐次线性微分方程
\item 高于二阶的某些常系数齐次线性微分方程
\item 简单的二阶常系数非齐次线性微分方程
\item 微分方程的简单应用
\end{enumerate}
\textbf{考试要求}

\begin{enumerate}
\item 了解微分方程及其阶、解、通解、初始条件和特解等概念。
\item 掌握变量可分离的微分方程及一阶线性微分方程的解法，会解齐次微分方程。
\item 会用降阶法解下列形式的微分方程：$y^{(n)}=f(x)$，$y''=f(x,y')$ 和 $y''=f(y,y')$。
\item 理解线性微分方程解的性质及解的结构。
\item 掌握二阶常系数齐次线性微分方程的解法，并会解某些高于二阶的常系数齐次线性微分方程。
\item 会解自由项为多项式、指数函数、正弦函数、余弦函数以及它们的和与积的二阶常系数非齐次线性微分方程。
\item 会用微分方程解决一些简单的应用问题。
\end{enumerate}
\end{quote}
\subsection*{本章星级总表}
\addcontentsline{toc}{subsection}{本章星级总表}

\begingroup\small\setlength{\tabcolsep}{4.0pt}
\begin{longtable}{@{}L{0.4471\linewidth}L{0.1118\linewidth}L{0.1118\linewidth}L{0.1118\linewidth}L{0.1118\linewidth}@{}}
\toprule
\textbf{考点} & \textbf{星级} & \textbf{考纲要求} & \textbf{难度} & \textbf{常见题型} \\
\midrule
\endhead
微分方程的基本概念（阶、通解、特解、初始条件） & $\bigstar$$\bigstar$ & 了解 & 基础 & 选择 \\
变量可分离的微分方程 & $\bigstar$$\bigstar$$\bigstar$$\bigstar$ & 掌握 & 基础 & 填空、解答 \\
齐次微分方程 & $\bigstar$$\bigstar$$\bigstar$ & 会 & 中等 & 解答 \\
一阶线性微分方程（公式法与常数变易法） & $\bigstar$$\bigstar$$\bigstar$$\bigstar$ & 掌握 & 中等 & 填空、解答 \\
可降阶的高阶微分方程（$y^{(n)}=f(x)$、$y''=f(x,y')$、$y''=f(y,y')$） & $\bigstar$$\bigstar$$\bigstar$ & 会用 & 中等 & 解答 \\
线性微分方程解的性质与解的结构 & $\bigstar$$\bigstar$$\bigstar$$\bigstar$ & 理解 & 中等 & 选择、解答 \\
二阶常系数齐次线性微分方程（特征方程三种情形） & $\bigstar$$\bigstar$$\bigstar$$\bigstar$$\bigstar$ & 掌握 & 基础 & 填空、解答 \\
高于二阶的某些常系数齐次线性微分方程 & $\bigstar$$\bigstar$ & 会 & 中等 & 填空 \\
二阶常系数非齐次线性微分方程（待定系数法，含和与积型自由项） & $\bigstar$$\bigstar$$\bigstar$$\bigstar$$\bigstar$ & 会 & 难 & 解答 \\
已知解反求微分方程或参数 & $\bigstar$$\bigstar$$\bigstar$ & 会 & 难 & 解答 \\
微分方程的简单应用（几何、物理、变化率） & $\bigstar$$\bigstar$$\bigstar$ & 会用 & 中等 & 解答、应用 \\
\bottomrule
\end{longtable}
\endgroup

本章主线只有一条：\textbf{先认类型，再套解法}。数二卷面上出现的方程几乎都落在四个口袋里——可分离、齐次、一阶线性、常系数线性。花 10 秒认类型，比盲目积分 10 分钟值钱得多。认类型的顺序建议固定为：先看能不能把 $y$ 与 $x$ 分开（可分离），再看 $\dfrac yx$ 是否整体出现（齐次），再看是否关于 $y$ 与 $y'$ 线性（一阶线性），最后看是否常系数线性（特征方程法）。

复习这一章还有一个效率上的建议：把四个"口袋"的识别标志默写成一张对照表，做题时按表逐条排除。绝大多数失分不是算不出来，而是类型认错之后一路硬算。凡是遇到陌生的方程，先问三句话：能不能把 $y$ 与 $x$ 分到等号两边？能不能把右端整体写成 $\dfrac yx$ 的函数？能不能整理成 $y'+P(x)y=Q(x)$？三句都否定，再看它是不是常系数线性方程，能不能用特征方程法。另外，解出通解后一定要回代检验一次——这是本章唯一的"自动纠错"手段，在分离变量、换元、待定系数这三类操作之后尤其必要，因为这三步都是最容易被一个小疏漏毁掉全题的地方。

\section{基本概念与一阶微分方程}

\subsection{微分方程的概念与解}

\subsubsection{微分方程的基本概念（阶、通解、特解、初始条件） $\bigstar$$\bigstar$}

\textbf{定义}：含有未知函数及其导数（或微分）的方程称为\textbf{微分方程}；未知函数是一元函数的微分方程称为\textbf{常微分方程}。方程中出现的未知函数最高阶导数的阶数，称为微分方程的\textbf{阶}。

$n$ 阶微分方程的一般形式是

\[F\left(x,y,y',\cdots,y^{(n)}\right)=0 .\]

若能从式中解出最高阶导数，就得到\textbf{已解出最高阶导数的形式} $y^{(n)}=f\left(x,y,y',\cdots,y^{(n-1)}\right)$，这是讨论解的存在唯一性时的标准写法。

\textbf{解、通解、特解}：若函数 $y=\varphi(x)$ 在区间 $I$ 上代入方程后成为恒等式，则称它是方程在 $I$ 上的解。含有 $n$ 个\textbf{相互独立}的任意常数的解

\[y=\varphi(x,C_1,C_2,\cdots,C_n)\]

称为 $n$ 阶方程的\textbf{通解}；在通解中给任意常数取定一组值所得的解称为\textbf{特解}。

"相互独立"这四个字要理解到位：它指的是\textbf{不能通过合并、重命名把常数个数减少}。例如 $C_1x+C_2x=(C_1+C_2)x$ 实际只有一个独立常数，$C_1e^{x}+C_2e^{x}$ 同理。反过来，$C_1e^{x}+C_2x$ 中的两个常数无法合并，是独立的。

\textbf{初始条件}：$n$ 阶方程的初始条件有 $n$ 个，即

\[y(x_0)=y_0,\quad y'(x_0)=y_1,\quad \cdots,\quad y^{(n-1)}(x_0)=y_{n-1}.\]

求满足初始条件的特解的问题称为\textbf{初值问题}。记住一条铁律：\textbf{初始条件的个数 = 方程的阶数 = 通解中独立任意常数的个数}。

\textbf{几何意义}：一阶方程的通解 $y=\varphi(x,C)$ 在平面上是一族曲线，称为\textbf{积分曲线族}；给定初始条件 $y(x_0)=y_0$ 就是从这族曲线中挑出过点 $(x_0,y_0)$ 的那一条。这也是"初始条件锁定唯一一条曲线"这一说法的来历，应用题的几何背景常常就藏在这里。

\textbf{线性与非线性}：若方程关于 $y,y',\cdots,y^{(n)}$ 都是一次的（不含它们的乘积、商与复合），则称为\textbf{线性}微分方程，其标准形为

\[y^{(n)}+a_1(x)y^{(n-1)}+\cdots+a_{n-1}(x)y'+a_n(x)y=f(x).\]

线性方程中 $f(x)\equiv 0$ 时称为\textbf{齐次}线性方程，$f(x)\not\equiv 0$ 时称为\textbf{非齐次}线性方程。判别时只需检查：各项中 $y$ 及其各阶导数是否都只以一次幂、且彼此不相乘地出现。例如 $y''+xy'=e^x$、$x^2y''+y=\sin x$ 都是线性方程（系数可以是 $x$ 的任意函数），而 $y''+yy'=0$、$y''+\sin y=0$、$(y')^2+y=0$ 都是非线性的。

\textbf{验证一个函数是解的一般步骤}：

\begin{enumerate}
\item 求所需阶数的导数（复合函数注意链式法则）；
\item 把 $y,y',y''$ 一起代入方程；
\item 化简，说明左右两端恒等；
\item 补充说明使等式成立的 \textbf{$x$ 的取值范围}。
\end{enumerate}
\textbf{例 1}：验证 $y=Cx^2$ 是方程 $y'=\dfrac{2y}{x}$ 的通解，并求满足 $y(1)=3$ 的特解。

解：$y'=2Cx$，右端 $\dfrac{2Cx^2}{x}=2Cx$，故对任意常数 $C$ 都成立，是通解（含一个任意常数，与一阶相符，区间取 $x>0$）。代入 $y(1)=3$ 得 $C=3$，特解为 $y=3x^2$。

\textbf{例 2}：验证 $y=C_1\cos x+C_2\sin x$ 是 $y''+y=0$ 的通解。

解：$y'=-C_1\sin x+C_2\cos x$，$y''=-C_1\cos x-C_2\sin x=-y$，故 $y''+y=0$ 恒成立。又 $C_1,C_2$ 相互独立（$\cos x$ 与 $\sin x$ 之比不是常数），所以它是二阶方程的通解。

\textbf{例 3（由通解反推方程）}：求以 $y=C_1e^{x}+C_2e^{2x}$ 为通解的微分方程。

解：$y'=C_1e^{x}+2C_2e^{2x}$，$y''=C_1e^{x}+4C_2e^{2x}$。用"消常数"的思路：$y'-y=C_2e^{2x}$，$y''-y'=2C_2e^{2x}$，两式相除比较得 $y''-y'=2(y'-y)$，即

\[y''-3y'+2y=0 .\]

（这类"由解反推方程"的题在后面的考点里还会系统讲。）

\begin{callkey}{必背}
\begin{itemize}
\item $n$ 阶方程的通解含 $n$ 个独立任意常数，初始条件也有 $n$ 个；
\item 线性方程的标准形：$y^{(n)}+a_1(x)y^{(n-1)}+\cdots+a_n(x)y=f(x)$，$f\equiv 0$ 为齐次；
\item 初始条件的个数永远等于阶数，这是检查答案的第一道关卡。
\end{itemize}
\end{callkey}

\begin{callwarn}{易错}
\begin{enumerate}
\item 验证解必须同时说明\textbf{成立的区间}。例如 $y=\ln x+C$ 只能写在 $x>0$ 上，不能写成"全体实数"。
\item 通解与特解的差别只在"任意常数是否取定"，不要把"含任意常数"当成特解。
\item 判断阶数只看\textbf{最高阶导数}，与方程中出现的多项式次数无关：$y''+y^3=0$ 是二阶方程。
\item 判断线性与否时不要被系数吓到：$x^3y''+e^{x}y'=\ln x$ 仍是\textbf{线性}方程，因为 $y,y',y''$ 都是一次的。
\item 常数"独立"是硬要求，做题时写完通解先数一遍独立常数个数。
\end{enumerate}
\end{callwarn}

\begin{callexam}{真题考法}
小题形式：给一个含参函数问它是否为某方程的解、由通解代初始条件求特解、判断方程是几阶、是否线性。陷阱常藏在"任意常数是否独立"，例如 $y=C_1x+C_2x$ 实际只含一个独立常数，不能作为二阶方程的通解。也有题目给出一个含参解，要求确定参数使该解满足给定方程，本质仍是代入比较。

\end{callexam}

\subsection{一阶微分方程}

\subsubsection{变量可分离的微分方程 $\bigstar$$\bigstar$$\bigstar$$\bigstar$}

\textbf{标准形式}：

\[\frac{dy}{dx}=f(x)g(y) .\]

\textbf{识别方法}：把方程整理成 $\dfrac{dy}{dx}=$（只含 $x$ 的式子）$\times$（只含 $y$ 的式子）。只要能写成"一边只含 $x$、一边只含 $y$"的乘积形式，就是变量可分离。若原方程是 $M(x)N(y)\,dx+P(x)Q(y)\,dy=0$，直接同除以 $N(y)P(x)$ 即可分离。

\textbf{解法（三步）}：

\begin{enumerate}
\item 分离变量：把 $y$ 放到一边、$x$ 放到另一边，
\end{enumerate}
\[\frac{dy}{g(y)}=f(x)\,dx\qquad (g(y)\ne 0).\]

\begin{enumerate}
\item 两边分别积分：
\end{enumerate}
\[\int \frac{dy}{g(y)}=\int f(x)\,dx+C .\]

\begin{enumerate}
\item 写出通解（可以是隐式）：$G(y)=F(x)+C$，其中 $G,F$ 分别是 $\dfrac1{g(y)}$ 与 $f(x)$ 的原函数。
\end{enumerate}
若 $g(y_0)=0$，则 $y\equiv y_0$ 也是解（常数解），它可能不包含在通解里，\textbf{必须单独写出}。判断"是否已经包含"的技巧是：看通解中令常数取某些值（含 $\pm\infty$ 的极限情形）能否得到该常数解。例如 $y'=2xy$ 的通解 $y=Ce^{x^2}$，令 $C=0$ 即得 $y\equiv 0$，已经包含；而 $y'=y^2$ 的通解 $y=-\dfrac{1}{x+C}$ 永远取不到 $0$，故 $y\equiv 0$ 是额外解。

\textbf{常用积分备查}（分离后几乎必用）：

\[\int\frac{dy}{y}=\ln\vert y\vert,\qquad \int y^{\alpha}\,dy=\frac{y^{\alpha+1}}{\alpha+1}\ (\alpha\ne-1),\]

\[\int\frac{dy}{y^{2}}=-\frac1y,\qquad \int\frac{dy}{1+y^{2}}=\arctan y,\qquad \int\tan y\,dy=-\ln\vert\cos y\vert .\]

\textbf{关于任意常数的写法}：积分后出现对数时，常把常数写成 $\ln\vert C\vert$ 或 $\ln C$（$C\ne 0$），这样能把两边的对数合并、去掉对数号，答案更简洁。例如

\[\ln\vert y\vert=x^{2}+C_1\ \Longrightarrow\ y=\pm e^{C_1}e^{x^{2}}\ \Longrightarrow\ y=Ce^{x^{2}},\]

最后一笔把 $\pm e^{C_1}$ 整体记作新的任意常数 $C$（这个 $C$ 只要求非零；若常数解 $y\equiv 0$ 也算在内，就放宽为任意常数）。

\textbf{例 1}：解方程 $y'=2xy$，并求满足 $y(0)=1$ 的特解。

解：分离变量得 $\dfrac{dy}{y}=2x\,dx$，积分得 $\ln\vert y\vert=x^2+C_1$，故 $y=Ce^{x^2}$；注意 $y\equiv 0$ 已含在 $C=0$ 中。由 $y(0)=1$ 得 $C=1$，特解为

\[y=e^{x^2}.\]

代回验证：$y'=2xe^{x^2}=2xy$，初始条件也满足。

\textbf{例 2}：求 $(1+x^{2})\,dy-(1+y^{2})\,dx=0$ 的通解。

解：整理成 $\dfrac{dy}{1+y^{2}}=\dfrac{dx}{1+x^{2}}$，两边积分得 $\arctan y=\arctan x+C$，即通解为

\[\arctan y-\arctan x=C .\]

若要化成显式，可用 $\tan(A-B)=\dfrac{\tan A-\tan B}{1+\tan A\tan B}$ 得 $\dfrac{y-x}{1+xy}=\tan C$，即 $y=\dfrac{x+C_0}{1-C_0x}$（$C_0=\tan C$）。两种写法都对，隐式解往往更省事，考场上不必强行显化。

\textbf{例 3}：求 $\dfrac{dy}{dx}=e^{x+y}$ 的通解。

解：右端可写成 $e^{x}\cdot e^{y}$，是标准的可分离形式。分离得 $e^{-y}\,dy=e^{x}\,dx$，积分得 $-e^{-y}=e^{x}+C$，即

\[e^{x}+e^{-y}=C .\]

注意这里常数必须满足 $C>0$（因为 $e^{x}+e^{-y}$ 恒正），这是变量可分离题中容易被忽略的取值限制。

\textbf{关于隐式解与显式解}：分离变量后得到的往往是隐式关系 $G(y)=F(x)+C$，只要它确定了 $y$ 是 $x$ 的函数，就是合法的通解。考试中不必强行解出 $y$：像 $\arctan y=\arctan x+C$、$e^{x}+e^{-y}=C$ 这样的写法既简洁又不易出错。只有当题目明确要求求出 $y$、或后续还要代入其他式子计算时，才考虑化成显式。

\begin{callwarn}{易错}
\begin{enumerate}
\item \textbf{除以 $g(y)$ 会丢解}。做完后回头检查 $g(y)=0$ 的常数解是否被漏掉。
\item 积分结果含 $\ln\vert y\vert$ 时不要丢绝对值，去绝对值时会引入符号分段。
\item 任意常数写一个即可，不要写成 $C_1+C_2$（那不是两个独立常数）。
\item 分离的前提是 $g(y)\ne 0$；若 $g$ 在区间内有零点，通解只能写在使 $g(y)\ne 0$ 的区间上。
\item $dy$ 与 $dx$ 前的因子要一起搬走，漏乘某个因子会让方程"看起来可分离，其实不是"。
\end{enumerate}
\end{callwarn}

\begin{callexam}{真题考法}
填空题与解答题第一问最常见。两种变式：一是方程本身含 $e^{x+y}$、$\cos y$ 这类"天然可分"的因子；二是先作因式分解再分离，例如 $(1+x^2)\,dy-(1+y^2)\,dx=0$。题目给了初始条件就一定要把 $C$ 定出来。另外，分离变量常常只是一道更大题目的第一步（后面接几何量计算、最值讨论），所以解出的通解要写得便于后续代入。

\end{callexam}

\subsubsection{齐次微分方程 $\bigstar$$\bigstar$$\bigstar$}

\textbf{标准形式}：能整理成

\[\frac{dy}{dx}=\varphi\!\left(\frac{y}{x}\right)\]

的方程称为齐次微分方程。判别标志有两个：

\begin{itemize}
\item \textbf{代入检验}：把 $x,y$ 同时换成 $tx,ty$，方程形式不变；
\item \textbf{次数检验}：写成 $M(x,y)\,dx+N(x,y)\,dy=0$ 后，$M,N$ 都是同一次数的齐次函数（例如 $x^{2}+y^{2}$、$xy$、$x^{2}-xy$ 都是二次齐次式）。
\end{itemize}
\textbf{解法（换元 $v=\dfrac{y}{x}$）}：

\begin{enumerate}
\item 令 $v=\dfrac{y}{x}$，即 $y=vx$，两边对 $x$ 求导得
\end{enumerate}
\[\frac{dy}{dx}=v+x\frac{dv}{dx}.\]

\begin{enumerate}
\item 代回原方程，得 $v+x\dfrac{dv}{dx}=\varphi(v)$，即
\end{enumerate}
\[x\frac{dv}{dx}=\varphi(v)-v .\]

\begin{enumerate}
\item 这是 $v$ 与 $x$ 的变量可分离方程，分离积分：
\end{enumerate}
\[\int \frac{dv}{\varphi(v)-v}=\int \frac{dx}{x}=\ln\vert x\vert+C .\]

\begin{enumerate}
\item 积出 $v$ 后把 $v=\dfrac{y}{x}$ \textbf{回代}，写成 $y$ 与 $x$ 的关系式。
\end{enumerate}
换元有效的道理很简单：齐次方程的解曲线在把图形按同一比例放大缩小时形状不变，因此描述"形状"的比值 $\dfrac yx$ 才是真正的变量，把它单独拎出来方程就降成了可分离。

\textbf{判别流程}：遇到 $\dfrac{dy}{dx}=$ 一个分式时，先看分子分母是不是同次齐次式。若是，同除分母中 $x$ 的最高次幂，把它整理成 $\dfrac yx$ 的函数；若分子分母不同次，优先考虑按一阶线性处理（把 $x$ 当未知函数往往能凑成线性）。写成微分形式 $M\,dx+N\,dy=0$ 时，若 $M,N$ 同次，则同除 $x^{n}$ 即可看出齐次结构；但若 $dx$ 前只含 $y$、$dy$ 前只含 $x$，那实际上是变量可分离而不是齐次，不要认错。这两条岔路的判别只需看一眼次数，却决定了后面是换元 $v=\dfrac yx$ 还是直接分离积分，值得养成先看次数的习惯。

\textbf{可化为齐次的两种情形}（考得较少，会认即可）：

\begin{itemize}
\item 若 $\dfrac{dy}{dx}=f\!\left(\dfrac{a_1x+b_1y+c_1}{a_2x+b_2y+c_2}\right)$ 且 $a_1b_2-a_2b_1\ne 0$，先解方程组 $\begin{cases}a_1x+b_1y+c_1=0,\\ a_2x+b_2y+c_2=0,\end{cases}$ 得交点 $(h,k)$，再作平移 $X=x-h,\;Y=y-k$，方程即化为关于 $X,Y$ 的齐次方程。
\item 若 $a_1b_2-a_2b_1=0$（两直线平行），令 $u=a_2x+b_2y$，方程化为可分离方程。
\end{itemize}
\textbf{例 1}：求 $\dfrac{dy}{dx}=\dfrac{x^{2}+y^{2}}{xy}$（$x>0$）的通解。

解：右端 $=\dfrac{x}{y}+\dfrac{y}{x}=\dfrac1v+v$。代入 $v+xv'=v+\dfrac1v$，得 $xv'=\dfrac1v$，分离得 $v\,dv=\dfrac{dx}{x}$，积分得 $\dfrac{v^{2}}{2}=\ln x+C$。回代 $v=\dfrac yx$：

\[y^{2}=2x^{2}\ln x+Cx^{2} .\]

验证：取 $C=1$，$y^{2}=2x^{2}\ln x+x^{2}$，两边求导得 $2yy'=4x\ln x+4x$，即 $y'=\dfrac{2x(\ln x+1)}{y}$；而原方程右端为 $\dfrac{x^{2}+y^{2}}{xy}=\dfrac{x^{2}+2x^{2}\ln x+x^{2}}{xy}=\dfrac{2x(\ln x+1)}{y}$，两者一致。

\textbf{例 2}：求 $\dfrac{dy}{dx}=\dfrac{y}{x}+\tan\dfrac yx$ 的通解。

解：设 $v=\dfrac yx$，则 $v+xv'=v+\tan v$，即 $x v'=\tan v$。分离得 $\dfrac{\cos v}{\sin v}\,dv=\dfrac{dx}{x}$，积分得 $\ln\vert\sin v\vert=\ln\vert x\vert+C_1$，即 $\sin v=Cx$。回代 $v=\dfrac yx$：

\[\sin\frac yx=Cx .\]

（这里用到了 $\int\cot v\,dv=\ln\vert\sin v\vert$，与前面备查表一致。）注意 $x\to 0$ 时该方程有奇性，通解只在 $\dfrac yx\ne k\pi$ 的区间上成立。

\begin{callwarn}{易错}
\begin{enumerate}
\item \textbf{"齐次"有两个完全不同的意思}：本节"齐次微分方程"指 $\frac{dy}{dx}=\varphi(\frac yx)$；线性方程中的"齐次"指 $f(x)\equiv 0$。看到"齐次方程"先看它是不是线性方程。
\item 换元后千万不要漏掉 $x\frac{dv}{dx}$ 这一项，这是本考点最高频的失分点。
\item $\varphi(v)-v=0$ 时，$v\equiv k$ 对应直线解 $y=kx$，需要单独检验是否漏掉。
\item 回代一步不能忘，答案里不应该再出现 $v$。
\item 换元后积出的是 $v$ 与 $x$ 的关系，回代时要把 $\ln\vert x\vert$ 这类项保留好，别在整理时把常数吞掉。
\item 若方程写成 $\frac{dx}{dy}=\varphi(\frac xy)$，对称地令 $v=\frac xy$ 即可，不要硬套 $\frac yx$。
\end{enumerate}
\end{callwarn}

\begin{callexam}{真题考法}
解答题中常作为第一步出现。《考试要求》只要求"会解齐次微分方程"，考法固定：给出分子分母同次的式子，换元 $v=\frac yx$ 化成可分离。陷阱是题干需要先移项、同除，把 $\frac{dy}{dx}$ 单独解出来才能看出齐次结构。也有题目把齐次方程写成微分形式 $(x^{2}+y^{2})\,dx-xy\,dy=0$，此时同样同除 $x^{2}$ 化为 $\frac yx$ 的函数。

\end{callexam}

\subsubsection{一阶线性微分方程（公式法与常数变易法） $\bigstar$$\bigstar$$\bigstar$$\bigstar$}

\textbf{标准形式}：

\[\frac{dy}{dx}+P(x)y=Q(x),\]

其中 $P(x),Q(x)$ 在考虑的区间上连续。$Q(x)\equiv 0$ 时称为一阶线性\textbf{齐次}方程，$Q(x)\not\equiv 0$ 时称为\textbf{非齐次}方程。\textbf{识别要点：$y$ 与 $y'$ 都是一次的，系数只能是 $x$ 的函数。}

\textbf{第一步：解对应的齐次方程} $y'+P(x)y=0$，它是可分离的：

\[\frac{dy}{y}=-P(x)\,dx\ \Longrightarrow\ \ln\vert y\vert=-\int P(x)\,dx+C_1\ \Longrightarrow\ y=Ce^{-\int P(x)\,dx}.\]

\textbf{第二步（积分因子法与通解公式）}：在方程两边同乘 $e^{\int P(x)\,dx}$（称为\textbf{积分因子}），左边恰好凑成乘积的导数：

\[\frac{d}{dx}\left[ye^{\int P(x)\,dx}\right]=Q(x)e^{\int P(x)\,dx}.\]

两边积分再解出 $y$，就得到\textbf{通解公式}

\[y=e^{-\int P(x)\,dx}\left[\int Q(x)e^{\int P(x)\,dx}\,dx+C\right].\]

注意括号内那个积分号\textbf{不再加常数}，通解中的任意常数 $C$ 已经写在括号外。

\textbf{第二步（常数变易法，与公式法等价）}：把齐次通解中的 $C$ 换成待定函数 $C(x)$，即设

\[y=C(x)e^{-\int P(x)\,dx}.\]

求导得 $y'=C'(x)e^{-\int P\,dx}-C(x)P(x)e^{-\int P\,dx}$，代入 $y'+P(x)y=Q(x)$，含 $C(x)$ 的两项正好抵消，剩下

\[C'(x)e^{-\int P\,dx}=Q(x)\ \Longrightarrow\ C'(x)=Q(x)e^{\int P(x)\,dx},\]

积分即得 $C(x)$。这个推导说明两件事：公式法不过是常数变易法的结果；含 $C(x)$ 的项一定抵消，这是"变易"能成立的原因。

\textbf{使用步骤}：

\begin{enumerate}
\item 把方程整理成 $y'+P(x)y=Q(x)$ 的标准形（$y'$ 的系数必须化到 $1$）；
\item 算出 $\int P(x)\,dx$（\textbf{不加常数}）；
\item 套公式，或按常数变易法走；
\item 有初始条件就定 $C$（注意：\textbf{先求通解再代初值}，不要一开始就把 $C$ 当常数用掉）。
\end{enumerate}
\textbf{例 1}：求 $y'-\dfrac{2}{x}y=x^{2}\;(x>0)$ 满足 $y(1)=1$ 的特解。

解：这里 $P(x)=-\dfrac2x$，$Q(x)=x^{2}$，

\[\int P(x)\,dx=-2\ln x,\qquad e^{-\int P\,dx}=x^{2},\qquad e^{\int P\,dx}=\frac{1}{x^{2}} .\]

代入公式：

\[y=x^{2}\left[\int x^{2}\cdot\frac{1}{x^{2}}\,dx+C\right]=x^{2}(x+C)=x^{3}+Cx^{2} .\]

由 $y(1)=1$ 得 $1=1+C$，故 $C=0$，特解为 $y=x^{3}$。

验证：$y'=3x^{2}$，$y'-\dfrac2xy=3x^{2}-2x^{2}=x^{2}$，初始条件 $y(1)=1$ 也满足。

\textbf{例 2（交换自变量）}：求 $\dfrac{dy}{dx}=\dfrac{1}{x+y^{2}}$ 的通解。

解：这个方程关于 $y$ 不是线性的，但把 $x$ 看成 $y$ 的函数就变成线性的：由 $\dfrac{dx}{dy}=x+y^{2}$ 得

\[\frac{dx}{dy}-x=y^{2}.\]

此时自变量是 $y$，$P=-1$，$Q=y^{2}$：

\[x=e^{y}\left[\int y^{2}e^{-y}\,dy+C\right]=e^{y}\left[-\left(y^{2}+2y+2\right)e^{-y}+C\right]=Ce^{y}-\left(y^{2}+2y+2\right).\]

验证：$x'=Ce^{y}-(2y+2)$，$x'-x=Ce^{y}-2y-2-Ce^{y}+y^{2}+2y+2=y^{2}$，成立。

\textbf{两个常见变形}：

\begin{itemize}
\item \textbf{交换自变量}（如例 2）：若方程关于 $y'$ 不是线性的、但关于 $\dfrac{dx}{dy}$ 是线性的，就把 $x$ 看成未知函数，化成 $\dfrac{dx}{dy}+P(y)x=Q(y)$，套同一公式。典型判据：$x$ 的最高次数高于 $y$ 时优先考虑这条路线。
\item \textbf{含变限积分的方程}：如 $y=x+\displaystyle\int_0^x y(t)\,dt$，先两边对 $x$ 求导脱去积分号，化为 $y'=1+y$ 这类方程，再由原式取 $x=0$ 得初始条件。
\end{itemize}
\textbf{例 3（积分方程）}：求 $y(x)=x+\displaystyle\int_0^x y(t)\,dt$ 的解。

解：两边求导得 $y'=1+y$，即 $y'-y=1$；由原式令 $x=0$ 得 $y(0)=0$。用公式：

\[y=e^{x}\left[\int 1\cdot e^{-x}\,dx+C\right]=e^{x}\left(-e^{-x}+C\right)=Ce^{x}-1 .\]

由 $y(0)=0$ 得 $C=1$，故 $y=e^{x}-1$。

\textbf{一阶线性方程的"解的结构"}：把通解 $y=Ce^{-\int P\,dx}+e^{-\int P\,dx}\displaystyle\int Qe^{\int P\,dx}dx$ 看成两项，第一项正是对应齐次方程的通解 $Y$，第二项（把常数取成 $0$）是非齐次方程的一个特解 $y^{*}$。也就是说，\textbf{一阶线性方程同样满足"非齐次通解等于齐次通解加非齐次特解"}，这与后面高阶线性方程的结构完全一致。看清这一点，公式就不必死记：先解齐次得 $Y$，再补一个特解即可。

\textbf{例 4（含参数的方程）}：求 $y'+\dfrac{a}{x}y=x^{b}$（$a,b$ 为常数，$x>0$）的通解。

解：$\int P\,dx=a\ln x$，故 $e^{-\int P\,dx}=x^{-a}$，$e^{\int P\,dx}=x^{a}$，代入公式：

\[y=x^{-a}\left[\int x^{b}\cdot x^{a}\,dx+C\right]=x^{-a}\left[\frac{x^{a+b+1}}{a+b+1}+C\right]\qquad (a+b\ne -1).\]

当 $a+b=-1$ 时内层积分为 $\ln x$，通解为 $y=x^{-a}\left(\ln x+C\right)$。这类题考查的是\textbf{分类讨论意识}：凡是积分结果中分母可能为零的参数取值，都必须单独说明。

\begin{callkey}{必背}
一阶线性方程的通解公式（$y'+P(x)y=Q(x)$）：

\[y=e^{-\int P(x)\,dx}\left[\int Q(x)e^{\int P(x)\,dx}\,dx+C\right]\]

\begin{itemize}
\item 外层永远是 $e^{-\int P}$，内层永远是 $e^{+\int P}$；
\item 齐次情形 $Q\equiv 0$ 时退化为 $y=Ce^{-\int P\,dx}$，也就是公式中取 $\int Qe^{\int P}dx=0$。
\end{itemize}
\end{callkey}

\begin{callwarn}{易错}
\begin{enumerate}
\item \textbf{必须先化成标准形}。像 $xy'+y=x^{2}$ 要同除 $x$ 得 $y'+\frac1xy=x$，此时 $P(x)=\frac1x$，而不是把 $y$ 的系数直接当作 $P$ 之外的东西。
\item 公式里 $e^{-\int P}$ 与 $e^{\int P}$ 不要写反，写反结果完全错误。记法：\textbf{外层是 $e^{-\int P}$}。
\item $\int P(x)\,dx$ 不加任意常数；任意常数只出现一次，在括号外。
\item 积分方程求导后要\textbf{回头取初始条件}，否则常数定不出来。取初始条件的办法是把 $x$ 取成积分下限。
\item 内层积分 $\int Q(x)e^{\int P\,dx}dx$ 常需要分部积分，算完后一定回代检查一步。
\end{enumerate}
\end{callwarn}

\begin{callexam}{真题考法}
填空（求通解或特解）与解答题（与几何、物理应用结合）都常见。识别陷阱：题干表面不是线性，例如含 $y^{2}$ 的项或形如 $y'=\frac{y}{x+y^{2}}$，此时应把 $x$ 当未知函数、$y$ 当自变量，整理成 $\frac{dx}{dy}-\frac1yx=y$ 后用公式。另一种考法是"由 $y'+P(x)y=Q(x)$ 及 $y(0)=y_0$ 求 $\lim_{x\to+\infty}y(x)$"，核心仍是先解出解析式再取极限。

\end{callexam}

\section{可降阶的高阶微分方程}

\subsection{降阶法}

\subsubsection{可降阶的高阶微分方程（$y^{(n)}=f(x)$、$y''=f(x,y')$、$y''=f(y,y')$） $\bigstar$$\bigstar$$\bigstar$}

考纲明确要求三类，判定口诀是：\textbf{缺谁就以谁为自变量}——方程里没有 $y$，就令 $p=y'$ 并把 $p$ 看成 $x$ 的函数；方程里没有 $x$，就令 $p=y'$ 并把 $p$ 看成 $y$ 的函数。降阶的实质是\textbf{换一个自变量，把二阶问题压成一阶问题}。

\textbf{类型 I：$y^{(n)}=f(x)$}。右端只含 $x$，逐次积分 $n$ 次，每次产生一个任意常数，共得 $n$ 个独立常数：

\[y=\underbrace{\int\cdots\int}_{n\ \text{次}}f(x)\,dx\cdots dx+\frac{C_1x^{n-1}}{(n-1)!}+\cdots+C_{n-1}x+C_n .\]

（后面那一串常数项之所以写成带阶乘的形式，是因为逐次积分时每个常数还要再积一次，$\frac{C_1x^{n-1}}{(n-1)!}$ 只是把"积出来以后长这个样子"写清楚；实际做题时从里往外逐次加常数即可。）

\textbf{类型 II：$y''=f(x,y')$（不显含 $y$）}。方程里没有 $y$ 本身，令

\[p=y',\qquad y''=p',\]

方程降为关于 $p$ 的一阶方程 $p'=f(x,p)$。解出 $p=p(x,C_1)$ 后，再由 $y'=p$ 积分一次得 $y$。更一般地，若方程是 $y^{(n)}=f\left(x,y^{(n-1)}\right)$，同样令 $p=y^{(n-1)}$ 即可。

\textbf{类型 III：$y''=f(y,y')$（不显含 $x$）}。方程里没有 $x$，令

\[p=y',\qquad y''=\frac{dp}{dx}=\frac{dp}{dy}\cdot\frac{dy}{dx}=p\frac{dp}{dy},\]

方程降为关于 $y$ 的一阶方程

\[p\frac{dp}{dy}=f(y,p).\]

解出 $p=g(y,C_1)$ 后代回 $\dfrac{dy}{dx}=g(y,C_1)$，这是变量可分离方程，再积一次得通解。

\textbf{例（类型 I）}：解 $y'''=e^{x}$。

解：逐次积分，$y''=e^{x}+C_1$，$y'=e^{x}+C_1x+C_2$，

\[y=e^{x}+\frac{C_1}{2}x^{2}+C_2x+C_3 .\]

（三个独立常数，与三阶相符。）

\textbf{例（类型 II）}：解 $y''=y'+x$。

解：令 $p=y'$，则 $p'-p=x$，这是一阶线性方程，$P=-1$，$Q=x$：

\[p=e^{x}\left[\int xe^{-x}\,dx+C_1\right]=e^{x}\left[-(x+1)e^{-x}+C_1\right]=C_1e^{x}-x-1 .\]

再积分：

\[y=C_1e^{x}-\frac{x^{2}}{2}-x+C_2 .\]

验证：$y'=C_1e^{x}-x-1$，$y''=C_1e^{x}-1$，而 $y'+x=C_1e^{x}-x-1+x=C_1e^{x}-1=y''$，成立。

\textbf{例（类型 III）}：解 $yy''=(y')^{2}$，且 $y(0)=1,\;y'(0)=1$。

解：令 $p=y'$，$y''=p\dfrac{dp}{dy}$，代入得 $y\,p\dfrac{dp}{dy}=p^{2}$。当 $p\not\equiv 0$ 时约去 $p$：

\[\frac{dp}{p}=\frac{dy}{y}\ \Longrightarrow\ \ln\vert p\vert=\ln\vert y\vert+C\ \Longrightarrow\ p=C_1y .\]

由 $y'(0)=1,\;y(0)=1$ 得 $C_1=1$，即 $y'=y$，解得 $y=C_2e^{x}$，再由 $y(0)=1$ 得 $C_2=1$，故

\[y=e^{x}.\]

验证：$y'=y''=e^{x}$，$yy''=e^{2x}=(y')^{2}$，成立。

\textbf{例（类型 II 的另一种走向）}：解 $y''=1+(y')^{2}$。

解：设 $p=y'$，则 $p'=1+p^{2}$，分离得 $\dfrac{dp}{1+p^{2}}=dx$，积分得 $\arctan p=x+C_1$，即 $p=\tan(x+C_1)$。再积分：

\[y=\int\tan(x+C_1)\,dx=-\ln\vert\cos(x+C_1)\vert+C_2 .\]

这个例子说明：降阶以后遇到的可能是反正切、反正弦等非初等类型的积分，只要算得出来就继续算，算不出来就写成隐式或定积分形式。

\textbf{降阶后常数如何分配}：降阶过程中每解一个一阶方程都会引入一个任意常数，因此用类型 II、III 处理二阶方程时，最终通解必定含两个独立常数。若题目给了初始条件，通常先用 $y'(x_0)=y_1$ 这个\textbf{导数条件}确定降阶后一阶方程里的常数，再用 $y(x_0)=y_0$ \textbf{函数值条件}确定最后一次积分产生的常数。顺序颠倒会让运算变乱，也容易把两个常数用混。

\begin{callwarn}{易错}
\begin{enumerate}
\item 类型 III 中 $y''=p\dfrac{dp}{dy}$，\textbf{不能写成 $p'$}。这是本考点第一大失分点（$p$ 是 $x$ 的函数还是 $y$ 的函数，必须先想清楚）。
\item 约去 $p$ 时注意 $p\equiv 0$（即 $y\equiv$ 常数）也是解，视题目要求检验是否舍去。
\item 每次积分产生一个常数，通解中常数个数必须等于阶数；中途用初始条件定常数时要按阶数分配。
\item 类型 II 解出 $p$ 后别忘了还有一次积分 $y=\int p\,dx$。
\item 类型 III 解出 $p=g(y,C_1)$ 后，$\dfrac{dy}{dx}=g(y)$ 还要再分离一次，不要以为已经做完。
\item 判断类型时先看方程里\textbf{缺哪个变量}，而不是先看阶数；同一个方程有时既能按类型 II 也能按类型 III 处理，选计算量小的那条路。
\end{enumerate}
\end{callwarn}

\begin{callexam}{真题考法}
解答题第二问高频，尤其是"不显含 $x$"型 $y''=f(y,y')$。常见设计是：第一问由几何或物理背景列出方程，第二问用降阶法求解并代初始条件。也有把类型 I 藏在 $y^{(4)}=f(x)$ 中，考查多次积分的耐心。另一类考法是给出 $y''=f(x,y')$ 并要求用 $p$ 表示，重在写清换元过程而不是最终积分。

\end{callexam}

\section{线性微分方程解的性质与解的结构}

\subsection{解的性质与结构定理}

\subsubsection{线性微分方程解的性质与解的结构 $\bigstar$$\bigstar$$\bigstar$$\bigstar$}

考虑 $n$ 阶线性微分方程

\[L[y]\equiv y^{(n)}+a_1(x)y^{(n-1)}+\cdots+a_{n-1}(x)y'+a_n(x)y=f(x),\]

其中 $a_i(x),f(x)$ 在区间 $I$ 上连续。$f(x)\equiv 0$ 时记作 $L[y]=0$，称为齐次方程。以下所有性质都来自同一个事实：\textbf{$L$ 是线性的}，即对任意常数 $c_1,c_2$ 有

\[L\left[c_1y_1+c_2y_2\right]=c_1L[y_1]+c_2L[y_2].\]

\textbf{性质一（齐次方程的叠加原理）}：若 $y_1,y_2$ 都是 $L[y]=0$ 的解，则对任意常数 $C_1,C_2$，$C_1y_1+C_2y_2$ 也是 $L[y]=0$ 的解。一般地，齐次方程的解集构成一个 $n$ 维线性空间。

\textbf{性质二（基本解组与齐次通解）}：若 $y_1,y_2,\cdots,y_n$ 是 $L[y]=0$ 的 $n$ 个\textbf{线性无关}的解，则齐次通解为

\[Y=C_1y_1+C_2y_2+\cdots+C_ny_n .\]

这样的 $n$ 个解称为\textbf{基本解组}。线性无关的意思是不存在不全为零的常数 $k_1,\cdots,k_n$ 使 $k_1y_1+\cdots+k_ny_n\equiv 0$；二阶情形最直观：$y_1,y_2$ 线性相关当且仅当 $\dfrac{y_2}{y_1}$ 是常数。

\textbf{性质三（非齐次解与齐次解的转化）}：若 $y_1,y_2$ 都是 $L[y]=f(x)$ 的解，则 $y_1-y_2$ 是 $L[y]=0$ 的解。更一般地，若非齐次方程的 $k$ 个解 $y_1,\cdots,y_k$ 的线性组合满足

\[\sum_{i=1}^{k}c_i=0 \ \Longrightarrow\ \sum_{i=1}^{k}c_iy_i \ \text{是齐次方程的解},\qquad
\sum_{i=1}^{k}c_i=1 \ \Longrightarrow\ \sum_{i=1}^{k}c_iy_i \ \text{是非齐次方程的解}.\]

原因只有一行：$L\!\left[\sum c_iy_i\right]=\sum c_iL[y_i]=f(x)\sum c_i$。\textbf{看到"非齐次方程若干解的线性组合"，先算系数和。}

\textbf{性质四（非齐次通解结构）}：

\[y=Y+y^{*},\]

其中 $Y$ 是齐次通解，$y^{*}$ 是非齐次方程的\textbf{某一个}特解。这条定理的作用是把"解非齐次方程"拆成两件独立的事：找齐次通解（纯代数）和找一个特解（待定系数或变易），互不干扰。

\textbf{性质五（叠加原理）}：若 $L[y_1]=f_1(x)$，$L[y_2]=f_2(x)$，则 $L[y_1+y_2]=f_1(x)+f_2(x)$。这是非齐次自由项"和型"可以拆开分求的依据。

\textbf{性质六（解的线性无关判别）}：二阶情形，设 $y_1,y_2$ 是 $y''+p(x)y'+q(x)y=0$ 的两个解，则它们线性无关的充要条件是\textbf{朗斯基行列式}

\[W(x)=\begin{vmatrix} y_1 & y_2\\ y_1' & y_2' \end{vmatrix}=y_1y_2'-y_2y_1'\not\equiv 0 .\]

进一步的结论（Abel 公式）：$W(x)=Ce^{-\int p(x)\,dx}$，因此 $W(x)$ 在区间 $I$ 上\textbf{要么恒为 $0$，要么处处不为 $0$}（因为指数函数永远不为零）。所以只需在任意一点验证 $W(x_0)\ne 0$，就能断定两个解线性无关；若 $y_2=ky_1$，则必有 $W\equiv 0$。这条结论在实际做题时非常省事：在 $x=0$（或任何方便的点）算一次行列式就够了。

\textbf{性质七（已知一个解求通解，降阶公式）}：设 $y_1$ 是 $y''+p(x)y'+q(x)y=0$ 的一个非零解，则与它线性无关的另一个解为

\[y_2=y_1\int\frac{e^{-\int p(x)\,dx}}{y_1^{2}}\,dx,\]

于是通解为 $y=C_1y_1+C_2y_2$。这个公式的来历是设 $y=y_1v$ 代入原方程，$v$ 满足一个可降阶的方程，解出 $v$ 就得到上式。做题时不必背推导，但要\textbf{记住必须先把 $y''$ 的系数化成 $1$}，$p(x)$ 才取对。

\textbf{例 1（验证基本解组）}：验证 $y_1=e^{x},\;y_2=xe^{x}$ 是 $y''-2y'+y=0$ 的基本解组，并写出通解。

解：$y_1'=e^{x},y_1''=e^{x}$，代入得 $e^{x}-2e^{x}+e^{x}=0$；$y_2'=(1+x)e^{x},\;y_2''=(2+x)e^{x}$，代入得 $(2+x)e^{x}-2(1+x)e^{x}+xe^{x}=0$。朗斯基行列式

\[W(x)=e^{x}\cdot(1+x)e^{x}-xe^{x}\cdot e^{x}=e^{2x}\ne 0,\]

故线性无关，通解为 $y=(C_1+C_2x)e^{x}$。

\textbf{例 2（系数和判据）}：设 $y_1=x,\;y_2=x+e^{x},\;y_3=x+e^{-x}$ 都是 $y''-y=-x$ 的解，求该方程的通解。

解：先找齐次解。$y_2-y_1=e^{x}$，$y_3-y_1=e^{-x}$，由性质三它们都是 $y''-y=0$ 的解，且 $W(e^{x},e^{-x})=-2\ne 0$，故齐次通解为 $C_1e^{x}+C_2e^{-x}$。再由性质四，取 $y^{*}=y_1=x$：

\[y=C_1e^{x}+C_2e^{-x}+x .\]

（也可直接看系数和：$2y_1-y_2-y_3=-e^{x}-e^{-x}$，系数和 $2-1-1=0$，确为齐次解。再验证 $y''-y=-x$：$y_1=x$ 时 $0-x=-x$，$y_2=x+e^{x}$ 时 $e^{x}-x-e^{x}=-x$，都对。）

\textbf{例 3（降阶公式）}：已知 $y_1=x$ 是 $x^{2}y''-2xy'+2y=0$ 的解，求通解。

解：化标准形得 $y''-\dfrac2xy'+\dfrac{2}{x^{2}}y=0$，故 $p(x)=-\dfrac2x$，$e^{-\int p\,dx}=e^{2\ln x}=x^{2}$，于是

\[y_2=x\int\frac{x^{2}}{x^{2}}\,dx=x\cdot x=x^{2}.\]

通解 $y=C_1x+C_2x^{2}$。验证 $y=x^{2}$：$2x^{2}-4x^{2}+2x^{2}=0$，成立。

\textbf{例 4（用 $W$ 反求系数）}：已知 $y_1=e^{x},\;y_2=e^{-x}$ 是 $y''+p(x)y'+q(x)y=0$ 的两个解，求 $p,q$。

解：$W=y_1y_2'-y_2y_1'=-e^{x}e^{-x}-e^{-x}e^{x}=-2$，为常数。由 $W=Ce^{-\int p\,dx}$ 及 $W$ 为常数可知 $p\equiv 0$；再把 $y_1=e^{x}$ 代入 $y''+q(x)y=0$ 得 $e^{x}+q(x)e^{x}=0$，即 $q(x)\equiv -1$。于是方程为 $y''-y=0$。

\textbf{为什么齐次解空间恰好是 $n$ 维}：线性方程的解集对加法与数乘封闭，所以构成线性空间；由存在唯一性定理，任给一组初始值 $\left(y_0,y_1,\cdots,y_{n-1}\right)$ 都有唯一解，而 $n$ 个初始值给出 $n$ 个自由度，因此这个空间的维数正好是 $n$。这一句话把"基本解组含 $n$ 个解""通解含 $n$ 个独立常数""初始条件有 $n$ 个"三件事串在了一起。做题时若发现常数个数不对，就知道某一步漏了或重了。

\textbf{用朗斯基行列式做判断的实战写法}：拿到两个函数先算 $W(x)$，化简到只剩一个明显非零的因子（如 $e^{2x}$、$-2$），说明它不为零，即可下结论"线性无关"。若 $W\equiv 0$，则在两个函数都连续可导的前提下它们线性相关，即一个是另一个的常数倍。需要注意的是，$W$ 恒为零也可能源于两个函数根本不满足同一个二阶方程，所以只有"已知二者是同一齐次方程的解"时，才能用 $W$ 判定线性相关性。

\begin{callkey}{必背}
解的结构三句话：

\begin{enumerate}
\item 齐次方程：$n$ 个线性无关解 $y_1,\cdots,y_n$ 的线性组合 $Y=\sum C_iy_i$ 就是全部解；
\item 非齐次方程：$y=Y+y^{*}$，齐次通解加一个非齐次特解；
\item 线性组合判据：$L\!\left[\sum c_iy_i\right]=f(x)\sum c_i$，看系数和是 $0$ 还是 $1$。
\end{enumerate}
\end{callkey}

\begin{callwarn}{易错}
\begin{enumerate}
\item 非齐次方程的\textbf{两个解之差}是齐次解，\textbf{两个解之和}一般什么都不是——不要想当然。
\item 齐次解加齐次解仍是齐次解；齐次解加非齐次解是非齐次解；两个非齐次解相加一般不是解。
\item 线性组合判据只看\textbf{系数和}：和为 $0$ 得齐次解，和为 $1$ 得非齐次解，其余情形两个结论都不成立。
\item 说"线性无关"必须有依据（朗斯基行列式非零，或两个解的比值不是常数），不能只因为"两个函数长得不一样"。
\item 降阶公式里的 $p(x)$ 是\textbf{标准形}中 $y'$ 的系数，用之前一定要把 $y''$ 的系数化成 $1$。
\item 由两个解反求 $p(x)$ 时要用 $p=-\dfrac{W'}{W}$，注意 $W$ 必须不为零，故只在使 $W\ne 0$ 的区间上有效。
\end{enumerate}
\end{callwarn}

\begin{callexam}{真题考法}
选择题最爱考三类判断：（a）给某非齐次方程的三个解，问哪个线性组合仍是解、哪个是齐次解；（b）给两个函数问是否线性无关；（c）给通解形式反推方程。解答题里，"已知齐次一个解，求通解"往往作为降低难度的台阶，后面再接初值或积分计算。这类题的计算量通常不大，得分点全在有没有把结构关系用对。

\end{callexam}

\section{常系数线性微分方程及其应用}

\subsection{常系数齐次线性方程}

\subsubsection{二阶常系数齐次线性微分方程（特征方程三种情形） $\bigstar$$\bigstar$$\bigstar$$\bigstar$$\bigstar$}

\textbf{标准形式}：

\[y''+py'+qy=0\qquad (p,q\ \text{为实常数}).\]

\textbf{特征方程法的由来}：方程里 $y,y',y''$ 的系数都是常数，说明"函数与其导数成固定比例"的结构最可能满足方程，而具有这种性质的只有指数函数，所以试设 $y=e^{rx}$。代入得

\[r^{2}e^{rx}+pre^{rx}+qe^{rx}=0\ \Longrightarrow\ e^{rx}\left(r^{2}+pr+q\right)=0 ,\]

约去恒不为零的 $e^{rx}$，就把微分方程化成了\textbf{代数方程}

\[r^{2}+pr+q=0,\qquad \Delta=p^{2}-4q .\]

按判别式 $\Delta$ 分三种情形写通解。

\begingroup\small\setlength{\tabcolsep}{4.0pt}
\begin{longtable}{@{}L{0.5550\linewidth}L{0.3906\linewidth}@{}}
\toprule
\textbf{特征方程的根} & \textbf{通解} \\
\midrule
\endhead
两个相异实根 $r_1\ne r_2$ & $y=C_1e^{r_1x}+C_2e^{r_2x}$ \\
二重实根 $r=-\frac p2$ & $y=(C_1+C_2x)e^{rx}$ \\
共轭复根 $r=\alpha\pm\beta i$，其中 $\alpha=-\frac p2,\allowbreak\ \beta=\frac{\sqrt{4q-p^{2}}}{2}$ & $y=e^{\alpha x}\left(C_1\cos\beta x+C_2\sin\beta x\right)$ \\
\bottomrule
\end{longtable}
\endgroup

速写口诀：\textbf{$y$ 求几次导就写几次方}——$y''\to r^{2}$，$y'\to r$，$y\to 1$，常数项直接搬下来。

三种情形的道理各有各的要点：

\begin{itemize}
\item \textbf{相异实根}：$e^{r_1x}$ 与 $e^{r_2x}$ 之比是 $e^{(r_1-r_2)x}$，不是常数，二者线性无关，直接线性组合即可。
\item \textbf{二重根}：$e^{rx}$ 只给出一个解，另一个解要用降阶公式补。由前面性质七算得 $y_2=xe^{rx}$，所以通解要写成 $(C_1+C_2x)e^{rx}$，\textbf{那个 $x$ 就是这样来的}。
\item \textbf{共轭复根}：形式上 $e^{(\alpha\pm\beta i)x}$ 都是解，但含虚数不便使用；由欧拉公式 $e^{\pm i\beta x}=\cos\beta x\pm i\sin\beta x$ 取实部与虚部，就得到两个实值解 $e^{\alpha x}\cos\beta x$、$e^{\alpha x}\sin\beta x$，它们的组合才是标准的实值通解。
\end{itemize}
\textbf{几个特殊情形}（考小题时很常见）：

\begin{itemize}
\item $q=0$：特征方程有零根 $r=0$，通解中含一个常数项，即 $y=C_1+C_2e^{-px}$（当 $p\ne 0$）；
\item $p=0$：特征根为纯虚根 $\pm\beta i$（$q>0$），解为简谐振动形式 $y=C_1\cos\beta x+C_2\sin\beta x$；
\item $p>0$ 且 $q>0$：特征根要么是两个负实根，要么是实部为负的共轭复根，于是\textbf{该方程的任意解都满足 $\lim\limits_{x\to+\infty}y(x)=0$}。这条性质常用在选择或证明题里。
\end{itemize}
\textbf{反问题}：若特征根是 $r_1,r_2$，则特征方程为 $(r-r_1)(r-r_2)=0$，即

\[r^{2}-(r_1+r_2)r+r_1r_2=0,\]

对应微分方程 $y''-(r_1+r_2)y'+r_1r_2y=0$。由解的形状反推方程时，把根和系数对应起来就能一次写出。

\textbf{例 1（相异实根）}：解 $y''-5y'+6y=0$。

解：$r^{2}-5r+6=0$ 得 $r_1=2,r_2=3$，通解 $y=C_1e^{2x}+C_2e^{3x}$。

\textbf{例 2（二重根 + 初值）}：解 $y''+4y'+4y=0$，$y(0)=1,\;y'(0)=2$。

解：$r^{2}+4r+4=(r+2)^{2}=0$，$r=-2$ 为二重根，通解 $y=(C_1+C_2x)e^{-2x}$，而

\[y'=C_2e^{-2x}-2(C_1+C_2x)e^{-2x}.\]

由 $y(0)=C_1=1$；由 $y'(0)=C_2-2C_1=2$ 得 $C_2=4$，故

\[y=(1+4x)e^{-2x}.\]

\textbf{例 3（共轭复根）}：解 $y''+2y'+5y=0$。

解：$r^{2}+2r+5=0$，$r=-1\pm 2i$，即 $\alpha=-1,\beta=2$，通解

\[y=e^{-x}\left(C_1\cos 2x+C_2\sin 2x\right).\]

\textbf{例 4（反问题）}：求以 $y=e^{2x}(C_1\cos 3x+C_2\sin 3x)$ 为通解的二阶常系数齐次方程。

解：特征根为 $2\pm 3i$，故 $r_1+r_2=4$，$r_1r_2=4+9=13$，特征方程 $r^{2}-4r+13=0$，所求方程为 $y''-4y'+13y=0$。

\textbf{例 5（含参数讨论）}：讨论 $y''+2by'+\left(b^{2}+1\right)y=0$ 的解的形态（$b$ 为实数）。

解：$\Delta=4b^{2}-4\left(b^{2}+1\right)=-4<0$，特征根恒为 $r=-b\pm i$。故无论 $b$ 取何值，解都是衰减或增长的振荡形式

\[y=e^{-bx}\left(C_1\cos x+C_2\sin x\right).\]

当 $b>0$ 时振幅衰减并趋于零，当 $b<0$ 时振幅随时间指数增长，$b=0$ 时为等幅振荡。

\textbf{解的性质在本章的三个用处}：一是判断极值——已知 $f$ 是某方程的解且 $f'(x_0)=0$，由方程直接读出 $f''(x_0)$ 的符号，无需解出 $f$；二是判断长期行为——由特征根实部的符号判断 $\lim_{x\to+\infty}y(x)$ 是零、是无穷还是持续振荡；三是反求方程——由通解的形状反推特征根，再由根写方程。这三类小题在选择题与填空题中反复出现，共同点是都不需要真的把方程解出来。

\begin{callkey}{必背}
二阶常系数齐次方程 $y''+py'+qy=0$ 的通解（特征方程 $r^{2}+pr+q=0$）：

\begin{itemize}
\item $\Delta>0$：$y=C_1e^{r_1x}+C_2e^{r_2x}$；
\item $\Delta=0$：$y=(C_1+C_2x)e^{rx}$，$r=-\dfrac p2$；
\item $\Delta<0$：$y=e^{\alpha x}\left(C_1\cos\beta x+C_2\sin\beta x\right)$，$\alpha=-\dfrac p2$，$\beta=\dfrac{\sqrt{4q-p^{2}}}{2}$。
\end{itemize}
\end{callkey}

\begin{callwarn}{易错}
\begin{enumerate}
\item \textbf{二重根时别漏 $x$}：写成 $y=C_1e^{rx}+C_2e^{rx}$ 等于只有一个独立常数，必然错。
\item 复根一定要化成\textbf{实值形式} $e^{\alpha x}(C_1\cos\beta x+C_2\sin\beta x)$；写成 $C_1e^{(\alpha+\beta i)x}+C_2e^{(\alpha-\beta i)x}$ 通常不给分。
\item $\beta$ 的值要算对：$\beta=\frac{\sqrt{4q-p^{2}}}{2}$，例如 $r=-1\pm2i$ 对应 $p=2,q=5$。
\item 判别式算错会导致根的类型判断错误，$p,q$ 是实常数时这一步必须先写出来。
\item 由通解反推方程时，要保证通解中独立常数的个数等于阶数（二阶必须恰好两个）。
\item 特征根是零根时，$e^{0\cdot x}=1$，通解里出现常数项，不要漏写。
\end{enumerate}
\end{callwarn}

\begin{callexam}{真题考法}
填空题几乎必有一道（求通解或求带初值的特解）。选择题的经典设计是：已知 $y=f(x)$ 是 $y''-2y'+4y=0$ 的解，且 $f(x_0)>0,\;f'(x_0)=0$，问 $f(x)$ 在 $x_0$ 处取极大还是极小。做法是由方程得 $f''(x_0)=-4f(x_0)<0$，故取极大值——判据来自方程本身，不需要解出 $f$。另一种设计是把解与广义积分结合：先由初值解出 $y$，再计算 $\int_0^{+\infty}y(x)\,dx$。还有一类是"含参数方程讨论解的有界性/衰减性"，关键是看特征根实部的符号。

\end{callexam}

\subsubsection{高于二阶的某些常系数齐次线性微分方程 $\bigstar$$\bigstar$}

\textbf{标准形式}：

\[y^{(n)}+a_1y^{(n-1)}+\cdots+a_{n-1}y'+a_ny=0\qquad (a_i\ \text{为常数}),\]

特征方程为

\[r^{n}+a_1r^{n-1}+\cdots+a_{n-1}r+a_n=0 .\]

\textbf{解的对应规则（"单根单写，复根复写，重根重写"）}：

\begingroup\small\setlength{\tabcolsep}{4.0pt}
\begin{longtable}{@{}L{0.4844\linewidth}L{0.4613\linewidth}@{}}
\toprule
\textbf{特征根} & \textbf{对应的解} \\
\midrule
\endhead
单实根 $r$ & $e^{rx}$ \\
$m$ 重实根 $r$ & $\left(C_1+C_2x+\cdots+C_mx^{m-1}\right)e^{rx}$ \\
单复根 $\alpha\pm\beta i$ & $e^{\alpha x}\left(C_1\cos\beta x+C_2\sin\beta x\right)$ \\
$m$ 重复根 $\alpha\pm\beta i$ & $e^{\alpha x}\left[P_{m-1}(x)\cos\beta x+Q_{m-1}(x)\sin\beta x\right]$，其中 $P_{m-1},Q_{m-1}$ 为 $m-1$ 次多项式 \\
\bottomrule
\end{longtable}
\endgroup

把每个特征根给出的解\textbf{全部相加}就是通解，其中独立任意常数的总数必定等于方程的阶数 $n$。重根情形可以这样记：\textbf{$m$ 重实根贡献 $m$ 个解，$m$ 重复根贡献 $2m$ 个解。}

解题步骤：

\begin{enumerate}
\item 写出特征方程，用有理根试探加分组分解（或提取公因式、平方差、完全平方）因式分解；
\item 求出全部特征根并标出重数；
\item 按上表逐根写解并相加；
\item 用"独立常数个数 $=n$"作最后检查。
\end{enumerate}
\textbf{因式分解的实战技巧}（本节唯一的难点）：

\begin{itemize}
\item 先试 $r=\pm 1,\pm 2,\pm 3$ 等有理根，凡是常数项 $a_n$ 的因数都有可能；
\item 找到一个根后立刻用综合除法或配凑降次，把问题化到二次；
\item 遇到 $r^{4}+2r^{2}+1$、$r^{4}-2r^{2}+1$ 这类"只含偶次幂"的式子，直接令 $t=r^{2}$ 或识别为完全平方；
\item 遇到 $r^{3}+r^{2}-r-1$ 这类式子，前两项与后两项分别提公因式即可。
\end{itemize}
\textbf{例 1（三阶、三个单实根）}：解 $y'''-6y''+11y'-6y=0$。

解：特征方程 $r^{3}-6r^{2}+11r-6=0$，试 $r=1$ 成立，分解得

\[r^{3}-6r^{2}+11r-6=(r-1)(r-2)(r-3),\]

故

\[y=C_1e^{x}+C_2e^{2x}+C_3e^{3x}.\]

\textbf{例 2（三阶、重根）}：解 $y'''+y''-y'-y=0$。

解：对特征方程分组分解：

\[r^{3}+r^{2}-r-1=r^{2}(r+1)-(r+1)=(r+1)\left(r^{2}-1\right)=(r+1)^{2}(r-1).\]

故 $r=-1$ 是二重根、$r=1$ 是单根：

\[y=(C_1+C_2x)e^{-x}+C_3e^{x}.\]

（三个独立常数，与三阶相符。）

\textbf{例 3（四阶、重根）}：解 $y^{(4)}-2y''+y=0$。

解：特征方程 $r^{4}-2r^{2}+1=(r^{2}-1)^{2}=(r-1)^{2}(r+1)^{2}$，$r=1$ 与 $r=-1$ 都是二重根，故

\[y=(C_1+C_2x)e^{x}+(C_3+C_4x)e^{-x}.\]

\textbf{例 4（四阶、四重根）}：解 $y^{(4)}-4y'''+6y''-4y'+y=0$。

解：特征方程为 $(r-1)^{4}=0$（这正是二项式系数 $1,4,6,4,1$ 的结构，见到这组系数就该立刻想到 $(r-1)^{4}$），$r=1$ 是四重根，故

\[y=\left(C_1+C_2x+C_3x^{2}+C_4x^{3}\right)e^{x}.\]

\textbf{例 5（四阶、重共轭复根）}：解 $y^{(4)}+2y''+y=0$。

解：特征方程 $r^{4}+2r^{2}+1=(r^{2}+1)^{2}$，$r=\pm i$ 是二重复根（$\alpha=0,\beta=1$），故

\[y=(C_1+C_2x)\cos x+(C_3+C_4x)\sin x .\]

\textbf{一点必要的提醒}：高于二阶的某些方程可以通过\textbf{换元降阶}的思路化为常系数方程再求解——遇到结构陌生的高阶方程，先想"能不能作一个变量代换把它变成熟悉的常系数方程"，这就是降阶思想的延伸。数二不需要深究具体换元形式，记住这个方向即可。

\begin{callwarn}{易错}
\begin{enumerate}
\item 因式分解是本节全部难点。先试 $r=\pm1,\pm2,\pm3$ 等有理根，再对余下的多项式用分组分解或公式。
\item 重根必须按次数补齐 $x$ 的各次幂：二重根写 $(C_1+C_2x)e^{rx}$，四重根要写到 $x^{3}$，不是 $(C_1+C_2)e^{rx}$。
\item 复根一定共轭成对出现，不会只剩一个纯虚根；漏掉共轭那一半会导致常数个数不足。
\item 最后一定数一遍独立常数个数，必须等于阶数，这是最省时的自检手段。
\item 特征方程只与各项系数有关，与原方程的阶数写法无关；先把它完整写下来再动手分解，不要跳步。
\end{enumerate}
\end{callwarn}

\begin{callexam}{真题考法}
填空题为主，给出三阶或四阶常系数齐次方程求通解。常见设计是特征方程能漂亮地分解，例如 $(r-1)^{2}(r+1)$、$(r^{2}+4)^{2}$、$(r-1)^{4}$ 这类结构；偶尔与"已知某个解求参数"结合，考查重数的判断。得分点很集中：特征方程写对、因式分解到位、重数对应的多项式次数不写错。

\end{callexam}

\subsection{常系数非齐次线性方程}

\subsubsection{二阶常系数非齐次线性微分方程（待定系数法，含和与积型自由项） $\bigstar$$\bigstar$$\bigstar$$\bigstar$$\bigstar$}

\textbf{标准形式}：

\[y''+py'+qy=f(x)\qquad (p,q\ \text{为实常数}),\]

其中自由项 $f(x)$ 是多项式、指数函数、正弦函数、余弦函数以及它们的和与积。

\textbf{待定系数法为什么可行}：当 $f(x)=P_m(x)e^{\lambda x}$ 时，把这种形式的函数代入常系数线性算子 $L$，指数因子 $e^{\lambda x}$ 会被整体提出、多项式仍化为同次多项式，即

\[L\left[Q(x)e^{\lambda x}\right]=e^{\lambda x}\tilde{Q}(x),\]

其中 $\tilde Q$ 与 $Q$ 次数相同（当 $\lambda$ 是特征根时次数会降低）。所以只要按 $f$ 的"形状"设一个带待定系数的同类函数，代入后比较同类项就一定能解出系数——这就是待定系数法的全部依据。也正因为如此，它\textbf{只对常系数、且自由项为上述类型时有效}。

\textbf{解题四步}：

\begin{enumerate}
\item \textbf{求齐次通解 $Y$}：解特征方程 $r^{2}+pr+q=0$，按三种情形写出 $Y=C_1y_1+C_2y_2$。
\item \textbf{设特解 $y^{*}$}：设 $k$ 为"共振重数"，即 $\lambda$（或 $\alpha+\beta i$）作为特征根的重数，取 $k=0,1,2$（不是特征根取 $0$，是单根取 $1$，是二重根取 $2$）。
\item \textbf{定系数}：代入原方程比较同类项。
\item \textbf{组装并定常数}：$y=Y+y^{*}$，有初始条件就定 $C_1,C_2$。
\end{enumerate}
\begingroup\small\setlength{\tabcolsep}{4.0pt}
\begin{longtable}{@{}L{0.4728\linewidth}L{0.4728\linewidth}@{}}
\toprule
\textbf{自由项 $f(x)$ 的形式} & \textbf{特解形式（$Q_m,R_m,S_m$ 为待定多项式）} \\
\midrule
\endhead
$P_m(x)e^{\lambda x}$ & $y^{*}=x^{k}Q_m(x)e^{\lambda x}$ \\
$P_m(x)$（即 $\lambda=0$） & $y^{*}=x^{k}Q_m(x)$ \\
$e^{\alpha x}\left[P_l(x)\cos\beta x+P_n(x)\sin\beta x\right]$ & $y^{*}=x^{k}e^{\alpha x}\left[R_m(x)\cos\beta x+S_m(x)\sin\beta x\right],\allowbreak\ m=\max\{l,n\}$ \\
$A\cos\beta x+B\sin\beta x$（即 $\alpha=0$） & $y^{*}=x^{k}\left(a\cos\beta x+b\sin\beta x\right)$ \\
\bottomrule
\end{longtable}
\endgroup

三条设定原则：

\begin{enumerate}
\item \textbf{多项式次数与 $f$ 中多项式的最高次数相同}，且\textbf{各项系数都要设出来}（即使 $f$ 中缺项）；
\item \textbf{指数因子 $e^{\lambda x}$ 照抄}，不能省；
\item \textbf{三角函数的角频率照抄}，$\cos\beta x$ 与 $\sin\beta x$ 要\textbf{成对出现}，即使 $f$ 中只有 $\cos\beta x$ 也要两个都设。
\end{enumerate}
\textbf{$k$ 的记忆法}：把 $\lambda$（或 $\alpha+\beta i$）拿去和特征根比对——\textbf{不是根就不乘 $x$，是单根乘一次 $x$，是二重根乘 $x^{2}$}。它只取决于特征根，与自由项里多项式的次数无关。

\textbf{自由项缺项与常数因子的处理}：$f(x)$ 中没有出现的项，在特解里也要照设。例如 $f(x)=x^{2}$ 时必须设 $ax^{2}+bx+c$，不能只设 $ax^{2}$；$f(x)=\sin 3x$ 时必须设 $a\cos 3x+b\sin 3x$，不能只设 $b\sin 3x$；$f(x)=e^{x}$ 时也要设 $ae^{x}$，待定系数 $a$ 的符号不能省。反过来，$f(x)$ 前面若有常数因子，直接并入待定系数即可，不必单独讨论。设全之后比较系数，本该为零的那些系数自然会解成零；而漏设则会让方程组出现矛盾，这正是"解不出系数"最常见的原因。

\textbf{若 $f$ 是和型}：用叠加原理拆成几项，分别求特解再相加，\textbf{每一项的 $k$ 要分别判断}。

\textbf{若 $f$ 是积型}：如 $x e^{x}\sin x$、$x^{2}e^{2x}$，按上表统一处理——"指数照抄、多项式取最高次、三角成对设、再看共振重数乘 $x^{k}$"。

\textbf{复数解法}（处理三角型自由项的一条捷径）：求 $y''+py'+qy=e^{\alpha x}P_m(x)\cos\beta x$ 时，可以先解

\[z''+pz'+qz=P_m(x)e^{(\alpha+\beta i)x},\]

求出复值特解 $z^{*}$ 后取其实部，即得原方程的一个特解 $y^{*}=\operatorname{Re}z^{*}$。这条路线把三角型化成指数型，设定形式时不必反复比较 $\cos,\sin$ 的系数，计算三角型大题时能省不少笔墨。

\textbf{例 1（多项式型）}：求 $y''-3y'+2y=x^{2}$ 的通解。

解：特征方程 $r^{2}-3r+2=0$，$r_1=1,r_2=2$，故 $Y=C_1e^{x}+C_2e^{2x}$。

这里 $\lambda=0$ 不是特征根，$k=0$，$m=2$，设 $y^{*}=ax^{2}+bx+c$。代入：

\[2a-3(2ax+b)+2(ax^{2}+bx+c)=2ax^{2}+(-6a+2b)x+(2a-3b+2c)=x^{2}.\]

比较系数：$2a=1\Rightarrow a=\dfrac12$；$-6a+2b=0\Rightarrow b=3a=\dfrac32$；$2a-3b+2c=0\Rightarrow 1-\dfrac92+2c=0\Rightarrow c=\dfrac74$。故

\[y=C_1e^{x}+C_2e^{2x}+\frac12x^{2}+\frac32x+\frac74 .\]

\textbf{例 2（共振，单重根）}：求 $y''-3y'+2y=e^{x}$ 的通解。

解：$Y=C_1e^{x}+C_2e^{2x}$。这里 $\lambda=1$ \textbf{恰好是特征单根}，所以 $k=1$，设 $y^{*}=axe^{x}$。代入：

\[{y^{*}}''-3{y^{*}}'+2y^{*}=e^{x}\left[a(2+x)-3a(1+x)+2ax\right]=-ae^{x}=e^{x},\]

得 $a=-1$，故 $y^{*}=-xe^{x}$，

\[y=C_1e^{x}+C_2e^{2x}-xe^{x}.\]

（直观理解：$f(x)=e^{x}$ 与齐次解"同频"，直接设 $ae^{x}$ 会被全部抵消，必须乘 $x$ 提升一档。）

\textbf{例 3（三角函数型）}：求 $y''+y=\cos 2x$ 的一个特解。

解：特征根 $r=\pm i$，而 $\pm 2i$ 不是特征根，$k=0$，设 $y^{*}=a\cos 2x+b\sin 2x$。由 ${y^{*}}''=-4y^{*}$ 得

\[{y^{*}}''+y^{*}=-3y^{*}=\cos 2x\ \Longrightarrow\ a=-\frac13,\ b=0,\qquad y^{*}=-\frac13\cos 2x .\]

\textbf{例 4（和型，用叠加原理拆开）}：求 $y''-y=e^{2x}+x$ 的一个特解。

解：拆成两个方程分别求。

\begin{itemize}
\item 对 $e^{2x}$：$\lambda=2$ 不是特征根（特征根为 $\pm1$），设 $ae^{2x}$，代入得 $4a-a=1$，$a=\dfrac13$；
\item 对 $x$：$\lambda=0$ 不是特征根，设 $bx+c$，代入得 $-(bx+c)=x$，故 $b=-1,\;c=0$。
\end{itemize}
\[y^{*}=\frac13e^{2x}-x .\]

验证：${y^{*}}'=\frac43e^{2x}-1$，${y^{*}}''=\frac43e^{2x}$，${y^{*}}''-y^{*}=\left(\frac43-\frac13\right)e^{2x}+x=e^{2x}+x$，成立。

\textbf{例 5（积型：多项式乘指数）}：求 $y''-y=xe^{x}$ 的一个特解。

解：特征根为 $\pm1$，$\lambda=1$ 是\textbf{单根}，故 $k=1$，设

\[y^{*}=x(ax+b)e^{x}=(ax^{2}+bx)e^{x}.\]

代入计算得

\[{y^{*}}''-y^{*}=\left(4ax+2a+2b\right)e^{x}=xe^{x},\]

比较系数：$4a=1\Rightarrow a=\dfrac14$；$2a+2b=0\Rightarrow b=-\dfrac14$。故

\[y^{*}=\frac14\left(x^{2}-x\right)e^{x}=\frac{x(x-1)}{4}e^{x}.\]

\textbf{例 6（积型：多项式乘三角函数，且共振）}：求 $y''+y=x\cos x$ 的一个特解。

解：特征根 $r=\pm i$，$\alpha+\beta i=i$ 是\textbf{特征单根}，故 $k=1$；又 $f$ 中多项式最高次数为 $1$，$m=1$，设

\[y^{*}=x\left[(ax+b)\cos x+(cx+d)\sin x\right].\]

代入 ${y^{*}}''+y^{*}=x\cos x$，比较 $\cos x$、$\sin x$、$x\cos x$、$x\sin x$ 的系数得

\[4c=1,\qquad 2a+2d=0,\qquad -4a=0,\qquad 2c-2b=0,\]

解得 $a=0,\;b=\dfrac14,\;c=\dfrac14,\;d=0$，故

\[y^{*}=\frac{x}{4}\cos x+\frac{x^{2}}{4}\sin x .\]

验证：${y^{*}}''+y^{*}=\frac14\left(3x\cos x-x^{2}\sin x\right)+\frac14\left(x\cos x+x^{2}\sin x\right)=x\cos x$，成立。

\textbf{例 7（指数乘三角且共振）}：求 $y''+2y'+5y=e^{-x}\cos 2x$ 的一个特解。

解：特征方程 $r^{2}+2r+5=0$ 的根为 $r=-1\pm 2i$，而自由项对应的 $\alpha+\beta i=-1+2i$ \textbf{正是特征根}，故 $k=1$。设

\[y^{*}=xe^{-x}\left(a\cos 2x+b\sin 2x\right).\]

记 $w=e^{-x}\left(a\cos 2x+b\sin 2x\right)$，则 $L\left[xw\right]=2\left(w'+w\right)$（含 $x$ 的项因 $w$ 是齐次解而消失），而

\[w'+w=e^{-x}\left(-2a\sin 2x+2b\cos 2x\right),\]

故 $L\left[y^{*}\right]=e^{-x}\left(-4a\sin 2x+4b\cos 2x\right)$。要求它等于 $e^{-x}\cos 2x$，得 $b=\dfrac14,\;a=0$，于是

\[y^{*}=\frac{x}{4}e^{-x}\sin 2x .\]

验证：$L\left[y^{*}\right]=e^{-x}\cdot 4\cdot\dfrac14\cos 2x=e^{-x}\cos 2x$，成立。

\begin{callwarn}{易错}
\begin{enumerate}
\item \textbf{$k$ 只看 $\lambda$（或 $\alpha+\beta i$）是不是特征根、是几重根}，与 $f$ 中多项式的次数无关；$f$ 中多项式只决定待定多项式的次数。
\item 比较系数时\textbf{不要因为 $f$ 缺项就不设该项}。例如 $f(x)=x^{2}$ 也要设 $ax^{2}+bx+c$；$f(x)=\cos 2x$ 也要设 $a\cos 2x+b\sin 2x$。
\item "和型"必须\textbf{拆开分别求}，拆开后每一项自己的共振重数要\textbf{分别判断}。
\item 待定系数法适用的前提是：\textbf{常系数}线性方程加自由项为多项式、指数、正余弦及其和与积。自由项是 $\ln x$、$\dfrac1x$、$\sec x$ 之类时不能硬套。
\item 最后写答案时不要只写 $y^{*}$，题目问"通解"就必须补上齐次通解 $Y$。
\item 求出的特解不要与齐次解"长得一样"：如果发现 $y^{*}$ 是齐次解的线性组合，说明 $k$ 取小了。
\item 定系数时方程组若矛盾，一定是设定形式错了（通常是漏了 $x^{k}$），回头检查比反复算积分快得多。
\end{enumerate}
\end{callwarn}

\begin{callexam}{真题考法}
解答题的必考题型（10 分上下），通常两问：先求通解，再求满足初值的特解；或与几何应用、体积面积结合。评分点上，齐次通解、特解设定形式、待定系数求解、初值定常数各占一块。最大的陷阱是自由项的 $\lambda$ 与特征根重合：漏乘 $x^{k}$ 会导致系数方程组无解，考场上发现自己"解不出系数"时，第一反应就应该是回头检查 $k$ 有没有取错。自由项为"和型"时，注意叠加原理只能用在线性方程上，拆开后不要漏加。

\end{callexam}

\subsubsection{已知解反求微分方程或参数 $\bigstar$$\bigstar$$\bigstar$}

这类题不给方程、只给解，反向把方程还原出来。常见四个入口。

\textbf{入口一：已知二阶常系数齐次方程的通解形式}。从通解读出特征根，再用根写特征方程：

\[(r-r_1)(r-r_2)=r^{2}-(r_1+r_2)r+r_1r_2=0\ \Longrightarrow\ y''-(r_1+r_2)y'+r_1r_2y=0 .\]

\textbf{例}：求以 $y=C_1e^{-x}+C_2e^{3x}$ 为通解的二阶常系数齐次方程。

解：$r_1=-1,\;r_2=3$，$r_1+r_2=2$，$r_1r_2=-3$，故方程为 $y''-2y'-3y=0$。

\textbf{入口二：已知某函数是解，求参数}。把解代入方程，比较同类项，解关于参数的方程。

\textbf{例}：设 $y=e^{kx}$ 是 $y''-3y'+2y=0$ 的解，求 $k$。

解：代入得 $(k^{2}-3k+2)e^{kx}=0$，即 $k^{2}-3k+2=0$，故 $k=1$ 或 $k=2$。

\textbf{入口三：已知解的参数形态，反求参数}。解中出现 $xe^{rx}$ 说明有重根；解中出现 $\cos\beta x,\sin\beta x$ 说明有共轭复根；解中只有两个不同的指数函数说明有两个相异实根。由形态反推判别式 $\Delta$ 的符号即可定出参数范围。

\textbf{例}：若方程 $y''+ay'+by=0$ 的通解中出现 $\cos 2x$ 与 $\sin 2x$（且不带 $e$ 的幂因子），求 $a,b$。

解：解中出现 $\cos 2x,\sin 2x$ 且无指数衰减因子，说明特征根是纯虚根 $\pm 2i$，故特征方程为 $r^{2}+4=0$，即 $a=0,\;b=4$。

\textbf{入口四：已知变系数方程的一个解，用降阶公式求通解}（见上一板块性质七）。若已知两个解 $y_1,y_2$，还可由

\[W(x)=y_1y_2'-y_2y_1'=Ce^{-\int p(x)\,dx}\]

反解出 $p(x)=-\dfrac{W'(x)}{W(x)}$，再把 $y_1$ 代入方程解出 $q(x)$。

\textbf{例}：已知 $y_1=e^{x},\;y_2=xe^{x}$ 是 $y''+p(x)y'+q(x)y=0$ 的两个解，求对应的常系数方程。

解：这两个解对应特征根 $r=1$（二重），特征方程为 $(r-1)^{2}=r^{2}-2r+1=0$，即 $p=-2,\;q=1$，故方程为 $y''-2y'+y=0$。

\textbf{入口五：已知非齐次方程的几个解}。先用"系数和为 $0$"的组合造齐次解，得到齐次通解，再用其中任意一个解作为特解，即得通解；若要进一步求方程，把通解形式与特征根对应起来即可。

\begin{callwarn}{易错}
\begin{enumerate}
\item 反求方程时必须保证通解中独立常数的个数等于阶数；给出 3 个解时，先判断它们之间是否线性相关。
\item 由通解反推时，$e^{\alpha x}(C_1\cos\beta x+C_2\sin\beta x)$ 对应的是一对共轭复根，$p=-2\alpha$，$q=\alpha^{2}+\beta^{2}$。
\item 变系数情形用降阶公式前，必须把 $y''$ 的系数化为 $1$，否则 $p(x)$ 取错。
\item 反求出的方程可以相差一个非零常数倍，通常约定首项系数为 $1$。
\item 已知"某个解"与已知"通解"完全不同：给一个解只能定出一个关系，必须配合其他条件才能定出全部系数。
\end{enumerate}
\end{callwarn}

\begin{callexam}{真题考法}
多以选择题或解答题的小问出现：给三个非齐次解求通解、给通解形式求方程、给一个解求通解。与解的结构定理联考时，关键是先由"系数和为 $0$"造出齐次解，再用"齐次通解加一个非齐次特解"组装答案。这类题常常不需要真正的积分计算，考的是结构关系是否清楚。

\end{callexam}

\subsection{微分方程的应用}

\subsubsection{微分方程的简单应用（几何、物理、变化率） $\bigstar$$\bigstar$$\bigstar$}

应用题的难点不在解方程，而在\textbf{把文字翻译成含导数的等式}。统一按四步走：

\begin{enumerate}
\item \textbf{设变量}：明确谁是自变量（常取 $x$ 或 $t$）、谁是未知函数，写清定义域（如 $x>0$）。
\item \textbf{找瞬时关系}：把"变化率""切线""速度""浓度"等翻译成 $y'$、$y''$、$\dfrac{dP}{dt}$ 等。关键心法是"看微小时间（或微小位移）内发生了什么变化，再除以这个微小量"。
\item \textbf{列方程并定初始条件}：一般是一阶方程（可分离或线性），初始条件由题意中某一时刻的状态给出。
\item \textbf{解方程并回到实际问题}：定出常数，回答题目所问的量（时间、温度、坐标、面积等），不要停在通解。
\end{enumerate}
\textbf{几何类常用翻译}：

\begin{itemize}
\item 切线斜率 $=y'$；过点 $(x,y)$ 的切线方程为 $Y-y=y'(X-x)$，它在 $y$ 轴上的截距为 $y-xy'$，在 $x$ 轴上的截距为 $x-\dfrac{y}{y'}$；
\item 法线斜率 $=-\dfrac{1}{y'}$（$y'\ne 0$），法线方程 $Y-y=-\dfrac1{y'}(X-x)$，它在 $y$ 轴上的截距为 $y+\dfrac{x}{y'}$；
\item 曲线上动点围成的面积 $A(x)$ 满足 $A'(x)=f(x)$（变限积分求导）；
\item 弧长 $s(x)=\displaystyle\int_{x_0}^{x}\sqrt{1+y'^{2}}\,dt$，对 $x$ 求导得 $s'(x)=\sqrt{1+y'^{2}}$；
\item 旋转体体积、侧面积这类量用积分表示后，对动端点求导即可得到含 $y$ 与 $y'$ 的方程。
\end{itemize}
\textbf{变化率类常用模型}：

\begin{itemize}
\item 牛顿冷却（加热）定律：$\dfrac{dT}{dt}=-k\left(T-T_a\right)$，$T_a$ 为环境温度；
\item 指数增长与衰变：$\dfrac{dP}{dt}=kP$（衰变时 $k<0$）；半衰期 $T$ 与衰变常数满足 $\lambda=\dfrac{\ln 2}{T}$；
\item Logistic 模型：$\dfrac{dP}{dt}=kP\left(1-\dfrac PM\right)$，$M$ 为环境容量；
\item 混合（盐水）问题：$\dfrac{dA}{dt}=$ 流入速率减流出速率，其中流入速率 $=c_{\text{in}}r_{\text{in}}$，流出速率 $=\dfrac{A(t)}{V(t)}\cdot r_{\text{out}}$；
\item 运动学：加速度 $a=\dfrac{dv}{dt}=v\dfrac{dv}{ds}$，质量为 $m$ 的物体受力 $F=m\dfrac{dv}{dt}$，阻力常写成 $-kv$ 或 $-kv^{2}$。
\end{itemize}
\textbf{例 1（冷却）}：物体温度 $T(t)$ 从 $100\,^{\circ}\mathrm{C}$ 开始下降，环境温度恒为 $20\,^{\circ}\mathrm{C}$，$10$ 分钟后温度为 $60\,^{\circ}\mathrm{C}$，求 $T(t)$。

解：由冷却定律 $\dfrac{dT}{dt}=-k(T-20)$，分离变量积分得

\[T-20=Ce^{-kt},\qquad T(0)=100\ \Longrightarrow\ C=80,\]

故 $T=20+80e^{-kt}$。由 $T(10)=60$ 得 $40=80e^{-10k}$，即 $e^{-10k}=\dfrac12$，所以

\[T(t)=20+80e^{-\frac{\ln 2}{10}t}\quad (^{\circ}\mathrm{C}).\]

若再问"何时降到 $30\,^{\circ}\mathrm{C}$"：由 $10=80e^{-kt}$ 得 $e^{-kt}=\dfrac18$，即 $kt=3\ln 2$，故 $t=30$ 分钟。

\textbf{例 2（几何，一阶线性）}：设曲线 $y=y(x)\;(x>0)$ 上任一点处的切线在 $y$ 轴上的截距等于该点横坐标的平方，求曲线方程。

解：切线在 $y$ 轴上的截距为 $y-xy'$，由题意

\[y-xy'=x^{2}\ \Longrightarrow\ y'-\frac1xy=-x .\]

这是一阶线性方程，$P(x)=-\dfrac1x$，$Q(x)=-x$，$\int P\,dx=-\ln x$：

\[y=e^{\ln x}\left[\int(-x)e^{-\ln x}\,dx+C\right]=x\left[\int(-1)\,dx+C\right]=x(C-x)=Cx-x^{2}.\]

验证：$y'=C-2x$，$y-xy'=Cx-x^{2}-Cx+2x^{2}=x^{2}$，符合题意。

\textbf{例 3（混合问题）}：容器内有 $100$ L 盐水，含盐 $10$ kg。以 $3$ L/min 的速度注入净水，同时以同样速度排出溶液（体积不变），求 $t$ 时刻容器内的含盐量 $A(t)$。

解：净水不含盐，故流入速率为 $0$；流出速率为 $\dfrac{A}{100}\cdot 3=\dfrac{3A}{100}$。于是

\[\frac{dA}{dt}=-\frac{3A}{100},\qquad A(0)=10 .\]

分离变量解得

\[A(t)=10e^{-0.03t}\ (\mathrm{kg}).\]

（可见含盐量按指数衰减，最终趋于 $0$，与实际相符。）

\textbf{例 4（运动学）}：质量为 $m$ 的物体在阻力 $F=-kv$ 作用下运动，初速度为 $v_0$，求速度 $v(t)$ 与位移 $s(t)$。

解：由牛顿第二定律 $m\dfrac{dv}{dt}=-kv$，即 $\dfrac{dv}{dt}=-\dfrac km v$，解得 $v(t)=v_0e^{-\frac km t}$。再积分得

\[s(t)=\int_0^{t}v_0e^{-\frac km \tau}\,d\tau=\frac{mv_0}{k}\left(1-e^{-\frac km t}\right),\]

当 $t\to+\infty$ 时位移趋于有限值 $\dfrac{mv_0}{k}$，符合"阻力与速度成正比"的实际情况。

\textbf{例 5（几何 + 最值，两问结构）}：设曲线 $y=y(x)\;(x>0)$ 过点 $(1,1)$，且其上任一点处的切线斜率为 $\dfrac{2x}{y}$，求该曲线方程。

解：由题意 $y'=\dfrac{2x}{y}$，即 $y\,dy=2x\,dx$，积分得 $\dfrac{y^{2}}{2}=x^{2}+C$。由 $y(1)=1$ 得 $\dfrac12=1+C$，即 $C=-\dfrac12$，故曲线为

\[y^{2}=2x^{2}-1\qquad (x>0,\ y>0).\]

这是"先解方程、再算几何量"的标准两问结构：第一问只要认准 $y'$ 的几何含义并分离变量；第二问往往把解出的表达式代入某个几何量（法线截距、面积、旋转体体积）求最值或极限。凡是几何应用题，第一步永远是把题中那句几何描述翻译成含 $y'$ 的等式，第二步才谈积分。

\begin{callwarn}{易错}
\begin{enumerate}
\item \textbf{符号方向}：冷却、衰变、流出都要写负号；"增长率与剩余量成正比"对应 $y'=ky$，若说"减少率"则 $y'=-ky$。
\item 切线截距公式别记反：由切线 $Y-y=y'(X-x)$ 令 $X=0$ 得截距 $y-xy'$；法线截距要用 $-\dfrac1{y'}$，结果为 $y+\dfrac{x}{y'}$。
\item 混合问题中"流出速率"依赖当前浓度 $\dfrac AV$，而体积 $V$ 可能随时间变化，不能当成常数。
\item 定义域要跟着题意走（如 $x>0$、$t\ge 0$），出现 $\ln$ 时尤其要说明。
\item 题目问的是实际量（时间、温度、坐标），\textbf{必须有明确的数值或表达式答案}，只给通解不得分。
\item 几何题里"截距""距离""面积"都要先用 $x,y,y'$ 表示出来再列方程，不要凭感觉写。
\end{enumerate}
\end{callwarn}

\begin{callexam}{真题考法}
解答题。两类原型反复出现：一类是"曲线满足某几何条件"（切线或法线截距、切线长、面积、弧长、旋转体体积），先列出含 $y'$ 的方程再用一阶解法；另一类是"变化率"背景（冷却、盐水浓度、人口、衰变、变力运动），先列一阶方程再解。第二问常接最值或极限，例如"求法线在 $y$ 轴上截距最小时的切点坐标"，此时要把第一问得到的 $y(x)$ 代进去求极值；也可能接 $\int_0^{+\infty}y(x)\,dx$ 这类广义积分。建模时只要抓住"某量对某变量的导数等于题中描述的变化率"，方程基本就出来了。

\end{callexam}

\subsection*{本章易错清单}
\addcontentsline{toc}{subsection}{本章易错清单}

\begin{enumerate}
\item 通解中\textbf{独立任意常数的个数必须等于方程的阶数}；$C_1x+C_2x$ 这类写法并不含两个独立常数，不能当作二阶方程的通解。
\item 分离变量时除以的因子可能取零，$g(y)=0$ 对应的\textbf{常数解必须单独检验}，漏掉就丢分。
\item 一阶线性方程必须先整理成 $y'+P(x)y=Q(x)$ 的\textbf{标准形}（$y'$ 的系数化为 $1$），否则 $P,Q$ 取错；$x$ 次数更高时要记得交换自变量。
\item 一阶线性通解公式里外层是 $e^{-\int P\,dx}$、内层是 $e^{\int P\,dx}$，\textbf{方向写反结果全错}；括号内的积分不再加常数。
\item 齐次方程换元 $v=\dfrac yx$ 后，$\dfrac{dy}{dx}=v+x\dfrac{dv}{dx}$ 中的 $x\dfrac{dv}{dx}$ \textbf{不能漏}；最后要把 $v$ 回代成 $\dfrac yx$。
\item "齐次"有两义：一阶"齐次方程"指 $\dfrac{dy}{dx}=\varphi\!\left(\dfrac yx\right)$，线性方程中的"齐次"指 $f(x)\equiv 0$，审题时先分清。
\item 不显含 $x$ 的方程 $y''=f(y,y')$ 必须用 $y''=p\dfrac{dp}{dy}$，\textbf{写成 $p'$ 是典型错误}；每积分一次补一个常数。
\item 二阶常系数齐次：二重根的解是 $(C_1+C_2x)e^{rx}$，\textbf{别漏 $x$}；共轭复根要写成 $e^{\alpha x}(C_1\cos\beta x+C_2\sin\beta x)$ 的\textbf{实值形式}。
\item 待定系数法中 $k$ 只看自由项的 $\lambda$（或 $\alpha+\beta i$）是特征方程的几重根，\textbf{与多项式的次数无关}；漏乘 $x^{k}$ 会导致系数方程组无解。
\item 设特解时多项式要\textbf{设全}（含缺项系数），三角函数要\textbf{成对设}（$\cos\beta x$ 与 $\sin\beta x$ 同时出现）；"和型"自由项要拆开分别求，各自的共振重数分别判断。
\item 解的结构判据：非齐次两个解之差是齐次解，两个解之和一般不是解；线性组合"系数和为 $0$ 是齐次解、和为 $1$ 是非齐次解"，其余情形结论均不成立。
\item 应用题要写清定义域与符号方向（衰减取负号、流出项用当前浓度），最后一定要用初始条件定出常数，并回到题目所问的实际量作答。
\end{enumerate}
\chapter*{附录一 五星考点冲刺路线}
\addcontentsline{toc}{chapter}{附录一 五星考点冲刺路线}
\markboth{附录一 五星考点冲刺路线}{附录一 五星考点冲刺路线}
\subsection{这份路线怎么用}

数学二的复习时间永远不够，所以本附录把全册考点按星级重新排了一遍，给出\textbf{先做什么、做到什么程度}。 用法：先看第一张总表确定优先级，再用"三轮时间表"安排进度，最后用"考前 7 天清单"收尾。

\subsection{冲刺优先级总表}

"达标标准"一栏写的是\textbf{能够独立做到的事}，做不到就说明这个考点还没过关，回到对应章节。

\begingroup\small\setlength{\tabcolsep}{4.0pt}
\begin{longtable}{@{}L{0.1272\linewidth}L{0.3052\linewidth}L{0.0862\linewidth}L{0.0890\linewidth}L{0.2865\linewidth}@{}}
\toprule
\textbf{优先级} & \textbf{考点} & \textbf{星级} & \textbf{建议投入} & \textbf{达标标准} \\
\midrule
\endhead
第 1 优先 & 未定式极限的计算（含泰勒与等价替换） & $\bigstar$$\bigstar$$\bigstar$$\bigstar$$\bigstar$ & 12 小时 & 任给一个 $\frac00$、$\infty-\infty$、$1^\infty$ 型极限，3 分钟内算对 \\
第 1 优先 & 一元函数积分学的计算（换元、分部、有理函数） & $\bigstar$$\bigstar$$\bigstar$$\bigstar$$\bigstar$ & 16 小时 & 不定积分与定积分常用题型不查表直接算 \\
第 1 优先 & 变限积分函数及其导数、牛顿-莱布尼茨公式 & $\bigstar$$\bigstar$$\bigstar$$\bigstar$$\bigstar$ & 8 小时 & 含 $x$ 的变限积分求导、求极限一步到位 \\
第 1 优先 & 微分中值定理与证明题（含辅助函数构造） & $\bigstar$$\bigstar$$\bigstar$$\bigstar$$\bigstar$ & 14 小时 & 见到"证明存在 $\xi$ 使…"能写出构造思路 \\
第 1 优先 & 导数定义与可导性判定 & $\bigstar$$\bigstar$$\bigstar$$\bigstar$$\bigstar$ & 6 小时 & 分段点、绝对值、抽象函数的可导性讨论不丢分 \\
第 1 优先 & 二重积分的计算（直角坐标、极坐标、换序） & $\bigstar$$\bigstar$$\bigstar$$\bigstar$$\bigstar$ & 12 小时 & 给区域能画出图并选对坐标系，5 分钟内算完 \\
第 1 优先 & 二阶常系数线性微分方程（齐次 + 非齐次） & $\bigstar$$\bigstar$$\bigstar$$\bigstar$$\bigstar$ & 8 小时 & 特征根三种情形与待定系数法全对 \\
第 2 优先 & 导数与微分的计算（隐函数、参数方程、高阶） & $\bigstar$$\bigstar$$\bigstar$$\bigstar$ & 8 小时 & 高阶导数与莱布尼茨公式能用 \\
第 2 优先 & 洛必达法则的适用条件与失效判断 & $\bigstar$$\bigstar$$\bigstar$$\bigstar$ & 4 小时 & 能说出什么时候不能用洛必达 \\
第 2 优先 & 单调性、极值、最值的判定 & $\bigstar$$\bigstar$$\bigstar$$\bigstar$ & 6 小时 & 含参数的单调性讨论完整不丢情况 \\
第 2 优先 & 定积分的性质、中值定理与定积分计算综合 & $\bigstar$$\bigstar$$\bigstar$$\bigstar$ & 8 小时 & 对称性、周期性、递推公式会用 \\
第 2 优先 & 反常积分的计算与敛散性判别 & $\bigstar$$\bigstar$$\bigstar$$\bigstar$ & 6 小时 & $p$ 积分结论与比较判别法熟练 \\
第 2 优先 & 定积分的几何与物理应用 & $\bigstar$$\bigstar$$\bigstar$$\bigstar$ & 6 小时 & 面积、体积、弧长、侧面积公式不出错 \\
第 2 优先 & 不等式证明与放缩 & $\bigstar$$\bigstar$$\bigstar$$\bigstar$ & 6 小时 & 常用不等式（$e^x\ge 1+x$、$\ln x\le x-1$ 等）能主动用出来 \\
第 2 优先 & 多元复合函数与隐函数求导 & $\bigstar$$\bigstar$$\bigstar$$\bigstar$ & 6 小时 & 抽象函数链式法则不漏项 \\
第 2 优先 & 多元函数极值与拉格朗日乘数法 & $\bigstar$$\bigstar$$\bigstar$$\bigstar$ & 5 小时 & 条件极值完整写出方程组并解出 \\
第 2 优先 & 数列极限（夹逼、单调有界、递推数列） & $\bigstar$$\bigstar$$\bigstar$$\bigstar$ & 6 小时 & 递推数列单调有界证明能独立完成 \\
第 2 优先 & 等价无穷小阶的确定与求参数 & $\bigstar$$\bigstar$$\bigstar$$\bigstar$ & 4 小时 & 已知极限定参数类题稳拿 \\
第 2 优先 & 积分中值问题与积分不等式 & $\bigstar$$\bigstar$$\bigstar$$\bigstar$ & 6 小时 & 能结合中值定理与变限积分构造证明 \\
第 3 优先 & 连续与间断点分类、闭区间连续函数性质 & $\bigstar$$\bigstar$$\bigstar$ & 4 小时 & 间断点分类不失分 \\
第 3 优先 & 微分概念、可微与可导的关系 & $\bigstar$$\bigstar$$\bigstar$ & 3 小时 & 概念判断题能排除错误选项 \\
第 3 优先 & 凹凸、拐点、渐近线 & $\bigstar$$\bigstar$$\bigstar$ & 4 小时 & 三类渐近线都会求 \\
第 3 优先 & 一阶微分方程（可分离、齐次、一阶线性） & $\bigstar$$\bigstar$$\bigstar$ & 5 小时 & 一阶线性公式法默写正确 \\
第 3 优先 & 可降阶的高阶方程与解的结构 & $\bigstar$$\bigstar$$\bigstar$ & 4 小时 & 三种可降阶形式识别正确 \\
第 3 优先 & 多元函数的最大值最小值与应用 & $\bigstar$$\bigstar$$\bigstar$ & 4 小时 & 有界闭区域上求最值步骤完整 \\
第 3 优先 & 柯西中值定理与泰勒法定理证明 & $\bigstar$$\bigstar$$\bigstar$ & 5 小时 & 能识别该用柯西还是泰勒 \\
第 4 优先 & 曲率、曲率圆与曲率半径 & $\bigstar$$\bigstar$ & 1 小时 & 记住公式会代入 \\
第 4 优先 & 二元函数极限与连续、可微判定 & $\bigstar$$\bigstar$ & 3 小时 & 知道判定可微的标准步骤 \\
第 4 优先 & 微分方程的应用 & $\bigstar$$\bigstar$ & 3 小时 & 能列方程并写出初始条件 \\
第 4 优先 & 函数图形的描绘 & $\bigstar$$\bigstar$ & 2 小时 & 知道作图五步 \\
\bottomrule
\end{longtable}
\endgroup

\subsection{三轮时间表（以 90 天为例，可按比例缩放）}

\textbf{第一轮：打地基（第 1--40 天，每天 3 小时）}

\begingroup\small\setlength{\tabcolsep}{4.0pt}
\begin{longtable}{@{}L{0.1159\linewidth}L{0.1787\linewidth}L{0.1159\linewidth}L{0.5008\linewidth}@{}}
\toprule
\textbf{顺序} & \textbf{内容} & \textbf{天数} & \textbf{说明} \\
\midrule
\endhead
1 & 函数、极限、连续 & 6 天 & 重点：极限计算与等价替换、间断点分类 \\
2 & 一元函数微分学 & 12 天 & 重点：导数定义、中值定理与证明题、导数应用 \\
3 & 一元函数积分学 & 12 天 & 重点：不定积分计算、变限积分、定积分应用、反常积分 \\
4 & 多元函数微积分学 & 6 天 & 重点：复合函数求导、极值、二重积分 \\
5 & 常微分方程 & 4 天 & 重点：一阶线性、二阶常系数、非齐次待定系数 \\
\bottomrule
\end{longtable}
\endgroup

第一轮要求：每个$\bigstar$$\bigstar$$\bigstar$$\bigstar$以上考点，\textbf{先看本册的"套路"，再自己做 3--5 道题}，错的记在本册对应考点的warning框旁。

\textbf{第二轮：补漏洞（第 41--75 天，每天 2.5 小时）}

\begin{itemize}
\item 只看每章的"星级总表"和"易错清单"：把$\bigstar$$\bigstar$$\bigstar$以下快速扫过，$\bigstar$$\bigstar$$\bigstar$$\bigstar$以上重新推一遍。
\item 每两天做一套真题（限时 180 分钟），对答案后\textbf{按考点归档错题}，而不是按套卷归档。
\item 出现连续两类题都错，说明是同一个考点的问题，回到本册对应章节重读。
\end{itemize}
\textbf{第三轮：考前冲刺（第 76--90 天，每天 2 小时）}

\begin{itemize}
\item 每三天过一遍本册"附录二：必背公式与结论速查"。
\item 重做真题中的错题，只做错题，不做新题。
\item 考前 7 天按下面的清单逐项打勾。
\end{itemize}
\subsection{考前 7 天清单}

\begin{enumerate}
\item 等价无穷小与泰勒展开表\textbf{默写一遍}（$\sin x,\ \tan x,\ \arcsin x,\ \arctan x,\ \ln(1+x),\ e^x,\ (1+x)^\alpha,\ \cos x$）。
\item 基本积分表\textbf{默写一遍}（$\frac{1}{x^2+a^2}$、$\frac{1}{\sqrt{a^2-x^2}}$、$\sec x$、$\csc x$、$\sqrt{a^2-x^2}$ 等）。
\item 定积分应用的六个公式（面积、旋转体体积、侧面积、弧长、形心、平均值）写一遍并给出适用条件。
\item 中值定理三条（罗尔、拉格朗日、柯西）的条件写全，辅助函数构造的四个常用原型默写。
\item 二阶常系数非齐次方程的特解设法表（自由项为 $P_m(x)$、$P_m(x)e^{\lambda x}$、$e^{\lambda x}(A\cos\beta x+B\sin\beta x)$ 三类）。
\item 二重积分的极坐标变换与换序步骤复述一遍。
\item 手算一遍：$\int_0^{\pi/2}\sin^n x\,dx$ 的递推（Wallis 公式）与 $\int_0^{+\infty}\frac{dx}{x^p}$ 的敛散结论。
\end{enumerate}
\chapter*{附录二 必背公式与结论速查}
\addcontentsline{toc}{chapter}{附录二 必背公式与结论速查}
\markboth{附录二 必背公式与结论速查}{附录二 必背公式与结论速查}
\subsection{等价无穷小（$x\to 0$ 时必须写清前提）}

\[\sin x\sim x,\quad \tan x\sim x,\quad \arcsin x\sim x,\quad \arctan x\sim x,\quad \ln(1+x)\sim x,\]

\[e^x-1\sim x,\quad a^x-1\sim x\ln a\ (a>0,a\ne1),\quad 1-\cos x\sim \frac{x^2}{2},\]

\[(1+x)^{\alpha}-1\sim \alpha x,\quad \sqrt[n]{1+x}-1\sim \frac{x}{n},\quad \log_a(1+x)\sim \frac{x}{\ln a}\ (a>0,a\ne1).\]

\textbf{三阶差（最容易漏，也最容易考）}：

\[x-\sin x\sim \frac{x^3}{6},\qquad \tan x-x\sim \frac{x^3}{3},\qquad x-\ln(1+x)\sim \frac{x^2}{2},\]

\[\arcsin x-x\sim \frac{x^3}{6},\qquad x-\arctan x\sim \frac{x^3}{3},\qquad \tan x-\sin x\sim \frac{x^3}{2}.\]

\begin{callwarn}{易错}
等价无穷小替换\textbf{只在乘除中无条件成立}；在加减中替换必须确认"两个无穷小不同阶"， 否则要用泰勒展开到足够阶数。典型反例：$\lim_{x\to0}\frac{\tan x-\sin x}{x^3}$， 若把 $\tan x$、$\sin x$ 都换成 $x$ 会得到 0，正确结果是 $\frac12$。

\end{callwarn}

\subsection{泰勒公式常用展开（皮亚诺余项）}

\[e^x=1+x+\frac{x^2}{2!}+\frac{x^3}{3!}+o(x^3),\qquad
\ln(1+x)=x-\frac{x^2}{2}+\frac{x^3}{3}+o(x^3),\]

\[\sin x=x-\frac{x^3}{3!}+\frac{x^5}{5!}+o(x^5),\qquad
\cos x=1-\frac{x^2}{2!}+\frac{x^4}{4!}+o(x^4),\]

\[\tan x=x+\frac{x^3}{3}+\frac{2x^5}{15}+o(x^5),\qquad
\arcsin x=x+\frac{x^3}{6}+o(x^3),\qquad \arctan x=x-\frac{x^3}{3}+o(x^3),\]

\[(1+x)^{\alpha}=1+\alpha x+\frac{\alpha(\alpha-1)}{2!}x^2+o(x^2),\qquad
\frac{1}{1-x}=1+x+x^2+o(x^2),\qquad \frac{1}{1+x}=1-x+x^2+o(x^2).\]

\subsection{求导公式}

\begingroup\small\setlength{\tabcolsep}{4.0pt}
\begin{longtable}{@{}L{0.2278\linewidth}L{0.2278\linewidth}L{0.2278\linewidth}L{0.2278\linewidth}@{}}
\toprule
\textbf{函数} & \textbf{导数} & \textbf{函数} & \textbf{导数} \\
\midrule
\endhead
$x^{\alpha}$ & $\alpha x^{\alpha-1}$ & $\tan x$ & $\sec^2 x$ \\
$a^x$ & $a^x\ln a$ & $\cot x$ & $-\csc^2 x$ \\
$\log_a x$ & $\frac{1}{x\ln a}$ & $\sec x$ & $\sec x\tan x$ \\
$\arcsin x$ & $\frac{1}{\sqrt{1-x^2}}$ & $\csc x$ & $-\csc x\cot x$ \\
$\arccos x$ & $-\frac{1}{\sqrt{1-x^2}}$ & $\ln\vert x\vert$ & $\frac1x$ \\
$\arctan x$ & $\frac{1}{1+x^2}$ & $\operatorname{arccot} x$ & $-\frac{1}{1+x^2}$ \\
\bottomrule
\end{longtable}
\endgroup

\textbf{莱布尼茨公式}：$(uv)^{(n)}=\sum_{k=0}^{n}C_n^k u^{(k)}v^{(n-k)}$。

\subsection{常用不定积分}

\[\int \frac{dx}{x^2+a^2}=\frac1a\arctan\frac xa+C,\qquad
\int \frac{dx}{\sqrt{a^2-x^2}}=\arcsin\frac xa+C,\]

\[\int \frac{dx}{x^2-a^2}=\frac{1}{2a}\ln\left|\frac{x-a}{x+a}\right|+C,\qquad
\int \frac{dx}{\sqrt{x^2\pm a^2}}=\ln\left|x+\sqrt{x^2\pm a^2}\right|+C,\]

\[\int \sqrt{a^2-x^2}\,dx=\frac x2\sqrt{a^2-x^2}+\frac{a^2}{2}\arcsin\frac xa+C,\]

\[\int \sec x\,dx=\ln|\sec x+\tan x|+C,\qquad \int \csc x\,dx=\ln|\csc x-\cot x|+C,\]

\[\int e^{ax}\sin bx\,dx=\frac{e^{ax}(a\sin bx-b\cos bx)}{a^2+b^2}+C,\qquad
\int e^{ax}\cos bx\,dx=\frac{e^{ax}(a\cos bx+b\sin bx)}{a^2+b^2}+C.\]

\subsection{定积分常用结论}

\textbf{对称性与周期性}：

\[\int_{-a}^{a}f(x)dx=\begin{cases}2\displaystyle\int_0^af(x)dx,& f\ \text{为偶函数},\\[4pt] 0,& f\ \text{为奇函数},\end{cases}
\qquad \int_a^{a+T}f(x)dx=\int_0^Tf(x)dx .\]

\textbf{区间再现}：$\displaystyle\int_a^bf(x)dx=\int_a^bf(a+b-x)dx$。

\textbf{华里士（Wallis）公式}：

\[\int_0^{\pi/2}\sin^n x\,dx=\int_0^{\pi/2}\cos^n x\,dx=
\begin{cases}\dfrac{(n-1)!!}{n!!}\cdot\dfrac{\pi}{2},& n\ \text{为偶数},\\[8pt]
\dfrac{(n-1)!!}{n!!},& n\ \text{为奇数}.\end{cases}\]

\textbf{积分中值定理}：$f$ 在 $[a,b]$ 上连续，则存在 $\xi\in[a,b]$ 使 $\int_a^bf(x)dx=f(\xi)(b-a)$。

\subsection{反常积分的敛散结论}

\begingroup\small\setlength{\tabcolsep}{4.0pt}
\begin{longtable}{@{}L{0.2954\linewidth}L{0.2954\linewidth}L{0.3376\linewidth}@{}}
\toprule
\textbf{积分} & \textbf{收敛条件} & \textbf{备注} \\
\midrule
\endhead
$\int_1^{+\infty}\frac{dx}{x^p}$ & $p>1$ & $p\le1$ 发散 \\
$\int_0^{1}\frac{dx}{x^p}$ & $p<1$ & $p\ge1$ 发散 \\
$\int_1^{+\infty}\frac{dx}{x(\ln x)^p}$ & $p>1$ & 用换元 $t=\ln x$ \\
$\int_0^{+\infty}e^{-\lambda x}dx$ & $\lambda>0$ & 值为 $\frac1\lambda$ \\
$\int_0^{+\infty}\frac{dx}{1+x^2}$ & 收敛 & 值为 $\frac{\pi}{2}$ \\
$\int_0^{+\infty}\frac{\sin x}{x}dx$ & 收敛（条件收敛） & 值为 $\frac{\pi}{2}$，数二只要求判断敛散性 \\
\bottomrule
\end{longtable}
\endgroup

\textbf{比较判别法}：若 $0\le f(x)\le g(x)$，则 $\int g$ 收敛 $\Rightarrow$ $\int f$ 收敛；$\int f$ 发散 $\Rightarrow$ $\int g$ 发散。 判敛时只需与 $\frac{1}{x^p}$（无穷区间）或 $\frac{1}{(x-a)^p}$（瑕点）比较。

\subsection{定积分的应用公式}

\begingroup\small\setlength{\tabcolsep}{4.0pt}
\begin{longtable}{@{}L{0.4708\linewidth}L{0.2288\linewidth}L{0.2288\linewidth}@{}}
\toprule
\textbf{量} & \textbf{公式} & \textbf{前提} \\
\midrule
\endhead
平面图形面积（直角坐标） & $S=\int_a^b\lvert f(x)-g(x)\rvert\,dx$ & $f,g$ 连续 \\
面积（极坐标） & $S=\frac12\int_{\alpha}^{\beta}r^2(\theta)d\theta$ & $r(\theta)\ge0$ \\
旋转体体积（绕 $x$ 轴） & $V=\pi\int_a^bf^2(x)dx$ & $y=f(x)\ge0$ \\
旋转体体积（柱壳法，绕 $y$ 轴） & $V=2\pi\int_a^bx\lvert f(x)\rvert dx$ & $0\le a<b$ \\
平行截面面积已知的立体 & $V=\int_a^bA(x)dx$ & $A(x)$ 为该处截面积 \\
弧长（直角坐标 / 参数 / 极坐标） & $\int_a^b\sqrt{1+y'^2}dx$；$\int_{\alpha}^{\beta}\sqrt{x'^2+y'^2}dt$；$\int_{\alpha}^{\beta}\sqrt{r^2+r'^2}d\theta$ & 曲线光滑 \\
旋转体侧面积（绕 $x$ 轴） & $S=2\pi\int_a^b\lvert y\rvert\sqrt{1+y'^2}dx$ & $y=f(x)$ 光滑 \\
形心横坐标 & $\bar x=\frac{\int_a^bx f(x)dx}{\int_a^bf(x)dx}$ & 平面薄片同理用二重积分 \\
函数平均值 & $\bar f=\frac{1}{b-a}\int_a^bf(x)dx$ & — \\
变力做功 & $W=\int_a^bF(x)dx$ & 力沿位移方向 \\
水压力 & $P=\int_a^b\rho g\,x\,l(x)dx$ & $x$ 为深度，$l(x)$ 为宽度 \\
引力 & $F=G\int_a^b\frac{m\,\rho\,dx}{x^2}$ & 细杆模型 \\
\bottomrule
\end{longtable}
\endgroup

\subsection{常用不等式与放缩}

\[e^x\ge 1+x\ (x\in\mathbb R),\qquad \ln x\le x-1\ (x>0),\]

\[\frac{x}{1+x}\le\ln(1+x)\le x\ (x>-1),\qquad \frac{2}{\pi}x\le\sin x\le x\ \left(0\le x\le\frac{\pi}{2}\right),\]

\[\lvert\sin x\rvert\le1,\qquad \lvert\arctan x\rvert<\frac{\pi}{2},\qquad \frac{x}{1+x^2}<\arctan x<x\ (x>0).\]

均值不等式：$\dfrac{a+b}{2}\ge\sqrt{ab}$；柯西不等式：$\left(\sum a_ib_i\right)^2\le\sum a_i^2\sum b_i^2$。

\subsection{微分方程解的形式}

\textbf{一阶线性} $y'+P(x)y=Q(x)$：

\[y=e^{-\int P(x)dx}\left(\int Q(x)e^{\int P(x)dx}dx+C\right).\]

\textbf{二阶常系数齐次} $y''+py'+qy=0$ 的特征方程 $r^2+pr+q=0$：

\begingroup\small\setlength{\tabcolsep}{4.0pt}
\begin{longtable}{@{}L{0.5138\linewidth}L{0.4319\linewidth}@{}}
\toprule
\textbf{特征根} & \textbf{通解} \\
\midrule
\endhead
两个不等实根 $r_1\ne r_2$ & $y=C_1e^{r_1x}+C_2e^{r_2x}$ \\
二重实根 $r$ & $y=(C_1+C_2x)e^{rx}$ \\
共轭复根 $\alpha\pm\beta i$ & $y=e^{\alpha x}(C_1\cos\beta x+C_2\sin\beta x)$ \\
\bottomrule
\end{longtable}
\endgroup

\textbf{二阶常系数非齐次} $y''+py'+qy=f(x)$ 的特解设法（$k$ 为 $\lambda$ 或 $\alpha\pm\beta i$ 作为特征根的重数，取 $0$ 或 $1$）：

\begingroup\small\setlength{\tabcolsep}{4.0pt}
\begin{longtable}{@{}L{0.4728\linewidth}L{0.4728\linewidth}@{}}
\toprule
\textbf{自由项 $f(x)$} & \textbf{特解形式 $y^*$} \\
\midrule
\endhead
$P_m(x)$ & $x^kQ_m(x)$ \\
$P_m(x)e^{\lambda x}$ & $x^kQ_m(x)e^{\lambda x}$ \\
$e^{\alpha x}(A\cos\beta x+B\sin\beta x)$ & $x^ke^{\alpha x}(a\cos\beta x+b\sin\beta x)$ \\
\bottomrule
\end{longtable}
\endgroup

\textbf{解的结构}：非齐次通解 $=$ 对应齐次通解 $+$ 一个非齐次特解； 非齐次的两个解之差是对应齐次方程的解；齐次解的任意线性组合仍是齐次解。

\begin{callkey}{必背}
微分方程题拿到手先做三件事：判断方程类型（可分离？线性？常系数？）、 写出通解形式再定常数、检查是否满足初始条件。 应用题别忘写初始条件——漏写初始条件是最常见的失分点。

\end{callkey}

\end{document}
